\documentclass[11pt,a4paper]{article}
\usepackage[T1]{fontenc}
\usepackage[utf8]{inputenc}
\usepackage{lmodern,microtype,amsmath,amssymb,amsthm,mathtools}
\usepackage[margin=25mm]{geometry}
\usepackage{booktabs,longtable,array,graphicx}
\usepackage[hidelinks]{hyperref}
\usepackage{enumitem}
\setlist{nosep}
\setlength{\emergencystretch}{2em}
\newtheorem{theorem}{Theorem}[section]
\newtheorem{proposition}[theorem]{Proposition}
\newtheorem{lemma}[theorem]{Lemma}
\newtheorem{corollary}[theorem]{Corollary}
\theoremstyle{definition}
\newtheorem{definition}[theorem]{Definition}
\newtheorem{example}[theorem]{Example}
\theoremstyle{remark}
\newtheorem{remark}[theorem]{Remark}
\DeclareMathOperator{\ord}{ord}
\DeclareMathOperator{\Disc}{Disc}
\DeclareMathOperator{\He}{He}
\DeclareMathOperator{\Tr}{Tr}
\DeclareMathOperator{\rank}{rank}
\DeclareMathOperator{\diag}{diag}
\DeclareMathOperator{\supp}{supp}
\DeclareMathOperator{\Rea}{Re}
\DeclareMathOperator{\Ima}{Im}
\newcommand{\workbenchurl}{https://github.com/KokunoYumeto/zeta-and-l-functions-workbench}
\newcommand{\recordurl}{https://github.com/KokunoYumeto/mathematics-workbenches}
\title{The simplicity question for the Riemann zeros after Alp\"oge's theorem}
\author{}
\date{23 November 2026}
\begin{document}
\maketitle
\begin{abstract}
Alp\"oge's zero-location theorem establishes the Riemann hypothesis and its Dirichlet generalization. Starting from that theorem, we express the remaining simplicity question through exact growth limits for finite Toeplitz and Jensen tests of the backward heat flow. On a fixed contour, Newton-sum determinants recover every endpoint multiplicity, with explicit coefficient and error bounds. We prove positive proportions of both horizontal Hurwitz directions, a common interval of simple real heat survival for every fraction below the Montgomery--Taylor proportion $0.67250\ldots$, and an isolated-pair collision expansion. In the circular unitary model the expansion transfers the extreme-gap process to a Weibull first-collision law and a Poisson early-collision process. Suzuki's inverse-compression convergence gives fixed-multiplicity recovery. Higher even trace powers give a complete finite arithmetic test at every fixed height, and a certified example reaches height $97.632$. Their integer inertia test realizes every fraction below $0.83625\ldots$ with logarithmic support. Independent contour and spectral probes agree through height $20$. A uniform pair estimate gives an all-window lower bound for the Weil form. A vanishing lower limit of the original trial ratio suffices for simplicity; its required quantitative estimate remains open. The mathematics and text were produced by GPT-6 Astra (ChatGPT Pro mode), with specified statements independently re-proved by separate instances of the same model; no human has checked this text.
\end{abstract}

\section{The question after the zero-location theorem}
Alp\"oge's theorem establishes the Riemann hypothesis and the generalized Riemann hypothesis for Dirichlet $L$-functions \cite{alpoge,lean,waterloo}. These zero-location conclusions make the earlier RH inputs available in Montgomery's two-thirds simple-zero theorem \cite[Corollary~2, p.~182]{montgomery1973} and in the $19/27$ theorem of Bui--Heath-Brown, which develops the mollified-derivative method of Conrey--Ghosh--Gonek \cite[p.~1]{bui-heathbrown2013}\cite[p.~497]{conrey-ghosh-gonek1998}. The heat family of de Bruijn and Newman \cite{debruijn1950,newman1976}, the Lehmer-pair method of Csordas--Smith--Varga \cite[Theorem~1, pp.~109--110]{csv}, Rodgers--Tao's proof of $\Lambda\ge0$ \cite[Theorem~1, p.~3]{rodgers-tao2020}, and Polymath~15's quantitative upper-bound method \cite[Theorems~1.1--1.2, pp.~2--3]{polymath2019} describe the boundary between real and nonreal heat zeros. Alp\"oge's theorem now gives $\Lambda=0$. Griffin--Ono--Rolen--Zagier prove fixed-degree Jensen hyperbolicity and its Hermite limit \cite[Theorems~1--3, p.~2]{griffin-ono-rolen-zagier2019}. We study the remaining simplicity question through finite tests, and prove quantitative statements about the zeros that those tests observe.

Before the zero-location theorem, Alp\"oge--Furman \cite[Theorems~A--B, pp.~1--2]{alpoge-furman2026} proved an unconditional simple-critical proportion $0.67250\ldots$ and distinct-zero proportion $0.83625\ldots$, also for each fixed primitive Dirichlet character. They credit discovery and writing of the argument to Claude, developed by Anthropic, and take responsibility for its verification; the earlier Claude manuscript gives the same stated proportions \cite[Theorems~A--B, pp.~1--2]{Claude2026}. Their method combines a finite Weil-form rank inequality with a prime-side second moment and an inertia calculation. The survival theorem here combines the restricted pair-correlation input with inverse-square separation estimates to obtain one common heat interval. The later gap work of Bondarenko--Heap explicitly acknowledges assistance by OpenAI's ChatGPT; its precise role and its arithmetic theorem are discussed in Section~\ref{sec:arithmetic-collisions}.

Lamzouri \cite[Theorem~1.1, p.~3; Section~2, pp.~4--8]{lamzouri-simple2026} obtains the same simple-critical and distinct proportions from a Hilbert-space inequality applied to the pair-correlation form. Knausg\aa rd \cite[Theorems~1.1 and~1.3, pp.~2--3]{knausgard2026} gives the respective distinct and simple-critical lower proportions $1645064/1960733$ and $1669159/2478195$, using inequalities for short consecutive blocks together with the arithmetic energy. The source specifies its finite-computation dependencies in Table~5 and Section~8, pp.~33--37, and credits OpenAI Codex and Anthropic Claude with parts of the argument and certificates on p.~37. These estimates refine the available density information; the absolute finite-packet thresholds below use the complete endpoint divisor. Cobeli--Zaharescu \cite[Theorem~1, p.~2]{cobeli-zaharescu2026} bound the proportion of zeros that are simple or on the critical line. Writing that union count as $N_s+N_0-N_s^0$, the zero-location theorem gives $N_0=N$ and $N_s^0=N_s$, so its value is now exactly $N$. This identity specifies how their horizontal-location question specializes here.

Mori's uniform interval form of the pair argument gives a lower
bound for the Weil form with exponent \(\mu^{6/5}\log\mu
(1+\log\log\mu)\), where \(\mu=e^{2L}\)
\cite[Theorem~4.3, pp.~11--13]{mori-rh-alone2026}.
Theorem~\ref{thm:cosine-lower-rate} applies the Montgomery--Taylor
cosine test to that uniform argument and obtains the exponent
\(\mu^{1/d_{\rm MT}}\), where \(d_{\rm MT}=0.8362503\ldots\).
The proof includes the weighted error estimate and the transfer through
Hall matching and Voss's sampling theorem. Its quantitative scale,
compared explicitly below with the prolate trial, leaves the simplicity
ratio open.

Throughout, write
\begin{align*}
 \xi(s)&=\tfrac12s(s-1)\pi^{-s/2}\Gamma(s/2)\zeta(s),
 &\Xi(z)&=\xi(\tfrac12-iz),\\
 H_0(x)&=\tfrac18\xi(\tfrac12+ix/2)=\tfrac18\Xi(x/2),
 &\partial_tH_t&=-\partial_x^2H_t.
\end{align*}
The coordinate map is $s=1/2+ix/2$; thus an ordinate $\gamma$ corresponds to $x=2\gamma$. A negative heat time is $t=-a$ with $a>0$. Every comparison below keeps these maps explicit.

\subsection{The exact size of a finite certificate}

This section expresses endpoint multiplicity by the exact asymptotic sizes of two finite tests.
We use
\[
 H_t(z)=\int_0^\infty e^{tu^2}\Phi(u)\cos(zu)\,du,
 \qquad H_0(z)=\tfrac18\xi(\tfrac12+\tfrac12iz),
 \qquad \partial_tH_t=-\partial_z^2H_t.
\]
The positive theta kernel is
\[
 \Phi(u)=\sum_{n\ge1}
 (2\pi^2n^4e^{9u}-3\pi n^2e^{5u})e^{-\pi n^2e^{4u}}.
\]
It makes the integral jointly entire in $t,z$.  For real $a\ge0$ put
\[
 F_a(w)=\frac{H_{-a}(i\sqrt w)}{H_{-a}(0)}
       =\sum_{k\ge0}b_k(a)w^k,
 \qquad
 b_k(a)=\frac{\int_0^\infty e^{-au^2}\Phi(u)u^{2k}\,du}
                   {(2k)!H_{-a}(0)}>0.
\]
The series defines $F_a$ without a choice of square root.  Write $b_k=0$ for $k<0$.
For increasing nonnegative integer lists $I=(i_1,\ldots,i_r)$ and
$J=(j_1,\ldots,j_r)$, its Toeplitz minor is
$\det[b_{j_v-i_u}]_{u,v=1}^r$.  A sequence is $PF_r$ when all such
minors of orders at most $r$ are nonnegative.  The least order of a
negative minor is $m_*(-a)$.  The least $R$ for which a negative minor
uses only indices $0,\ldots,R$ is $R_*(-a)$.

The Jensen polynomial and its detecting degree are
\[
 J_{N,a}(X)=\sum_{k=0}^N(N)_k b_k(a)X^k,
 \qquad
 j_*(-a)=\min\{N\ge1:J_{N,a}\hbox{ has a nonreal zero}\},
\]
where $(N)_k=N!/(N-k)!$.  An empty minimum is $+\infty$.
For positive real $x$ let $m_x$ denote the multiplicity of $x$ as a zero of $H_0$.
By Alp\"oge's theorem these are all its positive-half-plane zeros.
Let $h_{m,1},\ldots,h_{m,m}$ be the roots of the probabilists' Hermite
polynomial
\[
 \operatorname{He}_m(u)=
 \sum_{k=0}^{\lfloor m/2\rfloor}
 \frac{m!(-1)^k}{2^k k!(m-2k)!}u^{m-2k}.
\]
For $m\ge2$ define
\[
 D_m=\operatorname{diag}(h_{m,1},\ldots,h_{m,m}),\qquad
 \phi_m(X)=\int_{U(m)}e^{i\operatorname{Tr}(D_mUXU^*)}\,dU,
\]
where $X$ is Hermitian and Haar measure has total mass one, and set
\[
 \mathfrak r_m=\inf\{\lambda_{\max}(X)-\lambda_{\min}(X):\phi_m(X)<0\}.
\]
Set $\mathfrak r_1=+\infty$ and let $x_*$ be the least positive multiple
zero of $H_0$, with value $+\infty$ when the set is empty.

\begin{theorem}[Exact finite-detection limits]\label{thm:fd-limits}
The following limits hold in the extended positive real numbers:
\[
 \begin{split}
 \lim_{a\downarrow0}\sqrt a\,m_*(-a)
 &=\lim_{a\downarrow0}\sqrt a\,R_*(-a)
   =\min_{x:m_x\ge2}\frac{x\mathfrak r_{m_x}}{2\sqrt2},\\
 \lim_{a\downarrow0}a\,j_*(-a)&=\frac{x_*^2}{8}.
 \end{split}
\]
Every nonempty minimum in the first line is attained and is strictly
positive and finite.  Each detecting quantity is finite for every $a>0$.
The following assertions are consequently equivalent:
\[
 \begin{gathered}
 \text{every nontrivial zero of $\zeta$ is simple},\\
 \sqrt a\,m_*(-a)\longrightarrow\infty,\qquad
 \sqrt a\,R_*(-a)\longrightarrow\infty,\qquad
 a\,j_*(-a)\longrightarrow\infty.
 \end{gathered}
\]
Here each of the three limit assertions is separately equivalent to simplicity.
Moreover $\mathfrak r_2=\pi$ and $\mathfrak r_3=2\pi/\sqrt3$.
\end{theorem}

These formulas identify the arithmetic information needed to improve a
finite test beyond its universal scale.  The proof below establishes the
formulas; the assertion that their limits are infinite remains the simplicity question.
The use of coefficient positivity and the $PF$ characterization is classical
\cite{Edrei1953}.  Shi gives the heat coefficient equations and their
consecutive-minor formulation \cite{Shi2026}.  The asymptotic argument
here controls every row and column index at a prescribed minor order.


The Jensen normalization is related explicitly to that of
Griffin--Ono--Rolen--Zagier. Their coefficients $\gamma(n)$ are
specified by $8\xi(1/2+z)=\sum_{n\ge0}\gamma(n)z^{2n}/n!$.
Consequently $n!b_n(0)=\gamma(n)/(4^n\gamma(0))$ and
\[
 J_{N,0}(X)=\gamma(0)^{-1}J_\gamma^{N,0}(X/4).
\]
Their Hermite polynomial is
$2^{N/2}\operatorname{He}_N(X/\sqrt2)$.
Their fixed-degree, growing-shift limit
\cite[Theorem~3, p.~2]{griffin-ono-rolen-zagier2019} and the
negative-time, growing-degree limit here use these exact coefficient
and variable maps.


\subsection{Multiplicity on one contour}
Here is a complete local result. Let $U$ be a bounded conjugation-invariant Jordan domain with piecewise smooth boundary. Let $H(t,z)$ be jointly holomorphic near $\{0\}\times\overline U$, real for real $t,z$, and satisfy $H_t=-H_{zz}$. Its endpoint zeros in $U$ are distinct real numbers $c_1,\ldots,c_r$, with positive multiplicities $m_1,\ldots,m_r$, and it has no endpoint zero on the boundary. Set $M=\sum_i m_i$ and
\[
 S_j(t)=\frac1{2\pi i}\oint_{\partial U}z^j\frac{H_z(t,z)}{H(t,z)}\,dz,
 \qquad D_k(t)=\det[S_{u+v}(t)]_{0\le u,v<k},\quad D_0=1.
\]
\begin{theorem}[The complete local multiplicity partition]\label{thm:local-orders}
Sort the multiset $\bigsqcup_{i=1}^r\{0,1,\ldots,m_i-1\}$ into $b_1\le\cdots\le b_M$. Then
\[
 \ord_{t=0}D_k(t)=v_k:=\sum_{j=1}^k b_j\qquad(1\le k\le M).
\]
Each first difference $v_k-v_{k-1}$ is defined with $v_0=0$. The number of these differences equal to $h$ is $\#\{i:m_i\ge h+1\}$. All leading coefficients are positive and are given in Proposition~\ref{prop:leading-coefficient}.
\end{theorem}
Thus finitely many analytic germs determine every multiplicity enclosed by the contour. The statement concerns exactly that divisor. In the zeta application, Alp\"oge's theorem supplies its real centers. A theorem covering every growing contour would settle the simplicity question only after its conclusions have been established at every height.

\subsection{A complementary translation}
For $d\ge1$ let $\omega=e^{2\pi i/d}$ and put
\[
 A_d(s,q)=\frac1d\sum_{j=0}^{d-1}\zeta(s,1+\omega^j q).
\]
The Hurwitz branches use continuation from $q=0$.
\begin{theorem}[Multiplicity and monodromy of the rotation average]\label{thm:phase-orders}
At a nontrivial zero $\rho$ of multiplicity $m$, the monic polynomial of the nearby zeros of $A_d(\cdot,q)$ has discriminant of order $d(m-1)$ at $q=0$. A positive circuit of $q$ permutes the $m$ labels by $j\mapsto j+d\pmod m$. On disjoint disks containing multiplicities $m_1,\ldots,m_r$, the full discriminant has order $d(M-r)$.
\end{theorem}
For example, a triple zero and a simple zero give heat-minor orders $(0,0,1,3)$, whereas two double zeros give $(0,0,1,2)$. Their rotation-averaged full discriminants both have order $2d$. This is an exact comparison of the two statistics on the same two partitions.

\subsection{Arithmetic certificates for a bounded packet}
The inverse compressions studied by Suzuki also give a complete finite
test at a prescribed height. Here are the spaces and the exact test.
Use $\widehat f(z)=\int_{\mathbb R}f(x)e^{-izx}\,dx$ and let
$\Gamma$ be the set of distinct signed ordinates of the zeta zeros.
Its multiplicity at $\gamma$ is $m_\gamma$. The Weil form is
\[
 q(f,g)=\sum_{\gamma\in\Gamma}
 m_\gamma\widehat f(\gamma)\overline{\widehat g(\gamma)}.
\]
Sampling identifies its completion with $\ell^2(\Gamma,m)$.
Let $C$ multiply the $\gamma$ coordinate by $1/\gamma$.
For an integer $N\ge1$, let $V_N$ consist of the continuous
piecewise-linear functions supported on $[-N,N]$, with knots
at multiples of $2^{-N}$, value zero at the endpoints, and integral
zero. Let $P_N$ be its $q$-orthogonal projection and $C_N=P_NCP_N$.
Define, for $k\ge1$,
\[
 S_k=-2k[z^{2k}]\log\frac{\Xi(z)}{\Xi(0)},\quad
 D_{N,k}=S_k-\operatorname{tr}C_N^{2k},\quad
 H_{N,k}=\left(\frac2{D_{N,k}}\right)^{1/(2k)}.
\]
Theorem~\ref{thm:fd-arithmetic-complete} proves that these defects
are positive, decrease in $N$, and satisfy
\[
 D_{N,k}\downarrow2\sum_{\gamma>0}(m_\gamma-1)\gamma^{-2k},
 \qquad
 \sum_{0<\gamma\le T}(m_\gamma-1)\le\tfrac12T^{2k}D_{N,k}.
\]
With $\gamma_*=x_*/2$, its complete bounded-height assertion is
\[
 \begin{split}
 &\text{every zero with }0<\Im\rho\le T\text{ is simple}\\
 &\quad\Longleftrightarrow
   \text{some finite }N,k\text{ satisfies }D_{N,k}<2/T^{2k},
 \end{split}
 \qquad
 \sup_{N,k}H_{N,k}=\gamma_*.
\]
In particular the Jensen limit is exactly
$\tfrac12(\sup_{N,k}H_{N,k})^2$.
The matrix entries are finite prime sums and archimedean integrals,
and the scalar $S_k$ is a Taylor coefficient of $\xi$ at $1/2$.
The interval calculation in Section~\ref{sec:fd-arithmetic} gives
\[
       \text{every zero with }0<\Im\rho\le97.632
       \text{ is simple}.
\]
Proposition~\ref{prop:fd-arithmetic-separation} also gives
$|\gamma_1-\gamma_2|>6\cdot10^{-10}$ for every pair of positive
zero occurrences through height $90$.
Platt--Trudgian established simplicity through height
$3\cdot10^{12}$ by rigorous sign counting and Turing's method
\cite[abstract, p.~792; preprint Section~2, pp.~2--3]{platt-trudgian2021}.
The examples here give the arithmetic matrices and divided-mass error
budgets for the present finite criteria.
The global assertion still asks that these certificates reach every
height. For the first trace defect, the orbital estimate gives
\[
 \lim_{a\downarrow0}\sqrt a\,m_*(-a)
       \ge\frac{\pi}{\sqrt{D_{N,1}}},\qquad
 \lim_{a\downarrow0}a\,j_*(-a)\ge\frac1{D_{N,1}}.
\]

There is also an exact integer certificate. Put
\(\iota_N(T)=\operatorname{ind}_+(T H_N-G_N)\), where
\(G_N\) and \(H_N\) are the Gram and compact-primitive matrices
of \(V_N\). At a height with \(\Xi(T)\ne0\),
Proposition~\ref{prop:fd-inertia-complete} proves
\[
 \iota_N(T)\uparrow n_{\rm dist}(T),
 \qquad
 \text{the packet through }T\text{ is simple}
 \ \Longleftrightarrow\ \exists N:\ \iota_N(T)\ge n(T).
\]
Here \(n\) counts multiplicity and \(n_{\rm dist}\) counts
support points. For every \(0<\eta<d_{\rm MT}\), where
\(d_{\rm MT}=5/4-\cot(1/\sqrt2)/(2\sqrt2)=0.8362503\ldots\),
Theorem~\ref{thm:strong-inertia-density} gives constants
\(A_\eta,T_\eta>0\) such that
\[
 \iota_{\lceil A_\eta\log T\rceil}(T)
       \ge(d_{\rm MT}-\eta)n(T)\qquad(T\ge T_\eta).
\]
Its matrices have dimension
\(O_\eta(T^{A_\eta\log2}\log T)\) and prime range
\(n<e^2T^{2A_\eta}\). The exact relation to the trace test is
\(D_{N,k}=4k\int_0^\infty(n(t)-\iota_N(t))t^{-2k-1}dt\).
The required zero integer deficit at unbounded heights remains open.

The elementary multiplicity bound $m_\gamma<(3/4)\log\gamma$
for $\gamma\ge10$, combined with Platt--Trudgian's certified
height $3\cdot10^{12}$, gives
\[
 \lim_{a\downarrow0}\sqrt a\,m_*(-a)
 =\lim_{a\downarrow0}\sqrt a\,R_*(-a)>1.65\cdot10^{12},
 \qquad \lim_{a\downarrow0}a\,j_*(-a)>4.5\cdot10^{24}.
\]
Corollary~\ref{cor:fd-global-growth-numeric} also gives an explicit
two-sided comparison of these global constants. The arithmetic trace
example through $97.632$ gives $216.88$ and $4766.003712$ by the
same maps, using its own finite matrix data.
Section~\ref{sec:executed-probe} carries out the spectral projection
and integer-rounding certificate through height $20$.
The contour certificate in Section~\ref{sec:executed-contour}
independently proves the same packet simple and reaches eight simple
centers through height $45$, with a common real heat interval
$|t|\le10^{-257}$. These statements use their separately certified
error budgets.

\subsection{Directions and common continuation of Hurwitz zeros}
For a primitive Dirichlet character $\chi$ of conductor $q$, including
the conductor-one case, use the denominator translation
\[
 F_\chi(u,s)=\sum_{n\ge1}\chi(n)(n+u)^{-s}
   =q^{-s}\sum_{r=1}^q\chi(r)\zeta\left(s,\frac{r+u}{q}\right).
\]
Continuation from $u=0$ specifies the branches. At a simple zero
$\rho=1/2+i\gamma$ of $L_\chi$, its tangent is
\[
                  v_\rho=\frac{\rho L_\chi(\rho+1)}{L_\chi'(\rho)}.
\]
For zeta, let $\theta$ be the continuous Riemann--Siegel phase with
$\theta(0)=0$, set $Z(t)=e^{i\theta(t)}\zeta(1/2+it)$, and define
$S$ by the midpoint convention $N(t)=\theta(t)/\pi+1+S(t)$.
At the $n$th positive simple ordinate put
\[
 W_n=S(\gamma_n)-\pi^{-1}\arg\zeta(3/2+i\gamma_n),\qquad
 \delta_n=\arctan\frac1{2\gamma_n},
\]
where the argument is the Euler-product logarithm.
Theorem~\ref{thm:hurwitz-sign} proves the exact identity
\[
 \Re v_{\rho_n}
 =\frac{|\zeta(\rho_n+1)|}{|Z'(\gamma_n)|}
      \sqrt{\gamma_n^2+1/4}\sin(\pi W_n+\delta_n).
\]
The rightward sign is therefore exactly the positivity of this sine.
At $n=1169$ the uncorrected sine is negative and the corrected sine
positive; both signs have rigorous interval certificates.
The velocity formula is due to Frisk--de Gosson
\cite[Section~2, equation~(3), p.~2]{frisk-degosson}.
The phase identity here fixes the midpoint and exterior-argument
conventions needed to use it.

Let $N_{\chi,T}$ count the zeros in $T<\Im\rho\le2T$ with
multiplicity. The signed-density theorem proves that, for each fixed
primitive $\chi$, each strict sign of $\Re v_\rho$ has lower
proportion at least $10^{-14}$ relative to $N_{\chi,T}$.
These counts concern distinct simple centers.
There are also absolute $\kappa,c,c_1,T_0>0$ such that, for
$T\ge T_0$ and $q\le T^\kappa$, at least
$10^{-30}N_{\chi,T}$ centers of each sign have simple, pairwise
disjoint branches with ordinates in $(T,2T)$ throughout
\[
 (1+T|u|)e^{12T|\Im u|}\le c_1e^{s_T},
                  \qquad s_T=\sqrt{\tfrac12\log\log T}.
\]
The sections are holomorphic in the interior and continue locally
across the boundary; their strict sides persist for real
$0<u\le c/T$. This is Theorem~\ref{thm:uniform-deformation}.
Theorems~\ref{thm:gaussian-speed-component} and
\ref{thm:anisotropic-continuation} give separate lower counts for
logarithmic derivative and velocity sizes, and larger continuation
regions obtained by selecting larger derivatives.

For the heat family, the same arithmetic coordinate instead gives
\[
 x'(0)=\frac{H_0''(x)}{H_0'(x)},\qquad
 s'(0)=\frac{iH_0''(x)}{2H_0'(x)}.
\]
These formulas specify the exact comparison of the two tangents.
At multiple Hurwitz centers and at simple vertical tangencies,
Section~\ref{sec:motion-bending} gives the complete second jet,
including the first nonzero normal coefficient of every multiple
branch.

\subsection{Backward heat survival and the isolated-pair law}
Put $M_T=\#\{x\in(T,2T):H_0(x)=0\}$, counting multiplicity.
Put $c_{\rm MT}=3/2-\cot(1/\sqrt2)/\sqrt2=0.6725007\ldots$,
the Montgomery--Taylor simple-zero proportion.
For every $0<\eta<c_{\rm MT}$ there are $c_\eta,T_\eta>0$ such that,
for $T\ge T_\eta$, at least $(c_{\rm MT}-\eta)M_T$ distinct simple
centers have branches that remain distinct, real, simple and in
$(T,2T)$ for
\[
                    -c_\eta/(\log T)^2\le t\le0.
\]
This is Theorem~\ref{thm:two-thirds-survival}. Its arithmetic input
is the restricted pair-correlation theorem, applied after the
zero-location theorem. Corollary~\ref{cor:fd-density-detection} proves that every
selected packet has infinite Toeplitz order, window size and Jensen
detection degree throughout that common interval. Multiplication by
this packet gives explicit comparison maps from the complementary
factor to the full heat function.
At the same time, absolute $C,a_0>0$ give, for $0<a<a_0$, a simple
nonreal zero of $H_{-a}$ with
\[
 0<\Re z\le\exp(\exp(C/\sqrt a)),\qquad\Im z\le-1/2.
\]
Theorem~\ref{thm:nonreal-height} uses Dobner's uniform deformed
Dirichlet-series approximation and an explicit finite recurrence
argument.

The local collision statement applies to a real entire function $f$
of order at most one whose zeros are all real. For two adjacent
simple zeros $b\pm\delta/2$, set
\[
 G(z)=\frac{f(z)}{(z-b)^2-\delta^2/4},\quad
 L_b=\frac{G'(b)}{G(b)},\quad
 S_b=\sum_{f(\xi)=0,\,\xi\ne b\pm\delta/2}(b-\xi)^{-2},
 \quad \kappa_b=\delta^2S_b.
\]
For sufficiently small $\kappa_b$, the local continuation under
$e^{a\partial_z^2}f$ has a unique real double collision near
$(a,z)=(\delta^2/8,b-\delta^2L_b/4)$, at
\[
 a_c=\delta^2\left(\frac18+\frac{\kappa_b}{32}
                         +O(\kappa_b^2)\right),\qquad
 z_c=b-2a_cL_b+O(\delta\kappa_b^{3/2}).
\]
The constants and the neighborhood after scaling by $\delta$ are
absolute. Theorem~\ref{thm:isolated-pair} proves this statement and
the creation of a conjugate nonreal pair. The quantity $\kappa_b$
measures exactly how the other zeros enter the error.

\subsection{The circular unitary model}
Let $\theta_1<\cdots<\theta_N$ be the angles of a Haar unitary
matrix, with $\theta_{i+N}=\theta_i+2\pi$, and put
\[
 Z_N(z)=\prod_{i=1}^N2\sin\frac{z-\theta_i}{2},\quad
 Z_{N,-a}=e^{a\partial_z^2}Z_N,
 \quad \delta_N=\min_i(\theta_{i+1}-\theta_i).
\]
For $N\ge2$ its first loss of reality $A_N$ is almost surely
positive and finite. Theorem~\ref{thm:cue-collisions} proves
\[
 \frac{8A_N}{\delta_N^2}=1+O_{\mathbb P}(N^{-1/3}),\qquad
 \mathbb P\left\{\frac{8N^{8/3}A_N}{(72\pi)^{2/3}}>v\right\}
                 \longrightarrow e^{-v^{3/2}}\quad(v\ge0).
\]
The first pair starts at the smallest gap with probability tending
to one. The collision time and location process, under the same
time scaling, converges on bounded time intervals to a Poisson
process of intensity $(3/2)\sqrt v\,dv\,db/(2\pi)$.
All collisions there are double with probability tending to one,
and each creates a nonreal pair that persists for the rest of the
interval. The number of such pairs at time $v$ therefore tends to
$\operatorname{Poisson}(v^{3/2})$.
Ben Arous--Bourgade prove the required extreme-gap process
\cite[Theorem~1.1 and Corollary~1.2, pp.~2650--2651]{ben-arous-bourgade};
the heat calculation transfers their gaps to these collision times.

For the label statement fix $r\in[0,1)$ and $a\in\mathbb R$ and
choose an independent uniform index $I$. Define
\[
 \phi_{N,r}(x)=\Im\sum_l\log(1-re^{i(x-\theta_l)}),\qquad
 W_{N,i}=i-\frac{N\theta_i}{2\pi}-a-\frac{\phi_{N,r}(\theta_i)}\pi,
 \qquad \eta_{N,i}=\operatorname{sgn}\sin(\pi W_{N,i}),
\]
with the power-series logarithm, and let
$D_{j,N}=N(\theta_{I+j}-\theta_I)/(2\pi)$.
Theorem~\ref{thm:cue-labels} proves that for every fixed $j\ge1$,
$(W_{N,I}\bmod2,D_{j,N})$ tends to $(V,D_j)$, where $V$ is
uniform modulo two and independent of the $j$th positive point of
the reduced Palm sine process. Its kernel is
$k(x-y)-k(x)k(y)$, where $k(u)=\sin(\pi u)/(\pi u)$.
Thus the labels are asymptotically fair and their limiting
correlation is
$\mathbb E[1-2d_2(j-D_j)]$, with
$d_2(v)=\min_{m\in\mathbb Z}|v-2m|$.
The convergence is uniform on compact ranges of $r<1$.
The limiting correlation tends to zero as $j\to\infty$.
For adjacent labels the limiting agreement probability lies strictly
between
\[
 0.6665466018372107200719020708703
 \quad\text{and}\quad
 0.6665466018372107200719020748216.
\]
The two independent interval calculations are described in
Section~\ref{sec:cue-labels}.
Section~\ref{sec:arithmetic-collisions} states the exact gap and
isolation estimates that would transfer the collision limit to the
zeta divisor, and gives a real-zero model showing the limitation of
the proved fixed-gap bounds.

\subsection{The Weil ground-state programme}
Let $A_L$ be the positive operator of the closed Weil form on
$(-L,L)$ against the $L^2$ norm, with eigenvalues
$a_1(L)\le a_2(L)\le\cdots$ counted with multiplicity.
Write $a_e,a_o$ for its even and odd minima.
Proposition~\ref{prop:parity-comparison} proves that
$a_e<a_o$ gives a simple even ground state $e_L$ with nonzero
integral and a Fourier transform with simple real zeros. Along
every sequence $L\to\infty$ with
$\inf(1-a_e/a_o)>0$, its normalized transform tends locally
uniformly to
\[
 \frac{\widehat e_L(z)}{\widehat e_L(0)}
       \longrightarrow\prod_{\gamma>0,\,\gamma\in\Gamma}
                         (1-z^2/\gamma^2).
\]
Each distinct ordinate contributes one factor. Suzuki's
Theorem~3.3, equation~(4.1), and Remark~4.1
\cite[pp.~10--12]{suzuki} give the unshifted inverse-compression
convergence. The fixed-multiplicity recovery by a probe in
Section~\ref{sec:suzuki-probes} is a direct consequence of that
theorem: the limiting spectral mass is
$m_\gamma|\widehat\varphi(\gamma)|^2$, and division and rounding
recover $m_\gamma$. The explicit error includes the reciprocal
spectral separation and the probe's Fourier lower bound.

Let $K$ be the inverse Fourier transform of $\Xi$.
For every family of real even $L^2$-unit vectors $v_L$ in the
local form domains with $v_L\to K/\|K\|_2$, the exact implication
proved in Proposition~\ref{prop:liminf-rayleigh} is
\[
 \liminf_{L\to\infty}\frac{q(v_L)}{a_2(L)}=0
       \quad\Longrightarrow\quad
       \text{every nontrivial zeta zero is simple}.
\]
The even prolate trial of Connes--Consani--Moscovici is defined
explicitly in Section~\ref{sec:prolate-trial}; for $L=\log\lambda$
it satisfies, for every $\varepsilon>0$,
\[
 q(v_L)\ll_\varepsilon\lambda^{9+\varepsilon}e^{-4\pi\lambda^2}.
\]
An odd trial tending to $K'/\|K'\|_2$ also proves
$a_2(L)\ll_\varepsilon\lambda^{13+\varepsilon}e^{-4\pi\lambda^2}$.
Thus the remaining lower estimate has absolute polynomial exponent
$9+\alpha$, with $0<\alpha\le4$; the parameter $\alpha$ measures
the gain above the even trial's exponent $9$.
There is a weaker all-window lower estimate, proved in
Theorem~\ref{thm:cosine-lower-rate}: for some absolute \(C,L_0>0\),
\[
 a_1(L)\ge\exp\{-C e^{2L/d_{\rm MT}}L(1+\log L)\}
                 \qquad(L\ge L_0).
\]
It also bounds \(a_2(L)\) and the codimension-one minimum \(b_L\)
of Section~\ref{sec:theta-coercivity}. Its exponent
\(1/d_{\rm MT}=1.1958141\ldots\) is strictly between one and \(6/5\).
Substitution into the trial ratio leaves a growing positive term of
order \(e^{2L/d_{\rm MT}}L\log L\) in the upper-bound exponent.
Thus the lower estimate at the displayed prolate scale remains open.

The exact compact-operator example at the end of the spectral
argument shows why the commutator and compactness estimates alone
leave this quantitative step open.


\section{Proof overview}
The finite-detection problem has two simultaneous limits: the heat time
approaches zero while the minor order or polynomial degree grows.
A local Hermite expansion identifies the contribution of a multiple
endpoint. The difficulty is to prove a matching lower bound uniformly
over every row and column index, and over all the other zeros.
Canonical-product estimates control the remote factors; the
Harish--Chandra orbital integral gives the local sign at growing
Toeplitz order. For Jensen polynomials the falling-factorial operator
and its finite-degree stability symbol provide the corresponding
uniform lower bound. The resulting limits keep the multiplicity and
the height of the responsible endpoint explicit.

The arithmetic trace tests use the same divisor through another exact
map. The Taylor coefficients of $\log\Xi$ count every multiplicity,
whereas the inverse operator has one eigendirection at each distinct
ordinate. Compression decreases each even trace power. Their defect
therefore bounds the multiplicity excess. Letting the power grow
isolates the first possible multiple ordinate, proving completeness of
the finite test at every fixed height. The proof also bounds projection
loss for two nearby coordinate vectors and turns it into an explicit
separation estimate. This step uses only finite arithmetic matrices
once their entries have certified error bounds.

The integer version counts the compression eigenvalues above
\(1/T\). A graph of close zero labels supplies separated coordinates;
modulated compact tents and interpolation turn them into a positive
subspace for \(TH_N-G_N\). Keeping the test uniform over arbitrary
real intervals gives more information: a separated sampling set with a
controlled counting deficit. Hall matching transports an oversampled
grid into that set, and Voss's quantitative theorem bounds the local
Weil form from below. Its cost is determined by the interval scale
needed for the prime-side mean-square estimate.

The signed Hurwitz counts require a joint estimate for the exterior
phase and the mollified derivative. A square mollifier makes the
signed first moment real to leading order, while the first and second
absolute moments leave a strict numerical margin after the small-phase
set is removed. A further short-prime insertion proves Gaussian
restrictions of these moments. Curvature control and a logarithmic
comparison then select separated centers with large derivatives.
Rouch\'e's theorem applies to all their disks on one common parameter
region; uniform conductor estimates make the selection simultaneous
for conductors in a small power range.

The local difficulty is that individual zeros at a multiple endpoint need fractional powers of the heat time. The contour integrals are nevertheless holomorphic, because they are symmetric in the enclosed zeros. Rescaling at an endpoint produces a Hermite polynomial. Cauchy--Binet then writes each moment determinant as a sum of squared Vandermonde determinants. At positive heat time all their leading coefficients are positive. Selecting $k$ branches therefore becomes an integer minimization with no cancellation, and its first differences recover the partition.

For a finite-precision statement, one must bound both the signal and the error. Separation of the distinct endpoint centers bounds the leading coefficient from below. A zero-free contour in a complex time disk bounds the Newton sums and hence the determinant. Cauchy's formula then gives explicit coefficient and remainder errors. The same contour construction applies to a Hurwitz translation, whose first nonzero rotation-averaged coefficient never vanishes in the critical strip. Its local model is $y^m+cq^d$, from which the discriminant order and monodromy follow.

The two zero families have different defining differential identities. We transport both to $s=1/2+ix/2$ before comparing their tangent vectors. For the collision calculation, the complementary canonical product separates the close pair from all the remaining zeros. Removing its linear exponential is an exact conjugation of the heat operator; the quadratic term gives the first correction to the two-zero collision time.


In the circular unitary ensemble, a three-point determinant bound
makes extreme pairs isolated. The local collision expansion can then
be applied to the whole early-gap process of Ben Arous--Bourgade.
For the labels, the essential point is joint convergence with the gap:
a determinant estimate controls the nontrivial Fourier modes even
after a local count has been inserted. The sine-kernel Fredholm
determinants provide the adjacent-label value and its interval error.

The spectral argument compares the original and shifted Weil metrics.
An exact rank-one commutator identifies the quotient primitive, while
an even--odd relative gap bounds the change of metric. Suzuki's
inverse-compression convergence then identifies the limiting product
over distinct ordinates. A trial ratio tending to zero along one
sequence also forces the ground vector to converge to the inverse
Fourier transform of $\Xi$; comparison of these two limits gives
simplicity. The prolate estimates prove the required numerator bound.
The unresolved denominator is written as a finite matrix lower
certificate with an explicit projection error, so its remaining
arithmetic inequality is specified quantitatively.

\section{Proof of the finite-detection limits}
\label{sec:fd-main}
We prove Theorem~\ref{thm:fd-limits}, including the estimates uniform over every zero and every allowed minor index.

\subsection{The analytic estimate that controls all heights}

We first obtain a relative convolution estimate uniform near every ray.
The theta integral gives, uniformly for $t$ in a compact set,
\[
 \log\max_{|z|\le R}|H_t(z)|=O(R\log(R+3)).
\]
For example, split its integral at $u=\log(R+3)$; beyond that point
$\exp(-\pi e^{4u})$ dominates the other exponential factors.
It follows that $F_a$ has order at most $1/2$.  Hadamard factorization
therefore has genus zero and has constant zero-free factor.  At $a=0$,
\[
 H_0(z)=H_0(0)\prod_{x>0}(1-z^2/x^2)^{m_x},\qquad
 N_H(R)=\sum_{0<x\le R}m_x\ll R\log(R+3).
\]
In this subsection all sums over $x$ include multiplicity.
For nonreal $z$ define
\[
 L(z)=\frac{H_0'(z)}{H_0(z)},\qquad
 S(z)=\sum_{x>0}\left(|z-x|^{-2}+|z+x|^{-2}\right).
\]
The denominator is nonzero by Alp\"oge's theorem.  The entire remainder
\[
 Q_z(v)=\frac{H_0(z+v)}{H_0(z)e^{L(z)v}}
\]
satisfies
\begin{equation}\label{eq:fd-rem}
 |Q_z(v)|\le e^{S(z)|v|^2/2},\qquad
 |Q_z(v)-1|\le8S(z)|v|^2e^{2S(z)|v|^2}.
\end{equation}
Indeed it is the locally uniform product of
$(1+v/(z-\xi))\exp(-v/(z-\xi))$, for $\xi=\pm x$.
The inequality $|1+w|^2\le\exp(2\Re w+|w|^2)$ proves the first bound.
For the second, scale so $S=1$.  The constant and linear Taylor
coefficients are $1,0$, and Cauchy's estimate on the unit circle gives
$|Q(v)-1|\le2e^{1/2}|v|^2$ on $|v|\le1/2$.
On its complement the first bound gives
$|Q(v)-1|\le2e^{|v|^2/2}\le8|v|^2e^{2|v|^2}$.

The heat equation and the Fourier integral give the exact Gaussian identity
\[
 H_{-a}(z)=\frac1{\sqrt{4\pi a}}
             \int_{\mathbb R}e^{-u^2/(4a)}H_0(z+u)\,du.
\]
Translate the line of integration by $p=2aL(z)$.  The integrand is entire,
and its Gaussian decay on the two vertical ends dominates the preceding
order-one growth bound.  Completion of the square yields
\[
 \frac{H_{-a}(z)}{H_0(z)e^{aL(z)^2}}=\mathbb E Q_z(p+U),
 \qquad U\sim N(0,2a).
\]
Set $A=aS(z)$ and $B=a^2|L(z)|^2S(z)$.  Equation~\eqref{eq:fd-rem},
$|p+U|^2\le2|p|^2+2U^2$, and elementary Gaussian integration give
\begin{equation}\label{eq:fd-gauss}
 \left|\frac{H_{-a}(z)}{H_0(z)e^{aL(z)^2}}-1\right|
 \le16e^{16B}\left(
 \frac{4B}{\sqrt{1-16A}}+
 \frac{2A}{(1-16A)^{3/2}}\right).
\end{equation}
This estimate applies for $16A<1$.  In particular $A,B\le2^{-10}$
make its right side at most $7/16$.  To check that last rational bound,
use $e^{16B}<2$, $(1-16A)^{-1/2}<2$, and
$(1-16A)^{-3/2}<3$.

For $r=|z|$ and $|\Im z|\ge qr$, $0<q\le1$, partial summation of
$N_H(R)\ll R\log(R+3)$ gives
\begin{equation}\label{eq:fd-cones}
 S(z)\ll q^{-2}\frac{\log(r+3)}{r+1},\qquad
 |L(z)|\ll q^{-1}\log(r+3).
\end{equation}
Here is the estimate without a separation hypothesis.  By symmetry take
$\Re z\ge0$ and split the positive zeros at $r/2$ and $2r$.
Below $r/2$ both distances are at least $r/2$, so their contribution
to $S$ is at most $8N_H(r/2)/r^2$.  The paired expression
$L(z)=-2z\sum_x(x^2-z^2)^{-1}$ gives a contribution
$O(N_H(r/2)/r)$ to $L$.  In the middle interval
$|z-x|\ge qr$ and $|z+x|\ge r$, giving respectively
$2q^{-2}N_H(2r)/r^2$ and $2q^{-1}N_H(2r)/r$.
For the last interval, $\sum_{x>2r}x^{-2}\ll\log(r+3)/r$,
which bounds the corresponding contributions by a constant times this
sum and by a constant times $r$ times this sum.  A fixed neighborhood
of zero is handled directly by the convergent product.

Define
\[
 \alpha(a)=\sup_{H_{-a}(z)=0,\,\Re z>0}|\arg z|.
\]
There are no imaginary-axis zeros, since $H_{-a}(iy)>0$.
The bounds~\eqref{eq:fd-cones} and~\eqref{eq:fd-gauss} show that
$\alpha(a)=O(\sqrt a)$.  More precisely,
\begin{equation}\label{eq:fd-tail}
 \lim_{R\to\infty}\sup_{0<a\le1}
 \sup_{\substack{H_{-a}(z)=0,\,\Re z>0\\|z|\ge R}}
            \frac{|\arg z|}{\sqrt a}=0.
\end{equation}
For completeness, fix $\eta\in(0,1/2)$ and put $q=\eta\sqrt a$.
In that cone the parameters in~\eqref{eq:fd-gauss} are bounded by
\[
 C\eta^{-2}\frac{\log(r+3)}{r+1},\qquad
 C^3\eta^{-4}\frac{\log^3(r+3)}{r+1}.
\]
Both tend to zero uniformly in $0<a\le1$.  Thus at sufficiently large
radius no zero has $|\Im z|\ge\eta\sqrt a\,|z|$.
The inequality $\arcsin u\le2u$ for $0\le u\le1/2$ proves
\eqref{eq:fd-tail}.  The same proof with a sufficiently large fixed
multiple of $\sqrt a$ gives the global $O(\sqrt a)$ bound.

At an endpoint of multiplicity $m$, the substitution
$a=q^2$, $z=x+qu$ gives, after division by its leading coefficient
and $q^m$,
\[
 e^{D_u^2}u^m+O(q)
   =(i\sqrt2)^m\operatorname{He}_m(u/(i\sqrt2))+O(q).
\]
All lowest weighted Taylor terms follow from
$\partial_a^jH_{-a}=\partial_z^{2j}H_{-a}$.
The Hermite roots are distinct: their Jacobi matrix has positive
successive off-diagonal entries $\sqrt1,\ldots,\sqrt{m-1}$, so an
eigenvector is determined by its first coordinate and each eigenspace
has dimension one.  The matrix is symmetric and therefore all roots
are real.  The implicit function theorem and Rouch\'e's theorem give
\begin{equation}\label{eq:fd-localheat}
 z_j(-a)=x+i\sqrt{2a}\,h_{m,j}+O(a).
\end{equation}
Simple endpoint zeros have real simple continuations by conjugation.
Compact zero convergence and~\eqref{eq:fd-tail} now prove
\begin{equation}\label{eq:fd-anglelimit}
 \lim_{a\downarrow0}\frac{\alpha(a)}{\sqrt a}
   =\sup_x\frac{\sqrt2\,h_{m_x}}x,
 \qquad h_m=\max_j|h_{m,j}|.
\end{equation}
The supremum is a maximum whenever it is positive.  In fact
$h_m\le2\sqrt m$, whereas $m_x\le N_H(2x)\ll x\log(x+3)$,
so its terms tend to zero at infinity.

\subsection{From angles to all Toeplitz minors}

We next use the elementary Schur identities to keep control of arbitrary indices.
For the alphabet $be^{i\theta},be^{-i\theta}$,
\[
 h_d=b^d\frac{\sin((d+1)\theta)}{\sin\theta},
\]
with value $(d+1)b^d$ at $\theta=0$.  The dual Jacobi--Trudi identity
turns a Toeplitz minor of order $r$ in elementary coefficients into
a skew Schur function with diagram width at most $r$.
One direct index formula, for allowed lists $i_\ell\le j_\ell$, is
\[
 \lambda'_p=j_{r+1-p}+p-i_1-r,\qquad
 \mu'_p=i_{r+1-p}+p-i_1-r\quad(1\le p\le r).
\]
Transposing the minor and reversing both lists gives
$s_{\lambda/\mu}$.  These formulas also prove the stated width bound.
A two-variable tableau has zero weight sum when a column has more
than two boxes.  Each two-box column otherwise contributes $b^2$,
and the unpaired interval in each row contributes one of the displayed
$h_d$, with $d\le r$.  The unpaired intervals are independent because
paired boxes below occupy an initial row segment and paired boxes above
occupy a final row segment.  Thus
\begin{equation}\label{eq:fd-sector}
 (r+1)|\theta|<\pi
 \quad\Longrightarrow\quad
 (1+be^{i\theta}w)(1+be^{-i\theta}w)\text{ has }PF_r\text{ coefficients}.
\end{equation}
Coefficient convolution is Toeplitz multiplication, so Cauchy--Binet
preserves this property under products.  For a fixed minor its sum
is finite, since intermediate indices stay between its extreme row
and column indices.  Local uniform limits preserve the inequalities.

The reciprocal parameters of $F_a$ are $z^{-2}$, one for each
right-half-plane zero of $H_{-a}$, counted with multiplicity.  Their
arguments have absolute values at most $2\alpha(a)$.
The genus-zero product and~\eqref{eq:fd-sector} therefore prove
$PF_r$ whenever $2(r+1)\alpha(a)<\pi$.
Every allowed minor is then strictly positive.  Indeed there are
infinitely many linear or conjugate quadratic factors, each with
constant coefficient one and positive first coefficient.  In iterated
Cauchy--Binet, move the rightmost row index to its target one step at
a time, then the next, using a fresh factor for every step.  The
selected factor minor is its positive first coefficient; all stationary
factor minors are one.  This gives a positive term, and all other terms
are nonnegative.  Finitely many factors realize the required path,
so the same positive lower bound survives the infinite limit.
There must be infinitely many factors because otherwise $H_{-a}$
would be a polynomial, whereas its Fourier integral tends to zero on
the real axis and has positive value at zero.

It follows already that $m_*(-a),R_*(-a)\gg a^{-1/2}$.
Equation~\eqref{eq:fd-anglelimit} proves that simplicity of all endpoint
zeros gives $\sqrt a\,m_*(-a)\to\infty$.
Edrei's entire-function theorem \cite[Proposition A, pp.~86--87;
Section~3, p.~92]{Edrei1953} identifies $PF_\infty$ with all zeros of
$F_a$ being negative real.  The map $w=-z^2$ identifies this with
all-real zeros of $H_{-a}$.  Alp\"oge and Rodgers--Tao give
$\Lambda=0$, so $m_*(-a)$ is finite for $a>0$.

\subsection{The orbital constant and isolation of one cluster}

The exact constant requires a whole coalescing cluster to be treated at once.
For a nonzero traceless Hermitian $D$ with sign-symmetric spectrum,
$\phi_D(X)=\int e^{i\operatorname{Tr}(DUXU^*)}\,dU$ is real and takes
a negative value.  On the $m^2$-dimensional Euclidean space of Hermitian
matrices with norm squared $\operatorname{Tr}(X^2)$, Gaussian Fourier
integration gives
\[
 \int e^{-t\|X\|^2/2}\phi_D(X)\,dX
   =(2\pi/t)^{m^2/2}e^{-\|D\|^2/(2t)}.
\]
The right side tends to zero as $t\downarrow0$.  Nonnegativity of
$\phi_D$ everywhere would contradict its being at least $1/2$ on a
fixed ball about zero.  A negative value can be diagonalized, perturbed
to distinct eigenvalues, and translated by a scalar without changing its sign.
This proves $\mathfrak r_m<\infty$.

There is strict positivity on every compact spectral-diameter range
$c<\mathfrak r_m$.  Differentiation under the integral gives
$\Delta\phi_m=-\|D_m\|^2\phi_m$.  At a zero in a nonnegative
neighborhood, the first nonzero homogeneous Taylor term would be
nonnegative and harmonic.  Its average on a sphere is its value at
zero, hence zero, contradicting its nonvanishing.  Such a neighborhood
therefore contains no zero.  Scalar translation puts every spectrum of
diameter at most $c$ in $[0,c]$, a compact set, giving a positive minimum.
The same observation at zero gives $\mathfrak r_m>0$.

We use the classical Harish--Chandra--Itzykson--Zuber integral
\cite[(11)--(12), pp.~312--313]{Zuber2018}:
\[
 \mathcal I(A,B)=\int_{U(m)}e^{\operatorname{Tr}(AUBU^*)}\,dU
  =\left(\prod_{j=1}^{m-1}j!\right)
       \frac{\det[e^{a_i b_j}]}{\Delta(a)\Delta(b)}.
\]
Repeated eigenvalues are interpreted by analytic continuation.
For $\rho=(m-1,\ldots,0)$ the Weyl alternant gives the exact identity
\begin{equation}\label{eq:fd-weylhciz}
 \frac{s_\mu(e^{\varepsilon a_1},\ldots,e^{\varepsilon a_m})}
      {s_\mu(1^m)}
 =\frac{\mathcal I(\operatorname{diag}(a),
                   \operatorname{diag}(\varepsilon(\mu+\rho)))}
        {\mathcal I(\operatorname{diag}(a),
                   \operatorname{diag}(\varepsilon\rho))}.
\end{equation}
This follows by dividing the two alternants; their remaining Vandermonde
quotient is $s_\mu(1^m)=\Delta(\mu+\rho)/\Delta(\rho)>0$.

Suppose a fixed cluster has parameters
$b_\varepsilon e^{\varepsilon a_j(\varepsilon)}$, with positive
$b_\varepsilon$ and
$a_j(\varepsilon)=i\kappa h_{m,j}+O(\varepsilon)$, and is stable
under conjugation.  For $c<\mathfrak r_m/\kappa$, formula
\eqref{eq:fd-weylhciz} is uniformly positive for every partition of
at most $m$ rows and width $r$ with $\varepsilon r\le c$.
The numerator ranges over a compact set of matrices with eigenvalues
in $[0,c+(m-1)\varepsilon]$ and tends uniformly to the strictly
positive orbital value.  The denominator tends to one.  Conjugation
makes the quotient real.  Partitions of more than $m$ rows have value
zero.  The Littlewood--Richardson expansion of a skew Schur function
is a sum with nonnegative integer coefficients of ordinary Schur
functions whose width is no larger than that of the skew diagram;
it therefore gives $PF_r$ for the whole cluster.

In particular $\mathfrak r_m\ge\pi/h_m$.  Apply
\eqref{eq:fd-sector} to the pairs in $e^{i\varepsilon h_{m,j}}$ and
take partitions with $\varepsilon\mu$ approaching an arbitrary
spectrum of diameter less than $\pi/h_m$.  Formula~\eqref{eq:fd-weylhciz}
shows that every orbital value in this range is nonnegative.
For $m=2$ the same formula is
$\phi_2(X)=\operatorname{sinc}(\lambda_1-\lambda_2)$.
For $m=3$, factoring its exponential Vandermonde gives
\[
 \phi_3(X)=\prod_{i<j}\operatorname{sinc}
              \left(\tfrac{\sqrt3}{2}(\lambda_i-\lambda_j)\right).
\]
The first formula gives $\mathfrak r_2=\pi$.  The second is positive
below diameter $2\pi/\sqrt3$, and negative at the spectrum
$(0,L/2,L)$ for $2\pi/\sqrt3<L<4\pi/\sqrt3$.
Thus $\mathfrak r_3=2\pi/\sqrt3$.

We now justify isolation of the negative cluster from every other zero.
Let its parameters be $\xi_j=b e^{\varepsilon a_j}$, with
$a_j=id_j+O(\varepsilon)$, $\sum a_j=0$, and distinct real $d_j$.
Let $\zeta_1,\ldots,\zeta_p$ be a finite conjugation-stable prefix
with positive limiting values strictly larger than $b^0=\lim b$.
Allow repeated prefix limits, but require pairwise separation at least
$c\varepsilon$ for distinct prefix parameters.  Let the remaining
parameters $\eta_j$ have uniformly bounded $\ell^1$ norm, modulus at
most $c_0<b^0$, and positive endpoint limits.  Their generating product
$C(v)=\prod_j(1-\eta_jv)^{-1}$ then converges locally uniformly on
$|v|<1/c_0$ to the positive endpoint product.

Choose distinct $L_1>\cdots>L_m>0$ at a strictly negative orbital
value of $\operatorname{diag}(d)$ and choose $C_1>L_1$.
Put
\[
 \mu_i=\lfloor L_i/\varepsilon\rfloor,\quad
 n=\lceil C_1/\varepsilon\rceil,\quad
 \lambda=(n^p,\mu_1,\ldots,\mu_m).
\]
Both $\mu_m$ and $n-\mu_1$ exceed a positive multiple of
$1/\varepsilon$.  The cluster character divided by
$b^{|\mu|}s_\mu(1^m)$ tends to the chosen negative orbital value,
by~\eqref{eq:fd-weylhciz}.

Here are the determinant estimates that transfer its sign.
Put $q=p+m$, $\mathcal A=\{\zeta_i,\xi_j\}$ and choose
$c_0<c_1<b^0$.  Partial fractions on $|v|=1/c_1$ give the full
complete symmetric coefficients in the form
\[
 h_k=\sum_{\alpha\in\mathcal A}A_\alpha\alpha^k+g_k,
 \quad |A_\alpha|\le C_2\varepsilon^{-(q-1)},\quad
 |g_k|\le C_2\varepsilon^{-(q-1)}c_1^k.
\]
Indeed the only small denominators are pairwise top-pole differences,
each bounded below by a constant times $\varepsilon$; after removing
the poles the same circle bounds the analytic remainder.
In the $q$ by $q$ Jacobi--Trudi determinant for $s_\lambda$, the
term using all top poles is exactly
\[
 s_\lambda(\mathcal A)\prod_{\alpha\in\mathcal A}C(1/\alpha).
\]
To check this equality, factor the pole matrix as
$[\alpha^{\lambda_i-i}]_{i,\alpha}
\operatorname{diag}(A_\alpha)[\alpha^j]_{\alpha,j}$.
The Vandermonde determinants cancel the internal factors in the residues.
The displayed product is real and has a positive nonzero limit.

Write $P=\prod_i\zeta_i>0$ and $W=P^n b^{|\mu|}>0$.
Every term using at least one remainder column uses distinct top-pole
labels for the other columns, since repeated labels give identical
columns.  Sort the logarithms of their moduli and assign them to the
nonincreasing row exponents.  Replacing any top modulus by $c_1$
loses a fixed factor less than one raised to at least $\mu_m$.
The chosen cluster moduli differ from $b$ by $e^{O(\varepsilon^2)}$,
and all prefix row exponents are $n$; hence the full top weight is
$W e^{O(\varepsilon)}$.
The fixed number of residue factors contributes only a power of
$\varepsilon^{-1}$.  Thus the whole remainder is
$O(W\varepsilon^{-K}q_1^{\mu_m})$, with $q_1<1$.

Expand the Weyl numerator of $s_\lambda(\mathcal A)$ along its
first $p$ rows.  The assignment of every prefix column to these
rows gives exactly
\[
 P^n s_\mu(\xi_1,\ldots,\xi_m)
       \prod_{i,j}(1-\xi_j/\zeta_i)^{-1}.
\]
The prefix block is $P^{n+m}\Delta(\zeta)$; cancellation against
the full Vandermonde proves the formula even with coalescing prefix
parameters.  Its extra product has a positive limit.
Every other assignment exchanges a larger prefix modulus and a
smaller cluster modulus between row lengths differing by at least
$n-\mu_1$.  Sorting again bounds the error by
$O(W\varepsilon^{-K'}q_2^{n-\mu_1})$, $q_2<1$.
Division by the Vandermonde costs only another power of
$\varepsilon^{-1}$.  Both errors are exponentially small in
$1/\varepsilon$, whereas the leading negative term has magnitude
at least a positive constant times $W s_\mu(1^m)$.
Consequently $s_\lambda<0$ for all sufficiently small $\varepsilon$.
Dual Jacobi--Trudi realizes it as a Toeplitz minor of order
$\lambda_1\le n$ in the window $0,\ldots,n+q-1$.
By translating the negative orbital spectrum, its smallest eigenvalue
and the extra gap $C_1-L_1$ can both be chosen arbitrarily small.
Thus the upper limit of $\varepsilon\lambda_1$ is at most the
negative spectral diameter plus an arbitrary positive error.

\begin{proof}[Proof of the Toeplitz limits in Theorem~\ref{thm:fd-limits}]
Put $\varepsilon=\sqrt a$.  Equation~\eqref{eq:fd-localheat} gives
near each endpoint $x$, after removing the positive geometric mean,
\[
 z_j(-a)^{-2}=b_x(a)e^{\varepsilon a_{x,j}},\qquad
 a_{x,j}=-\frac{2i\sqrt2}{x}h_{m_x,j}+O(\varepsilon).
\]
Let $\mathcal L=\min_{m_x\ge2}x\mathfrak r_{m_x}/(2\sqrt2)$.
Since $\mathfrak r_m\ge\pi/h_m$ and $h_{m_x}/x\to0$, these
local constants tend to infinity with $x$.  Every nonempty minimum
is therefore attained among finitely many endpoints.

Fix $c<\mathcal L$.  Choose a large circle so that~\eqref{eq:fd-tail}
makes every exterior parameter angle smaller than
$\pi\varepsilon/(2(c+1))$.  All exterior quadratic factors are
$PF_{\lfloor c/\varepsilon\rfloor}$ for small $\varepsilon$.
There are finitely many interior endpoint clusters, and the uniform
orbital argument gives the same property for each one; a simple
cluster is a positive linear factor.  Cauchy--Binet and the positive
path argument give strict positivity of every allowed minor of that
order.  This proves both lower limits, since $R_*\ge m_*-1$.
It also proves their divergence when $\mathcal L=\infty$.

For a multiple endpoint attaining $\mathcal L$, choose a negative
orbital spectrum within any prescribed positive error of its radius.
All preceding endpoint clusters form the finite prefix in the
isolation argument; Hermite simplicity gives their
$\varepsilon$ separation.  The complementary tail hypotheses hold
uniformly: the theta growth estimate and Jensen's formula give
$N_t(R)\ll R\log(R+3)$ on compact time intervals, and partial
summation makes $\sum_{|z_j(t)|>R}|z_j(t)|^{-2}$ tend uniformly
to zero.  Compact zero convergence handles the finite intervening
annulus and gives the required strict modulus separation.
Cluster isolation now produces negative minors of order $n$ in a
window of size $n+O_x(1)$, with
$\limsup\sqrt a\,n\le\mathcal L$ plus the arbitrary error.
This proves the upper limits and the first line of the theorem.
\end{proof}

\subsection{The exact Jensen degree}

We finish the proof by controlling both the finite endpoint clusters and the remaining zeros of every Jensen polynomial.
For an entire series $F(w)=\sum b_kw^k$ with $b_k\ge0$, its normalized
Jensen approximation obeys
\begin{equation}\label{eq:fd-japprox}
 \sup_{|w|\le R}|J_N(w/N)-F(w)|
 \le\frac{R^2F''(R)}{2N}\le\frac{F(2R)}{2N}.
\end{equation}
For $k\le N$, the coefficient multiplier is
$t_{N,k}=\prod_{j=0}^{k-1}(1-j/N)$ and
$0\le1-t_{N,k}\le k(k-1)/(2N)$.
For $k>N$ the same bound holds with $t_{N,k}=0$.
Sum the coefficient bounds and use $k(k-1)\le2^k$.
This proves~\eqref{eq:fd-japprox}.
Every $F_a$ with $a>0$ has a nonreal zero by the already established
Newman boundary.  Rouch\'e on a small disk around one such zero,
disjoint from the real axis, proves that $j_*(-a)<\infty$.

We need a finite-degree stability statement, for which the exact
input is Borcea--Br\"and\'en's Theorem~2(b), printed p.~470
\cite{BB2009}.  A real operator on polynomials of degree at most $N$
preserves hyperbolicity whenever its symbol $T[(z+w)^N]$ is nonzero
for $\Im z,\Im w>0$.
For $N\ge2$ and a conjugate parameter pair $be^{\pm i\theta}$, the symbol of
$1+2b\cos\theta D+b^2D^2$ is
\[
 (z+w)^{N-2}\big((z+w)^2+2Nb\cos\theta(z+w)+N(N-1)b^2\big).
\]
The quadratic discriminant is $4Nb^2(1-N\sin^2\theta)$.
Thus $N\sin^2\theta\le1$ and $|\theta|<\pi/2$ make its roots
negative real and its symbol stable.  A positive linear factor
$1+bD$ has stable symbol $(z+w)^{N-1}(z+w+Nb)$.
For $N=1$ the quadratic-factor operator has symbol
$z+w+2b\cos\theta$, so the same conclusion holds directly.
Their successive application to $z^N$ gives
\begin{equation}\label{eq:fd-operatorj}
 F(D)z^N=z^N J_N(1/z).
\end{equation}
Finite products preserve hyperbolicity, and coefficientwise limits of
monic hyperbolic polynomials remain hyperbolic.  The latter assertion
follows directly from Rouch\'e on a disk about any proposed nonreal
limit root.  We have proved that every parameter pair satisfying
$N\sin^2\theta\le1$ guarantees hyperbolicity of the degree-$N$
Jensen polynomial of the full product.

To examine a multiple endpoint, set
\[
 G(a,w)=H_{-a}(i\sqrt w)=\sum_{k\ge0}c_k(a)w^k,
 \qquad
 \mathcal J(a,\epsilon,w)=\sum_{k\ge0}c_k(a)w^k
                  \prod_{j=0}^{k-1}(1-j\epsilon).
\]
This interpolation is jointly entire.  Indeed the theta integral gives,
on $|a|\le A$,
$|G(a,w)|\le\exp(C_A(1+|w|^{3/4}))$.
Splitting the integral at $u=\log(\sqrt{|w|}+3)$ proves the stronger
bound $\exp(C_A(1+\sqrt{|w|})\log(\sqrt{|w|}+3))$.
Cauchy's estimate consequently gives
$|c_k(a)|\le C_A^k k^{-4k/3}$.
On $|\epsilon|\le E$ the multiplier product is at most
$(1+E)^k k!$, so the interpolating series converges normally on
compact sets.  In particular
\begin{equation}\label{eq:fd-jinterp}
 \mathcal J(a,1/N,w)=H_{-a}(0)J_{N,a}(w/N).
\end{equation}

The coefficient of $\epsilon^s$ in this interpolation is a differential
operator $Q_s$ of order at most $2s$, with principal part
\[
 \frac{(-1)^s w^{2s}}{2^s s!}D_w^{2s}.
\]
To verify it, let $e_s(k)$ be the elementary symmetric function of
$0,1,\ldots,k-1$.  Its recursion
$e_s(k+1)-e_s(k)=k e_{s-1}(k)$ shows inductively that it is a
polynomial of degree $2s$ with leading coefficient $1/(2^s s!)$.
Expand it in falling factorials of $k$; these are exactly the
coefficient multipliers of $w^jD_w^j$.
The heat equation becomes
\begin{equation}\label{eq:fd-wheat}
 \partial_aG=\mathcal L G,\qquad
 \mathcal L=-4wD_w^2-2D_w.
\end{equation}
Its $r$th Taylor coefficient is $\mathcal L^rG(0,w)/r!$ and has
principal differential term $(-4w)^rD_w^{2r}/r!$.

Suppose $G(0,w)=A(w-w_0)^m+O((w-w_0)^{m+1})$,
where $w_0=-x^2$ and $A\ne0$.
Substitute $a=\gamma q^2$, $\epsilon=q^2$, $w=w_0+qu$.
A term of bidegree $(r,s)$ in $a,\epsilon$ has differential order
at most $2(r+s)$, so no power of $q$ below $m$ occurs.
At degree $m$, only its principal differential coefficient contributes.
Summing over $r+s=k$ gives
\[
 \sum_{r+s=k}\frac{(-4\gamma w_0)^r}{r!}
                 \frac{(-w_0^2/2)^s}{s!}
       =\frac{(4\gamma x^2-x^4/2)^k}{k!}.
\]
Thus the quotient extends analytically to $q=0$ and
\begin{equation}\label{eq:fd-jlocal}
 \frac{\mathcal J(\gamma q^2,q^2,w_0+qu)}{Aq^m}
   =e^{c_x(\gamma)D_u^2}u^m+O(q),\qquad
 c_x(\gamma)=4\gamma x^2-x^4/2.
\end{equation}
The error is uniform on compact sets of $u,\gamma$ by joint analyticity.
When $m\ge2$ and $\gamma>x^2/8$, the leading polynomial has a
simple nonreal pair, since it is
$(i\sqrt{2c_x(\gamma)})^m
\operatorname{He}_m(u/(i\sqrt{2c_x(\gamma)}))$.
Rouch\'e preserves this pair.  Set $N=\lceil\gamma/a\rceil$,
$q=N^{-1/2}$, and replace $\gamma$ by $aN\to\gamma$ in
\eqref{eq:fd-jlocal}.  Equation~\eqref{eq:fd-jinterp} gives
\begin{equation}\label{eq:fd-jupper}
 \limsup_{a\downarrow0}a j_*(-a)\le x_*^2/8
 \quad\text{when }x_*<\infty.
\end{equation}

The matching lower bound needs more than compact Jensen approximation.
Fix $0<\gamma_0<x_*^2/8$ and choose
$\kappa=1/(4\sqrt{\gamma_0})$.
By~\eqref{eq:fd-tail}, choose $R_0$ so that all heat zeros of modulus
at least $R_0$ have $|\arg z|\le\kappa\sqrt a$ for $0<a\le1$.
Choose $R>R_0^2$ with no endpoint zero on $|w|=R$ and at least one
endpoint zero inside.  For small complex $a$ this contour contains
a fixed number $d$ of zeros of $F_a$.
The integrals
\[
 p_k(a)=\frac1{2\pi i}\int_{|w|=R}
             w^k\frac{\partial_wF_a(w)}{F_a(w)}\,dw
 \quad(1\le k\le d)
\]
are holomorphic and are their power sums.  Newton identities give a
monic polynomial with these zeros.  Divide by its nonzero constant
coefficient to obtain $P_a(0)=1$.
Then $F_a=P_a T_a$, with $P_a$ of degree $d$ and $T_a$ holomorphic
and nonvanishing inside the circle.  This factorization is jointly
holomorphic: inside, use the Cauchy integral of the boundary quotient;
outside, division is already holomorphic.  Both factors have real
coefficients for real $a$.

The finite polynomial $P_a$ has the same lowest weighted jets as $G$.
To see this, locally write $P_a=U(a,w)G(a,w)$, where $U$ is a
nonvanishing holomorphic unit.  At an endpoint $w_0$ of multiplicity
$m$, let $A_0$ be the leading coefficient of $P_0$ there.
For $2r\le m$, the coefficient of $a^r$ in $P_a$ is
\[
 \frac{A_0(-4w_0)^r m!}{r!(m-2r)!}(w-w_0)^{m-2r}
       +O((w-w_0)^{m-2r+1}).
\]
Indeed a term using a positive power $a^j$ from $U$ has vanishing
order at least $m-2r+2j$; the $j=0$ term has exactly the displayed
principal derivative from~\eqref{eq:fd-wheat}.  Variation of its
coefficient or of $U(0,w)$ adds at least one order.

Apply the finite interpolation
$\mathcal J_P(a,\epsilon,w)=\sum_{k=0}^d b_k^{P}(a)w^k
\prod_{j<k}(1-j\epsilon)$ to $P_a$.
The preceding weighted-jet argument proves~\eqref{eq:fd-jlocal}
for $\mathcal J_P$, with $A_0$ replacing $A$, uniformly for
$0\le\gamma\le\gamma_0$.
At a multiple endpoint, $x\ge x_*$ gives
$c_x(\gamma)<0$ uniformly on this compact interval; its limiting
polynomial is a real Hermite polynomial with simple separated roots.
At a simple endpoint the limiting polynomial is $u$.
Uniform Rouch\'e disks about these real roots each contain exactly
one root; real coefficients force that root to be real.  Their
$w$ coordinates remain negative.  The finitely many endpoint
multiplicities sum to $d$, so all roots of the degree-$d$ finite
interpolation have been counted.  Hence there is $N_1>d$ such that
\[
 N\ge N_1,\quad 0\le aN\le\gamma_0
 \quad\Longrightarrow\quad
 J_N(P_a)\text{ has only real negative zeros}.
\]

For the remaining finitely many $N<N_1$, strictness at $a=0$ is
needed.  Factor $P_0(w)=\prod_{i=1}^d(1+\beta_iw)$ with
$\beta_i>0$ and apply $1+\beta_iD$ successively to $z^N$.
The nonzero roots at each stage are negative and simple, with a
possible root at zero.  In fact, for a polynomial whose zero root
has multiplicity $v$ and whose other roots are simple negative
$\xi_1<\cdots<\xi_s$,
\[
 \frac{p'}p(z)=\frac vz+\sum_{j=1}^s\frac1{z-\xi_j}
\]
is strictly decreasing on each real complementary interval.
The equation $p'/p=-1/\beta$ has one simple solution to the left
of the smallest root and between consecutive distinct roots,
including the interval ending at zero when $v>0$.
No solution lies to the right of the largest root.  Old nonzero
roots cease to be roots; zero's multiplicity decreases by one.
This proves the induction and shows that every $J_N(P_0)$ has
only simple negative roots.  Coefficient continuity preserves them
for small real $a$.  Intersect these finitely many neighborhoods.
We have proved head hyperbolicity for every $N\le\gamma_0/a$.

Finally every conjugate parameter pair of $T_a$ has argument
$|\theta|\le2\kappa\sqrt a$ and therefore
\[
 N\sin^2\theta\le4\gamma_0\kappa^2=1/4
 \quad (N\le\gamma_0/a).
\]
Every real tail parameter is positive.  The already proved
finite-degree operator criterion shows that $T_a(D)$ preserves
hyperbolicity in degree at most $N$, first for finite tail products
and then by coefficientwise limits.  Apply it to the hyperbolic
polynomial $P_a(D)z^N$ and use~\eqref{eq:fd-operatorj}.
This gives hyperbolicity of $J_{N,a}$ for every $N\le\gamma_0/a$.
Thus
$\liminf_{a\downarrow0}a j_*(-a)\ge\gamma_0$.
Let $\gamma_0\uparrow x_*^2/8$ and combine with
\eqref{eq:fd-jupper}; for $x_*=\infty$ let $\gamma_0$ grow without
bound.  This proves the Jensen assertion and completes the proof
of Theorem~\ref{thm:fd-limits}.

\subsection{An explicit small-time certificate}

We give rational constants for the preceding lower bounds.
A direct ball evaluation of the entire completed function gives
\[
 0<p_0:=F_0'(0)=\frac{\xi''(1/2)}{8\xi(1/2)}<\frac12,
 \qquad
 p_0=0.00577624827885474269723345260758\ldots .
\]
The decimal is accompanied by an outward ball enclosure in the
\href{\workbenchurl}{\emph{Detection constants} certificate}.
The inequality, together with $\cos v\ge1-v^2/2$, proves
$H_0(x)\ge H_0(0)(1-p_0x^2)>0$ for $|x|\le1$.
Thus $x_1>1$ and $\sum_xx^{-2}=p_0<1/2$.

Hasanalizade--Shen--Wong's bound
\cite[Corollary~1.2, (1.5), p.~2]{HSW2021} is
\[
 \left|N_\zeta(T)-\frac{T}{2\pi}\log\frac{T}{2\pi e}\right|
 \le0.1038\log T+0.2573\log\log T+9.3675\quad(T\ge e).
\]
It gives $N_\zeta(e)<10$ and $N_\zeta(T)\le6T\log(T+3)$
for $T\ge e$.  Since $N_H(X)=N_\zeta(X/2)$, one obtains the
convenient bound $N_H(X)\le100X\log(X+3)$ for every $X>0$.
For $X<1$ the count is zero; for $1\le X<2e$ it is less than ten;
for $X\ge2e$ the preceding estimate is stronger.

For $r\ge1/2$, the three ranges in the proof of~\eqref{eq:fd-cones}
give the following respective upper bounds; write $\ell=\log(r+3)$:
\[
\begin{array}{c|cc}
 \text{range}&S&|L|\\ \hline
 x<r/2&400\ell/r&(400/3)\ell\\
 r/2\le x\le2r&800q^{-2}\ell/r&800q^{-1}\ell\\
 x>2r&3200\ell/r&(3200/3)\ell.
\end{array}
\]
For the last row, partial summation gives
$\sum_{x>2r}x^{-2}\le400\ell/r$.
Thus $S\le13200q^{-2}\ell/(r+1)$ and
$|L|\le2000q^{-1}\ell$.
For $r\le1/2$, the product and $p_0<1/2$ give
$S<4$ and $|L|\le4r/3$.
Both cases fit~\eqref{eq:fd-cones} with $C=2^{14}$.
The elementary inequalities
$\ell/(r+1)\le2$ and $\ell^3/(r+1)\le16$ give
$S\le2^{15}q^{-2}$ and $|L|^2S\le2^{46}q^{-4}$.
For the second inequality, use $r+1\ge(r+3)/3$ and the maximum
$\max_{y>0}(\log y)^3/y=27/e^3$ on $y\ge1$.
Taking $q=2^{16}\sqrt a\le1/2$ in~\eqref{eq:fd-gauss} proves
\[
 \alpha(a)\le2^{17}\sqrt a\qquad(0<a\le2^{-34}).
\]
Consequently, for every integer $r\ge1$, all allowed Toeplitz minors
of orders at most $r$ are strictly positive on
\[
 -\frac{2^{-36}}{(r+1)^2}\le t\le0.
\]
For $0<a\le2^{-38}$ this yields
\[
 m_*(-a)\ge\left\lfloor\frac{2^{-18}}{\sqrt a}\right\rfloor,
 \qquad
 R_*(-a)\ge\left\lfloor\frac{2^{-18}}{\sqrt a}\right\rfloor-1.
\]
The Jensen operator criterion gives, on the same time range,
\[
 j_*(-a)>\left\lfloor\frac{2^{-36}}a\right\rfloor.
\]
Indeed $|\theta|\le2\alpha(a)\le2^{18}\sqrt a\le1/2$,
so $N\sin^2\theta\le1$ throughout the displayed degree range.

\subsection{Computing the Hermite orbital radii}

The orbital constants can be reduced to finitely many one-variable equations.
Put
\[
 P_m(z)=i^m\operatorname{He}_m(z/i),\qquad
 S_m(x)=\sum_{j=1}^m\frac{e^{ih_{m,j}x}}{P_m'(ih_{m,j})},
 \qquad
 W_{m,k}(x)=\det[S_m^{(p+q)}(x)]_{p,q=0}^k.
\]
Then
\begin{equation}\label{eq:fd-wradius}
 \mathfrak r_m=
 \min_{0\le k\le\lfloor(m-2)/2\rfloor}
        \inf\{x>0:W_{m,k}(x)=0\}.
\end{equation}
Here a first-zero infimum is $+\infty$ when there is no positive zero.
To prove the formula, let $E_m$ be the real solution space of
$P_m(D)f=0$.  Lagrange interpolation gives
$S_m^{(j)}(0)=0$ for $j<m-1$ and $S_m^{(m-1)}(0)=1$.
Thus $u_j=S_m^{(m-1-j)}$, $0\le j\le m-1$, is a canonical basis
at zero: $u_j^{(i)}(0)=0$ for $i<j$ and $u_j^{(j)}(0)=1$.
The space is invariant under translations and reflections.
The confluent HCIZ formula says that $\phi_m$ vanishes at a spectrum
with distinct locations $t_\ell$ and multiplicities $\nu_\ell$ precisely
when the matrix with columns
\[
 \left((ih_{m,j})^r e^{ih_{m,j}t_\ell}/r!\right)_{j=1}^m,
 \qquad 0\le r<\nu_\ell,
\]
is singular.  This is the evaluation map from $E_m$ to the corresponding
Hermite jets.  Singularity is equivalent to a nonzero real solution
having at least $m$ zeros there, counted with multiplicity.
Consequently the critical interval length for the extended Chebyshev
property of $E_m$ is exactly $\mathfrak r_m$: below it all these
orbital values are strictly positive; beyond it a radial path from
zero to a negative value crosses a singular evaluation matrix.
At the endpoint take a convergent sequence of spectra with negative
values and decreasing diameters; its limit is a zero value, since a
negative limiting value could be radially contracted.

Carnicer--Mainar--Pe\~na's critical-length theorem
\cite[Section~3, Proposition~3, p.~6]{Carnicer2023} now applies.
It identifies the critical length of a reflection-invariant space with
the first positive zero among $W(u_j,\ldots,u_{m-1})$,
$j>(m-1)/2$.  Reversing its columns and putting $k=m-1-j$
gives exactly~\eqref{eq:fd-wradius}.
Also $\sin(h_mx)\in E_m$ has $m$ zeros on
$[0,(m-1)\pi/h_m]$, so
\[
 \frac\pi{h_m}\le\mathfrak r_m\le\frac{(m-1)\pi}{h_m}.
\]
For example $S_2=\sin x$ and
$S_3=(1-\cos(\sqrt3x))/3$.  The latter first zero is tangential,
which is why a sign-change search on the scalar functions alone
would be insufficient.

There is a rigorous algorithm for arbitrary rational accuracy in
\eqref{eq:fd-wradius}.  At zero,
\[
 W_{m,k}(x)=c_{m,k}x^{d}(1+O(x)),\quad
 d=(k+1)(m-1-k),\quad
 c_{m,k}=(-1)^{k(k+1)/2}\prod_{j=0}^k\frac{j!}{(m-1-j)!}.
\]
This follows by inserting the leading monomials in the determinant
and evaluating their Vandermonde.  Divide by $c_{m,k}x^d$.
The resulting entire function has value one at zero and computable
Taylor coefficients.  If
$W_{m,k}(x)=\sum_J A_Je^{b_Jx}$ is its Cauchy--Binet expansion,
then $J$ ranges over $(k+1)$-subsets,
\[
 A_J=\frac{\Delta((ih_{m,j})_{j\in J})^2}
             {\prod_{j\in J}P_m'(ih_{m,j})},\qquad
 b_J=\sum_{j\in J}ih_{m,j}.
\]
Writing $B=\sum_J|A_J|$ and $M=(k+1)h_m$, the degree-$N$
Taylor remainder of the normalized function on $|x|\le R$ is at most
\[
 \frac{BM^d}{|c_{m,k}|}
      \frac{(MR)^{N+1}e^{MR}}{(N+d+1)!},
\]
and its derivative is bounded by
$BM^{d+1}e^{MR}/|c_{m,k}|$.
A rational mesh therefore certifies positivity on every rational
$[0,c]$ with $c<\mathfrak r_m$.
Independently enumerate distinct rational spectra and certify a
negative orbital value by HCIZ.  Every rational
$C>\mathfrak r_m$ admits such a spectrum of diameter less than $C$,
by continuity.  Dovetail the two searches until their interval width
is smaller than the requested error.  Both searches terminate, including
when a scalar Wronskian has a tangential first zero.

Novak's uniform nonvanishing theorem supplies the dependence on $m$
\cite[Theorem~1.2, p.~6]{Novak2022}.  It gives an absolute
$\epsilon_0>0$ such that
$\int_{U(m)}\exp(zm\operatorname{Tr}(AUBU^*))\,dU\ne0$
for $|z|\le\epsilon_0$ and spectra of $A,B$ in the unit disk.
For a Hermitian $X$ of diameter $R$, take
$A=2(X-cI)/R$, $B=D_m/h_m$ and $z=iRh_m/(2m)$,
where $c$ is the spectral midpoint.  The integral is exactly $\phi_m(X)$.
It is real, starts at one on the path $uX$, and never vanishes when
$R\le2\epsilon_0m/h_m$.  Therefore
\[
 \frac{2\epsilon_0m}{h_m}\le\mathfrak r_m
                 \le\frac{(m-1)\pi}{h_m}.
\]
The Hermite Jacobi matrix gives
$\sqrt{m-1}\le h_m\le2\sqrt{m-1}$, the lower bound following
from its mean squared eigenvalue.
Its last $k$ coordinates have off-diagonal weights at least
$\sqrt{m-k}$.  The positive sine eigenvector of the unweighted
$k$-vertex path gives
$h_m\ge2\sqrt{m-k}\cos(\pi/(k+1))$.
Taking $k\to\infty$ with $k/m\to0$ proves $h_m/(2\sqrt m)\to1$.
In particular
\[
 \epsilon_0\sqrt m\le\mathfrak r_m\le\pi\sqrt{m-1},\qquad
 \epsilon_0\le\liminf\frac{\mathfrak r_m}{\sqrt m}
 \le\limsup\frac{\mathfrak r_m}{\sqrt m}\le\frac\pi2.
\]
The source theorem gives no numerical value of $\epsilon_0$ used here;
the rational small-time certificate above is independent of that constant.

\section{The finite unitary character associated with the coefficients}
\label{sec:fd-unitary}
This section gives the exact change of variables that turns coefficient positivity into positivity on a fixed unitary group.
For \(a\ge0\) put
\[
 p(a)=F_a'(0),\qquad E_a(w)=F_a(w/p(a)),\qquad
 \Psi_a(z)=E_a(z-1)=\sum_{k\ge0}a_k(a)z^k.
\]
Thus \(E_a(0)=E_a'(0)=1\) and \(\Psi_a(1)=1\).
For a unitary \(d\) by \(d\) matrix \(U\) with eigenvalues
\(z_1,\ldots,z_d\), define
\[
 \chi_{a,d}(U)=\prod_{\nu=1}^d\Psi_a(z_\nu).
\]
A normalized character here means a continuous central function of
positive type with value one at the identity: for every finite list
\(U_i\), its matrix \([\chi(U_i^{-1}U_j)]_{i,j}\) is positive
semidefinite.

\begin{proposition}[Positivity at every fixed unitary rank]
\label{prop:fd-unitary}
For each \(d\ge1\) there is \(\eta_d>0\) such that \(\chi_{a,d}\)
is a normalized character for \(0\le a\le\eta_d\). Its expansion is
\[
 \chi_{a,d}(U)=\sum_{\ell(\lambda)\le d}
   c_{\lambda,d}(a)s_\lambda(U),\qquad
 c_{\lambda,d}(a)=\det[a_{\lambda_i-i+j}(a)]_{i,j=1}^d>0.
\]
Here \(\lambda\) is a partition padded with zeros, \(s_\lambda\)
is its irreducible polynomial character, and \(a_k=0\) for \(k<0\).
The expansion converges absolutely and uniformly, and
\[
 \sum_{\ell(\lambda)\le d}c_{\lambda,d}(a)s_\lambda(1^d)=1.
\]
For each \(a>0\), some finite rank fails positive type. If \(d_*(-a)\)
is the first such rank, then \(d_*(-a)\to\infty\) as \(a\downarrow0\).
Near zero it is exactly the least negative Toeplitz-minor order
of the translated coefficients \((a_k(a))_{k\ge0}\).
\end{proposition}

\begin{proof}
The genus-zero product takes the form
\[
 E_a(w)=\prod_j(1+\beta_j(a)w),\qquad
 \beta_j(a)=\frac1{p(a)z_j(-a)^2},\qquad
 \sum_j\beta_j(a)=1.
\]
There is one parameter for each right-half-plane heat zero, with
multiplicity. At \(a=0\) these parameters are positive and there
are infinitely many. Their sum is one, so their maximum is
strictly less than one. Choose a fixed \(b<1\) larger than that
maximum. The function \(E_0\) has no zero on \(|w|\le1/b\);
compact convergence gives the same property for small \(a\).
Thus every \(|\beta_j(a)|<b\).

Translation now gives an exact parameter map:
\[
 \Psi_a(z)=E_a(-1)\prod_j(1+\alpha_j(a)z),\qquad
 \alpha_j(a)=\frac{\beta_j(a)}{1-\beta_j(a)}.
\]
The real parameters \(\beta_j\) are positive near zero, their
factors \(1-\beta_j\) are positive, and conjugate parameters
give the positive product \(|1-\beta_j|^2\). Therefore
\(E_a(-1)>0\). The \(\alpha_j\) are summable in absolute value.
For \(\beta=re^{i\theta}\), \(r\le b\), differentiation of the
argument along the circular arc gives
\[
 \left|\arg\frac{\beta}{1-\beta}\right|
 \le\frac{|\theta|}{1-b}.
\]
Indeed the derivative of \(\log(\beta/(1-\beta))\) with respect
to \(\theta\) is \(i/(1-re^{i\theta})\), of modulus at most
\((1-b)^{-1}\). The angle at \(\theta=0\) is zero.
The uniform heat-angle estimate makes these translated angles
tend uniformly to zero. For each fixed \(d\), the pairwise
Schur criterion~\eqref{eq:fd-sector} and Cauchy--Binet now give
\(PF_d\) for \((a_k(a))\) on a small closed interval. The
positive path term proves strict positivity at every allowed
index, since there are still infinitely many factors.

We give the precise character calculation. For any entire
\(\Psi(z)=\sum a_kz^k\), multiply \(\prod_i\Psi(z_i)\) by the
Weyl alternant \(\det[z_i^{d-j}]\). Cauchy--Binet expands the
resulting determinant in alternants with decreasing exponents
\(\lambda_j+d-j\). Division by the original alternant gives
\[
 \prod_i\Psi(z_i)=
 \sum_{\ell(\lambda)\le d}
       \det[a_{\lambda_i-i+j}]_{i,j=1}^d s_\lambda(z_1,\ldots,z_d).
\]
The Weyl interpretation is
\cite[Theorem~5.22.1, p.~129]{Etingof2011}; the compact-group
positive-type convention is \cite[Section~4.1, p.~56]{Faraut2007}.
This identity is initially a degree-wise identity of power series.
To justify it uniformly on the unit torus, choose \(R>1\) and
let \(M_R=\max_{|z|=R}|\Psi(z)|\). Cauchy's estimate gives
\[
 \left|\det[a_{\lambda_i-i+j}]\right|
 \le d!M_R^d R^{-|\lambda|}.
\]
The sum of the indices in each determinant term is \(|\lambda|\);
terms with negative indices vanish. The Weyl dimension formula
bounds \(s_\lambda(1^d)\) by a polynomial in \(|\lambda|\), and
there are at most \((n+1)^d\) partitions of \(n\) with at most
\(d\) parts. Hence the series converges absolutely and uniformly.
Its value at coincident eigenvalues follows by continuity.

For the heat coefficients the determinant is an allowed Toeplitz
minor: its row indices are \(i-\lambda_i\), its column indices
are \(j\), and a common integer shift makes them nonnegative.
It is consequently positive. Evaluation at \(U=I\) gives the
displayed sum of probability weights, and a positive mixture
of normalized irreducible characters is of positive type.

We also need the converse at a fixed rank. For an entire real
\(\Psi\) with \(a_0>0\) and \(\Psi(1)=1\), positive type makes
every irreducible coefficient in the preceding expansion
nonnegative, by compact-group orthogonality, equivalently by
\cite[Section~4.1, p.~56]{Faraut2007}. Padding a partition
of length at most \(k\le d\) by zeros multiplies its determinant
by \(a_0^{d-k}\). Thus the specialization \(h_n\mapsto a_n/a_0\)
is nonnegative on every ordinary Schur function of length at most
\(d\).

An allowed minor of order \(k\), with row indices \(i_\ell\)
and column indices \(j_\ell\), becomes the skew Schur value
\(a_0^k s_{\lambda/\mu}\) under this specialization, where
\[
 \lambda_p=j_{k+1-p}+p-i_1-k,\qquad
 \mu_p=i_{k+1-p}+p-i_1-k,\qquad 1\le p\le k.
\]
Transposition and reversal of both lists prove the identity by
ordinary Jacobi--Trudi. The Littlewood--Richardson expansion
of \(s_{\lambda/\mu}\) has nonnegative coefficients and only
partitions contained in \(\lambda\), hence of length at most \(k\).
All these values are nonnegative. Forbidden minors vanish by
upper triangularity. This proves the exact equivalence
\[
 \chi_d\text{ is of positive type on }U(d)
 \quad\Longleftrightarrow\quad (a_k)_{k\ge0}\in PF_d
 \qquad(a_0>0).
\]
It also proves the minimum-order assertion, since \(E_a(-1)>0\)
near zero.

To finish, suppose that all ranks were of positive type for
some \(a>0\). Rank one gives nonnegative Taylor coefficients.
If \(k\) is the first nonzero index, write \(\Psi_a(z)=z^k B(z)\).
The products of \(B\) remain of positive type at every rank,
because they are multiplied by the one-dimensional unitary
character \(\det(U)^{-k}\). Now \(B(0)>0\), so the proved
equivalence makes its coefficient sequence \(PF_\infty\).
Edrei's theorem makes every zero of \(B\), and therefore every
zero of \(\Psi_a\), nonpositive real. But \(H_{-a}\) has a
nonreal zero \(z\) with \(\Re z\ne0\), and
\(1-p(a)z^2\) is a nonreal zero of \(\Psi_a\). This contradiction
proves finiteness of \(d_*(-a)\). Positivity at rank \(d+1\)
restricts to positivity at rank \(d\), since
\(\chi_{a,d+1}(U\oplus1)=\chi_{a,d}(U)\). The intervals already
proved for every fixed \(d\) then give its divergence.
\end{proof}

The formulas specify the full relation with the earlier Toeplitz
test: reciprocal endpoint parameters first undergo
\(z^{-2}\mapsto z^{-2}/p(a)\), and then
\(\beta\mapsto\beta/(1-\beta)\). The character determinant uses
the coefficients after both operations. All claims about the
untranslated detecting order use \(F_a\), as in
Theorem~\ref{thm:fd-limits}.

\section{What derivative moments would have to prove}
\label{sec:fd-moments}
This section turns derivative moments into an exact finite inequality and records the quantifiers needed for simplicity at every height.
For a finite multiset $\mathcal Z$ of nontrivial zeta zeros, with each zero
repeated according to its analytic multiplicity, put
\[
 u_\rho=|\zeta'(\rho)|^2,\quad
 s_k=\sum_{\rho\in\mathcal Z}u_\rho^k\ (k\ge1),\quad
 s_0=|\mathcal Z|.
\]
Let $M_0$ be the sum of the multiplicities of the multiple zeros in
$\mathcal Z$.  It is zero or an integer at least two.
Let $r$ be the number of distinct positive values among the $u_\rho$.
Define the Christoffel variational quantity
\[
 \lambda_d=\min_{P\in\mathbb R[X],\,\deg P\le d,\,P(0)=1}
                  \sum_{\rho\in\mathcal Z}P(u_\rho)^2.
\]
This is the classical Christoffel function at the origin for the finite
measure $\sum_\rho\delta_{u_\rho}$.
The variational and inverse-moment-matrix descriptions are classical
\cite{LasserrePauwels2019}; the same operation appears in the discussion
of moment-based inertia counts in Claude's simple-zero paper
\cite[Section~7.5(d),(f)]{Claude2026}.
The finite derivative measure here permits the following exact endpoint formula.

\begin{proposition}[A finite moment certificate for simplicity]
\label{prop:fd-moment}
For $0\le d\le r$, with the determinant of an empty matrix taken as one,
\[
 \lambda_d=
 \frac{\det[s_{i+j}]_{i,j=0}^{d}}
      {\det[s_{i+j}]_{i,j=1}^{d}},\qquad
 \lambda_r=M_0.
\]
The denominator is positive.  The sequence $\lambda_d$ is nonincreasing,
and $\lambda_d\ge M_0$ for all $d$.  The zeros in $\mathcal Z$ are
all simple precisely when one of these finite tests has $\lambda_d<2$.
At degree one,
\[
 \lambda_1=s_0-\frac{s_1^2}{s_2},
\]
whenever $r\ge1$; this is the usual Cauchy--Schwarz simple-zero count.
\end{proposition}

\begin{proof}
Write the measure as
$\mu=M_0\delta_0+\sum_{j=1}^r w_j\delta_{v_j}$,
where the $v_j>0$ are distinct and $w_j$ are positive integers.
Every admissible polynomial has squared integral at least $M_0$.
The polynomial $P(X)=\prod_{j=1}^r(1-X/v_j)$ attains this lower
bound and has degree $r$.  Enlarging the admissible polynomial space
makes the minima nonincreasing.
For $d\le r$, the lower moment block
$C=[s_{i+j}]_{i,j=1}^d$ is positive definite: its quadratic form is
$\sum_jw_j(\sum_{i=1}^d c_i v_j^i)^2$.
A nonzero polynomial $XQ(X)$ of degree at most $d$ cannot vanish at
$r\ge d$ distinct positive points in addition to zero.
For $P=1+\sum_{i=1}^d c_iX^i$, the objective is
$s_0+2b^Tc+c^TCc$, where $b=(s_1,\ldots,s_d)^T$.
Its minimum is $s_0-b^TC^{-1}b$.
The block determinant formula gives the displayed quotient, including
the case where the full moment matrix is singular.
At $d=0$ the formula reads $\lambda_0=s_0$.
When $r=0$ all the mass is at zero and this already equals $M_0$.
Finally a multiple analytic zero contributes at least two to $M_0$.
Thus $\lambda_d<2$ forces $M_0=0$, and $M_0=0$ gives
$\lambda_r=0$.  The empty multiset has $\lambda_0=0$ and also
satisfies the conclusion.
\end{proof}

The moments can also be evaluated on the same fixed contours used
for the heat minors.  Let $f(z)=H_0(z)$ and let a positively oriented
contour $\Gamma$ enclose a finite set of positive endpoint zeros,
with no zero on its boundary.  Alp\"oge's theorem makes those zeros
real, so $f'(x)^2\ge0$.  Put
\begin{equation}\label{eq:fd-contourmom}
 s_k^{\Gamma}=\frac1{2\pi i}\int_\Gamma
            f'(z)^{2k}\frac{f'(z)}{f(z)}\,dz,\qquad k\ge0.
\end{equation}
The residue at a zero $x$ of multiplicity $m_x$ is
$m_x f'(x)^{2k}$.  Hence~\eqref{eq:fd-contourmom} gives exactly
the moments of the positive finite measure
$\sum_xm_x\delta_{f'(x)^2}$.  The proposition applies unchanged,
and counts exactly the same multiple endpoints.
This is an explicit relation between the moment certificate and
contour integration, through the holomorphic observable $(f')^2$.
The two derivative normalizations are related at
$\rho=1/2+i\gamma$, $x=2\gamma$, by
\[
 f'(2\gamma)=\frac{i}{16}A(\rho)\zeta'(\rho),\qquad
 A(s)=\tfrac12s(s-1)\pi^{-s/2}\Gamma(s/2),
\]
so
\[
 f'(2\gamma)^2=|f'(2\gamma)|^2
             =\frac{|A(\rho)|^2}{256}|\zeta'(\rho)|^2.
\]
The factor is positive at every nontrivial zero.  Both measures
therefore have the same atom at zero, with its full multiplicity,
although the positive atoms have the displayed different locations.

A rigorous numerical certificate need not invert a nearly singular
moment matrix.  Choose a specific real polynomial
$P(X)=\sum_{i=0}^d c_iX^i$, with $c_0=1$.
Suppose approximate moments obey
$|s_k-\widetilde s_k|\le\epsilon_k$ for $0\le k\le2d$.
Then
\begin{equation}\label{eq:fd-momentcert}
 \sum_{\rho\in\mathcal Z}P(u_\rho)^2
 \le\sum_{i,j=0}^dc_ic_j\widetilde s_{i+j}
       +\sum_{i,j=0}^d|c_ic_j|\epsilon_{i+j}.
\end{equation}
This follows by expanding the square and applying the triangle
inequality to each moment error.  An upper bound less than two
certifies simplicity throughout the contour.
For the derivative values of $f$, every $s_k$ on the right is the
contour integral~\eqref{eq:fd-contourmom}.
Thus contour separation and bounds for $f,f'$ and their approximation
errors yield a finite certificate using the same boundary data as the
heat-minor method.  Its numerical threshold is an absolute number.

The exact all-height problem in this formulation is
\[
 \text{for every }T>0\text{ there are }d\ge0\text{ and }P(0)=1
 \text{ with }
 \sum_{0<\Im\rho\le T}P(|\zeta'(\rho)|^2)^2<2.
\]
A proved simple-zero proportion supplies an upper bound of order
$N(T)$ for the corresponding count of remaining multiplicities.
The displayed criterion requires a bound below two, at each height.
Equation~\eqref{eq:fd-momentcert} specifies the error precision that
would turn a moment estimate into that integer-valued certificate.
No estimate reaching this threshold at every height is proved here.

\section{Finite zero surgery and the information in leading moments}
\label{sec:fd-surgery}
This section gives an exact transformation preserving the leading positive-moment data while creating a prescribed multiple endpoint.
Choose two known simple positive ordinates $u<v$, for example the
first two, and set
\[
 R(s)=\frac{(s-1/2)^2+v^2}{(s-1/2)^2+u^2},\qquad
 \widetilde\zeta(s)=R(s)\zeta(s),\qquad
 \widetilde\xi(s)=R(s)\xi(s).
\]
The apparent poles at $1/2\pm iu$ are removable by simplicity.
The map removes those two zeros and adds one further copy of each
zero at $1/2\pm iv$.
For every $T>v$, the positive-ordinate zero counts, including
multiplicity, satisfy
\[
 \widetilde N(T)=N(T),\qquad
 \widetilde N_{\rm simple}(T)=N_{\rm simple}(T)-2.
\]
The completed function is entire, of the same order, and obeys
$\widetilde\xi(1-s)=\widetilde\xi(s)$ and the same conjugation law,
because $R(1-s)=R(s)$ and $R(\overline s)=\overline{R(s)}$.
Its nontrivial zeros all lie on the critical line by Alp\"oge's theorem
and the explicit zero replacement.  In the variable $z$ of $H_0$,
\[
 \frac{\widetilde H_0(i\sqrt w)}{\widetilde H_0(0)}
 =\frac{H_0(i\sqrt w)}{H_0(0)}
       \frac{1+w/(4v^2)}{1+w/(4u^2)}.
\]
The right side cancels one positive real product parameter and
repeats another; its endpoint Taylor coefficients remain positive.
This identity states exactly which coefficient object the
transformation produces.

For every other positive-ordinate zero,
\begin{equation}\label{eq:fd-surgeryder}
 \widetilde\zeta'(1/2+i\gamma)
  =\frac{\gamma^2-v^2}{\gamma^2-u^2}\,
                          \zeta'(1/2+i\gamma).
\end{equation}
This includes multiple zeros, where both derivatives vanish.
For $p>0$ put $M_p(T)=\sum_{0<\Im\rho\le T}|\zeta'(\rho)|^p$
and define $\widetilde M_p$ with the changed zero multiset.
For every $Y>2v$ and every $T\ge Y$,
\begin{equation}\label{eq:fd-momentsurgery}
 |\widetilde M_p(T)-M_p(T)|
 \le C_{p,Y}+\delta_p(Y)M_p(T),\qquad
 \delta_p(Y)=1-\left(1-\frac{v^2-u^2}{Y^2-u^2}\right)^p.
\end{equation}
Here $C_{p,Y}$ is the finite sum of absolute moment differences below
$Y$, including the two replaced endpoints.
To prove the inequality, for $\gamma\ge Y$ the positive ratio in
\eqref{eq:fd-surgeryder} lies between
$1-(v^2-u^2)/(Y^2-u^2)$ and one.
Raise it to the power $p$, multiply by $|\zeta'(\rho)|^p$, and
sum.  Since $\delta_p(Y)\to0$, every divergent leading positive
moment asymptotic is transported with relative error tending to zero.
The exact counting identities likewise transport every asymptotic
simple-zero proportion.  Nonetheless the zeros at $v$ are double.
These are explicit limitations of that collection of asymptotic data.

The transformation also identifies an arithmetic property that it changes.
As $\sigma\to+\infty$ along the real axis,
\[
 \widetilde\zeta(\sigma)
 =1+\frac{v^2-u^2}{\sigma^2}+O(\sigma^{-3})+O(2^{-\sigma}).
\]
It therefore has no ordinary Dirichlet series with leading coefficient
one that is absolutely convergent in a right half-plane.
Indeed, absolute convergence at $\sigma_0$ would imply
\[
 \left|\sum_{n\ge2}a_n n^{-\sigma}\right|
 \le2^{-(\sigma-\sigma_0)}
        \sum_{n\ge2}|a_n|n^{-\sigma_0}=O(2^{-\sigma}),
\]
contradicting the nonzero $\sigma^{-2}$ term.
Thus the exact transformation preserves the stated endpoint and
leading-moment information while changing the displayed Dirichlet-series
property.  A proof of simplicity using arithmetic information must use
more than the transported data.

\subsection{The positive kernel under finite zero surgery}
\label{sec:fd-positive-surgery}

The same transformation has an explicit positive kernel. We first
give the kernel estimate needed for its construction. The classical
properties of the Riemann kernel are collected in
\cite[Theorem~4.1, p.~10]{Csordas2015}; the monotonicity estimate is
proved here so that the positivity argument has a complete local input.

\begin{lemma}
\label{lem:fd-kernel-decreasing}
The Riemann kernel \(\Phi\) is strictly decreasing on \((0,\infty)\).
It is even and analytic in \(|\Im u|<\pi/8\), and every derivative
is \(O(\exp(-c e^{4|u|}))\) on the real axis for every \(0<c<\pi\).
\end{lemma}

\begin{proof}
For \(u\ge0\), write \(y_n=\pi n^2e^{4u}\). The defining series is
\[
 \Phi(u)=\sum_{n\ge1}e^u(2y_n^2-3y_n)e^{-y_n}.
\]
It and its derivatives converge locally uniformly in the strip,
since \(\Re(e^{4u})>0\) there. The stated decay follows by
absorbing the polynomial-exponential factors into
\(\exp((\pi-c)e^{4u})\). To verify evenness, let
\(\theta(t)=\sum_{n\in\mathbb Z}e^{-\pi n^2t}\).
The Gaussian Poisson identity gives
\(\theta(t)=t^{-1/2}\theta(1/t)\). Hence
\(\Theta(u)=e^u\theta(e^{4u})\) is even, and termwise
differentiation gives
\(\Phi=(\Theta''-\Theta)/16\). In particular \(\Phi'(0)=0\).

For the individual summands \(\Phi_n\), direct differentiation gives
\[
 \Phi_n'(u)=e^u y_ne^{-y_n}p(y_n),\qquad
 \Phi_n''(u)=e^u y_ne^{-y_n}Q(y_n),
\]
where
\[
 p(y)=-8y^2+30y-15,\qquad
 Q(y)=32y^3-224y^2+330y-75.
\]
The larger root of \(p\) is
\((15+\sqrt{105})/8<79/25\).
For \(u\ge1/500\), the elementary bounds
\(\pi>157/50\) and \(e^{1/125}>126/125\) give
\(y_n>79/25\). Every summand of \(\Phi'\) is then negative.

It remains to consider \(0\le u\le1/500\).
Here \(3<y_1<16/5\), using \(\pi<22/7\) and
\(e^{1/125}<125/124\).
On \([3,16/5]\),
\[
 Q'(y)=96y^2-448y+330
 \le96(16/5)^2-448\cdot3+330<0.
\]
Thus \(Q(y_1)\le Q(3)=-237\).
Since \(e<3\), we have \(y_1e^{-y_1}>1/27\), so
\(\Phi_1''<-79/9\).
For \(n\ge2\), one has \(y_n\ge12\),
\(y_n<(16/5)n^2\), \(e^u<101/100\), and
\(Q(y_n)<35y_n^3\). Therefore
\[
 \sum_{n\ge2}\Phi_n''(u)
 <\frac{101}{100}\,35(16/5)^4
                 \sum_{n\ge2}n^8e^{-3n^2}.
\]
Successive terms of the last sum have ratio at most
\((3/2)^8e^{-15}<1/10000\). Using \(e>8/3\), its first term
and geometric tail give
\[
 \sum_{n\ge2}n^8e^{-3n^2}
 <256(3/8)^{12}\frac{10000}{9999}<\frac1{500}.
\]
Finally
\(\frac{101}{100}35(16/5)^4/500<79/9\).
Thus \(\Phi''<0\) on this short interval.
Together with \(\Phi'(0)=0\), this proves strict decrease everywhere.
\end{proof}

\begin{proposition}[A positive kernel for the exact replacement]
\label{prop:fd-positive-surgery}
Let \(U=2u\) and \(V=2v\) be the two simple positive endpoints
used in the finite surgery, so \(0<U<V\). Define, for \(x\ge0\),
\[
 \Psi_U(x)=U\int_x^\infty\sin(U(t-x))\Phi(t)\,dt,
 \qquad
 \widetilde\Phi(x)=\Phi(x)+
                   \left(\frac{V^2}{U^2}-1\right)\Psi_U(x),
\]
and extend both functions evenly to the real axis.
Then \(\Psi_U\) and \(\widetilde\Phi\) are strictly positive,
even, smooth, and analytic in \(|\Im x|<\pi/8\).
Every real-axis derivative has superexponential decay of the
form in Lemma~\ref{lem:fd-kernel-decreasing}. Their cosine transforms
satisfy the exact identities
\[
 \int_0^\infty\Psi_U(x)\cos(zx)\,dx
       =\frac{H_0(z)}{1-z^2/U^2},\qquad
 \int_0^\infty\widetilde\Phi(x)\cos(zx)\,dx
       =\frac{V^2-z^2}{U^2-z^2}H_0(z).
\]
All apparent singularities are removable. In particular the entire
heat family
\[
 \widetilde H_t(z)=\int_0^\infty
                e^{t x^2}\widetilde\Phi(x)\cos(zx)\,dx
\]
is defined for every complex \(t,z\), and its endpoint is precisely
the completed-function surgery already described.
For real \(a\ge0\), every coefficient of
\(\widetilde H_{-a}(i\sqrt w)/\widetilde H_{-a}(0)\) is positive.
\end{proposition}

\begin{proof}
Put \(b=\pi/U\). Pairing successive half periods gives
\[
 \frac{\Psi_U(x)}U
 =\int_0^b\sin(Ut)\sum_{j\ge0}
  \bigl[\Phi(x+t+2jb)-\Phi(x+t+(2j+1)b)\bigr]dt>0.
\]
The kernel's strict decrease proves the sign; its decay justifies
the pairing and exchange of sum and integral.
Differentiating the defining integral gives
\[
 \Psi_U'(x)=-U^2\int_x^\infty\cos(U(t-x))\Phi(t)\,dt,
 \qquad
 \Psi_U''+U^2\Psi_U=U^2\Phi.
\]
At the origin, \(\Psi_U'(0)=-U^2H_0(U)=0\).
The differential equation with even analytic forcing therefore
has an even analytic solution in the same strip. This proves the
claimed regularity of the even extension.

On the real axis, \(|\sin(Uv)|\le Uv\) gives
\[
 |\Psi_U(x)|\le U^2\int_0^\infty v\Phi(x+v)\,dv.
\]
The defining theta series gives
\(\Phi(y)\le C e^{9y-\pi e^{4y}}\) for \(y\ge0\).
Since \(e^{4v}-1\ge4v\), it follows that
\[
 |\Psi_U(x)|\le
 \frac{CU^2 e^{9x-\pi e^{4x}}}{(4\pi e^{4x}-9)^2}
       =O_U(e^{x-\pi e^{4x}}).
\]
The first derivative has the analogous estimate from its integral;
the differential equation and induction give the bound for every
higher derivative. Thus integration by parts has no boundary
contribution at infinity, and its contribution at zero vanishes
because \(\Psi_U'(0)=0\).
If \(K_U\) is the cosine transform of \(\Psi_U\), the equation
therefore gives
\((U^2-z^2)K_U(z)=U^2H_0(z)\).
Both transforms are entire by the decay just established, proving
the first identity and removability at \(\pm U\).
Adding the two kernels gives the second identity because
\[
 1+\frac{V^2/U^2-1}{1-z^2/U^2}
       =\frac{V^2-z^2}{U^2-z^2}.
\]
Its coefficient \(V^2/U^2-1\) is positive, so
\(\widetilde\Phi>0\). The same decay permits arbitrary complex
heat time. Finally, expanding \(\cosh(x\sqrt w)\) shows that its
normalized coefficients are the strictly positive integrals
\[
 \frac{1}{(2j)!\widetilde H_{-a}(0)}
   \int_0^\infty x^{2j}e^{-a x^2}\widetilde\Phi(x)\,dx.
\]
\end{proof}

This proves preservation of the explicitly stated positive-kernel
heat class. The stronger definition of an admissible kernel in
\cite[Definition~1.2, p.~2]{Csordas2015} also requires strict decrease
of the transformed kernel; that additional property is not used here.
The exact Dirichlet-series obstruction and the transported positive
derivative-moment asymptotics remain those proved above for the same
rational replacement.

\begin{proposition}[The exact heat threshold after replacement]
\label{prop:fd-surgery-threshold}
The real times at which every zero of \(\widetilde H_t\) is real
are exactly \([0,\infty)\). Thus this transformed heat family has
Newman threshold zero and a double endpoint zero at \(V\).
\end{proposition}

\begin{proof}
We first verify forward preservation directly for the functions in
this family. Fix a real time \(r\) whose endpoint
\(f=\widetilde H_r\) has only real zeros. The positive kernel gives
\(f(0)>0\), and the preceding tail estimates imply
\[
 \log\max_{|z|\le R}|f(z)|=O_r(R\log(R+2)).
\]
For example, split the defining integral at a constant multiple of
\(\log(R+2)\); on the remaining tail the negative
\(-c e^{4x}\) term dominates \(Rx+|r|x^2\).
Evenness and the Hadamard factorization therefore give
\[
 \frac{f(z)}{f(0)}=\prod_{j\ge1}(1-z^2/x_j^2),
 \qquad \sum_jx_j^{-2}<\infty,
\]
where positive real zeros are listed with multiplicity.
Let \(p_M\) be a finite initial product. For every complex \(z\),
\[
 |p_M(z)|\le\prod_{j\ge1}(1+|z|^2/x_j^2)
             =\frac{f(i|z|)}{f(0)}.
\]

Every polynomial \(p+c p'\), with \(c\) real, has only real zeros
when \(p\) does. Indeed, away from the roots the equation is
\(p'/p=-1/c\); its left side is strictly decreasing between
successive distinct real roots and has the required infinite
endpoint limits. A root of multiplicity \(m\) retains
multiplicity \(m-1\), and counting the remaining roots proves
the assertion. The case \(c=0\) is immediate.
Consequently
\[
 (1-tD^2/n)^n p_M
 = (1-\sqrt{t/n}\,D)^n(1+\sqrt{t/n}\,D)^n p_M
\]
has only real zeros for \(t\ge0\).
Its coefficient limit as \(n\to\infty\) is
\(e^{-tD^2}p_M\), so this polynomial also has only real zeros.

For a standard real Gaussian \(Z\), the polynomial heat identity is
\[
 e^{-tD^2}p_M(z)=\mathbb E\,p_M(z+i\sqrt{2t}\,Z).
\]
The product bound and the growth estimate dominate the integrand
by a Gaussian-integrable function, uniformly for \(z\) in a compact
set and \(t\) in a bounded nonnegative interval. Passing to the
limit \(M\to\infty\) therefore gives
\[
 \lim_{M\to\infty}e^{-tD^2}p_M(z)
 =\frac{\mathbb E f(z+i\sqrt{2t}\,Z)}{f(0)}
 =\frac{\widetilde H_{r+t}(z)}{f(0)}.
\]
The last identity follows by integrating the positive-kernel
representation against the Gaussian; its expectation multiplies
\(\cos(zx)\) by \(e^{t x^2}\), and the same decay justifies Fubini.
The limit is nonzero, since \(\widetilde H_{r+t}(0)>0\).
Hurwitz's theorem on each half-plane consequently proves forward
preservation of real zeros. At \(r=0\), all zeros are real by the
exact endpoint replacement, so every \(t\ge0\) has that property.

At \(V\), the endpoint has multiplicity exactly two. Write
\(g=\widetilde H_0''(V)/2\ne0\).
Joint analyticity of the heat family and
\(\partial_t\widetilde H_t=-\partial_z^2\widetilde H_t\) give
\[
 \frac{\widetilde H_{-a}(V+\sqrt a\,w)}{a g}
          =w^2+2+O(\sqrt a),\qquad a\downarrow0,
\]
uniformly on compact \(w\)-sets.
Rouch\'e's theorem near \(w=\pm i\sqrt2\) gives nonreal zeros
for every sufficiently small negative time.
An earlier negative time with only real zeros would propagate
forward, by the argument just proved, to one of these small
negative times. This contradiction excludes every negative time.
\end{proof}

\subsubsection*{The exact change in the trace certificate}

The replacement also has an explicit action on the spectral and
arithmetic data. In the notation of
Theorem~\ref{thm:fd-arithmetic-complete}, its counted product moments
and distinct-coordinate trace obey
\[
 \widetilde S_k=S_k+2v^{-2k}-2u^{-2k},\qquad
 \operatorname{tr}(\widetilde C^{2k})
     =\operatorname{tr}(C^{2k})-2u^{-2k}.
\]
The first identity counts the removed pair and the added copy of
the pair at \(v\). In the second, the coordinate at \(v\) was
already present, while the coordinate at \(u\) disappears.
Subtracting gives the exact limiting-defect change
\[
 \widetilde D_{\infty,k}=D_{\infty,k}+2v^{-2k}.
\]
Here the two Hilbert spaces retain their own weighted coordinate
norms; this calculation compares the displayed traces directly.

On their common compactly supported test-function space, write
\(F=\widehat f\), \(G=\widehat g\). The corresponding Weil forms
satisfy
\[
 \widetilde q(f,g)-q(f,g)
 =\sum_{\varepsilon\in\{\pm1\}}
   \bigl(F(\varepsilon v)\overline{G(\varepsilon v)}
        -F(\varepsilon u)\overline{G(\varepsilon u)}\bigr).
\]
For \(h=f*\widetilde g\), Fourier multiplication therefore gives
the equivalent arithmetic correction
\[
 \widetilde W(h)-W(h)
   =2\int_{\mathbb R}h(x)
                \bigl(\cos(vx)-\cos(ux)\bigr)\,dx.
\]
Thus each finite Gram matrix changes by a sum of four explicitly
given evaluation tensors, of rank at most four. The correction
survives in the limiting multiplicity defect with the positive
value just computed. These identities specify the finite arithmetic
information changed by the surgery, alongside the positive-kernel
and leading-moment information that it preserves.


\subsection{The density quantifier for local finite tests}

For a positive simple endpoint $x$, choose a conjugation-invariant disk
containing no other zero of $H_0$.  Its unique zero remains real and
simple for a sufficiently short real time interval about zero.
The normalized polynomial of the corresponding squared-variable
contour has degree one, $1+\beta_x(a)w$, with $\beta_x(a)>0$.
Every one of its Toeplitz minors is nonnegative and every Jensen
polynomial is hyperbolic throughout that interval.  This follows
from its positive linear factor, or from the exact differential
identity $(1+\beta_x(a)D)z^N=z^{N-1}(z+N\beta_x(a))$.
Consequently every proved positive-density set of simple endpoints
provides the same-density family of local tests with infinite detecting
order on their respective persistence intervals.

A quantitative persistence theorem makes this uniform on the set of
labels to which it applies.  Specifically, a real simple zero branch
on $[-\tau,0]$ defines the same positive degree-one local factor for
all $a\in[0,\tau]$; this identity transfers the entire persistence
interval to the local detecting test.  The persistence theorem elsewhere
in this paper supplies its stated common value $\tau=c/\log^2T$.

The global quantities of Theorem~\ref{thm:fd-limits} use the product
over every endpoint.  Their full scaled limits give an exact quantifier
check: any sequence $a_n\downarrow0$ on which
$\sqrt{a_n}m_*(-a_n)\to\infty$, or
$a_nj_*(-a_n)\to\infty$, already proves simplicity of every endpoint.
The local degree-one factors enter that global product by multiplication;
a single multiple factor contributes the finite limiting constant
computed in Theorem~\ref{thm:fd-limits}.  This product map is the
precise relation between the positive-density local conclusion and
the all-endpoint growth problem.

\section{The trace determinants on a fixed contour}\label{sec:local}
This section proves the local partition formula and makes its finite-precision requirements explicit. All notation is as in Theorem~\ref{thm:local-orders}. By compactness of the boundary, $H(t,z)$ remains nonzero there for all sufficiently small complex $t$. The argument principle fixes the enclosed degree $M$, and proves that every $S_j$ is holomorphic in $t$.

\subsection{Hermite splitting and the leading coefficient}
At a zero $c$ of multiplicity $m$, write $H(0,c+y)=a_my^m+O(y^{m+1})$, where $a_m\ne0$. The heat equation recursively determines every time derivative by $\partial_t^nH=(-1)^n\partial_z^{2n}H$. Consequently, on putting $t=u^2/2$ and $z=c+uy$,
\[
 \frac{H(u^2/2,c+uy)}{a_mu^m}
 =\sum_{\ell=0}^{\lfloor m/2\rfloor}
 \frac{(-1)^\ell m!}{2^\ell\ell!(m-2\ell)!}y^{m-2\ell}+O(u)
 =\He_m(y)+O(u).
\]
The series is convergent locally uniformly in $y$. To see the divisibility, the coefficient of $u^{2n+j}$ comes from $\partial_z^{2n+j}H(0,c)$ and vanishes for $2n+j<m$; joint holomorphy supplies convergence. The monic probabilists' Hermite polynomial $\He_m$ has $m$ distinct real roots $h_{m,j}$. Its three-term recurrence is the characteristic-polynomial recurrence of the symmetric tridiagonal matrix with off-diagonals $\sqrt1,\ldots,\sqrt{m-1}$. That matrix has simple eigenvalues: an eigenvector is determined by its first coordinate through the tridiagonal equations, and a zero first coordinate forces all coordinates to vanish. The implicit function theorem now gives all the branches
\[
 z_{ij}(t)=c_i+\sqrt{2t}\,h_{m_i,j}+O(t).
\]
For small positive $t$ they are real, by conjugation and uniqueness.

For $0\le a\le m$ define
\[
 C(m,0)=1,\qquad
 C(m,a)=m^a\prod_{\ell=1}^{a-1}(m-\ell)^{a-\ell}\quad(a\ge1).
\]
\begin{proposition}[The exact leading coefficient]\label{prop:leading-coefficient}
With $v=v_k$ from Theorem~\ref{thm:local-orders}, the coefficient $c_k=[t^v]D_k(t)$ is
\[
 c_k=2^v\!\sum_{\substack{0\le a_i\le m_i,\ \sum_i a_i=k\\
                  \sum_i\binom{a_i}{2}=v}}
       \prod_i C(m_i,a_i)\prod_{i<j}(c_j-c_i)^{2a_i a_j}>0.
\]
\end{proposition}
\begin{proof}
For the $M$ enclosed roots $z_1,\ldots,z_M$, repeated with multiplicity, let $V_{u,j}=z_j^u$ for $0\le u<k$. Then the moment matrix is $VV^{\mathsf T}$. Cauchy--Binet gives
\[
 D_k=\sum_{|I|=k}\prod_{\substack{p<q\\p,q\in I}}(z_q-z_p)^2,
\]
the classical subdiscriminant identity; see Roy \cite[Proposition~1]{roy}. A selection containing $a_i$ roots from the $i$th center has order $e(a)=\sum_i\binom{a_i}{2}$. Each internal squared difference contributes one power of $t$; cross-center differences have nonzero limits. The successive costs at center $i$ are $0,1,\ldots,m_i-1$. Choosing the $k$ least costs respects each center's preceding costs, so the minimum of $e(a)$ is $v_k$. Every term attaining the minimum has a positive leading coefficient. This proves the order formula without any cancellation.

It remains to evaluate the sum of the internal squared Vandermonde factors on $a$ Hermite roots. Let $J_m$ be the tridiagonal matrix just described. Cramer's rule at its last vertex gives the resolvent entry $\He_{m-1}(z)/\He_m(z)$. Differentiating the Hermite recurrence, or its displayed coefficients, gives $\He_m'=m\He_{m-1}$. The spectral measure at that vertex therefore puts mass $1/m$ at each root. Reverse the vertices. The monic orthogonal-polynomial recurrence now has coefficients $m-1,m-2,\ldots,1$. The squared norm of its degree-$j$ monic polynomial is $\prod_{\ell=1}^j(m-\ell)$: apply that polynomial in $J_m$ to the first basis vector, obtaining the product of the first $j$ off-diagonal entries times the $(j+1)$st basis vector. Changing from monomials to these monic polynomials has determinant one. Thus the degree-$a$ moment determinant for the counting measure is
\[
 m^a\prod_{j=0}^{a-1}\prod_{\ell=1}^j(m-\ell)=C(m,a).
\]
Substitution into Cauchy--Binet proves the coefficient formula. Finally, the multiplicity of the integer $h$ among the costs is exactly the number of centers with $m_i\ge h+1$, proving the reconstruction assertion.
\end{proof}

At zero time the moment form is $\sum_i m_iP(c_i)Q(c_i)$, of rank $r$. In particular
\[
 D_r(0)=\prod_i m_i\prod_{i<j}(c_j-c_i)^2>0,
 \qquad D_k(0)=0\quad(k>r).
\]
\begin{proposition}[The first negative local minor]\label{prop:first-negative}
Set $Q(z)=\prod_i(z-c_i)$. For $r<M$,
\[
 D_{r+1}'(0)=2D_r(0)\sum_i m_i(m_i-1)Q'(c_i)^2>0.
\]
Hence $D_{r+1}(-a)<0$ for all sufficiently small $a>0$.
\end{proposition}
\begin{proof}
Replace the degree-$r$ monomial in the moment basis by the monic polynomial $Q$. This basis change has determinant one. Its last row and column vanish at $t=0$, so differentiating the determinant leaves $D_r(0)$ times the derivative of the moment of $Q^2$. For a polynomial $P$, differentiate its contour moment, use $H_t=-H_{zz}$, and integrate by parts on the closed contour to obtain
\[
 \left.\frac d{dt}\sum_zP(z)\right|_{t=0}
 =\sum_i\operatorname*{Res}_{z=c_i}
      P'(z)\frac{H_{zz}(0,z)}{H(0,z)}.
\]
Writing $H(0,z)=(z-c_i)^{m_i}g_i(z)$, the double-pole coefficient is $m_i(m_i-1)$ and the simple-pole coefficient is $2m_i g_i'(c_i)/g_i(c_i)$. For $P=Q^2$, $P'(c_i)=0$ and $P''(c_i)=2Q'(c_i)^2$. The residue is therefore the summand asserted. At least one $m_i$ exceeds one.
\end{proof}

\subsection{Two partitions worked out}
For $P(z)=z^3(z-1)$, its heat evolution is
\[
 e^{-tD^2}P=z^4-z^3-12tz^2+6tz+12t^2.
\]
Newton's identities give
\[
 D_1=4,\quad D_2=3+96t,\quad
 D_3=36t+864t^2+9216t^3,
\]
and the leading term of $D_4$ is $864t^3$. The costs are $0,0,1,2$, giving orders $(0,0,1,3)$. For $P(z)=z^2(z-1)^2$, the costs are $0,0,1,1$, giving orders $(0,0,1,2)$; its leading coefficients for $D_1,D_2,D_3,D_4$ are respectively $4,4,32,64$. The complete determinant polynomials are computed by exact rational arithmetic in the workbench. The implementations compare all coefficients, as well as orders and leading terms.

The closeness of different centers affects the signal. For $P(z)=z^2-\delta^2$,
\[
 D_2(t)=4\delta^2+8t,
\]
so a simple endpoint pair gives a negative minor at times $t<-\delta^2/2$. For $P(z)=z^2(z-\delta)$,
\[
 D_3(t)=8t(\delta^4+18\delta^2t+108t^2).
\]
Its detecting derivative is $8\delta^4$, tending to zero with the distance to the other center. These exact models explain why a negative observation at a prescribed nonzero time needs endpoint control before it can certify multiplicity at time zero.

\subsection{Explicit contour and coefficient errors}
Choose a complex time radius $R>0$ on which the boundary is zero-free. Let $\ell$ be the boundary length and put
\[
 Z=\max(1,\max_{z\in\partial U}|z|),\quad
 A=\frac{\ell}{2\pi}\max_{\substack{|t|\le R\\z\in\partial U}}
       \left|\frac{H_z(t,z)}{H(t,z)}\right|,
 \qquad B_k=k^{k/2}A^kZ^{k(k-1)}.
\]
These are finite numbers for this contour. The boundary integral gives $|S_j(t)|\le AZ^j$. Divide row $u$ and column $v$ of the moment matrix by $Z^u$ and $Z^v$ and apply Hadamard's determinant inequality. Thus
\begin{equation}\label{eq:det-bound}
 |D_k(t)|\le B_k\qquad(|t|\le R).
\end{equation}
For $r>1$ put $\delta=\min_{i<j}|c_i-c_j|>0$; for $r=1$ put $\delta=1$. Proposition~\ref{prop:leading-coefficient} immediately gives
\begin{equation}\label{eq:lead-lower}
 c_k\ge 2^{v_k}\delta^{k(k-1)-2v_k}.
\end{equation}
Indeed, every factor $C(m_i,a_i)$ is a positive integer, at least one minimizing selection exists, and the cross-center exponent is $k(k-1)-2v_k$.

\begin{proposition}[A coefficient certificate]\label{prop:coefficient-certificate}
For $0<\varrho<R$, take $N$ equally spaced points $t_l=\varrho e^{2\pi il/N}$. Let $\widetilde D_l$ approximate $D_k(t_l)$ with absolute error at most $\varepsilon$. For $0\le j<N$ define
\[
 \widetilde d_j=\frac1{N\varrho^j}\sum_{l=0}^{N-1}
           \widetilde D_l e^{-2\pi ilj/N}.
\]
Then
\[
 |\widetilde d_j-[t^j]D_k(t)|
 \le \varepsilon\varrho^{-j}
   +B_kR^{-j}\frac{(\varrho/R)^N}{1-(\varrho/R)^N}.
\]
Furthermore, for $|t|\le R/2$,
\[
 |D_k(t)-c_kt^{v_k}|
 \le 2B_k|t|^{v_k+1}/R^{v_k+1}.
\]
\end{proposition}
\begin{proof}
Write $D_k(t)=\sum_{n\ge0}d_nt^n$. Cauchy's estimate from \eqref{eq:det-bound} is $|d_n|\le B_kR^{-n}$. The discrete Fourier sum of exact samples equals $d_j+\sum_{h\ge1}d_{j+hN}\varrho^{hN}$, since $j<N$. Summing the geometric majorant proves the aliasing bound, and averaging the sample errors proves the first term. The lower Taylor coefficients vanish through degree $v_k-1$ by Theorem~\ref{thm:local-orders}. Summing the coefficient bounds above degree $v_k$ proves the remainder estimate.
\end{proof}

For example, a lower bound $c\le c_k$ and the last estimate certify the sign of $D_k(-a)$ for
\[
 0<a<\min\{R/2,\ cR^{v_k+1}/(2B_k)\}.
\]
Exact recovery of an order from approximate coefficients uses the uniform lower bound on its first nonzero coefficient in \eqref{eq:lead-lower}, scanning $0\le j\le\binom{k}{2}$ with an error below one third of the minimum of those lower bounds. Coefficients preceding the first nonzero coefficient have absolute value at most that error; the first nonzero one has absolute value at least twice that error.

There is also a direct route to the time disk. Let $\eta=\min_{\partial U}|H(0,z)|>0$ and let $V$ bound $|H_t|$ on $|t|\le R$, $z\in\partial U$. The estimate $RV<\eta$ proves boundary nonvanishing by integrating in $t$. A bound for $H_z$ then bounds $A$ above. This identifies the precise analytic estimates required as the contour grows.

\subsection{Two executed contour certificates}
\label{sec:executed-contour}
The error bounds above can be evaluated using the defining zeta
function and whole-edge complex enclosures. We give two such
certificates, at positive ordinate heights $20$ and $45$. The
coordinate is still $z=2\gamma$.

We first give a uniform elementary bound for the heat perturbation.
For $u\ge0$ the kernel in Section~\ref{sec:fd-positive-surgery} satisfies
\[
 0<\Phi(u)
 \le 2\pi^2 e^{9u-\pi e^{4u}}
       \sum_{n\ge1}n^4e^{-\pi(n^2-1)e^{4u}}
 <40e^{9u-3e^{4u}}.
\]
Indeed, the sum is less than $2$. Its terms with $n\ge2$ are
bounded by $n^4e^{-3(n^2-1)}$; the first is less than $1/32$,
and successive ratios are less than
$(3/2)^4e^{-15}<1/1000$. Since
$e^{4u}\ge1+4u+8u^2$, it follows that
\[
 u^j e^{u^2+u/4}\Phi(u)
 <40e^{-3}u^j e^{-23u^2-11u/4},\qquad 0\le j\le3.
\]
Each integral on the right is less than $1$: drop its last
exponential and use, respectively,
\[
 \int_0^\infty u^j e^{-23u^2}\,du
 =\frac{\sqrt\pi}{2\sqrt{23}},\quad
 \frac1{46},\quad
 \frac{\sqrt\pi}{4\,23^{3/2}},\quad
 \frac1{2\,23^2},
\]
together with $40e^{-3}<5/2$, $\sqrt\pi<2$ and
$\sqrt{23}>4$. Differentiating the cosine integral therefore proves
\begin{equation}\label{eq:uniform-contour-heat}
 |H_t(z)-H_0(z)|\le |t|,\qquad |\partial_zH_t(z)|<1
 \quad (|t|\le1,\ |\Im z|\le1/4).
\end{equation}

For $T=20,45$, let $U_T$ be the rectangle
$0<\Re z<2T$, $|\Im z|<1/4$. The certified data are
\[
\begin{array}{c|c|c|c|c}
 T&\min_{\partial U_T}|H_0|\text{ exceeds}&R_T&A_T\text{ is below}
       &S_0(0)\\ \hline
 20&10^{-6}&10^{-8}&14\cdot10^6&1\\
 45&10^{-14}&10^{-16}&3\cdot10^{15}&8
\end{array}
\]
Here $A_T$ has exactly the definition in the preceding subsection.
The boundary length is $4T+1$, so
\[
 A_T\le\frac{4T+1}{2\pi(\eta_T-R_T)}
\]
follows directly from \eqref{eq:uniform-contour-heat}. All the
lower bounds cover complete edge segments. They are obtained by
subdividing each edge into dyadic intervals and evaluating
\[
 H_0(z)=\frac18(s-1)\pi^{-s/2}\Gamma(1+s/2)\zeta(s),
 \qquad s=\tfrac12+iz/2,
\]
as a complex ball on the entire segment. No numerical list of zeros
enters this computation.

The Newton integral is evaluated by certified complex quadrature,
using the exact logarithmic derivative
\[
 \frac{H_0'(z)}{H_0(z)}=\frac i2\left\{
 \frac1{s-1}-\frac12\log\pi
 +\frac12\psi(1+s/2)+\frac{\zeta'(s)}{\zeta(s)}\right\}.
\]
Its enclosures lie within $10^{-6}$ of the displayed integers;
the argument principle makes those integers exact. The real
endpoint values have opposite strict signs on the heat-coordinate
intervals
\[
 (28,29),\ (42,43),\ (50,51),\ (60,61),\
 (65,66),\ (75,76),\ (81,82),\ (86,87).
\]
The first rectangle uses just the first interval. For the second
rectangle the eight sign changes account for its entire degree,
counted with multiplicity. Thus every enclosed zero is simple, and
the distance between successive centers is greater than $4$.
These sign enclosures and the Newton integral also give a direct
verification of reality and simplicity on this particular rectangle.

For the second rectangle, $Z<91$ and $B_k<10^{240}$ for
$1\le k\le8$. Its leading coefficients satisfy
$D_k(0)\ge4^{k(k-1)}\ge1$. The Taylor remainder bound gives
\[
 |D_k(t)-D_k(0)|<\tfrac15
       \qquad (|t|\le10^{-257},\ 1\le k\le8).
\]
For real $t$ in this interval all the leading Hankel determinants
are positive. This implies that the eight roots are real and simple.
To verify that implication directly, a repeated root makes the
moment matrix singular. For distinct roots, a nonreal conjugate
pair admits a real polynomial of degree below eight taking the
values $i,-i$ there and zero at the other roots, by interpolation
and conjugation. Its moment square equals $-2$, contradicting
positive definiteness. For $T=20$ the constant enclosed degree is
one throughout $|t|\le10^{-8}$; the unique zero is therefore simple
throughout that complex disk and real at real heat times.

There is also an executed contour comparison for the ordinary
Hurwitz translation at height $20$. On the corresponding
$s$-rectangle, $3/8\le\Re s\le5/8$, $0\le\Im s\le20$,
so $|s|<21$. For $|u|\le1/4$, integration of each summand in $u$
gives
\[
 |\zeta(s,1+u)-\zeta(s)|
 \le 2|s|e^7|u|\sum_{n\ge1}n^{-11/8}
 <200000|u|.
\]
Here $|n+w|\ge3n/4$ and $|\arg(n+w)|\le1/(3n)$ on
the straight segment from $0$ to $u$; also
$\zeta(11/8)<1+8/3<4$. These estimates justify the convergent
difference series throughout this rectangle. The completion factor
$(s-1)\pi^{-s/2}\Gamma(1+s/2)/8$ has modulus less than $3$:
the Euler integral bounds the gamma modulus by
$\Gamma(1+\Re s/2)\le1$, the latter inequality following from
logarithmic convexity between $1$ and $2$. Hence the completed
translation differs from $H_0$ by less than $600000|u|$.
For $|u|\le10^{-12}$ this is smaller than the certified boundary
minimum. Rouch\'e's theorem gives exactly one simple enclosed zero
throughout that complex parameter disk. Averaging the same estimate
over any finite set of rotations gives the identical conclusion
for every rotation average in Theorem~\ref{thm:phase-orders}.

The computations use FLINT's rigorous complex-ball integration
\cite[section ``Integration'']{flint-acb-calc}. Their editable
program and interval receipts are in the workbench cited in the
closing note. The height-$20$ contour will be compared with an
executed probe-mass certificate in Section~\ref{sec:suzuki-probes}.

\section{The Hurwitz comparison on the same divisor}\label{sec:phase}
This section proves Theorem~\ref{thm:phase-orders} and the exact map from a spectral jet to prime periods. The derivative formula for the Hurwitz zeta function \cite[25.11.17]{dlmf} gives, locally normally for $|q|<1$,
\[
 \zeta(s,1+q)=\sum_{n\ge0}b_n(s)q^n,
 \qquad b_n(s)=\frac{(-1)^n(s)_n}{n!}\zeta(s+n).
\]
Here $(s)_n=s(s+1)\cdots(s+n-1)$. One obtains the identity first by an absolutely convergent binomial expansion for $\Re s>1$, then by continuation. Near a nontrivial zero $\rho$, it is normally convergent directly: $(s)_n/n!$ grows at most polynomially on a compact $s$-set, whereas $|q|^n$ decays geometrically and $\zeta(s+n)$ is uniformly bounded for $n\ge1$. Since $0<\Re\rho<1$, the Euler product proves $\zeta(\rho+n)\ne0$ for $n\ge1$, and no factor of $(\rho)_n$ vanishes.

The rotation average cancels exactly the coefficients whose degree is not a multiple of $d$. Thus
\[
 A_d(s,q)=\zeta(s)+b_d(s)q^d+O(q^{2d}),\qquad b_d(\rho)\ne0.
\]
Write $a=\zeta^{(m)}(\rho)/m!\ne0$. With $q^d=w^m$ and $s=\rho+wy$, division by $w^m$ produces $ay^m+b_d(\rho)+O(w)$. The leading polynomial has $m$ distinct roots. The implicit function theorem gives
\[
 s_j(q)=\rho+c e^{2\pi ij/m}q^{d/m}+O(q^{2d/m}),
 \qquad c^m=-b_d(\rho)/a.
\]
Every squared root difference has order $2d/m$ and a nonzero coefficient. Their product has order $d(m-1)$. A positive circuit in $q$ advances each label by $d$ modulo $m$, yielding $\gcd(m,d)$ cycles of length $m/\gcd(m,d)$. The polynomial coefficients and discriminant are holomorphic, because the contour Newton sums are holomorphic and Newton's identities express the coefficients in them. Between different disks, root differences have nonzero limiting values. Multiplication proves the full order $d\sum_i(m_i-1)$ and completes the proof of Theorem~\ref{thm:phase-orders}.

\subsection{The full local jet and its inverse}
For $p>1$, let $a=\log p$, fix a phase $\omega$ of modulus one, and set $u=\omega p^{\rho-1/2}$. The following explicit isomorphism compares the entire local multiplicity data:
\begin{equation}\label{eq:jet-isomorphism}
 \begin{split}
 \mathbb C[U,U^{-1}]/(U-u)^m&\longrightarrow
                 \mathbb C[X]/(X-\rho)^m,\\
 U&\longmapsto u\sum_{j=0}^{m-1}\frac{a^j(X-\rho)^j}{j!}.
 \end{split}
\end{equation}
Its inverse sends
\[
 X\longmapsto\rho+\frac1a\sum_{j=1}^{m-1}
       \frac{(-1)^{j+1}}j\left(\frac{U-u}{u}\right)^j.
\]
The two composites are the identity because the exponential and logarithm power series are inverse modulo the $m$th power of their nilpotent variable. The image of $U$ is a unit, so Laurent inversion is allowed. On the $j$th associated grade, $(U-u)^j$ maps to $(ua)^j(X-\rho)^j$. In particular, every local nilpotent order is preserved by this map.

Two distinct rational primes separate all complex centers jointly. Indeed, agreement of the two periods at $\rho$ and $\sigma$ gives
\[
 (\rho-\sigma)\log p=2\pi ik,\qquad
 (\rho-\sigma)\log q=2\pi il
\]
for integers $k,l$. Hence $p^l=q^k$, and prime valuations give $k=l=0$, then $\rho=\sigma$. Together with \eqref{eq:jet-isomorphism} and the Chinese remainder theorem, this yields the full finite algebra. Explicitly, the local kernel in $\mathbb C[U^{\pm1},V^{\pm1}]$ is
\[
 \left((U-u_i)^{m_i},\ V-v_i\sum_{j=0}^{m_i-1}
      \binom{\log q/\log p}{j}\left(\frac{U-u_i}{u_i}\right)^j\right).
\]
The kernels for distinct centers are comaximal, since their pairs of constant coordinates are distinct. Their intersection is the global kernel, and the quotient is $\prod_i\mathbb C[X]/(X-\rho_i)^{m_i}$.

For a single prime, two distinct centers can collide for at most one choice of the rational prime, by the same valuation argument. Therefore among any $\binom M2+1$ rational primes at least one separates every enclosed center. The minimum, over those primes, of the discriminant orders of the period images is exactly $d(M-r)$: local orders are preserved by the nonzero derivative $(\log p)p^{s-1/2}$, and cross-center collisions can only add nonnegative order.

\subsection{Two prime periods at finite precision}
Here the zero-location theorem makes the periods unit-modulus coordinates. On a fixed contour enclosing ordinates $\gamma_j$ with multiplicities $m_j$, define
\[
 S_{kl}=\sum_jm_je^{i(k\log p+l\log q)\gamma_j}
 =\frac1{2\pi i}\oint e^{(s-1/2)(k\log p+l\log q)}
                         \frac{\zeta'(s)}{\zeta(s)}\,ds.
\]
Let a target interval be $|\gamma-\gamma_0|\le h$, and let every other enclosed ordinate have distance in $[\Delta,H]$, where $0\le h<\Delta\le H$. Put $a=\log p$, $b=\log q$, $M=\sum_jm_j$, and
\[
 \sigma=\min_{\Delta\le x\le H}
       \max\{|e^{iax}-1|,|e^{ibx}-1|\}>0.
\]
Positivity follows from compactness and the two-prime separation just proved. For $K\ge2$ set $w_k=(K-|k|)/K^2$ for $|k|<K$. Given errors $|\widetilde S_{kl}-S_{kl}|\le\varepsilon$, form
\[
 B=\Re\sum_{|k|,|l|<K}w_kw_l
        e^{-i(ka+lb)\gamma_0}\widetilde S_{kl}.
\]
\begin{proposition}[An integer mass certificate]\label{prop:prime-rounding}
For the total multiplicity $m$ in the target interval,
\[
 |B-m|\le\varepsilon+
       \frac{M(K^2-1)(a^2+b^2)h^2}{12}
       +\frac{4M}{K^2\sigma^2}.
\]
An error bound below $1/2$ certifies $m$ by nearest-integer rounding.
\end{proposition}
\begin{proof}
The normalized Fej\'er kernel is
\[
 F_K(x)=K^{-2}\left|\sum_{j=0}^{K-1}e^{ijx}\right|^2
        =\sum_{|k|<K}w_ke^{ikx}.
\]
It lies between zero and one, and the geometric sum gives
$F_K(x)\le4/(K^2|1-e^{ix}|^2)$ away from its period. Averaging
$1-\cos((j-l)x)\le(j-l)^2x^2/2$ gives
$1-F_K(x)\le(K^2-1)x^2/12$. The exact statistic is
$\sum_jm_jF_K(a(\gamma_j-\gamma_0))F_K(b(\gamma_j-\gamma_0))$.
The deficit on the target is bounded by the sum of the two quadratic deficits. At every other point at least one phase distance is at least $\sigma$, so its product is bounded by $4/(K^2\sigma^2)$. The weights are nonnegative and sum to one in each coordinate, bounding the moment-error contribution by $\varepsilon$.
\end{proof}

This estimate can be made constructive on a known interval: its defining separation function is Lipschitz with constant $\max(a,b)$. Certified midpoint evaluations on a mesh, minus this constant times the half mesh size, give a lower bound for $\sigma$. Refinement eventually proves a positive bound for any fixed $[\Delta,H]$.

The finite models also quantify a precision requirement. A double atom at $\gamma$ and two unit atoms at $\gamma\pm\tau$ have moment difference
\[
 2|1-\cos((ka+lb)\tau)|\le(K-1)^2(a+b)^2\tau^2,
 \qquad |k|,|l|<K.
\]
The map between these two measures and their displayed finite moment arrays proves the bound directly. A fixed positive tolerance therefore admits arbitrarily close such arrays. To read the multiplicity of a single distinct ordinate from the rounded interval mass requires a certified interval containing just that distinct ordinate. The heat determinant certificate makes the same dependence visible through its cross-center powers in \eqref{eq:lead-lower}.

\section{The two zero motions and the corrected Hurwitz sign}
\label{sec:motion-local}

We compute the tangent vectors of the heat family and the Hurwitz translation in the same arithmetic coordinate. This identifies precisely which motion is measured by the phase statistic. The Hurwitz zero-velocity equation was already written by Frisk and de Gosson \cite[Section~2, equation~(3)]{frisk-degosson}; Garunk\v stis and Steuding studied its trajectories and their endpoint classifications \cite[Section~2]{garunkstis-steuding}. We combine the exact velocity with the midpoint phase convention below.
Write
\[
 \xi(s)=\tfrac12s(s-1)\pi^{-s/2}\Gamma(s/2)\zeta(s),
 \qquad H_0(x)=\tfrac18\xi(\tfrac12+ix/2).
\]
The heat family satisfies $\partial_tH_t=-\partial_x^2H_t$.
The second family is $F(q,s)=\zeta(s,1+q)$, where the last function is the Hurwitz zeta function. At $q=0$ it equals $\zeta(s)$. The identity
\[
 \partial_qF(q,s)=-s\zeta(s+1,1+q)
\]
follows first by differentiating the absolutely convergent series for $\Re s>1$ and then by meromorphic continuation; see \cite[25.11.17]{dlmf}.

\begin{proposition}[Exact tangent vectors]\label{prop:two-tangents}
At a simple real zero $x$ of $H_0$, the heat zero and its arithmetic coordinate satisfy
\[
 x'(0)=\frac{H_0''(x)}{H_0'(x)}\in\mathbb R,
 \qquad
 \left.\frac{d}{dt}\left(\tfrac12+ix(t)/2\right)\right|_{t=0}
       =\frac{iH_0''(x)}{2H_0'(x)}.
\]
The heat branch stays real for sufficiently small real $t$ of either sign, so its arithmetic coordinate has real part $1/2$ on that interval.
At a simple zero $\rho$ of $\zeta$, the Hurwitz branch has derivative
\[
 v_\rho:=s'(0)=\frac{\rho\zeta(\rho+1)}{\zeta'(\rho)}.
\]
\end{proposition}
\begin{proof}
The holomorphic implicit function theorem applies because the respective derivatives with respect to the zero coordinate are nonzero. Differentiating the two zero equations gives the displayed formulas. The function $H_t$ is real on the real axis for real $t$. Conjugating its unique local zero section produces another section with the same initial value; uniqueness identifies the two sections. Thus $x(t)$ is real throughout a sufficiently small real parameter interval. The affine map $x\mapsto1/2+ix/2$ gives the asserted arithmetic coordinate and its derivative.
\end{proof}

Let $\theta(u)$ denote the continuous Riemann--Siegel theta function with $\theta(0)=0$, and let
\[
 Z(u)=e^{i\theta(u)}\zeta(1/2+iu).
\]
At a positive zero ordinate, use the midpoint convention in
\[
 N(u)=\theta(u)/\pi+1+S(u).
\]
Under the zero-location input, enumerate the positive ordinates as $\gamma_n$, including every multiplicity. At a simple ordinate define
\[
 \phi_n=\arg\zeta(3/2+i\gamma_n),\qquad
 W_n=S(\gamma_n)-\phi_n/\pi,\qquad
 \delta_n=\arctan\frac1{2\gamma_n}.
\]
Here $\phi_n$ uses the logarithm specified by the absolutely convergent Euler product.

\begin{theorem}[The corrected Hurwitz sign]\label{thm:hurwitz-sign}
For every simple positive-ordinate zero $\rho_n=1/2+i\gamma_n$,
\[
 \Re v_{\rho_n}
 =\frac{|\zeta(\rho_n+1)|}{|Z'(\gamma_n)|}
       \left(\gamma_n\sin\pi W_n+\tfrac12\cos\pi W_n\right).
\]
In particular the strict rightward sign is exactly
$\sin(\pi W_n+\delta_n)>0$, and tangency occurs exactly when
$W_n+\delta_n/\pi\in\mathbb Z$.
At the indexed zero $n=1169$, the sine without $\delta_n$ is negative and the corrected sine is positive.
\end{theorem}
\begin{proof}
At a zero, differentiating $Z$ yields
$Z'(\gamma_n)=ie^{i\theta(\gamma_n)}\zeta'(\rho_n)$.
Since $Z(0)<0$, the sign at the crossing of the $n$th simple zero is
$\operatorname{sgn}Z'(\gamma_n)=(-1)^{n+1}$. This parity remains valid when preceding zeros have higher multiplicity: an even multiplicity changes neither the sign nor the parity of the count, while an odd multiplicity changes both.
The midpoint count reads
$\theta(\gamma_n)/\pi=n-3/2-S(\gamma_n)$. Hence
\[
 \frac1{\zeta'(\rho_n)}
 =\frac{ie^{i\theta(\gamma_n)}}{Z'(\gamma_n)}
 =\frac{e^{-i\pi S(\gamma_n)}}{|Z'(\gamma_n)|}.
\]
Insert $\zeta(\rho_n+1)=|\zeta(\rho_n+1)|e^{i\phi_n}$ into
Proposition~\ref{prop:two-tangents} and take real parts. Finally,
\[
 \gamma_n\sin\pi W_n+\tfrac12\cos\pi W_n
 =\sqrt{\gamma_n^2+1/4}\sin(\pi W_n+\delta_n),
\]
and the prefactor is strictly positive.

The interval computation described in the workbench uses the certified indexed zeta-zero routine and complex-ball power series for $\zeta'$. It gives
\begin{align*}
 \gamma_{1169}&=1613.0788554510274235\ldots,\\
 W_{1169}&=-0.0000606299724187599509\ldots,\\
 \delta_{1169}/\pi&=0.0000986653178528222938\ldots,\\
 \Re v_{\rho_{1169}}&=0.025925651967163814\ldots.
\end{align*}
The saved enclosures exclude zero for the last quantity and for each of the two sine values, with the signs stated in the theorem. The argument of the exterior value can be evaluated by its principal branch here: the Euler-product logarithm obeys
$|\Im\log\zeta(3/2+iu)|\le\log\zeta(3/2)<\pi$, so the two branches agree.
\end{proof}

\subsection{The exact periodic information required by the sign}
For every odd integer $k$, the midpoint convention gives
\[
 e^{i\pi kW_n}
 =e^{-3\pi ik/2}(-1)^n
       e^{-ik\theta(\gamma_n)}e^{-ik\phi_n}.
\]
Put $Y_n=W_n+\delta_n/\pi$. Its strict direction is
$g(Y_n)$, where $g(y)=\operatorname{sgn}(\sin\pi y)$ and
$g(y)=0$ at the integers. We give an exact finite estimate for the
mean of this sign, and a model that preserves Gaussian moments and
all even Fourier modes while keeping the sign positive.

\begin{proposition}[Periodic sign and Gaussian moments]
\label{prop:periodic-gaussian-gate}
There are real random variables $X_\sigma$, $\sigma\to\infty$,
whose normalized moments converge to every standard Gaussian moment,
and for which, for every nonzero integer $\ell$ and integer $k\ge0$,
\[
 \mathbb E\left[e^{2\pi i\ell X_\sigma}
                     (X_\sigma/\sigma)^k\right]\longrightarrow0,
 \qquad g(X_\sigma)=1\quad\hbox{almost surely}.
\]
For arbitrary real $y_1,\ldots,y_M$, $M\ge1$, put
$\langle f\rangle=M^{-1}\sum_{j=1}^M f(y_j)$ and
$b_\epsilon=M^{-1}\#\{j:\operatorname{dist}(y_j,\mathbb Z)<\epsilon\}$.
For $1/2\le r<1$, $0<\epsilon\le1/2$, and integer $J\ge0$,
\begin{align*}
 |\langle g\rangle|
 \le{}&\frac4\pi\sum_{l=0}^J
      \frac{r^{2l+1}}{2l+1}
       |\langle\sin((2l+1)\pi y)\rangle|+b_\epsilon\\
 &+\frac{3(1-r)}{\pi\sin(\pi\epsilon)}
 +\frac{4r^{2J+3}}{\pi(2J+3)(1-r^2)}.
\end{align*}
\end{proposition}
\begin{proof}
Let $G$ be standard normal and let $U$ be independent and uniform
on $(0,1)$. Set
\[
 X_\sigma=2\lfloor\sigma G/2\rfloor+U.
\]
Then $|X_\sigma-\sigma G|\le2$, so
$X_\sigma/\sigma\to G$ in every fixed $L^p$, $p>0$.
Also $g(X_\sigma)=\operatorname{sgn}(\sin\pi U)=1$.
Writing $A_\sigma=2\lfloor\sigma G/2\rfloor/\sigma$, the binomial
expansion of $(A_\sigma+U/\sigma)^k$ has leading term
$A_\sigma^k$. Its product with $e^{2\pi i\ell U}$ has expectation
zero, by independence and $\mathbb E e^{2\pi i\ell U}=0$.
Every remaining term has a factor $\sigma^{-1}$ or smaller,
and the moments of $A_\sigma$ are bounded. This proves all the
weighted limits. The exact map on periodic phases is
$\mathbb R/2\mathbb Z\to\mathbb R/\mathbb Z$, $y\mapsto y$:
the uniform measure on the half-interval $(0,1)$ pushes forward to
the uniform measure on the target circle. Its lift still has sign
one. Replacing $U$ by $U+1$ gives sign minus one with the same
limiting moments and even Fourier modes.

For the finite estimate, the absolutely convergent Abel sum is
\[
 g_r(y)=\frac4\pi\sum_{l\ge0}
       \frac{r^{2l+1}\sin((2l+1)\pi y)}{2l+1}
 =\frac2\pi\arctan\frac{2r\sin\pi y}{1-r^2}.
\]
The identity follows from the power series for
$\tfrac12\log((1+z)/(1-z))$ at $z=re^{i\pi y}$.
The relevant quotient has positive real part, which fixes the
arctangent branch. Hence $g_r$ has the sign of $g$, with
$|g_r|\le1$, and, away from the $\epsilon$-neighborhood of the
integers,
\[
 |g(y)-g_r(y)|
 =\frac2\pi\arctan\frac{1-r^2}{2r|\sin\pi y|}
 \le\frac{3(1-r)}{\pi\sin(\pi\epsilon)}.
\]
On the remaining set the difference is at most one.
Finally, the omitted Fourier tail has absolute value at most
$4r^{2J+3}/[\pi(2J+3)(1-r^2)]$. Averaging, applying the triangle
inequality, and combining these three estimates proves the claim.
\end{proof}

For a finite collection of simple zeta zeros, apply the estimate with
$y_j=Y_{n_j}$. Its left side is exactly the absolute difference
between the fractions of strict rightward and strict leftward
branches. The displayed odd-frequency means and boundary counts
are therefore explicit unresolved quantities for half-balance.
The correction from $W_n$ to $Y_n$ changes a sign only inside a
neighborhood of an integer of radius $\delta_n/\pi$; its counting
cost is exactly bounded by the number in that neighborhood.

\section{Normal coefficients at multiple zeros and simple tangencies}
\label{sec:motion-bending}

The denominator translation has an explicit second jet. We retain
that jet because the first velocity can be vertical. Throughout
this section, a normal coefficient means the real part of a
coefficient in the coordinate $s-\rho$, with the positive real
parameter oriented by
\[
 F(s,u)=\sum_{n\ge1}\frac{\chi(n)}{(n+u)^s}
 =L(s)-usL(s+1)+\frac{u^2s(s+1)}2L(s+2)+O(u^3).
\]
The conductor-one case is $L=\zeta$. The equality first holds
in the absolutely convergent half-plane; its normally convergent
rising-factorial expansion continues it to a neighborhood of
each critical-line zero. The Euler product, primitive functional
equation and endpoint nonvanishing used below have the exact
normalizations in \cite[Section~25.15, equations (2), (5), (6), (9)]{dlmf}.

\begin{lemma}[Reflected logarithmic derivative]
\label{lem:reflected-positive-new}
Let $\Lambda$ be an entire function of order at most one, with
infinitely many zeros in $0\le\Re s\le1$, satisfying
$\Lambda(s)=\epsilon\overline{\Lambda(1-\bar s)}$, $|\epsilon|=1$.
At a zero $\rho=1/2+i\gamma$ of multiplicity $m\ge1$,
\begin{equation}\label{eq:reflected-positive-new}
 \Re\frac{\Lambda'}{\Lambda}(\rho+1)
 =m+\sum_{\rho'\ne\rho}
 \frac{3/2-\Re\rho'}{(3/2-\Re\rho')^2+(\gamma-\Im\rho')^2}>m.
\end{equation}
The sum counts multiplicity and converges absolutely.
\end{lemma}

\begin{proof}
Set $G(z)=\Lambda(\rho+z)$ and write
$G(z)=b_mz^m+b_{m+1}z^{m+1}+\cdots$. Reflection gives
$\Re(b_{m+1}/b_m)=0$. Hadamard factorization yields
\[
 G(z)=z^m e^{az+b}\prod_{\alpha\ne0}
       (1-z/\alpha)e^{z/\alpha},\qquad
 \alpha=\rho'-\rho,\quad a=b_{m+1}/b_m.
\]
Here $\sum|\alpha|^{-2}<\infty$. At $z=1$, the real part
of the logarithmic derivative is
$m+\sum\Re((1-\alpha)^{-1}+\alpha^{-1})$.
Since $|\Re\alpha|\le1/2$, the real parts of
$\alpha^{-1}$ are absolutely summable. Reflection pairs
$\alpha$ with $-\bar\alpha$, so their sum is zero.
Each remaining term is positive, and its tail is
$O(|\alpha|^{-2})$. This proves the formula and its strict sign.
\end{proof}

For a nonprincipal primitive character of parity $a\in\{0,1\}$,
put
\[
 Q_\chi(s)=(q/\pi)^{(s+a)/2}\Gamma((s+a)/2),\qquad
 \Lambda_\chi=Q_\chi L.
\]
The theta integral proves that $\Lambda_\chi$ is entire of
order at most one: after splitting its Mellin integral at one,
Poisson summation replaces the small-variable part by a conjugate
theta series that decays exponentially. The resulting integral
has growth $\exp(O_q(R\log(R+2)))$ on $|s|\le R$.
The functional equation cancels each trivial zero against its
gamma pole; at the endpoint this uses $L(1,\chi)\ne0$.
Thus its remaining zeros are precisely the nontrivial ones,
and the reflected lemma applies also for complex $\chi$.

\begin{proposition}[Exact second normal coefficient]
\label{prop:motion-second-jet}
At a critical-line zero of $L$ of order $m\ge1$, write
\[
 L(\rho+w)=A_mw^m+A_{m+1}w^{m+1}+O(w^{m+2}),\qquad
 D_\rho=\frac1\rho+\frac{L'}L(\rho+1)-\frac{A_{m+1}}{A_m}.
\]
For every nonprincipal primitive character,
\begin{equation}\label{eq:dirichlet-bending-positive-new}
 \Re D_\rho>m+h_a(\gamma)>m,
 \qquad
 \begin{cases}
 h_0(\gamma)=\dfrac1{2(\gamma^2+1/4)}
                  -\dfrac\pi{2\cosh(\pi\gamma)},\\[4pt]
 h_1(\gamma)=\dfrac\pi{2\cosh(\pi\gamma)}.
 \end{cases}
\end{equation}
For $L=\zeta$, the corresponding bound is
\begin{equation}\label{eq:zeta-bending-positive-new}
 \Re D_\rho>m-\frac3{2(\gamma^2+9/4)}
                    -\frac\pi{2\cosh(\pi\gamma)}>0.
\end{equation}
Every multiple zero ($m\ge2$) has convergent local sections
after $u=\tau^m$ of the form
\begin{equation}\label{eq:multiple-bending-jets-new}
 s_j(u)=\rho+c_j u^{1/m}
              +\frac{D_\rho}{m}c_j^2u^{2/m}+O(u^{3/m}),
 \qquad c_j^m=\frac{\rho L(\rho+1)}{A_m}\ne0.
\end{equation}
A vertical first coefficient $c_j=i\beta_j$ has strictly
negative second normal coefficient
$-\beta_j^2\Re D_\rho/m$. Consequently every branch has
a nonzero normal coefficient by the second Puiseux order.
When vertical first coefficients occur, the numbers on the
left and right are respectively $(m/2+1,m/2-1)$ for even $m$,
and $((m+1)/2,(m-1)/2)$ for odd $m$.
\end{proposition}

\begin{proof}
Reflection of the completed germ gives
\[
 \Re\frac{A_{m+1}}{A_m}=-\Re\frac{Q_\chi'}{Q_\chi}(\rho).
\]
Substitution expresses $\Re D_\rho$ as the left side of
\eqref{eq:reflected-positive-new} plus
$\Re(1/\rho+Q_\chi'/Q_\chi(\rho)-Q_\chi'/Q_\chi(\rho+1))$.
The conductor cancels. Digamma reflection and recurrence
\cite[Section~5.5, equations (2), (4)]{dlmf} evaluate this last
term as $h_a(\gamma)$. For even parity use
\[
 \Re\{\psi(\rho/2)-\psi((\rho+1)/2)\}
       =-\frac\pi{\cosh(\pi\gamma)};
\]
the recurrence $\psi(1+\rho/2)=\psi(\rho/2)+2/\rho$
gives the odd case. Both corrections are positive, since
$\cosh(\pi\gamma)\ge1+\pi^2\gamma^2/2>
\pi(\gamma^2+1/4)$.

For zeta use the entire completion
$\xi(s)=s(s-1)\pi^{-s/2}\Gamma(s/2)\zeta(s)/2$.
Its pole-removal factors contribute
\[
 \Re\{1/\rho+1/(\rho-1)-1/(\rho+1)\}
             =-3/[2(\gamma^2+9/4)],
\]
giving the first inequality in
\eqref{eq:zeta-bending-positive-new}. For the second inequality we use the following
small-height certificate.
The positive theta representation and the normalization
$H_0(x)=\xi(1/2+ix/2)/8$ give
\[
 \frac{H_0(x)}{H_0(0)}\ge1-c_1x^2,\qquad
 c_1=\frac{\xi''(1/2)}{8\xi(1/2)}.
\]
Ball power-series arithmetic at 80 and 120 decimal digits,
using the regular expression
$\xi(s)=(s-1)\pi^{-s/2}\Gamma(1+s/2)\zeta(s)$,
certifies
\[
 0.00577624827885474269<c_1<0.00577624827885474270<1/128.
\]
Thus $H_0(x)>H_0(0)/2$ for $|x|\le8$, and every
critical-line zeta zero has $|\gamma|>4$. The two subtracted
terms in \eqref{eq:zeta-bending-positive-new} are then less
than $(3/2+1/\pi)/16<11/96<1$. The certificate consists
only of finite ball inequalities; the passage to all these
heights uses $\cos v\ge1-v^2/2$ in the positive integral.

Finally, substitute $w=c\tau+d\tau^2+\cdots$, $u=\tau^m$,
in the displayed deformation. The coefficients of $\tau^m$
and $\tau^{m+1}$ give
\[
 A_mc^m=\rho L(\rho+1),\qquad
 mA_mc^{m-1}d+A_{m+1}c^{m+1}
       -\{L(\rho+1)+\rho L'(\rho+1)\}c=0.
\]
The $u^2$ term has degree $2m\ge m+2$ and does not enter.
The distinct nonzero roots $c$ and the implicit function
theorem give all the sections and their convergent expansions.
Their regular polygon has two vertical vertices for even $m$
and one for odd $m$ whenever it has any; the other vertices
split equally by the sign of their real part. Adding the
strictly left-bending vertices gives the asserted counts.
\end{proof}

\begin{lemma}[Exterior phase bound]
\label{lem:exterior-sector-new}
For every Dirichlet character, with the continuous Euler
argument,
\[
 \left|\arg\frac{L(5/2+it)}{L(3/2+it)}\right|<3/4.
\]
Writing $R=L(\rho+2)/L(\rho+1)$, one consequently has
\begin{equation}\label{eq:exterior-positive-new}
 P_\rho:=\Im((\rho+1)R)
       >|R|\frac{23\gamma-36}{32}>0\qquad(\gamma\ge2).
\end{equation}
\end{lemma}

\begin{proof}
The absolute Euler logarithm is at most
$\log(\zeta(3/2)/\zeta(5/2))$. The integral tail after five,
and the first four terms respectively, give the rational bounds
\[
 \begin{split}
 \zeta(3/2)&<\frac98+\frac5{14}+\frac7{36}+\frac{99}{100}
                =\frac83-\frac1{12600},\\
 \zeta(5/2)&>\frac{33}{32}+\frac7{40}+\frac5{78}
                >\frac{127}{100}.
 \end{split}
\]
These follow from $\sqrt2>7/5$, $\sqrt3>12/7$,
$\sqrt5>20/9$, $\sqrt2<10/7$, and $\sqrt3<26/15$,
each checked by squaring. Also
$e^{3/4}>\sum_{j=0}^4(3/4)^j/j!=4331/2048>800/381$.
This proves the phase bound. For $|\theta|<3/4$,
$\cos\theta>23/32$ and $\sin\theta>-3/4$;
expanding $\Im((3/2+i\gamma)|R|e^{i\theta})$
proves \eqref{eq:exterior-positive-new}.
\end{proof}

\begin{proposition}[The second jet at a simple tangent]
\label{prop:simple-tangent-new}
At every simple critical-line zero, let
$v=\rho L(\rho+1)/L'(\rho)$. The simple section has
\begin{equation}\label{eq:simple-bending-jet-new}
 s(u)=\rho+vu+
 \left(D_\rho v^2-\frac{(\rho+1)L(\rho+2)}{2L(\rho+1)}v\right)u^2
 +O(u^3).
\end{equation}
At a vertical velocity $v=i\beta$, put $d=\Re D_\rho$.
For a primitive nonprincipal character with $\gamma\ge2$,
and for every positive-height zeta zero, $d>0$, $P_\rho>0$,
and the second normal coefficient is exactly
\[
       -\beta^2d+\frac{\beta P_\rho}{2}
       =\beta d\left(\frac{P_\rho}{2d}-\beta\right).
\]
It is strictly negative for $\beta<0$. For $\beta>0$ it is
positive below $P_\rho/(2d)$ and negative above that value;
it vanishes at equality.
\end{proposition}

\begin{proof}
For $m=1$ the $u^2$ deformation term must be retained.
Substitute $s=\rho+vu+bu^2$ and match the coefficient of
$u^2$:
\[
 L'(\rho)b+\tfrac12L''(\rho)v^2
 -\{L(\rho+1)+\rho L'(\rho+1)\}v
 +\tfrac12\rho(\rho+1)L(\rho+2)=0.
\]
Division by $L'(\rho)$ proves the formula. The previous
proposition proves $d>0$ also at multiplicity one, and the
exterior phase lemma proves $P_\rho>0$ in the stated ranges.
Taking real parts with $v=i\beta$ proves every sign assertion.
In the notation of the corrected phase, differentiation gives
$d=\psi_\chi'(\gamma)$ exactly: the completed local germ
contributes $\theta_\chi'$, the exterior logarithm contributes
$\Re(L'/L)(\rho+1)$, and
$(-\arctan(1/(2\gamma)))'=\Re(1/\rho)$.
\end{proof}

\paragraph{The central real-character specialization.}
For a real nonprincipal primitive character, the completed
functional equation has root number $+1$; its exact sign is
proved through the quotient of the quadratic-field and rational
zeta completions in \cite[Section~1, equations (1.3)--(1.6)]{gross-signs}.
The central multiplicity $m$ is therefore even. The cited
zero-location theorem and positivity of the Euler product on
$(1,\infty)$ make the leading central coefficient $A_m$
positive, including $m=0$. The paired product gives
\[
 C_\chi:=2+\frac{L'}L(3/2)+\frac{Q_\chi'}{Q_\chi}(1/2)
 =m+\sum_{\gamma>0}\frac2{1+\gamma^2}
 +\begin{cases}2-\pi/2,&a=0,\\ \pi/2,&a=1.\end{cases}
\]
The sum includes multiplicity. At a central zero,
\eqref{eq:multiple-bending-jets-new} gives
\[ c_j^m=L(3/2)/(2A_m)>0,\qquad D_\rho=C_\chi. \]
For $m\equiv2\pmod4$ its branches split equally left and
right; for $m\equiv0\pmod4$, $m\ge4$, the two vertical
vertices give $(m/2+1,m/2-1)$. These formulas assign the
local branches of the central divisor without asserting that
this divisor has positive multiplicity for any character.

\paragraph{The exact Dedekind map.}
For the norm shift $\sum_{\mathfrak a}(N\mathfrak a+u)^{-s}$,
the same rising-factorial expansion replaces $L$ by $\zeta_K$.
The entire completion has pole-removal factors $s(s-1)$,
$r_1$ gamma factors $\Gamma(s/2)$ and $r_2$ factors $\Gamma(s)$.
The preceding calculation therefore gives the explicit identity
\[
 \Re D_\rho=\Re\frac{\xi_K'}{\xi_K}(\rho+1)
 -\frac3{2(\gamma^2+9/4)}
 -\frac{r_1\pi}{2\cosh(\pi\gamma)}
 -\frac{r_2}{2(\gamma^2+1/4)}.
\]
Its first term exceeds $m$, so $\Re D_\rho>0$ whenever
$\gamma^2\ge ((r_2+3)/2+r_1/\pi)/m$.
For a quadratic field the exact factorization
$\zeta_K=\zeta L(\,\cdot\,,\chi_K)$ gives an exterior argument
of absolute value below $3/2$. The alternating cosine
polynomial gives $\cos(3/2)>7/100$, whence
$P_\rho>|R|(7\gamma/100-3/2)>0$ for $\gamma\ge24$.
Thus the same second-jet sign formula holds there. The displayed completion and the quadratic factorization
give the corresponding maps for these coefficient identities.

\section{Positive proportions of both horizontal directions}
\label{sec:signed-density}
The local velocity formula becomes a counting theorem by combining a
mixed phase average with a square mollifier. The phase average controls
where the horizontal component can be small; the mollifier controls
how strongly the derivative mass can concentrate there.

Fix a primitive character $\chi$ of conductor $q$, including $q=1$ for
zeta, and write $L_\chi(s)=L(s,\chi)$. All limits in this section keep
$\chi$ fixed. The translated family whose tangent is counted is
\[
 F_\chi(u,s)=\sum_{n\ge1}\chi(n)(n+u)^{-s}
 =q^{-s}\sum_{r=1}^q\chi(r)\zeta\left(s,\frac{r+u}{q}\right).
\]
The first identity holds for $\Re s>1$ and $|u|<1$; the finite
Hurwitz sum continues it near every nontrivial zero. At $u=0$ it
is $L_\chi(s)$ and $\partial_uF_\chi(0,s)=-sL_\chi(s+1)$,
so its simple-zero derivative is the velocity displayed below.
We use the zero-location theorem stated in the introduction,
including its assertion for the auxiliary Dirichlet characters in the
proof of Lemma~\ref{lem:dirichlet-additive}. Let $N_{\chi,T}$ count the
zero occurrences $\rho=1/2+i\gamma$, $T<\gamma\le2T$, with multiplicity.

\begin{theorem}\label{thm:positive-directions}
For either choice of sign, the number $n_{\chi,\pm}(T)$ of distinct
simple zeros in this interval satisfying
\[
 \pm\Re\frac{\rho L_\chi(\rho+1)}{L_\chi'(\rho)}>0
\]
obeys
\[
 \liminf_{T\to\infty}\frac{n_{\chi,\pm}(T)}{N_{\chi,T}}
       \ge10^{-14}.
\]
The numerical constant is independent of the fixed character. The height
from which a bound with any smaller constant holds may depend on it.
\end{theorem}

\subsection{A mixed phase average}
Let $\epsilon_\chi$ be the root number, $e_\chi\in\{0,1\}$ its parity,
and
\[
 Q_\chi(s)=(q/\pi)^{(s+e_\chi)/2}\Gamma((s+e_\chi)/2),\qquad
 \Delta_\chi(s)=\epsilon_\chi Q_\chi(1-s)/Q_\chi(s).
\]
Choose a continuous phase $\theta_\chi$ for which
$Z_\chi(t)=e^{i\theta_\chi(t)}L_\chi(1/2+it)$ is real. Then
$\Delta_\chi(1/2+it)^{-1}=e^{2i\theta_\chi(t)}$.
For $\sigma>1$, let $\phi_{\chi,\sigma}(t)$ be the imaginary part
of the Euler logarithm of $L_\chi(\sigma+it)$.

\begin{lemma}\label{lem:mixed-phase}
With $\delta_t=\arctan(1/(2t))$ and
$\psi_\chi(t)=\theta_\chi(t)+\phi_{\chi,3/2}(t)-\delta_t$, one has
\[
 \sum_{T<\gamma\le2T}\cos^2\psi_\chi(\gamma)
       =(\tfrac12+o_\chi(1))N_{\chi,T}.
\]
At a simple zero the absolute horizontal component of its velocity,
divided by the absolute velocity, is $|\cos\psi_\chi(\gamma)|$.
\end{lemma}

\begin{proof}
For every fixed $x>0$ we first prove
\begin{equation}\label{eq:twisted-hardy}
 \sum_{0<\gamma\le T}e^{2i\theta_\chi(\gamma)}x^{i\gamma}
 =-\frac{i^{e_\chi}\epsilon_\chi^{-1}}{\sqrt{qx}}
       \sum_{n\le qxT/(2\pi)}\chi(n)\Lambda(n)e^{-2\pi in/(qx)}
       +O_{\chi,x}(\sqrt T\log^2T)=O_{\chi,x}(T).
\end{equation}
Here and below zero sums count multiplicity. Integrate
$(L_\chi'/L_\chi)(s)\Delta_\chi(s)^{-1}x^{s-1/2}/(2\pi i)$
on a rectangle with vertical sides $-3/2$ and $c=1+1/\log T$,
fixed bottom ordinate above $2\pi$, and top within one unit of $T$
at distance $\gg_\chi1/\log T$ from all ordinates. The zero count in
any unit interval is $O_\chi(\log T)$, from the completed logarithmic
derivative at $2+it$; its partial fractions give
$L_\chi'/L_\chi=O_\chi(\log^2T)$ on the top. Stirling gives
\[
 \Delta_\chi(\sigma+it)^{-1}
 =\epsilon_\chi^{-1}(qt/(2\pi))^{\sigma-1/2}
 e^{i(t\log(qt/(2\pi))-t+\pi(2e_\chi-1)/4)}(1+O_\chi(t^{-1})).
\]
The gamma factor on the left edge is $O(t^{-2})$, so that edge
has a bounded integral. The top contributes $O(\sqrt T\log^2T)$.
Expanding the right edge gives the coefficient $-\chi(n)\Lambda(n)$ and the phase
$2\pi u\log(qxu/(en))$, $u=t/(2\pi)$.

For clarity, the stationary-phase error is summable here. Put
$r=n/(qx)$. Outside $[r/2,2r]$, the first derivative of
$u\log(u/(er))$ has absolute value at least $\log2$;
integration by parts gives $O(\sqrt T\log T)$ after summation.
The range $r\le1$ is treated by splitting at $u=2$: the bounded
lower interval and the first-derivative test above it contribute
$O_{\chi,x}(\sqrt T\log T)$. Thus the stationary formula below is
used only for $r>1$.
On $[A,B]=[1,T/(2\pi)]\cap[r/2,2r]$, the stationary-phase
formula of Arias de Reyna \cite[Lemmas~2.2 and~2.5, p.~5]{arias}
is
\[
 \mathbf1_{A\le r\le B}r^c e^{i\pi/4-2\pi ir}
 +O\left(r^{c-1/2}+
 \frac{r^{c+1/2}}{|A-r|+\sqrt r}+
 \frac{r^{c+1/2}}{|B-r|+\sqrt r}\right).
\]
The first and lower-end errors are $O(\sqrt T\log T)$ using
$\sum_{n\le X}\Lambda(n)n^{-1/2}\ll\sqrt X$. Near the upper
endpoint the error is bounded by
$C\sqrt T\log T\sum_{n\asymp T}(|n-qxT/(2\pi)|+\sqrt T)^{-1}
\ll\sqrt T\log^2T$.
The phase constants multiply to $i^{e_\chi}$, proving the first
identity in \eqref{eq:twisted-hardy}; Chebyshev's prime bound proves
the second. Moving the top back to $T$ changes at most $O(\log T)$
unit-modulus summands.

The absolutely convergent series
\[
 \frac{L_\chi(\sigma+it)}{\overline{L_\chi(\sigma+it)}}
 =\sum_{m,n\ge1}\frac{\chi(n)\mu(m)\overline{\chi(m)}}{(mn)^\sigma}
                e^{it\log(m/n)}
\]
has absolute coefficient sum at most $\zeta(\sigma)^2$.
Apply \eqref{eq:twisted-hardy} at each fixed $x=m/n$, divide by
$N_\chi(T)\sim T\log T/(2\pi)$, and use dominated convergence.
This proves that the mixed mode
$e^{2i\theta_\chi(\gamma)}L_\chi(\sigma+i\gamma)/
\overline{L_\chi(\sigma+i\gamma)}$ has mean zero. Multiplication
by $e^{-2i\delta_\gamma}$ changes the mean by $o(1)$.
Subtracting the formulas at $2T$ and $T$, and taking real parts,
gives the assertion about $\cos^2\psi_\chi$.
At a simple zero,
$L_\chi'(\rho)=-ie^{-i\theta_\chi(\gamma)}Z_\chi'(\gamma)$
with real nonzero $Z_\chi'$. Since $\arg(i\rho)=\pi-\delta_\gamma$,
the velocity formula gives the last assertion exactly.
\end{proof}

\begin{corollary}[The size of the normal component]
\label{cor:unsigned-transversality}
Let $N_{\chi,s,T}$ count simple zeros in $(T,2T]$, and at
each such zero put $r_\rho=|\Re v_\rho|/|v_\rho|$.
Then
\[
 \liminf_{T\to\infty}\frac1{N_{\chi,T}}
       \sum_{\rho\text{ simple},\ T<\gamma\le2T}r_\rho^2
       \ge\frac16.
\]
For $0<\eta<1/2$, the number $G_{\chi,\eta}(T)$ of these
zeros with $r_\rho\ge\eta$ obeys
\[
 \liminf\frac{G_{\chi,\eta}(T)}{N_{\chi,T}}
 \ge\frac{1-4\eta^2}{6(1-\eta^2)},\qquad
 \liminf\frac{G_{\chi,\eta}(T)}{N_{\chi,s,T}}
 \ge\frac{1-4\eta^2}{4(1-\eta^2)}.
\]
In particular, the corresponding bounds at $\eta=1/4$ are
$2/15$ and $1/5$. For every fixed $\varepsilon>0$, the zeros
counted at a positive threshold also satisfy
$|\Re v_\rho|\gg_{\chi,\eta,\varepsilon}T^{1-\varepsilon}$.
The same limiting bounds hold with the interval $(0,T]$.
\end{corollary}

\begin{proof}
The degree-one pair-correlation theorem
\cite[Theorems~1.1--1.2]{rudnick-sarnak} gives
$N_{\chi,s,T}\ge(2/3-o(1))N_{\chi,T}$.
Here is its exact counting consequence. Take nonnegative
Schwartz tests $g$ with $g(0)=1$, Fourier support in $(-1,1)$,
and sine-kernel integrals tending to $1/3$, constructed in
Section~\ref{sec:survival}. A center of multiplicity $m\ge2$
contributes $m(m-1)\ge m$ to the ordered off-diagonal sum.
Subtracting the global pair sums at $2T$ and $T$ bounds the
nonsimple occurrences by $(1/3+o(1))N_{\chi,T}$; all other
terms are nonnegative.

Abbreviate $N=N_{\chi,T}$, $N_s=N_{\chi,s,T}$ and $G=G_{\chi,\eta}(T)$.
Lemma~\ref{lem:mixed-phase} gives
\[
 \sum_{\rho\text{ simple}}r_\rho^2
       \ge N_s-N/2+o(N),\qquad
 \sum_{\rho\text{ simple}}r_\rho^2
       \le G+\eta^2(N_s-G).
\]
Thus $G\ge N_s-N/[2(1-\eta^2)]+o(N)$, which proves
both fractions and the mean-square assertion. The absolutely
convergent Euler product bounds $|L(\rho+1)|$ below by
$\zeta(3/2)^{-1}$. The character growth bound and Cauchy's
formula in Section~\ref{sec:dirichlet-moments} give
$|L'(\rho)|\ll_{\chi,\varepsilon}T^\varepsilon$;
the exact velocity formula proves the speed estimate.
Finally sum over finitely many dyadic intervals and then let
the number of intervals tend to infinity. The omitted initial
segment has arbitrarily small relative zero count, proving
the versions with cutoff $T$.
\end{proof}

\subsection{The square mollifier and its two constants}
Let $d_{-1/2}(n)$ be the coefficients of $\zeta(s)^{-1/2}$ in
$\Re s>1$. They are real, multiplicative and bounded in modulus by
one. Fix $0<\nu<1/4$ and a real polynomial $P$ satisfying $P(1)=1$
and having a zero of order at least sixteen at zero. Define
\[
 B_{\chi,T}(s)=\sum_{n\le T^\nu}d_{-1/2}(n)\chi(n)n^{-s}
                 P\left(\frac{\log(T^\nu/n)}{\nu\log T}\right).
\]
Write $M=T^\nu$ and
$b_T(n)=d_{-1/2}(n)P(\log(M/n)/\log M)\mathbf1_{n\le M}$.
For $k\ge1$, the binomial expansion gives
\[
 d_{-1/2}(p^k)=-\frac{\binom{2k}{k}}{4^k(2k-1)}.
\]
With $C_P=\max(1,\sup_{[0,1]}|P|)$, the exact squared coefficient
and its bound are therefore
\[
 a_T(n)=\sum_{uv=n}b_T(u)b_T(v),\qquad
 |a_T(n)|\le C_P^2\tau_2(n),\qquad a_T(n)=0\quad(n>M^2).
\]
Thus the square has length $T^{2\nu}<T^{1/2}$. The contour estimates
below apply this divisor envelope to the literal coefficients.
Write $L=\log T$.

\begin{proposition}\label{prop:square-moments}
There are real constants $a=a(P,\nu)$ and $b=b(P,\nu)$ such that
\begin{align}
 \sum_{T<\gamma\le2T}L_\chi'(\rho)B_{\chi,T}(\rho)^2
       &=(a+o_\chi(1))N_{\chi,T}L,\label{eq:square-first}\\
 \sum_{T<\gamma\le2T}|L_\chi'(\rho)B_{\chi,T}(\rho)^2|^2
       &=(b+o_\chi(1))N_{\chi,T}L^2.\label{eq:square-second}
\end{align}
They are independent of the fixed character. In addition, putting
\[
 K_{\rho,\chi}=\frac{\Re\{\rho L_\chi(\rho+1)
                  \overline{L_\chi'(\rho)}\}}{\gamma},
\]
one has
\begin{equation}\label{eq:square-cancel}
 \sum_{T<\gamma\le2T}K_{\rho,\chi}|B_{\chi,T}(\rho)|^2
       =o_\chi(N_{\chi,T}L).
\end{equation}
If $I_P=\int_0^1P(t)dt$ and $J_P=\int_0^1(1-t)P(t)dt$, then
\[
 a(P,\nu)=\tfrac12+\nu I_P+\tfrac{\nu^2}{4}(I_P^2-J_P).
\]
The constant $b$ is the finite integral defined below.
\end{proposition}

The discrete moment formulas used to prove this proposition are given,
with their contour proofs for all fixed primitive characters and the
displayed divisor envelopes, in Section~\ref{sec:dirichlet-moments}.
For zeta, the main terms agree with Benli--Elma--Ng
\cite[Theorem~1.1, equations~(9)--(14), pp.~3--4]{benli-elma-ng}
and Heap--Li--Zhao \cite[Theorem~2, equations~(5)--(9), pp.~3--5]{heap-li-zhao}.
The coefficient hypotheses of the first reference include its
submultiplicativity condition~(7); the contour derivation here proves
the needed formula directly from the divisor envelope above.
The second moment counts zero occurrences with multiplicity.
We give the kernel evaluation and
the rational inequality here, since they determine the positive counting
proportion.

Introduce variables $s_1,s_2,t_1,t_2$ and formal series
\[
 E_L(x)=\prod_{i=1}^2(1+x/s_i)^{1/2}=\sum l_kx^k,
 \qquad E_R(x)=\prod_{j=1}^2(1+x/t_j)^{1/2}=\sum r_kx^k.
\]
Define the following two finite polynomials:
\begin{align*}
 j={}&l_1r_1+l_1r_2+\tfrac12l_1+l_2r_1-l_3
                       +\tfrac12r_1-r_3+\tfrac13,\\
 i={}&l_1r_1^2+l_1r_1r_2+l_1r_3-\tfrac13l_1
       +l_2r_1^2-l_2r_1-l_2r_2-\tfrac12l_2-2l_3r_1+l_4\\
    &+\tfrac12r_1^2+r_1r_2-r_1r_3+\tfrac16r_1+r_2^2
                  +\tfrac12r_2-r_4-\tfrac18.
\end{align*}
Expand
\[
 j+2i=\sum_{\mathbf k}c_{\mathbf k}s_1^{-k_1}s_2^{-k_2}
                                      t_1^{-k_3}t_2^{-k_4}.
\]
This specifies sixty-two rational coefficients, of total degree at most
four. For $P(z)=\sum p_jz^j$ set
\[
 \Psi_k^P(z)=\sum_{j\ge1}p_j
        \frac{\Gamma(j+1)}{\Gamma(j+k+1/2)}z^{j+k-1/2}.
\]
On the domain
\[
 \mathcal D=\{x_{ij}\ge0:\ \sum_jx_{ij}\le1,\ \sum_ix_{ij}\le1\}
\]
put $z_1=1-x_{11}-x_{12}$, $z_2=1-x_{21}-x_{22}$,
$z_3=1-x_{11}-x_{21}$ and $z_4=1-x_{12}-x_{22}$. Then
\begin{equation}\label{eq:square-b}
 b(P,\nu)=\frac1{\Gamma(1/4)^4}
 \int_{\mathcal D}\prod_{i,j}x_{ij}^{-3/4}
   \sum_{\mathbf k}c_{\mathbf k}\nu^{|\mathbf k|-1}
            \prod_{l=1}^4\Psi_{k_l}^P(z_l)\,\prod_{i,j}dx_{ij}.
\end{equation}
All singularities of the weight are integrable.

\begin{proof}[Evaluation of the constants in Proposition~\ref{prop:square-moments}]
The first-moment residue calculation in Section~\ref{sec:dirichlet-moments}
gives, writing $a(n)$ for the coefficients of $B_{\chi,T}^2$, the leading
expression
\[
 \tfrac12L^2+\sum_p\frac{a(p)}p
       (-L\log p+\tfrac12\log^2p)
       +\sum_{p<r}\frac{a(pr)}{pr}\log p\log r+O_{\chi,P}(L).
\]
For $u=\log p/\log(T^\nu)$ and $v=\log r/\log(T^\nu)$, the
coefficients are
\[
 a(p)=-P(1-u)\mathbf1_{u\le1},\qquad
 a(pr)=\tfrac12\{P(1-u-v)\mathbf1_{u+v\le1}
       +P(1-u)P(1-v)\mathbf1_{u\le1,\ v\le1}\}.
\]
The measure $(\log Y)^{-1}\sum_{p\le Y}(\log p)/p\,
\delta_{\log p/\log Y}$ tends weakly to Lebesgue measure on $[0,1]$,
by the prime number theorem and partial summation. Diagonal prime
pairs and repeated prime powers contribute $o(L^2)$. The prime
integral is $\nu I_P-\nu^2J_P/2$; the pair integral is
$\nu^2(J_P+I_P^2)/4$. Their sum gives $a$.

We describe every limiting operation for the second constant. Use four
inverse Mellin variables for the four factors of $B$. Scaling each
variable by $L$ gives the weight
\[
 W(s)=e^{\nu s}\sum_j\frac{p_jj!}{\nu^j s^{j+1}},\qquad\Re s=1/\nu.
\]
The gcd sum in the continuous part of the discrete moment has Euler
product
\begin{equation}\label{eq:square-euler-F}
 C(u,v;\alpha,\beta)
 \frac{\prod_{i,j}Z_q(1+u_i+v_j)^{1/4}}
 {\prod_i Z_q(1+u_i+\beta)^{1/2}
       \prod_j Z_q(1+v_j+\alpha)^{1/2}},
\end{equation}
where $Z_q(s)=\zeta(s)\prod_{p\mid q}(1-p^{-s})$. The shifted
right-edge part has Euler product
\begin{align}\label{eq:square-euler-G}
 D(u,v;\alpha,\beta,\gamma)&
 \frac{Z_q(1+\alpha+\beta)Z_q(1+\beta+\gamma)}{Z_q(1+\beta)}
 \frac{\prod_{i,j}Z_q(1+u_i+v_j)^{1/4}}
       {\prod_iZ_q(1+u_i+\beta)^{1/2}}\nonumber\\
 &\hspace{10mm}\cdot\prod_j
 \frac{Z_q(1+v_j)^{1/2}}
 {Z_q(1+v_j+\alpha)^{1/2}Z_q(1+v_j+\gamma)^{1/2}}.
\end{align}
To verify these identities, pair one prime from the left coefficient
with one from the right. For the first product the linear terms are
$-\frac12\sum_i p^{-1-u_i-\beta}
-\frac12\sum_jp^{-1-v_j-\alpha}
+\frac14\sum_{i,j}p^{-1-u_i-v_j}$.
For the second, use the balanced relation $hm=kn$, with coefficient
$\mu*n^{-\alpha}*n^{-\gamma}$ on $m$ and $n^{-\beta}$ on $n$;
the same pairing gives precisely the displayed factors.
After removing these linear terms, the local errors are uniformly
$O(p^{-3/2})$ when the real parts are sufficiently near zero.
Thus $C,D$ converge absolutely there. At the origin the two half-power
series on each coefficient side combine to $\mu$. The first local
gcd sum is $1-1/p$; in the second, the combined left and right
coefficients are both $\delta_1$. These observations prove
$C(0)=D(0)=1$, including the prime powers. Primes dividing $q$ are
omitted from both sides.

Put $\alpha=A/L$, $\beta=B/L$, $\gamma=C_1/L$, on distinct fixed
small Cauchy circles in $A,B,C_1$. Every coefficient-variable zeta
argument remains to the right of one. On the complete Mellin lines,
$|Z_q(1+\epsilon+it)|+|Z_q(1+\epsilon+it)^{-1}|=O_q(L)$ for
$\epsilon\asymp1/L$. The two Euler expressions above are therefore
$O_q(L^3)$ and $O_q(L^6)$, respectively. The four weights decay at
least as $(1+|\Im s|)^{-17}$. Removing any region
$|\Im s|>\sqrt L$ costs $O_q(L^{-5})$ and $O_q(L^{-2})$, respectively;
these are smaller than the respective scales $L^{-1}$ and one.
On the remaining box use
$Z_q(1+z)=E_qz^{-1}(1+O_q(z))$, $E_q=\varphi(q)/q$.
The first expression has total exponent $-1$ and its explicit pole
in the continuous part cancels $E_q^{-1}$; the second expression has
total exponent zero. Thus no conductor factor remains.

Write
\[
 F_0(s,t)=\frac{(s_1s_2t_1t_2)^{1/2}}
                    {\prod_{i,j}(s_i+t_j)^{1/4}}.
\]
The two scaled limits are $F_0E_L(B)E_R(A)$ and
\[
 R(A,B,C_1)=F_0\frac{B}{(A+B)(B+C_1)}E_L(B)E_R(A)E_R(C_1).
\]
They have polynomial growth, while $s_i+t_j$ stays in a right
half-plane. The Mellin weights give an integrable majorant;
dominated convergence and Cauchy's formula justify every shift
derivative. The continuous part has normalized kernel
\[
 F_0\frac{E_L(B)E_R(A)-e^{-A-B}E_L(-A)E_R(-B)}{A+B},
\]
whose $AB$ coefficient is $F_0j$. The right-edge part is
$R(A,B,C_1)+e^{-A-B}R(-B,-A,C_1)+e^{-B-C_1}R(A,-C_1,-B)$.
Dividing by $F_0$, its numerator over
$(A+B)(B+C_1)(A-C_1)$ is
\begin{align*}
 &B(A-C_1)E_L(B)E_R(A)E_R(C_1)\\
 &\quad-A(B+C_1)e^{-A-B}E_L(-A)E_R(-B)E_R(C_1)\\
 &\quad+C_1(A+B)e^{-B-C_1}E_L(-C_1)E_R(A)E_R(-B).
\end{align*}
Each of the three denominator factors divides this numerator as a
convergent power series. Expansion through total degree six gives
$ABC_1$ coefficient $i$. Reflection gives the second copy, so the
full kernel is $F_0(j+2i)$.

Finally write
$(s_i+t_j)^{-1/4}=\Gamma(1/4)^{-1}\int_0^\infty
 x^{-3/4}e^{-x(s_i+t_j)}dx$. Absolute convergence justifies Fubini.
The single-variable inverse Laplace integral is
\[
 \frac1{2\pi i}\int W(s)s^{1/2-k}e^{-xs}ds
 =\nu^{k-1/2}\Psi_k^P(1-x/\nu)\mathbf1_{x\le\nu}.
\]
Scaling the four cross variables by $\nu$ gives a factor $\nu$ from
the measures and $\nu^{|\mathbf k|-2}$ from the four vertices.
This is exactly \eqref{eq:square-b}. The zero-location theorem
identifies the differentiated factors with a positive squared modulus,
which proves \eqref{eq:square-second}. The contour argument proving
\eqref{eq:square-cancel} appears in Section~\ref{sec:dirichlet-moments}.
\end{proof}

\subsection{A rational strict inequality}
We now prove a strict inequality for the constants, including the
smoothing required by the preceding Mellin argument. Put
$P_0(z)=2z-z^2$ and $\nu_0=1/4$. The finite functional gives
$a_0=1543/2304$ and
\[
 \Psi_k^{P_0}(z)=\frac{2z^{k+1/2}}{\Gamma(k+3/2)}
                         \left(1-\frac z{k+3/2}\right).
\]
After factoring $\sqrt{z_1z_2z_3z_4}/\pi^2$ out of the kernel
sum in \eqref{eq:square-b}, denote its remaining polynomial by
$Q(S,D,E)$, where
$S=z_1+z_2=z_3+z_4$, $D=(z_1-z_2)^2$, $E=(z_3-z_4)^2$.
This defines $Q$ exactly through the sixty-two rational coefficients.
Its constant and linear-$D$ coefficients at $E=0$ are
\begin{align*}
 q_0(S)={}&\frac8{245581875}\bigl(776S^8+75340S^7+287460S^6
 -5504100S^5-11151525S^4\\
 &\hspace{23mm}+151420500S^3+43659000S^2
 -1746360000S+2619540000\bigr),\\
 q_D(S)={}&-\frac8{245581875}\bigl(5207S^6-140S^5+948840S^4
 +1451100S^3\\
 &\hspace{23mm}-49075950S^2+30492000S+261954000\bigr).
\end{align*}
For $0\le S\le2$ and $0\le D,E\le\min(S,2-S)^2$, one has
\begin{equation}\label{eq:kernel-majorant}
 Q(S,D,E)\le U(S,D,E):=q_0(S)+q_D(S)D-(2-3S/4)E.
\end{equation}
Here $q_D<0$, and $U\ge0$ on the larger domain
$0\le D\le\min(S,2-S)^2$, $0\le E\le(2-S)^2$.
We give a finite rational certificate. Set
\[
 A=-[D^2]Q(S,D,0)-[D^3]Q(S,D,0)\,D,\qquad
 B=-(2-3S/4)-\frac{Q(S,D,E)-Q(S,D,0)}E.
\]
The identity $U-Q=D^2A+EB$ is exact. After the substitutions
$(S,D,E)=(u,u^2v,u^2w)$ and $(2-u,u^2v,u^2w)$, all tensor
Bernstein coefficients of $A$ and $B$ on the unit cube are positive.
Their minimum values are
\[
 \begin{array}{c|cc}
 &S=u&S=2-u\\\hline
 A&5060984/81860625&64264/3898125\\
 B&4133699/13097700&38706253/491163750
 \end{array}
\]
The coefficients of $q_D(2u)$ are negative. The minimum Bernstein
coefficients of
$q_0(S)+q_D(S)\min(S,2-S)^2-(2-3S/4)(2-S)^2$
on the two halves are $167144539/6682500$ and
$1492032128/245581875$. These checks use the exact conversion
\[
 [B_{\mathbf k}]p=\sum_{\mathbf j\le\mathbf k}a_{\mathbf j}
       \prod_l\frac{\binom{k_l}{j_l}}{\binom{n_l}{j_l}},
 \qquad p(x)=\sum a_{\mathbf j}x^{\mathbf j},
\]
where $n_l$ is the degree in variable $l$. Bernstein basis functions
are nonnegative and sum to one, so these rational inequalities prove
\eqref{eq:kernel-majorant} on its entire domain.

Set $r_i=x_{i1}+x_{i2}$ and $x_{11}=r_1v_1$, $x_{21}=r_2v_2$.
Then $\sqrt{z_3z_4}\le S/2$, and positivity of $U$ allows the
column constraints to be dropped after replacing $Q$ by $U$.
The $v_i$ have normalized beta law $\operatorname{Beta}(1/4,1/4)$,
with $\mathbb EE=2(r_1^2+r_2^2)/3$. The remaining normalized row
variables $R_i$ have beta law $\operatorname{Beta}(1/2,3/2)$.
Their moments are $(1/2)_k/(2)_k$. Direct polynomial integration yields
\begin{align*}
 b(P_0,1/4)&\le\frac1{4\pi}\mathbb E\left[
 \frac S2\left(q_0(S)+q_D(S)(R_1-R_2)^2
 -(2-3S/4)\frac23(R_1^2+R_2^2)\right)\right]\\
 &=\frac{1505515846753}{536481792000\pi},\qquad S=2-R_1-R_2.
\end{align*}
The program \texttt{certify\_square\_mollifier} constructs the rational
kernel, records every Bernstein coefficient, and repeats this integration.
It uses exact rational arithmetic throughout.

Machin's identity and its alternating arctangent series give
$\pi>314159/100000$. Thus, writing
$B_0=1505515846753/[536481792000(314159/100000)]$,
\[
 2a_0^2-B_0=\frac{5046955580047}{1348324666343424}>0.
\]
To meet the high-order endpoint condition take
$P_K(z)=(2z-z^2)(1-(1-z)^K)^{15}$.
It vanishes to order sixteen, satisfies $P_K(1)=1$, and
$0\le P_K'\le32$. The derivative converges in $L^4[0,1]$ to
$P_0'$, by dominated convergence. The identities
\[
 \Psi_0^P(z)=\frac1{\sqrt\pi}\int_0^z\frac{P'(t)}{\sqrt{z-t}}dt,
 \qquad
 \Psi_k^P(z)=\frac1{\Gamma(k-1/2)}\int_0^zP(t)(z-t)^{k-3/2}dt
 \quad(k\ge1)
\]
and H\"older's inequality prove uniform convergence of every vertex
needed in \eqref{eq:square-b}. In particular
$\|\Psi_0^{P_K}-\Psi_0^{P_0}\|_\infty
\le3^{3/4}\|P_K'-P_0'\|_4/\sqrt\pi$.
The kernel weight is integrable, so $a$ and $b$ converge as well.
They are continuous in $\nu>0$. We may therefore fix once and for
all a rational $\nu<1/4$ sufficiently close to $1/4$, and then a
sufficiently large integer $K$, for which
\[
 a\ge a_*:=1543/2304-1/2000,\qquad
 b\le b_*:=B_0+1/2000<1.
\]
Both parameters remain fixed as $T$ tends to infinity.

\begin{proof}[Proof of Theorem~\ref{thm:positive-directions}]
Choose the fixed mollifier just constructed. Let $\eta=1/50$ and
$p_\eta=1/[2(1-\eta^2)]$. Lemma~\ref{lem:mixed-phase} shows that
the occurrences with $|\cos\psi_\chi|<\eta$ number at most
$(p_\eta+o(1))N_{\chi,T}$: their contribution to the average
$\cos^2\psi_\chi$ is at most $\eta^2$, and every other contribution
is at most one. Put $A_\rho=L_\chi'(\rho)B_{\chi,T}(\rho)^2$.
The first moment and the triangle inequality give
$\sum|A_\rho|\ge(a_*+o(1))N_{\chi,T}L$.
Cauchy and the second moment bound the mass on the excluded set by
$(\sqrt{b_*p_\eta}+o(1))N_{\chi,T}L$. Thus its complement carries
at least $(d+o(1))N_{\chi,T}L$, where
$d=a_*-\sqrt{b_*p_\eta}$.

The Euler product bounds
$\zeta(3/2)^{-1}\le|L_\chi(3/2+it)|\le\zeta(3/2)$ and the exact
phase formula imply
\[
 |K_{\rho,\chi}|\,|B_{\chi,T}(\rho)|^2
  =\frac{|\rho|}{\gamma}|L_\chi(\rho+1)|\,|A_\rho|
                  |\cos\psi_\chi(\gamma)|.
\]
The total absolute weighted mass is therefore at least
$(\eta d/\zeta(3/2)+o(1))N_{\chi,T}L$.
Equation~\eqref{eq:square-cancel} divides it equally between the
two signs, up to $o(N_{\chi,T}L)$. On either sign set the second
moment of this weighted statistic is at most
$(\zeta(3/2)^2b_*+o(1))N_{\chi,T}L^2$. A further application of
Cauchy gives for each sign the lower proportion
\[
 \frac{\eta^2d^2}{4\zeta(3/2)^4b_*}.
\]
Every nonzero summand belongs to a simple center: at a multiple zero
$L_\chi'(\rho)=0$. At a simple center its sign equals that of the
velocity, because multiplication by $|L_\chi'(\rho)|^2/\gamma$ is
positive. The exact rational computation gives
$a_*^2-b_*p_\eta>1/2000$ and $0<a_*,b_*<1$, hence $d>1/4000$.
Finally $\zeta(3/2)<3$, by the integral test. The last lower
proportion exceeds
$1/12960000000000>10^{-14}$, as claimed.
\end{proof}

\section{The discrete moments for a fixed primitive character}
\label{sec:dirichlet-moments}
This section proves the contour formulas used in
Proposition~\ref{prop:square-moments}. Keeping the character phases in
the additive transform is essential: their cancellation, followed by
character orthogonality, gives the same leading constants for every
fixed primitive character.

Throughout, $q$ is fixed, $Z_q(s)=\zeta(s)\prod_{p\mid q}(1-p^{-s})$,
$E_q=\varphi(q)/q$, and coefficient polynomials have the form
$A_\chi(s)=\sum_{n\le y}a(n)\chi(n)n^{-s}$, where the real coefficients
satisfy $|a(n)|\le C\tau_r(n)(\log(2n))^r$ for fixed $C,r$.
Their supports satisfy $y\le T^\vartheta$ for a fixed $\vartheta<1/2$.
In particular, the squared mollifier has
$|a_T(n)|\le C_P^2\tau_2(n)$. In the character averages below its
subsequence $a_T(eh)$ is kept exactly and bounded by
$C_P^2\tau_2(e)\tau_2(h)$. This follows from the divisor envelope
and $\tau_2(eh)\le\tau_2(e)\tau_2(h)$; the argument requires no
product inequality for the coefficients $a_T$ themselves.
All shift variables have modulus at most a fixed constant divided by
$\log T$; all assertions below are uniform on slightly larger shift
circles, so Cauchy differentiation is permitted.

\subsection{Bounds for auxiliary characters}
We first record explicitly the consequence of the zero-location input
used for the growing auxiliary moduli. For every fixed $\eta,\epsilon>0$,
\begin{equation}\label{eq:uniform-character-growth}
 |L(\sigma+it,\psi)|+|L(\sigma+it,\psi)^{-1}|
 \ll_{\eta,\epsilon}\{Q(|t|+3)\}^{\epsilon},
 \qquad \tfrac12+\eta\le\sigma\le2,
\end{equation}
for nonprincipal characters modulo $Q$. For principal characters the
same bound holds on the line $\sigma=1/2+\eta<1$, and on the horizontal
parts of the strip when $|t|\ge2$.

Here is a proof including uniformity in $Q$. For primitive nonprincipal
$\psi$, the real part of the completed logarithmic derivative at $2+it$
is the positive zero sum
\[
 \sum_\rho\frac{3/2}{(3/2)^2+(t-\gamma)^2}
     =\Re\frac{\Lambda'}\Lambda(2+it,\psi)
     \ll\log\mathcal C,\qquad \mathcal C=Q(|t|+3).
\]
In the Hadamard product centered at $1/2$, the linear coefficient is
purely imaginary by conjugate reflection, and the centered zero
reciprocals have zero real part; this proves the equality. The Euler
series at two and the gamma logarithmic derivative prove the bound.
Subtracting logarithmic derivatives at $\sigma+it$ and $2+it$ then gives
$|L'/L(\sigma+it,\psi)|\ll_\eta\log\mathcal C$.
For $x>1$, contour integration of $-(L'/L)(s+w)x^w/w^2$ gives
\begin{align*}
 -\frac{L'}L(s,\psi)={}&
 \sum_{n\le x}\frac{\Lambda(n)\psi(n)}{n^s}
                     \frac{\log(x/n)}{\log x}
 +\frac1{\log x}\left(\frac{L'}L\right)'(s,\psi)\\
 &+\frac1{\log x}\sum_\rho\frac{x^{\rho-s}}{(\rho-s)^2}
 +\frac1{\log x}\sum_{\rho_{\rm triv}}
                       \frac{x^{\rho_{\rm triv}-s}}{(\rho_{\rm triv}-s)^2}.
\end{align*}
The residues at zeros are counted with multiplicity. The functional
equation makes a left edge tend to zero; horizontal edges can be
chosen away from ordinates using the positive zero sum, which gives
$L'/L=O(\log^2(QH))$ there. The quadratic denominators make both zero
sums absolutely convergent.
Integrate this identity in the real coordinate from $\sigma$ to infinity,
and choose $x=\max(3,\log\mathcal C)$. The prime term is $O(\sqrt x)$;
the derivative term is $O_\eta(\log\mathcal C/\log x)$; the nontrivial
zero term is
$O_\eta(x^{-\eta}\log\mathcal C/(\log x)^2)$; the trivial-zero term
is bounded. It follows that
\[
 |\log L(s,\psi)|\ll_\eta\sqrt{\log\mathcal C}
       +\frac{\log\mathcal C}{\log\log(\mathcal C+3)}
       =o_\eta(\log\mathcal C).
\]
Exponentiation proves \eqref{eq:uniform-character-growth}. Passing to
an imprimitive character adds deleted Euler factors and their inverses,
whose product is $O_{\eta,\epsilon}(Q^\epsilon)$: absorb each sufficiently
large prime factor into $p^\epsilon$ and bound the finite remaining
product separately. For the principal character the pole at one adds
a bounded or $O(\sqrt x)$ term in the indicated regions. Fixed-order
derivatives follow from Cauchy's formula on a smaller strip.

For a fixed character these estimates also imply the critical-line
bound $L_\chi^{(j)}(1/2+it)\ll_{\chi,j,\epsilon}(|t|+3)^\epsilon$.
To see this, use \eqref{eq:uniform-character-growth} on
$\Re s=1/2+\eta$, reflect it to $1/2-\eta$ with the functional
equation, and apply the three-lines theorem in that strip. Its
reflection costs at most $T^{\eta+\epsilon}$; choose $\eta$ smaller
than the desired exponent, and apply Cauchy on a fixed smaller
strip. This also supplies the bounds on the horizontal contour edges
used below.

\subsection{Averaging the additive transform}
Write $e(u)=e^{2\pi iu}$. The stationary transform of the functional
equation leads to the following expression, for $T\le U\le2T$:
\begin{equation}\label{eq:additive-transform}
 \frac{\tau(\bar\chi)}q
 \sum_{\substack{k\le y\\(k,q)=1}}
   \frac{a_2(k)\bar\chi(k)}{k^{1-\beta}}
 \sum_{n\le qUk/(2\pi)}\frac{\chi(n)b(n)}{n^\beta}
                  e\left(-\frac n{qk}\right).
\end{equation}
Here $b=c*a_1$, where $c$ is a convolution of at most four shifted
constant or M\"obius sequences, with at most one M\"obius factor.
Fixed-order shift derivatives are allowed.

\begin{lemma}\label{lem:dirichlet-additive}
All nonprincipal character components of
\eqref{eq:additive-transform} have total
$O_{q,\epsilon}(T^{1/2+\epsilon}y)$. Its principal component is
exactly
\begin{equation}\label{eq:principal-transform}
 \frac1{\varphi(q)}
 \sum_{\substack{k\le y\\(k,q)=1}}\frac{a_2(k)}{k^{1-\beta}}
 \sum_{\substack{n\le qUk/(2\pi)\\(n,q)=1}}
     \frac{b(n)}{n^\beta}
       \frac{\mu(k/(n,k))}{\varphi(k/(n,k))}.
\end{equation}
Replacing its inner sums by the residues to the right of
$\Re s=1/2+\eta$, for a fixed sufficiently small $\eta>0$, has
the same total error bound.
\end{lemma}

\begin{proof}
Put $d=(n,k)$, $k'=k/d$, $n=dm$ and $Q=qk'$. On the units modulo
$Q$, expand
\[
 \chi(m)e(-m/Q)=\sum_{\psi\bmod Q}u_\psi\psi(m),
 \qquad \sum_\psi|u_\psi|^2=1.
\]
Chinese remaindering and the primitive Gauss sum give
\[
 u_{\psi_0}=\frac{\mu(k')\chi(k')\tau(\chi)\bar\chi(-1)}
                      {\varphi(q)\varphi(k')}.
\]
Using $\tau(\chi)\tau(\bar\chi)=\chi(-1)q$ and
$\chi(d)\chi(k')=\chi(k)$ cancels every character phase in the
principal scalar, leaving $\mu(k')/[\varphi(q)\varphi(k')]$.
This proves \eqref{eq:principal-transform}.

For the error, factor the relevant Dirichlet series exactly as
\begin{align*}
 \sum_{(m,Q)=1}\frac{b(dm)\psi(m)}{m^s}
 =\sum_{e\mid d}&\left(
 \sum_{\substack{h\le y/e\\(h,Qd/e)=1}}
              \frac{a_1(eh)\psi(h)}{h^s}\right)\left(
 \sum_{(m,Q)=1}\frac{c((d/e)m)\psi(m)}{m^s}\right).
\end{align*}
For a divisor $h_0\mid dm$, this is the bijection
$e=(h_0,d)$, $h_0=eh$, $(h,d/e)=1$, followed by $m=hm'$.
The second factor is a product of shifted $L$-functions and one
possible reciprocal, times local modifications at primes dividing
$d/e$. Before a valuation shift, the local series is
$C_p(z)=A_p(z)/\prod_{j=1}^{d_p}(1-b_jz)$, where $d_p\le4$,
$A_p$ is either $1$ or $1-a z$, and in the latter case $d_p\le3$.
All $a,b_j$ have modulus $p^{O(1/\log T)}$.
The case consisting of a single shifted M\"obius factor is handled
separately: $C_p(z)=1-a z$, its valuation-one shifted series is
$-a$, and every higher shifted series is zero. Its nonzero ratio
to $C_p$ is $-a/(1-a z)$ and satisfies the bound below. An empty
convolution has zero shifted series. In every remaining case
$d_p\ge1$.
For $v\ge1$, multiply the shifted series
$\sum_{n\ge0}c(p^{n+v})z^n$ by the same denominator.
The coefficient recurrence for $C_p$ makes every coefficient of
degree at least $d_p$ vanish. The remaining numerator therefore has
degree at most $d_p-1$; it need not be linear.
Its coefficients are bounded by
$O((v+1)^3p^{O((v+1)/\log T)})$, by expanding the at most four
geometric factors. Dividing by $C_p$ adds only the possible factor
$(1-a z)^{-1}$. Thus shifting by $v$ costs
$O((v+1)^3p^{O(v/\log T)})$; on
$\Re s\ge1/2+\eta$, every added denominator stays away from zero.
Thus \eqref{eq:uniform-character-growth} bounds its multiplier by
$\tau_4(d/e)\{dQ(|\Im s|+3)\}^\epsilon$.

Orthogonality and Cauchy in residue classes bound the finite polynomial
by
\begin{align*}
 \sum_{\psi\bmod Q}\left|
 \sum_{\substack{h\le y/e\\(h,Qd/e)=1}}
       a_1(eh)\psi(h)h^{-s}\right|^2
 &\le\varphi(Q)\left(1+\frac y{eQ}\right)
                      \sum_{h\le y/e}|a_1(eh)|^2h^{-2\Re s}\\
 &\ll_{\eta,r}\tau_r(e)^2(\log(2y))^{2r}(Q+y/e).
\end{align*}
Here $\tau_r(eh)\le\tau_r(e)\tau_r(h)$ and
$\sum\tau_r(h)^2h^{-1-2\eta}<\infty$. Together with
$\sum|u_\psi|^2=1$, this bounds the combined integrand by
\[
 \{dQ(|\Im s|+3)T\}^\epsilon\tau_{r+4}(d)\sqrt{Q+y}.
\]
Use truncated Perron inversion at $x=qUk/(2\pi d)$, replacing the
cutoff temporarily by $\lfloor x\rfloor+1/2$, with height $H=T^{10}$.
The initial line is $1+C/\log T$, beyond all shifts. The elementary
kernel is
\[
 \frac1{2\pi i}\int_{c-iH}^{c+iH}\frac{u^s}{s}ds
 =\mathbf1_{u>1}+O\left(u^c\min\{1,(H|\log u|)^{-1}\}\right).
\]
It follows by closing the infinite vertical contour on the appropriate
side and integrating its tails by parts. Half-integrality and the
divisor coefficient bound make the near-cutoff error
$O(d^\epsilon T^\epsilon x/H)$; the remaining absolute Dirichlet
norm is a fixed power of $\log T$ and is harmless. Apply the kernel
after combining characters: the original multiplier has modulus one,
and its nonprincipal remainder has modulus at most two, so no number
of characters enters the truncation error.

Shift to $1/2+\eta$. Nonprincipal characters contribute no poles;
record the explicit shifted poles for the principal one. The preceding
integrand bound, and $x^s/s$, give
\[
 O_{q,\epsilon}\left(T^\epsilon(Tk/d)^{1/2}
                       \tau_{r+4}(d)\sqrt{qk/d+y}\right).
\]
Choose the preliminary exponents and $\eta$ small enough to absorb
$H^\epsilon$, $x^\eta$ and logarithms. Horizontal edges have the
factor $H^{-1}$ and satisfy the same bound. The residues change by
at most $T^\epsilon d^\epsilon$ times fixed divisor factors when the
cutoff is restored. Multiplying by the outer $1/k$ and summing gives
\[
 T^{1/2+\epsilon}\sqrt y\sum_{k\le y}k^{-1/2}|a_2(k)|
                         \sum_{d\mid k}d^{-1/2}\tau_{r+4}(d)
 \ll T^{1/2+\epsilon}y.
\]
Slightly enlarging the shift circles allows differentiation; endpoint
partial summation allows fixed logarithmic weights.
\end{proof}

\subsection{The three residues and the discrete second moment}
Choose $c=1+C_0/L$, with the fixed $C_0$ larger than all scaled
shift bounds, and a fixed lower ordinate $t_0>2\pi$. Define
\[
 \mathcal K(\alpha,\beta,\gamma;U)
 =\frac1{2\pi i}\int_{c+it_0}^{c+iU}
 \frac{L_\chi(s+\gamma)}{L_\chi(s)}L_\chi(s+\alpha)
 L_{\bar\chi}(1-s+\beta)A_\chi(s)A_{\bar\chi}(1-s)\,ds.
\]
Applying $L_{\bar\chi}(1-s+\beta)=
\Delta_\chi(s-\beta)^{-1}L_\chi(s-\beta)$ and expanding absolutely
on the right line gives \eqref{eq:additive-transform}, with
$c=\mu*n^{-\gamma}*n^{-\alpha}*n^\beta$ in its coefficient notation,
and error $O_{q,\epsilon}(T^{1/2+\epsilon}y)$.
For complex $\beta$, make the change of variable $z=s-\beta$ before
applying the stationary-phase formula. For each pair of coefficient
indices,
\[
 (m/k)^{-s}\Delta_\chi(s-\beta)^{-1}
   =(m/k)^{-\beta}(m/k)^{-z}\Delta_\chi(z)^{-1}.
\]
The new vertical line has real part $c-\Re\beta$, which belongs
to a fixed compact interval about one and remains to the right of
every Dirichlet-series convergence boundary. Thus the amplitude
exponent in the cited formula is the real number
$c-\Re\beta-1/2$. Restoring the endpoint ordinates from
$t_0-\Im\beta,U-\Im\beta$ to $t_0,U$ adds segments of length
$O(L^{-1})$. The absolute Dirichlet norm and Stirling's bound give
total contribution $O(\sqrt{qT}\,yL^{C_1})$, with the bottom
contributing only $O(\sqrt q\,yL^{C_1})$, uniformly for $q\le T$.
Both lie within the stated errors. The factor $(m/k)^{-\beta}$
gives exactly $m^{-\beta}k^\beta$ in
\eqref{eq:additive-transform}.

To verify the error explicitly, the stationary point for inner index
$m$ and outer index $k$ is $t=2\pi m/(qk)$; the gamma quotient
contributes scalar $\tau(\bar\chi)/q$ and phase $e(-m/(qk))$.
The $O(1/t)$ relative Stirling error is $O_q(y\sqrt T\log^C T)$.
Outside a factor-two neighborhood of the stationary point, integration
by parts has the same bound. Inside it, the stationary formula stated
in the proof of Lemma~\ref{lem:mixed-phase} is applied only when
$m/(qk)>1$. The remaining terms are treated on a bounded lower
interval and by the first-derivative test above it, with the same
total error. For $m/(qk)>1$ the stationary formula gives errors bounded,
for each outer coefficient, by
\[
 C_q k^{-1/2}\sum_{m\le C_qkT}\frac{\tau_R(m)}{\sqrt m}
 \ll\sqrt T(kT)^\epsilon
\]
and
\[
 C_q\sqrt T(kT)^\epsilon\sum_{m\asymp qkT}
       (|m-qkU/(2\pi)|+qk\sqrt T)^{-1}
 \ll\sqrt T(kT)^\epsilon\log T.
\]
Connecting the shifted contour endpoints costs the same amount.
Summing the divisor-bounded outer coefficients proves the assertion.

For coprime $h,k$ put $f_{\alpha,\gamma}=\mu*n^{-\alpha}*n^{-\gamma}$
and
\[
 Z^{(q)}_{\alpha,\beta,\gamma,h,k}
 =h^{-\beta}\frac{Z_q(1+\alpha+\beta)Z_q(1+\beta+\gamma)}{Z_q(1+\beta)}
 \prod_{p^v\parallel k}
 \frac{\sum_{m\ge0}f_{\alpha,\gamma}(p^{m+v})p^{-m(1+\beta)}}
      {\sum_{m\ge0}f_{\alpha,\gamma}(p^m)p^{-m(1+\beta)}}.
\]
Then $\mathcal K$ equals
\begin{equation}\label{eq:dirichlet-three-residue}
 \frac1{2\pi}\int_0^U\sum_{\substack{(h,k)=1\\(ghk,q)=1}}
 \frac{a(gh)a(gk)}{ghk}
 \left\{\begin{aligned}
 &Z^{(q)}_{\alpha,\beta,\gamma,h,k}\\
 &+(qt/(2\pi))^{-\alpha-\beta}Z^{(q)}_{-\beta,-\alpha,\gamma,h,k}\\
 &+(qt/(2\pi))^{-\beta-\gamma}Z^{(q)}_{\alpha,-\gamma,-\beta,h,k}
 \end{aligned}\right\}dt
 +O_{q,\epsilon}(T^{1/2+\epsilon}y).
\end{equation}
Here is the residue computation. In the principal transform unfold
$b=c*a$, and express the two coefficient indices as $gh,gk$ with
$(h,k)=1$. The remaining Dirichlet series is
\[
 \sum_{(n,q)=1}c(n)n^{-s-\beta}
                \frac{\mu(k/(n,k))}{\varphi(k/(n,k))}.
\]
Its unrestricted product is
$Z_q(s)Z_q(s+\beta+\gamma)Z_q(s+\beta+\alpha)/Z_q(s+\beta)$.
At $p^v\parallel k$ its weight is zero for exponents below $v-1$,
$-1/(p-1)$ at $v-1$, and one thereafter. There are precisely the
three possible poles $1$, $1-\alpha-\beta$, $1-\beta-\gamma$.
For the first, let $F_p(X)=\sum f_{\alpha,\gamma}(p^j)X^j$,
$C_p(X)=F_p(X)/(1-p^\beta X)=\sum c_jX^j$, and $X=p^{-1-\beta}$.
The recurrence $c_j=f_{\alpha,\gamma}(p^j)+p^\beta c_{j-1}$ gives
\[
 (1-1/p)\left(\sum_{j\ge v}c_jX^j-
                         \frac{c_{v-1}X^{v-1}}{p-1}\right)
             =\sum_{j\ge v}f_{\alpha,\gamma}(p^j)X^j.
\]
This proves the displayed local factor after restoring the cutoff.
The residue $E_q$ cancels $q/\varphi(q)$ exactly.
The substitutions $(\alpha,\beta,\gamma)\mapsto
(-\beta,-\alpha,\gamma)$ and $(\alpha,-\gamma,-\beta)$ move each
other pole to the first and yield the remaining terms by
$x^{1-z}/(1-z)=\int_0^x t^{-z}dt$. Coinciding-shift values are
removable, since the original finite-coefficient integral is holomorphic.

For use in \eqref{eq:square-euler-F}--\eqref{eq:square-euler-G}, we
state the resulting discrete moment explicitly. Put
\begin{align*}
 \mathcal J_q(\alpha,\beta;U)
 =\frac1{2\pi}\int_0^U&\log\frac{qt}{2\pi}
 \sum_{\substack{(h,k)=1\\(ghk,q)=1}}\frac{a(gh)a(gk)}{ghk}\\
 &\cdot\left\{\frac{Z_q(1+\alpha+\beta)}{h^\beta k^\alpha}
 +(qt/(2\pi))^{-\alpha-\beta}
              \frac{Z_q(1-\alpha-\beta)}{h^{-\alpha}k^{-\beta}}\right\}dt.
\end{align*}
Let $\mathcal I_q(\alpha,\beta;U)$ be the $\gamma$ derivative at zero
of the main term in \eqref{eq:dirichlet-three-residue}. Then
\begin{align}\label{eq:dirichlet-discrete}
 \sum_{T<\gamma\le2T}&L_\chi(\rho+\alpha)
 L_{\bar\chi}(1-\rho+\beta)A_\chi(\rho)A_{\bar\chi}(1-\rho)\nonumber\\
 &=\bigl[\mathcal J_q(\alpha,\beta;U)+\mathcal I_q(\alpha,\beta;U)
        +\mathcal I_q(\beta,\alpha;U)\bigr]_{U=T}^{2T}
       +O_{q,\epsilon}(T^{1/2+\epsilon}y).
\end{align}
Indeed, integrate $(L_\chi'/L_\chi)(s)$ times the four displayed
factors over a rectangle. The right edge is
$\partial_\gamma\mathcal K|_{\gamma=0}$.
On the left use $L_\chi'/L_\chi=\Delta_\chi'/\Delta_\chi-
L_{\bar\chi}'/L_{\bar\chi}(1-s)$ and reflect. Conjugation and
reflection transpose the shifts; the principal expression in
\eqref{eq:dirichlet-three-residue} has real coefficients, which yields
$\mathcal I_q(\beta,\alpha)$. The gamma logarithmic derivative is
$-\log(qt/(2\pi))+O_q(1/t)$. Its main integral can be evaluated by
setting $\gamma=0$ in \eqref{eq:dirichlet-three-residue}; the quotient
$L(s+\gamma)/L(s)$ becomes one, and the third term vanishes through
$1/Z_q(1)$. Endpoint partial summation supplies the logarithm and gives
$\mathcal J_q$.

For completeness the horizontal edges satisfy
$L_\chi'/L_\chi=O_q(\log^2T)$ at heights avoiding zeros, and
\[
 |L_\chi(s+\alpha)L_{\bar\chi}(1-s+\beta)|
       \ll_{q,\epsilon}T^{|\Re s-1/2|+\epsilon},\qquad
 |A_\chi(s)A_{\bar\chi}(1-s)|\ll y^{1+\epsilon}
\]
through the strip. Their contribution is
$O_{q,\epsilon}(yT^{1/2+\epsilon})$.
Restoring the endpoints $T,2T$ costs
$O_{q,\epsilon}(yT^\epsilon\log T)$, using the local zero count and
the critical-line bound. All residues count multiplicity. Thus the
formula and its differentiated versions apply even when there are
multiple zeros. Taking $A=B_{\chi,T}^2$ and differentiating both
shifts proves the moment interface used in the kernel calculation.

\subsection{The first moment and the signed mean}
To prove \eqref{eq:square-first}, integrate
$L_\chi'(s)^2A_\chi(s)/L_\chi(s)$ with
$A_\chi=B_{\chi,T}^2$. At a multiple zero the integrand is holomorphic
and the target summand vanishes. This pole calculation for
$L_\chi'^2/L_\chi$ is also the simple-zero transform in
Banks \cite[Section~2, p.~3]{banks-simple-transform}; the polynomial
factor here inserts the specified mollifier. The right-edge series has no constant
term and contributes $O(y^{1+\epsilon}\log^C T)$. Reflection of the
left edge, followed by the stationary transform, gives a negative
principal transform with inner coefficient
\[
 c_k(n)=(\log k)^2-(\log n)^2+(\Lambda*\log)(n).
\]
This follows directly from
\[
 \Delta_\chi(s)^{-1}\{D_\chi(s)^2L_\chi(s)
       -2D_\chi(s)L_\chi'(s)+L_\chi'(s)^2/L_\chi(s)\},
 \qquad D_\chi=\Delta_\chi'/\Delta_\chi,
\]
after replacing $D_\chi$ by $-\log(qt/(2\pi))$ and evaluating the
logarithm at the stationary point. The reflected expression is
conjugated; its principal residues are real, so their values are
unchanged. All remaining components have the error in
Lemma~\ref{lem:dirichlet-additive}.

The principal series is
\[
 W_k(s)=\sum_{(n,q)=1}c_k(n)n^{-s}
              \frac{\mu(k/(n,k))}{\varphi(k/(n,k))}.
\]
For its residues write
\[
 F_k(s)=Z_q(s)A_k(s),\quad
 A_k(s)=\prod_{p^v\parallel k}
             \frac{p^{-(v-1)s}(p^{1-s}-1)}{p-1}.
\]
The coefficient identity
$(\Lambda*\log)(n)-(\log n)^2=-\sum_{p^j\mid n}j(\log p)^2$
is obtained from $Z_q'^2/Z_q-Z_q''=-Z_q(Z_q'/Z_q)'$.
Locally its insertion multiplies the coefficient of $X^j$ by
$j(j+1)/2$, that is, applies $(X^2\partial_X^2+2X\partial_X)/2$.
Let $z=s-1$, $Z_q(s)=E_q/z+C_q+O_q(z)$. Direct expansion gives
\begin{align*}
 W_1(s)&=-E_qz^{-3}+O_q(z^{-2}),\\
 W_p(s)&=\frac{E_q\log p}{p-1}z^{-2}
 +\left\{-E_q\left(\frac{3\log^2p}{2(p-1)}
        +\frac{\log^2p}{(p-1)^2}\right)
        +\frac{C_q\log p}{p-1}\right\}z^{-1}+O_{p,q}(1),\\
 W_{pr}(s)&=-\frac{E_q\log p\log r}{(p-1)(r-1)}z^{-1}+O_{p,r,q}(1)
 \quad(p\ne r).
\end{align*}
These are obtained by applying the displayed differential operator
to the rational local geometric factor. If $\omega(k)\ge3$, its
$\omega(k)$ zeros at one cancel every possible pole. Terms involving
repeated prime powers contribute $O_{q,P}(L)$ after division by
$U/(2\pi)$: their residues are bounded by
\[
 C|a(p^v)|p^{-v}(Lv^2\log p+v^3\log^2p),\quad v\ge2,
\]
and by $C|a(p^vr^w)|p^{-v}r^{-w}v^2w^2\log p\log r$ for two
primes. A repeated-prime sum converges absolutely; the other costs
at most $O(L)$. Replacing $p-1$ by $p$ also costs $O(L)$.
Taking residues of $W_k(s)x^s/s$ with $x=qUk/(2\pi)$ and its
negative outer scalar yields
\[
 \frac U{2\pi}\left\{\tfrac12\log^2U
 +\sum_p\frac{a(p)}p(-\log U\log p+\tfrac12\log^2p)
 +\sum_{p<r}\frac{a(pr)}{pr}\log p\log r+O_{q,P}(L)\right\}.
\]
The conductor cancels because $qE_q/\varphi(q)=1$; deleting its
finitely many primes affects only the error. This proves the first
residue expression used in Proposition~\ref{prop:square-moments}.

We finally justify the signed cancellation. For two polynomials with
real underlying coefficients and lengths at most $y<T^{1/2}$,
\begin{equation}\label{eq:real-bilinear}
 \Im\sum_{T<\gamma\le2T}L_\chi'(\rho)
             X_\chi(\rho)Y_{\bar\chi}(1-\rho)
       =O_{q,\epsilon}(T^{1/2+\epsilon}y).
\end{equation}
Integrate $L_\chi'^2/L_\chi$ times these polynomials on the same
rectangle. On the right, the diagonal relation $nh=k$ has character
factor $\chi(n)\chi(h)\bar\chi(k)=1$ and real coefficient.
The off-diagonal integral is $O(y^{1+\epsilon}\log^C T)$: group
$m=nh$, split $m$ outside or inside $[k/2,2k]$, and in the latter
range use $|\log(m/k)|\gg|m-k|/k$ and its harmonic distance sum.
The reflected left edge is the three-term expression just displayed,
with the outer and inner polynomials exchanged. Its stationary
principal component is real by the exact scalar
\eqref{eq:principal-transform}; at the stationary point each gamma
logarithm is the real number $\log(m/k)$. The nonprincipal components,
horizontal edges and restored endpoints have the error above. This
proves \eqref{eq:real-bilinear}, retaining complex characters throughout.

Choose a fixed $\delta>0$ with $\nu+\delta<1/2$, and apply
\eqref{eq:real-bilinear} to
\[
 X_\chi=B_{\chi,T},\qquad
 Y_{\bar\chi}(s)=B_{\bar\chi,T}(s)
            \sum_{m\le T^\delta}\bar\chi(m)m^{-1-s}.
\]
The exterior tail is at most $2T^{-\delta/2}$. By the second moment,
its total contribution is bounded by
\[
 2T^{-\delta/2}
 \left(N_{\chi,T}\sum_{T<\gamma\le2T}
             |L_\chi'(\rho)B_{\chi,T}(\rho)^2|^2\right)^{1/2}
       =o(N_{\chi,T}L).
\]
If $L_\chi'(\rho)\overline{L_\chi(\rho+1)}=u+iv$, then exactly
$K_{\rho,\chi}=v+u/(2\gamma)$. The imaginary part is controlled
by \eqref{eq:real-bilinear}; the last term has weighted sum
$O(N_{\chi,T}L/T)$ by Cauchy and the Euler product bound.
This proves \eqref{eq:square-cancel}, and completes the proof of
Proposition~\ref{prop:square-moments}.

\section{Gaussian restrictions of the direction moments}
\label{sec:gaussian-moments}
We strengthen the counting argument by retaining zeros at which a short
prime sum is positive. This supplies a lower bound for the derivative
on a set that still has both strict velocity signs. Throughout this
section the fixed mollifier, its constants $a,b$, and the phase
$\psi_\chi$ are those of Section~\ref{sec:signed-density}. Write
$N=N_{\chi,T}$, $L=\log T$, $V=\log L$, and put
\[
 d_T=e^{-(\log V)^2},\quad z_T=d_T^{-1},\quad x_T=T^{d_T},\quad
 S_T=\sum_{\substack{z_T<p\le x_T\\p\nmid q}}\frac1p,
 \quad \sigma_T^2=S_T/2.
\]
The prime number theorem, followed by partial summation, gives
\[
 S_T=V-(\log V)^2-2\log\log V+O_\chi(1)\sim V.
\]
Only this asymptotic is needed. It also follows from the uniform
character-growth estimate \eqref{eq:uniform-character-growth} by Perron
applied to $-\zeta'/\zeta$: shifting to $1/2+\epsilon$ gives
$\sum_{n\le u}\Lambda(n)=u+O_\epsilon(u^{1/2+2\epsilon})$;
partial summation and the convergent higher-prime-power contribution
then give $\sum_{p\le u}p^{-1}=\log\log u+C+o(1)$.
Define
\begin{align*}
 P_T(s)&=\sum_{z_T<p\le x_T}\chi(p)p^{-s},&
 X_\rho&=\Re P_T(\rho)/\sigma_T,\\
 A_\rho&=L_\chi'(\rho)B_{\chi,T}(\rho)^2,&
 C_\rho&=K_\rho|B_{\chi,T}(\rho)|^2.
\end{align*}
Here $K_\rho$ is the real quantity in \eqref{eq:square-cancel}.
Let $m_k$ denote the standard Gaussian moments.

\begin{theorem}\label{thm:gaussian-direction-moments}
For each fixed nonnegative integer $k$, as $T\to\infty$,
\begin{align*}
 N^{-1}\sum X_\rho^k&\longrightarrow m_k,&
 (NL)^{-1}\sum A_\rho X_\rho^k&\longrightarrow a m_k,\\
 (NL^2)^{-1}\sum|A_\rho|^2X_\rho^k&\longrightarrow b m_k,&
 N^{-1}\sum e^{2i\psi_\chi(\gamma)}X_\rho^k&\longrightarrow0,\\
 (NL)^{-1}\sum C_\rho X_\rho^k&\longrightarrow0.
\end{align*}
Every sum counts all zero occurrences in $(T,2T]$. The same five
limits hold with $X_\rho^k$ replaced by $\mathbf1_I(X_\rho)$,
and $m_k$ replaced by the Gaussian mass of any fixed finite interval $I$.
\end{theorem}

\subsection{Uniform rational sampling}
We record the two sampling calculations needed in the proof. They
retain the multiplicities and the arithmetic coefficients.
At endpoints a bounded distance from $T,2T$ and a distance
$\gg_\chi L^{-1}$ from the ordinates, the contour for $L_\chi'/L_\chi$
gives, for rational $x>1$ of polynomial size,
\begin{equation}\label{eq:rational-landau-new}
 \sum_{T<\gamma\le2T}x^\rho
 =-\frac T{2\pi}\chi(x)\Lambda(x)
 +O_\chi\left(xL^2+
       \frac{x\log(2x)}{\operatorname{dist}_*(x,\mathbb Z)}
       +\frac L{\log x}\right).
\end{equation}
The main coefficient is zero unless $x$ is an integer prime power;
at an integer the distance means the distance to the other integers.
To prove the estimate, integrate $(L_\chi'/L_\chi)(s)x^s$ with right
edge $1+1/L$. The term $n=x$ of its absolutely convergent series is
the displayed main term. The closest remaining integer costs
$O(x\log(2x)/\operatorname{dist}_*)$; splitting the other integers
into $n<x/2$, $x/2\le n\le2x$ and $n>2x$, and using
$|\log(n/x)|\gg |n-x|/x$ in the middle range, gives $O(xL^2)$.
On the reflected left edge, integration by parts in the gamma
logarithm gives $O(L/\log x)$; the absolutely convergent reflected
series is bounded by the same estimate. The good horizontal edges
give $O(xL^2)$ by partial fractions and the local zero count.
These estimates prove \eqref{eq:rational-landau-new}. Moving the
endpoints changes $O_\chi(L)$ summands, and will always be estimated
separately below.

We also need the even Hardy phase at rational frequencies. For coprime
$a,b\le T$ and $T^{-1}\le b/a\le T$, its uniform form is
\begin{align}\label{eq:rational-phase-new}
 \sum_{0<\gamma\le T}e^{2i\theta_\chi(\gamma)}(b/a)^{i\gamma}
 ={}&-\frac{qT}{2\pi\varphi(q)}\sqrt{b/a}\,
       \frac{\chi(b)\overline{\chi(a)}\mu(b)}{\varphi(b)}\nonumber\\
 &+O_{\chi,\epsilon}\left(T^{1/2+\epsilon}\sqrt b+
                         (a/b+\sqrt{a/b})L^2\right).
\end{align}
Here any fixed $\epsilon>0$ is allowed. A direct proof keeps the
weights in \eqref{eq:twisted-hardy}. Move its left edge to $-1/2$.
Stirling and the functional equation bound that edge by
$O_\chi(x^{-1}L^2)$ and the top by $O_\chi(\sqrt{xT}L^2)$.
In the stationary-phase calculation isolate the nearest integers to
$qxT/(2\pi)$; the bound
\[
 \sum_{n\asymp B}\frac1{|n-B|+b_0\sqrt T}
 \ll\frac1{b_0\sqrt T}+1+\log(2+T),\qquad b_0=qx,
\]
gives the remaining $x^{-1/2}L^2$ error. Thus the prime sum in
\eqref{eq:twisted-hardy} is uniform in this whole range, with error
$O_\chi((\sqrt{xT}+x^{-1}+x^{-1/2})L^2)$.

For $x=b/a$, let $Q=qb$, $Y=qbT/(2\pi a)$. The exact induced Gauss
sum is
\[
 \sum_{\substack{u\bmod Q\\(u,Q)=1}}\chi(u)e^{-2\pi iau/Q}
       =\mu(b)\chi(b)\overline{\chi(-a)}\tau(\chi).
\]
Indeed insert $\sum_{d\mid(u,b)}\mu(d)$, put $u=dv$ and sum over
the lifts modulo $q$. Their additive sum vanishes unless $b/d$ divides
$a$, so coprimality forces $d=b$. The remaining primitive Gauss sum
is the asserted one, including its zero value when $(ab,q)>1$.
Expand $\chi(u)e^{-2\pi iau/Q}$ in characters modulo $Q$.
Parseval says the squared sum of its coefficients is one, so their
absolute sum is at most $\sqrt{\varphi(Q)}$. For each character the
Perron argument just used, now with
\eqref{eq:uniform-character-growth} and Cauchy's formula for its
logarithmic derivative, gives
\[
 \sum_{n\le Y}\psi(n)\Lambda(n)
   =\mathbf1_{\psi\ {\mathrm{principal}}}Y
          +O_\epsilon(Y^{1/2+\epsilon}Q^\epsilon).
\]
Extra prime powers for an induced character cost only
$O(\omega(Q)\log(2Y))$. Apply Parseval, divide by $\sqrt{qb/a}$,
and decrease the chosen growth exponent if necessary. The resulting
error is the first term in \eqref{eq:rational-phase-new}.
The root number $\tau(\chi)/(i^{e_\chi}\sqrt q)$ cancels the Gauss
phase. The remaining main coefficient is the one displayed there.
When $(b,q)>1$ it vanishes; otherwise
$\varphi(qb)=\varphi(q)\varphi(b)$. This proves all cases.

\subsection{The first moment with two coefficient polynomials}
For the prime insertions we need a bilinear version of the first
moment, which we state explicitly. Let
$X_\chi(s)=\sum x(n)\chi(n)n^{-s}$ and
$Y_{\bar\chi}(s)=\sum y(n)\bar\chi(n)n^{-s}$ have real divisor-bounded
coefficients and length $y\le T^\theta$, with fixed $\theta<1/2$.
The moment $\sum_{0<\gamma\le U}L_\chi(\rho+\alpha)
X_\chi(\rho)Y_{\bar\chi}(1-\rho)$ equals, up to
$O_{q,\epsilon}(T^{1/2+\epsilon}y)$,
\begin{align}\label{eq:bilinear-first-shift}
 &\frac1{2\pi}\int_0^U\log\frac{qt}{2\pi}
       \sum_{(n,q)=1}\frac{(n^{-\alpha}*x)(n)y(n)}n\,dt
 -\frac U{2\pi}\sum_{(n,q)=1}
       \frac{(\Lambda*n^{-\alpha}*x)(n)y(n)}n\nonumber\\
 &\quad+\frac1{2\pi}\int_0^U\left.\partial_\gamma
       \sum_{\substack{(h,k)=1\\(ghk,q)=1}}
       \frac{y(gh)x(gk)}{ghk}
       \left\{Z^q_{\alpha,\gamma,h,k}
       +(qt/(2\pi))^{-\alpha-\gamma}
                    Z^q_{-\gamma,-\alpha,h,k}\right\}
                              \right|_{\gamma=0}dt,
\end{align}
where
\[
 Z^q_{\alpha,\gamma,h,k}
 =h^{-\alpha}k^{-\gamma}
   \frac{Z_q(1+\alpha+\gamma)}{Z_q(1+\alpha)}
        \prod_{p\mid k}\frac{1-p^\gamma}{1-p^{-1-\alpha}}.
\]
Shifts are $O(L^{-1})$ and coinciding-pole values are removable.
For zeta these are exactly the two equivalent published main terms
of Benli--Elma--Ng \cite[Theorem~1.1, equations~(9)--(14),
pp.~3--4]{benli-elma-ng}. The divisor-bounded version used here,
including the coefficients with prime insertions, follows from the
contour derivation below and Lemma~\ref{lem:dirichlet-additive}.

Here is the contour derivation for the fixed character, using the
already proved error machinery. Integrate
$(L_\chi'/L_\chi)(s)L_\chi(s+\alpha)X_\chi(s)Y_{\bar\chi}(1-s)$.
The right-edge diagonal is the negative convolution term above, since
its indices satisfy $nh=k$ and their character factors cancel.
The harmonic off-diagonal estimate in \eqref{eq:real-bilinear} bounds
the remainder. Reflect the left edge. Before the outside minus sign,
its upper-right expression is the conjugate of
\[
 \frac1{2\pi i}\int
 \Delta_\chi(s-\bar\alpha)^{-1}
 \left(\frac{\Delta_\chi'}{\Delta_\chi}(s)
                         -\frac{L_\chi'}{L_\chi}(s)\right)
 L_\chi(s-\bar\alpha)Y_\chi(s)X_{\bar\chi}(1-s)\,ds.
\]
Replace the parentheses temporarily by $L_\chi(s+\gamma)/L_\chi(s)$,
rename the barred shift, and put $v=s-\alpha$. The inner convolution
is $1*n^{-\alpha-\gamma}*(\mu(n)n^{-\alpha})$.
Lemma~\ref{lem:dirichlet-additive} applies with three factors and the
two different coefficient polynomials. Its principal series is
$Z_q(s)Z_q(s+\alpha+\gamma)/Z_q(s+\alpha)$.
At $p^v\parallel k$, the same principal weights as in that lemma are
zero below $v-1$, $-1/(p-1)$ at $v-1$, and one thereafter.
The local numerator has coefficients $f_0=1$ and
$f_j=(1-p^\gamma)p^{-j(\alpha+\gamma)}$ for $j\ge1$.
Summing its tail gives the local ratio
\[
 \frac{(1-p^\gamma)p^{-v(1+\alpha+\gamma)}}{1-p^{-1-\alpha}}.
\]
Restoring the outer cutoff and the $h$ factor gives $Z^q$.
The residue $\prod_{p\mid q}(1-p^{-1})$ cancels $q/\varphi(q)$.
The pole $1-\alpha-\gamma$ gives the other term by the substitution
$(\alpha,\gamma)\mapsto(-\gamma,-\alpha)$.
At $\gamma=0$ the first term has only $k=1$, while the second
vanishes through $1/Z_q(1)$. This gives the diagonal with
$n^{-\alpha}*x$ when the gamma logarithm is restored.
Differentiation in $\gamma$ gives the last term of
\eqref{eq:bilinear-first-shift}, with its positive sign.
The horizontal and Perron errors are those already proved in
Section~\ref{sec:dirichlet-moments}; uniform Cauchy differentiation
adds only fixed logarithmic powers. This proves the interface.

\subsection{Proof of the five limits}
\begin{proof}[Proof of Theorem~\ref{thm:gaussian-direction-moments}]
First expand $P_T(\rho)^j\overline{P_T(\rho)}^{\,l}$ for fixed
$j,l$. Equal products of primes require $j=l$, and unique
factorization gives the diagonal $j!S_T^j+O_j((1+S_T)^{j-1}/z_T)$.
For the off-diagonal main terms in \eqref{eq:rational-landau-new},
an integer prime-power ratio adds copies of just one prime.
Their absolute contribution is at most
$C_{j,l}T(1+V)^{j+l}\sum_{p\le x_T}\log p/p
\ll T d_TL(1+V)^{j+l}$.
After normalization this tends to zero, since $d_TV^r\to0$ for
every fixed $r$. All remaining indices are $T^{o(1)}$ and their
nonzero rational distances and logarithms are $T^{-o(1)}$.
The contour errors, even summed absolutely, are
$T^{o(1)}L^{O_{j,l}(1)}=o(N)$, as are the endpoint strips.
Expanding the real part proves the unweighted Gaussian moments.

For the two mollifier moments, fix $K\ge k$, work through degree $K$
in a formal variable $z$, and put
\[
 E_K(w)=\sum_{j=0}^K w^j/j!,\qquad
 Q_z(s)=B_{\chi,T}(s)^2E_K(zP_T(s)/(2\sigma_T)).
\]
Use $X=Q_z$, $Y=E_K(zP_T/(2\sigma_T))$ in
\eqref{eq:bilinear-first-shift}, and use $Q_z$ on both sides in
\eqref{eq:dirichlet-discrete}. The coefficient of $z^k$ is the
desired insertion divided by $k!$. The largest length is
$T^{2\nu+Kd_T}<T^\theta$ for one fixed $\theta<1/2$, eventually.
The coefficients are uniformly divisor-bounded: a $j$th prime
power has coefficient at most $j!\tau_j(n)$, and
$\sigma_T^{-1}\le1$. Interpolation at $2K+1$ real bounded values of
$z$ supplies each coefficient while preserving all error bounds.

We give the Euler-product and tail checks. At a marked prime introduce
insertion variables $u,v$, and put $x=-1+u$, $y=-1+v$ for the
two squared-mollifier coefficients. The coefficients of $p^{-1}$ in
the four relevant products, with all scaled shifts zero, are
\[
 \begin{array}{c|c}
 \text{first diagonal}&v(1+x)=uv\\
 \text{first two-residue product}&xv+v=uv\\
 \text{second gcd product}&xy+x+y=-1+uv\\
 \text{second full three-residue product}&1+xy+x+y=uv.
 \end{array}
\]
These follow by listing $g=p$, $h=p$, $k=p$ in the exact gcd sums.
In the first two-residue product the $k$-only term vanishes through
$1-p^\gamma$; the final row includes its explicit zeta prefactor.
Thus every ratio to its uninserted product has first coefficient
$uv/p$. Higher coefficients through fixed formal degree are
$O_K(p^{-2})$. The uninserted first two-residue product itself is
\[
 C_1\frac{Z_q(1+\alpha+\gamma)}{Z_q(1+\alpha)}
       \prod_{i=1}^2
       \frac{Z_q(1+u_i)^{1/2}}{Z_q(1+u_i+\gamma)^{1/2}},
\]
where $C_1$ converges absolutely near zero and equals one there.
The two second-moment products are
\eqref{eq:square-euler-F}--\eqref{eq:square-euler-G}.

Use the scaled Mellin weights from the square-moment proof. Their
decay is $O((1+|\Im s|)^{-J-1})$, $J\ge16$.
On $|\Im s_i|\le R=d_T^{-1/2}$, replacing the local shifted powers
by one costs $O(d_TR)$ per marked first coefficient. Their product,
with $u=v=z/(2\sigma_T)$, is therefore, coefficientwise,
\[
 \frac{\text{inserted product}}{\text{uninserted product}}
    =\exp(z^2S_T/(4\sigma_T^2))+o(1)=e^{z^2/2}+o(1).
\]
The error is a fixed power of $1+V$ times
$d_TRV+z_T^{-1}$, which tends to zero. This is uniform on separated
physical shift circles of radius $O(L^{-1})$.
On all Mellin lines the insertion coefficients are $O_K((1+V)^K)$.
The original global bounds for the second-moment products are
$O_q(L^3)$ and $O_q(L^6)$, and the first two-residue product is
$O_q(L^2)$. Removing $|\Im s_i|>\sqrt L$ costs respectively
$O(V^KL^{3-J/2})$, $O(V^KL^{6-J/2})$ and
$O(V^KL^{2-J/2})$, negligible at the required scales.
On the intermediate region $R<\max|\Im s_i|\le\sqrt L$ the Laurent
approximation gives normalized kernels bounded by
$C\prod_i(1+|\Im s_i|)^2$. Its integral is
$O(V^KR^{-14})=o(1)$. The absolutely convergent Euler remainders
tend uniformly to one in the central box.

The first diagonal has the same Gaussian factor and uninserted value
one. In its $\Lambda$ convolution every index supplied by $Y$ is a
product of at most $K$ marked primes, so its coefficient is
$O_K(\log x_T(1+V)^K)=o(L)$; it remains negligible after physical
shift differentiation. The errors $T^{1/2+\epsilon}T^{2\nu+Kd_T}$
are $o(T)$ for a fixed sufficiently small $\epsilon>0$.
All conductor residues cancel as in the square-moment proof.
Cauchy differentiation in the physical shifts and extraction of
$z^k/k!$ prove the limits with weights $A_\rho$ and $|A_\rho|^2$.

For the even phase expand the exterior Euler ratio and truncate its
two indices at $M=T^{1/2}$. Its tail is $O(M^{-1/2})$, while the
fixed prime insertion is $T^{o(1)}$, so this costs $o(N)$.
The arctangent correction also costs $o(N)$.
Apply \eqref{eq:rational-phase-new} to each resulting reduced ratio.
All its numerator and denominator indices are at most $MT^{o(1)}$.
Here is an absolute bound for the main coefficient. At a prime let
$i,j$ be its exponents in the two inserted products and $r,s$ its
exponents in the exterior numerator and denominator. After the
square-root frequency factor and $1/\varphi(b)$ are included, its
absolute weight is at most
\[
                    2p^{-\max(i+r,j+s)-r-s}.
\]
At an unmarked prime the sum is $1+O(p^{-2})$; at a marked prime
it is $O_k(p^{-\max(i,j)})$. Hence the full main coefficient is
$O_k((1+V)^k)$, giving $O_{\chi,k}(T(1+V)^k)=o(N\sigma_T^k)$.
The $\sqrt b$ error has only one harmonic exterior sum, since
$\sqrt b\le\sqrt{m h}$ with $h$ an inserted prime product.
The other exterior sums converge or cost $O(\sqrt M)$.
Thus, for fixed sufficiently small $\epsilon$, all errors are
$T^{1/2+\epsilon+o(1)}=o(N)$. This proves the phase limit.

For the real signed weight expand $X_\rho^k$ into its monomials.
Apply \eqref{eq:real-bilinear} to $BP_T^j$ and $BP_T^{k-j}$, with
the exterior series truncated at $T^{\delta_0}$, where
$\nu+\delta_0<1/2$. The imaginary part is $o(NL)$.
Restoring the exterior tail costs at most
\[
 C T^{-\delta_0/2}
 \left(\sum|A_\rho|^2\sum|X_\rho|^{2k}\right)^{1/2}=o(NL).
\]
The additional real term divided by $2\gamma$ in the exact formula
for $K_\rho$ is $O(NL/T)$ by the same inequality. This proves the
fifth polynomial-moment limit.

Finally consider the positive measures of masses $1/N$ and
$|A_\rho|^2/(NL^2)$ at $X_\rho$. Taylor's formula for their
characteristic functions has limiting remainder at degree $2j$
bounded by $C|t|^{2j}m_{2j}/(2j)!=C|t|^{2j}/(2^jj!)$.
Their characteristic functions therefore tend to those of Gaussian
measures of masses one and $b$. Fourier inversion gives convergence
on smooth compactly supported functions; upper and lower smooth
approximations give convergence on intervals.
For the complex first measure, the Taylor remainder is bounded by
\[
 \frac{|t|^k}{k!}
 \left(\frac1{NL^2}\sum|A_\rho|^2\right)^{1/2}
 \left(\frac1N\sum|X_\rho|^{2k}\right)^{1/2},
\]
whose limiting bound tends to zero. Its variation on any set is
bounded by the geometric mean of the two positive masses there.
The phase measure is dominated by the first positive measure, and
$|C_\rho|\le6|A_\rho|$ dominates the signed measure. The same Taylor,
Fourier and boundary-approximation argument gives all three remaining
interval limits. Gaussian boundary masses vanish, completing the proof.
\end{proof}

\section{Common disks for positive proportions of branches}
\label{sec:common-deformation}
The direction theorem concerns derivatives at the initial zeros. We
now continue fixed positive proportions of the actual branches over a
single parameter disk. A stronger derivative estimate also enlarges
the disk on which these branches remain simple and disjoint.

\begin{theorem}\label{thm:common-deformation}
There are absolute constants $c_0,c_1>0$ with the following property.
For every fixed primitive character $\chi$ and all sufficiently large
$T$ depending on $\chi$, each strict normal-velocity sign has at least
$10^{-24}N_{\chi,T}$ distinct simple centers $\rho$ in $(T,2T)$.
Their branches in the family $F_\chi(u,s)$ of
Section~\ref{sec:signed-density} are holomorphic, simple and pairwise
disjoint throughout
\[
                   |u|\le c_1\frac{\sqrt{\log\log T}}T.
\]
For every positive real $u\le c_0/T$, the same two sets of branches
lie on their respective strict sides of the critical line. All these
branches have ordinates in $(T,2T)$ throughout the complex disk.
\end{theorem}

\subsection{Weighted curvature}
Let $B=B_{\chi,T}$, $D_\rho=L_\chi'(\rho)$, and put
\[
 R=\frac1{1000L},\qquad
 M_\rho=\max_{|z|\le R/2}|L_\chi''(\rho+z)|.
\]
There are absolute constants $C_Q,C_M$ such that, eventually for each
fixed character,
\begin{equation}\label{eq:weighted-curvature-new}
 \sum|B(\rho)|^4|L_\chi(\rho+z)|^2\le C_QN
       \quad(|z|\le R),\qquad
 \sum|B(\rho)|^4M_\rho^2\le C_MNL^4.
\end{equation}
To prove the first assertion use the undifferentiated shifted moment
\eqref{eq:dirichlet-discrete}, with its square-mollifier coefficients.
Take outer physical shift circles of radii $2/(1000L)$ and
$3/(1000L)$, and a smaller auxiliary circle. The separate pole
denominators stay apart on these circles. The estimates in
\eqref{eq:square-euler-F}--\eqref{eq:square-euler-G} give $O(TL)$:
the first coefficient product is $O(L^{-1})$, its explicit pole is
$O(L)$ and the gamma logarithm is $O(L)$; the second product is
$O(1)$ and its auxiliary derivative is $O(L)$.
The Mellin tails were bounded in the proof of
Proposition~\ref{prop:square-moments}. The conductor residues cancel,
so the limiting constants are absolute; the lower-order errors have a
character-dependent starting height. The combined finite shifted sum
is holomorphic. The maximum principle extends the bound to both
interior disks. Set the shifts to $z,\bar z$; reflection on the
critical line makes this sum precisely the positive one in
\eqref{eq:weighted-curvature-new}. Since $N\sim TL/(2\pi)$ this
proves the assertion. Cauchy's formula gives
\[
 M_\rho^2\le256R^{-4}\frac1{2\pi}
                  \int_0^{2\pi}|L_\chi(\rho+Re^{it})|^2dt.
\]
Multiplying by $|B|^4$, summing and using the first assertion proves
the second.

\subsection{A logarithm sampled at separated centers}
Fix $0<\nu_0<1/100$, $x=T^{\nu_0}$ and $\eta>0$. Let
\[
 w_x(n)=\begin{cases}
 1,&n\le x,\\
 \log(x^2/n)/\log x,&x<n\le x^2,\\
 0,&n>x^2,
 \end{cases}
 \qquad
 P^*_{\chi,T}(t)=\sum_{2\le n\le x^2}
 \frac{\Lambda(n)w_x(n)\chi(n)}{n^{1/2+\eta/L+it}\log n}.
\]
The needed logarithmic estimate, with every ambient multiplicity, is
\begin{equation}\label{eq:log-kernel-new}
 \left|\log|L_\chi(1/2+\eta/L+it)|-\Re P^*_{\chi,T}(t)\right|
 \le C_\eta E_T(t)+1,\qquad
 E_T(t)=\sum_{\gamma'}\frac1{1+L^2(t-\gamma')^2}.
\end{equation}
The sum includes all positive, negative and zero ordinates of the
same character. One can take
$C_\eta=\max(4/(\nu_0\eta),3/\nu_0^2)$.

For completeness, integrate
\[
 \frac1{2\pi i\log x}\int_{(2)}
 -\frac{L_\chi'}{L_\chi}(s+w)
                    \frac{x^{2w}-x^w}{w^2}\,dw,
       \qquad \Re s>1/2.
\]
Perron's elementary $w^{-2}$ integral gives the smoothed prime sum.
Shift to the left along lines midway between the trivial zeros. The
functional equation bounds the remaining integral by
$O_\chi(x^{-A}(\log(A+|t|+3)+1)/A)$ on a line at distance $A$
to the left. Good horizontal edges tend to zero by the local zero
count and the squared denominator. The zero residues converge
absolutely, giving
\begin{align*}
 -\frac{L_\chi'}{L_\chi}(s)
 ={}&\sum_n\Lambda(n)w_x(n)\chi(n)n^{-s}
 +\frac1{\log x}\sum_{\omega\ {\mathrm{zero}}}
          \frac{x^{2(\omega-s)}-x^{\omega-s}}{(\omega-s)^2}\\
 &-\frac{\mathbf1_{q=1}}{\log x}
          \frac{x^{2(1-s)}-x^{1-s}}{(1-s)^2}.
\end{align*}
Integrate horizontally from $\Re s=1/2+\eta/L$ to infinity.
The Euler logarithm determines the branch; the zero-location input
makes the ray zero-free. For a nontrivial zero put
$v=L(\sigma-1/2)$ and $y=L(t-\gamma')$. Its integrated absolute
contribution is at most
\[
 \frac1{\nu_0}\int_\eta^\infty
       \frac{e^{-\nu_0v}+e^{-2\nu_0v}}{v^2+y^2}\,dv
                \le\frac{C_\eta}{1+y^2}.
\]
For $|y|\le1$ use the integral of $2v^{-2}$; for $|y|\ge1$
integrate the exponentials divided by $y^2$. The trivial-zero sum
is $O(x^{-1/2}/(T^2(\log x)^2))$ and the zeta pole costs at most
$2x/(T^2(\log x)^2)$. Both together are eventually below one.
Taking real parts proves \eqref{eq:log-kernel-new} and also
justifies the interchanges by nonnegative domination.

We will use two elementary sampling consequences. First,
\begin{equation}\label{eq:prime-comparison-new}
 \sum_{T<\gamma\le2T}
    |P^*_{\chi,T}(\gamma)-P_T(\rho)|^2=o(NV),
\end{equation}
where $P_T$ is the marked prime polynomial of
Section~\ref{sec:gaussian-moments}. To check this, split the difference
into omitted primes, the change $p^{-\eta/L}-1$ at included primes,
and higher prime powers. Their squared coefficient norms with weight
$1/n$ are, respectively,
$O_\chi(1+(\log V)^2)$, $O(d_T^2V)$, and $O(1)$.
For a prime-supported piece, two distinct primes have no integer
prime-power ratio, so its off-diagonal main terms in
\eqref{eq:rational-landau-new} vanish. For the higher prime powers
the coefficients are nonnegative; the integer-ratio main term in the
expanded square is
\[
             -\frac T\pi\frac{a_na_m\Lambda(m/n)}m\le0
             \quad(m>n).
\]
The character factors cancel when $(m,q)=1$, and otherwise the
term is zero. This proves the sign also for complex characters.
The absolute errors, with $Y=x^2$, are bounded by
\[
 L^2\sum_{n<m\le Y}\frac1n\ll YL^3,\quad
 \sum_{n<m\le Y}\frac{\log(2m/n)}
                  {n\operatorname{dist}_*(m/n,\mathbb Z)}
                       \ll Y^2L,\quad
 L\sum_{n<m\le Y}\frac1{m\log(m/n)}\ll YL^2.
\]
Use $\operatorname{dist}_*\ge1/n$ and
$\log(m/n)\ge(m-n)/m$. They are $o(N)$; the endpoint strips cost
$O_\chi(YL)=o(N)$. Apply the square inequality to the three pieces.
Since $(\log V)^2=o(V)$, this proves
\eqref{eq:prime-comparison-new}.

Second, for every set $\mathcal S$ of distinct centers in $(T,2T)$
with mutual separation at least $b/L$,
\begin{equation}\label{eq:reversed-kernel-new}
                     \sum_{\gamma\in\mathcal S}E_T(\gamma)
                          \le C(b)N.
\end{equation}
For each real $u$, separation gives
$\sum_{\gamma\in\mathcal S}(1+L^2(\gamma-u)^2)^{-1}\le C(b)$
by comparison with the convergent lattice sum.
Reverse the sum over ambient zeros in $[T/2,3T]$, whose multiplicity
is $O(N)$. The remaining zeros with $|\gamma'|\le4T$ contribute
$O(\#\mathcal S/(TL))$, since their distance is at least $T/2$.
On the shells $2^jT<|\gamma'|\le2^{j+1}T$, $j\ge2$, the count
is $O(2^jT(L+j))$ and the distance at least $2^{j-1}T$.
Their contribution is bounded by
\[
 \frac{C\#\mathcal S}{TL^2}\sum_{j\ge2}2^{-j}(L+j)
                 \ll\frac{\#\mathcal S}{TL}.
\]
These counts have absolute eventual constants for each fixed
character; negative ordinates are counted using its conjugate.
Tonelli's theorem applies throughout. This proves
\eqref{eq:reversed-kernel-new} without removing any complementary
multiple zero.

\subsection{Localization for every fixed power of the conductor}

The logarithmic comparison itself admits every fixed polynomial
conductor range. We record its quantifiers separately from the
moment estimates and the selection of signed sets. All definitions of
$P_T,\sigma_T,X_\rho,R,M_\rho$ remain unchanged.

\begin{proposition}\label{prop:uniform-separated-log-localization}
Fix $A,h,K>0$, a bounded interval $J$, and $\epsilon>0$.
Uniformly over primitive characters of conductor $q\le T^A$,
including conductor one, and all sets $\mathcal S_T$ of distinct
simple centers in $(T,2T)$ satisfying
\[
 |\gamma-\gamma'|\ge h/L\quad(\rho\ne\rho'),\qquad
 M_\rho\le KL|D_\rho|\quad(\rho\in\mathcal S_T),
\]
one has
\begin{align}
 &\#\left\{\rho\in\mathcal S_T:X_\rho\in J,
 \left|\frac{\log|D_\rho|-\log L}{s_T}-X_\rho\right|>\epsilon
 \right\}\nonumber\\
 &\hspace{25mm}\ll_{A,h,K,J,\epsilon}
 N_{\chi,T}\left(V^{-1/4}+\frac{(\log V)^2}{V}\right),
 \qquad s_T=\sqrt{V/2}.\label{eq:polynomial-conductor-localization}
\end{align}
The starting height is uniform in the character and in the choice of
$\mathcal S_T$. Every ambient zero count includes multiplicity.
\end{proposition}

\begin{proof}
Put $Q=\log(qT)$, so $L\le Q\le(1+A)L$, and set
\[
 W=V^{1/4},\qquad y=T^{1/W},\qquad a=\frac1{\log y}=\frac WL,
 \qquad c=\min(1/4000,1/(4K)),\qquad v=c/L.
\]
Write
\[
 F_u(t)=\sum_{\gamma'}\frac{u}{u^2+(t-\gamma')^2}.
\]
The centered Hadamard identity proved in
Section~\ref{sec:dirichlet-moments}, with the gamma factor removed,
gives uniformly for $0<u\le1$ and $T\le t\le2T$
\begin{equation}\label{eq:localization-real-log-derivative}
 \Re\frac{L_\chi'}{L_\chi}(1/2+u+it)
     =F_u(t)-\frac12\log\frac{qt}{2\pi}+O(T^{-1}).
\end{equation}
For complex characters the exact reflection is
$\Lambda_\chi(s)=\varepsilon_\chi
\overline{\Lambda_\chi(1-\bar s)}$; thus the centered linear
coefficient has zero real part. For conductor one use the entire
completion with its factors $s(s-1)$. This proves the displayed
identity for the same character without pairing its positive and
negative ordinates.

At $u=1$ the Euler logarithmic derivative is absolutely bounded,
so $F_1(t)=\tfrac12\log(q(|t|+2))+O(1)$ for all real $t$.
Consequently each unit interval at height $t$ contains
$O(\log(q(|t|+2)))$ occurrences. These bounds also give
\begin{equation}\label{eq:localization-occurrence-count}
 N:=N_{\chi,T}\asymp TQ,\qquad
 \#\{\gamma':T/2\le\gamma'\le3T\}\ll N.
\end{equation}
Here is a direct lower bound for the first assertion. For
$5T/4\le t\le7T/4$, dyadic summation of the local count gives
\[
 \sum_{\gamma'\notin(T,2T]}
       \frac1{1+(t-\gamma')^2}\ll Q/T.
\]
Indeed a distance shell of size $2^jT$ contains
$O(2^jT(Q+j+1))$ occurrences. Integrating $F_1$ over this middle
interval gives $TQ/4+O(T)$; the outside occurrences contribute
$O(Q)$ and each inside occurrence contributes at most $\pi$.
Thus $\pi N\ge TQ/4-O(T+Q)$. The upper bound and the second
assertion follow by covering with unit intervals. This proves
\eqref{eq:localization-occurrence-count} uniformly in the present
range.

List the ordinates in $\mathcal S_T$ as $t_j$. For an even
nonnegative function decreasing on the positive ray, separation
$\delta=h/L$ gives
\[
 \sum_j f(t_j-b)\le2f(0)+\delta^{-1}\int_{\mathbb R}f(t)\,dt.
\]
Retain the nearest point on each side of $b$, and bound each
remaining value by the integral over the preceding interval of length
$\delta$. Apply this to $f(t)=u/(u^2+t^2)$ and sum over the
occurrences in $[T/2,3T]$. The complementary occurrences have
$\sum(t_j-\gamma')^{-2}\ll Q/T$ by the same dyadic argument.
Since $\#\mathcal S_T\le N$, we obtain
\begin{equation}\label{eq:localization-summed-kernel}
 \sum_jF_u(t_j)\ll
 N\left(u^{-1}+L/h\right)+NuQ/T\qquad(0<u\le1).
\end{equation}

Define the auxiliary polynomial
\[
 \mathcal Q_y(t)=\sum_{2\le n\le y}
 \frac{\Lambda(n)\chi(n)}{n^{1/2+a+it}\log n}
                       \frac{\log(y/n)}{\log y}.
\]
The smoothed explicit formula follows the method of
\cite[Lemma~1 and its integrated form, pp.~985--987]{soundararajan-moments2009}
and \cite[equations~(2.2)--(2.6), pp.~111--112]{sono-dirichlet2013};
we retain an absolute error in both directions.
The $w^{-2}$ contour used in the proof of
\eqref{eq:uniform-character-growth} gives the following integrated
identity, with $s_a=1/2+a+it$:
\begin{align}
 \log|L_\chi(s_a)|
 ={}&\Re\mathcal Q_y(t)
       -a\Re\frac{L_\chi'}{L_\chi}(s_a)+\Re Z_y(t)+R_y(t),
       \label{eq:localization-smoothed-log}\\
 Z_y(t)={}&\frac1{\log y}\sum_{\gamma'}\int_a^\infty
 \frac{y^{-u+i(\gamma'-t)}}{(-u+i(\gamma'-t))^2}\,du.
 \nonumber
\end{align}
To check every term, integrate
$-(L_\chi'/L_\chi)(s+w)y^w/w^2$ on a right vertical line.
The residue at zero is
$-(L_\chi'/L_\chi)(s)\log y-(L_\chi'/L_\chi)'(s)$;
a zero of multiplicity $m$ contributes
$-m y^{\rho-s}/(\rho-s)^2$, and the zeta pole contributes
$y^{1-s}/(1-s)^2$. Shift left as in that earlier contour proof,
divide by $\log y$, and integrate horizontally from $s_a$ to
infinity. The derivative term gives the second term in
\eqref{eq:localization-smoothed-log}. The trivial-zero and pole
integrals give
\[
 |R_y(t)|\ll
 \frac{y^{-1/2-a}+\mathbf1_{q=1}y^{1/2-a}}
                 {T^2(\log y)^2}\ll T^{-1}.
\]
This follows by summing the parity progression of trivial zeros,
using denominators at least $T^2$, and integrating the exponential
numerators. All constants are absolute. The nontrivial-zero integral
satisfies
\[
 |Z_y(t)|\le
 \frac{y^{-a}}{(\log y)^2}
       \sum_{\gamma'}\frac1{a^2+(t-\gamma')^2}
             =e^{-1}aF_a(t).
\]
This absolute majorant justifies the integration and interchange of
the zero sum. The zero-free horizontal ray is supplied by the
admitted zero-location theorem.
Equations~\eqref{eq:localization-real-log-derivative} and
\eqref{eq:localization-summed-kernel} now give
\begin{equation}\label{eq:localization-auxiliary-log-mean}
 \sum_j\left|\log|L_\chi(1/2+a+it_j)|
                      -\Re\mathcal Q_y(t_j)\right|
                      \ll_{A,h}NW.
\end{equation}

The bridge to $v$ uses the exact integral of
\eqref{eq:localization-real-log-derivative}. In particular,
\[
 \int_v^aF_u(t)\,du
 =\frac12\sum_{\gamma'}
       \log\frac{a^2+(t-\gamma')^2}{v^2+(t-\gamma')^2}.
\]
The integrands are nonnegative. Integrating
\eqref{eq:localization-summed-kernel} therefore shows
\begin{align*}
 &\sum_j\left|\log|L_\chi(1/2+a+it_j)|
                  -\log|L_\chi(1/2+v+it_j)|\right|\\
 &\quad\ll N\{\log(a/v)+L(a-v)/h+(a^2-v^2)Q/T+aQ\}
            \ll_{A,h,K}NW.
\end{align*}
Only the displacement $v$ is used for Taylor's formula. It lies
inside $R/2$, and the stated curvature condition gives
\[
 L_\chi(\rho+v)=vD_\rho(1+e_\rho),\qquad
 |e_\rho|\le vM_\rho/(2|D_\rho|)\le Kc/2\le1/8.
\]
Thus, with
$U_\rho=\log|D_\rho|-\log L-\Re\mathcal Q_y(\gamma)$,
\begin{equation}\label{eq:localization-derivative-log-mean}
             \sum_{\mathcal S_T}|U_\rho|\ll_{A,h,K}NW.
\end{equation}
The bounded extra logarithmic factors are
$-\log c-\log|1+e_\rho|$.

We next compare the literal coefficients of $\mathcal Q_y$ and
$P_T$. Eventually $z_T<x_T<y$. For a prime $p\le y$, put
$r_p=\log p/\log y$ and $w_p=e^{-r_p}(1-r_p)$. The difference
coefficient is exactly
\[
 \frac{\chi(p)}{\sqrt p}
       (w_p-\mathbf1_{z_T<p\le x_T}).
\]
Use $0\le w_p\le1$, $|w_p-1|\le2r_p$, and the
reciprocal-prime estimates of Section~\ref{sec:gaussian-moments}.
Partial summation also gives
$\sum_{p\le u}(\log p)^2/p\ll(\log u)^2$.
The full coefficient square sum is consequently at most
\begin{align*}
 &\sum_{p\le z_T}\frac1p
 +\frac4{(\log y)^2}\sum_{p\le x_T}\frac{(\log p)^2}p
 +\sum_{x_T<p\le y}\frac1p
 +\sum_p\sum_{j\ge2}\frac1{j^2p^j}\\
 &\quad\ll 1+\log\log z_T+(d_TW)^2+\log(1/(d_TW))
       \ll(\log V)^2.
\end{align*}
The prime-power coefficient is bounded by $1/(jp^{j/2})$.
These bounds retain $\chi(p^j)=\chi(p)^j$ and use only
$|\chi(n)|\le1$, uniformly in the conductor.

Here a direct separated sampling inequality suffices. For
$H(t)=\sum_{n\le y}b_nn^{-it}$ and $E=\sum|b_n|^2$,
expansion on any interval of length at most $2T$ gives
\[
 \int|H|^2\ll(T+y\log(2y))E,\qquad
 \int|H'|^2\ll(T+y\log(2y))(\log y)^2E.
\]
For the off-diagonal terms use
$2|b_nb_m|\le|b_n|^2+|b_m|^2$ and
$\sum_{m\ne n}|\log(m/n)|^{-1}\ll y\log(2y)$,
which follows from
$|\log(m/n)|\ge|m-n|/\max(m,n)$.
The intervals $I_j=[t_j-\delta/2,t_j+\delta/2]$ have disjoint
interiors. Averaging the fundamental theorem of calculus on each
one gives
\[
 |H(t_j)|^2\le\delta^{-1}\int_{I_j}|H|^2
                             +2\int_{I_j}|HH'|.
\]
Sum and apply Cauchy's inequality. Since $y\log(2y)=o(T)$ and
$\log y\le L$, the result is
$\sum_j|H(t_j)|^2\ll_h TLE\ll_h NE$.
Applying it to the coefficient difference proves
\begin{equation}\label{eq:localization-polynomial-square-mean}
 \sum_{\rho\in\mathcal S_T}
 |\mathcal Q_y(\gamma)-P_T(\rho)|^2\ll_hN(\log V)^2.
\end{equation}

Finally, the variance has its stated normalization uniformly in
this larger conductor range. Set $H_q=\max(3,\log(q+3))$.
The prime divisors larger than $H_q$ contribute at most
$\log q/(H_q\log H_q)$ to $\sum_{p\mid q}p^{-1}$;
the others contribute at most $\log\log H_q+O(1)$.
Thus the omitted primes contribute $O_A(1+\log V)$, and
\[
 2\sigma_T^2
 =V-(\log V)^2-2\log\log V+O_A(1+\log V),\qquad
 \frac{\sigma_T}{s_T}=1+O_A((\log V)^2/V).
\]
Put $H_\rho=\mathcal Q_y(\gamma)-P_T(\rho)$. There is the
exact equality
\[
 \frac{\log|D_\rho|-\log L}{s_T}-X_\rho
 =\frac{U_\rho+\Re H_\rho}{s_T}
              +X_\rho(\sigma_T/s_T-1).
\]
On $X_\rho\in J$ the last term is at most $\epsilon/2$ eventually,
uniformly. Each exception therefore has
$|U_\rho|>\epsilon s_T/4$ or $|H_\rho|>\epsilon s_T/4$.
Markov's inequality, \eqref{eq:localization-derivative-log-mean},
and \eqref{eq:localization-polynomial-square-mean} bound their number by
\[
 \frac4{\epsilon s_T}\sum|U_\rho|
 +\frac{16}{\epsilon^2s_T^2}\sum|H_\rho|^2
 \ll_{A,h,K,\epsilon}
 N\left(V^{-1/4}+\frac{(\log V)^2}{V}\right).
\]
This proves the proposition with all the asserted uniformities.
\end{proof}

\subsection{Selection with a large derivative}
\begin{lemma}\label{lem:large-derivative-selection}
There are absolute $K_0\ge1$, $h>0$, $\eta>0$, and two sets of
simple centers, each containing at least $10^{-24}N$ points, such that
\begin{gather*}
 T+1<\Im\rho<2T-1,\qquad
 M_\rho\le K_0L|D_\rho|,\qquad
 |D_\rho|\ge\frac L{2\eta}e^{\sigma_T/2-1},\\
 \frac{|\Re v_\rho|}{|v_\rho|}\ge10^{-12},\qquad
 v_\rho=\frac{\rho L_\chi(\rho+1)}{D_\rho}.
\end{gather*}
The two sets have opposite strict signs of $\Re v_\rho$, and all
their centers together have mutual separation at least $4h/L$.
\end{lemma}

\begin{proof}
Restrict the five interval limits to $I=[1,2]$. Its Gaussian mass
$p$ exceeds $1/24$, since $p\ge e^{-2}/\sqrt{2\pi}$.
The first two mollifier masses become $apNL+o(NL)$ and
$bpNL^2+o(NL^2)$; the phase mass has
$\sum_I\cos^2\psi=(p/2+o(1))N$.
With $\eta_0=1/50$, the bad phase set has size at most
$(p/(2(1-\eta_0^2))+o(1))N$. The same exact numerical margin
as in the direction theorem gives
\[
 \sum_{I,\,|\cos\psi|\ge\eta_0}|A_\rho|
     \ge\left(p\left(a_*-\sqrt{b_*/(2(1-\eta_0^2))}\right)
                                                   -o(1)\right)NL
     >(p/4000-o(1))NL.
\]
The Euler bounds $1/3<|L_\chi(3/2+i\gamma)|<3$ and the exact
phase identity give total absolute $C_\rho$ mass greater than
$(p/600000-o(1))NL$. Its signed sum is $o(NL)$ by the fifth interval
limit. Each strict sign therefore has mass exceeding $NL/40000000$
eventually. Also $\sum|A_\rho|^2\le NL^2$ and
$|C_\rho|\le6|A_\rho|$ eventually.

Put $\kappa=10^{-9}$ and $\lambda=10^{-12}$.
Discard $|\cos\psi|<\lambda$, which costs at most $6\lambda NL$.
The unit endpoint strips have $O_\chi(L)$ occurrences and therefore
cost $o(NL)$ by Cauchy. Choose
$K_0\ge100\sqrt{C_M}/\kappa$. On the set
$M_\rho>K_0L|D_\rho|$, the absolute signed mass is at most
\[
 \frac6{K_0L}\sum |B(\rho)|^2M_\rho
 \le\frac{6\sqrt{C_M}}{K_0}NL<\kappa NL/8.
\]
Multiple centers carry zero mass and are removed. Choose
$h\le\min(1/16000,1/(100K_0))$.
Taylor's formula at any remaining center shows there is no other
such center within $4h/L$: otherwise
$|D_\rho|\le2K_0h|D_\rho|$, a contradiction. The segment lies in
the defining curvature disk.

Fix $\eta\le\min(1/4000,1/(4K_0))$ and apply
\eqref{eq:log-kernel-new}. By \eqref{eq:prime-comparison-new},
$|P^*_{\chi,T}-P_T|\le\sigma_T/4$ fails at only $o(N)$ occurrences.
By \eqref{eq:reversed-kernel-new},
$C_\eta E_T\le\sigma_T/4$ fails at only $O(N/\sigma_T)=o(N)$
remaining centers. Their total absolute signed mass is $o(NL)$ by
Cauchy. Since $X_\rho\in[1,2]$ at every center under consideration,
\[
          \log|L_\chi(\rho+\eta/L)|\ge\sigma_T/2-1.
\]
Taylor's formula gives
$L_\chi(\rho+\eta/L)=(\eta/L)D_\rho(1+e_\rho)$ with
$|e_\rho|\le K_0\eta/2\le1/8$. Hence the displayed exponential
lower bound for $|D_\rho|$ holds.
Each sign still has absolute mass at least $\kappa NL$.
Cauchy and $\sum C_\rho^2\le36NL^2$ give at least
$(\kappa^2/36)N>10^{-24}N$ distinct centers of each sign.
All other asserted conditions were retained in the deletions.
\end{proof}

\subsection{The two continuation scales}
\begin{proof}[Proof of Theorem~\ref{thm:common-deformation}]
The exact denominator-shift series is
\[
 F_\chi(u,s)=L_\chi(s)+
        \sum_{k\ge1}\frac{(-u)^k(s)_k}{k!}L_\chi(s+k).
\]
It follows first from the binomial series in $\Re s>1$, then by
analytic continuation; its terms with $k\ge1$ are locally normally
convergent near the high centers for $|u|<1$.
Put $G(s)=sL_\chi(s+1)$. On fixed disks of radius less than $1/100$
around those centers, absolute Dirichlet series give
\[
 c_GT\le|G(\rho)|\le C_GT,\qquad |G'(s)|\le C_GT.
\]
All constants are absolute. The lower bound uses the Euler product
at real part $3/2$. For $v=|s|\le3T$,
$|(s)_k|\le(v)_k$ and
$\sum_{k\ge0}(v)_k|u|^k/k!=(1-|u|)^{-v}$.
Thus, when $|u|\le1/10$, the whole perturbation satisfies
\begin{equation}\label{eq:shift-exponential-bound}
 |F_\chi(u,s)-L_\chi(s)|+
       |\partial_sF_\chi(u,s)-L_\chi'(s)|
       \le C T|u|e^{C T|u|}.
\end{equation}
Cauchy's formula on a slightly larger fixed $s$ disk gives the
derivative bound. The terms of order at least two separately give
\begin{equation}\label{eq:shift-quadratic-bound}
 F_\chi(u,s)=L_\chi(s)-uG(s)+R(s,u),\qquad
 |R|+|R_s|\le C_RT^2|u|^2
       \quad(T|u|\le1/10).
\end{equation}

Keep the sets in Lemma~\ref{lem:large-derivative-selection} and the
pairwise disjoint center disks $|s-\rho|<r=h/L$.
On their boundaries,
\[
 |L_\chi(s)-D_\rho(s-\rho)|
       \le\tfrac12M_\rho r^2\le(K_0h/2)|D_\rho|r.
\]
Choose an absolute $d>0$ with $Cd<1/4$, where $C$ is from
\eqref{eq:shift-exponential-bound}. Uniformly for $|u|\le d\sigma_T/T$,
the ratio of the perturbation to $|D_\rho|r$ is at most
\[
 C'\sigma_T\exp\{(Cd-1/2)\sigma_T\}\longrightarrow0.
\]
Rouch\'e's theorem gives exactly one zero in each disk for every
such complex parameter. Moreover throughout each disk,
\[
 |\partial_sF_\chi(u,s)-D_\rho|
 \le K_0h|D_\rho|+CT|u|e^{CT|u|}<|D_\rho|/2
\]
eventually. The zero is simple. Its local implicit sections patch by
uniqueness to a holomorphic section throughout the parameter disk.
Since $\sigma_T\sim\sqrt{\log\log T/2}$, an absolute $c_1>0$
gives the claimed disk. The fixed center disks are disjoint and lie
between the two discarded ordinate strips, proving the simultaneous
distinctness and height assertions on that whole disk.

We establish the uniform normal direction on the stated positive
real interval. Let $s_\rho(u)=\rho+z(u)$. For $T|u|\le c_0$ with
$c_0$ small and absolute, the zero equation and
\eqref{eq:shift-quadratic-bound} first give
$|z(u)|\le C|u|T/|D_\rho|$. Subtract the linear displacement
$uv_\rho=uG(\rho)/D_\rho$ from that equation. For $u\ne0$,
\[
 \left|\frac{z(u)}{uv_\rho}-1\right|
 \le C'\left(
       \frac{M_\rho|u|T}{|D_\rho|^2}
       +\frac{|u|T}{|D_\rho|}+T|u|\right)
 \le C'((K_0+1)c_0+c_0/L),
\]
where we used the weaker consequence $|D_\rho|\ge L$.
Choose the same $c_0$ small enough that the final bound is less than
$10^{-12}/2$. For positive real $u$ the margin
$|\Re v_\rho|\ge10^{-12}|v_\rho|$ then forces $\Re z(u)$ to have
the prescribed strict sign. This establishes the second scale and
completes the proof.
\end{proof}

Every selected zero defines a degree-one analytic cover over the
common complex disk: the derivative in the zero coordinate is
nonzero everywhere and the unique section has just been constructed.
The two normal orientations are the signs of the real differential
of this same section on the specified positive parameter ray.

\section{Logarithmic speeds and continuation near the real axis}
\label{sec:common-regions}
The full logarithmic comparison gives counting information on the
sizes of the velocities. An absolutely convergent difference formula
then enlarges the continuation region in directions close to the real
parameter axis. Put $s_T=\sqrt{(\log\log T)/2}$.

\begin{theorem}\label{thm:gaussian-speed-component}
Fix a primitive character and a finite interval $I$ of positive length.
For either strict sign of $\Re v_\rho$, the lower proportion, relative
to $N_{\chi,T}$, of distinct simple centers for which any one of
\[
 \frac{\log|L_\chi'(\rho)|-\log\log T}{s_T},\qquad
 \frac{\log|v_\rho|-\log(T/\log T)}{s_T},\qquad
 \frac{\log|\Re v_\rho|-\log(T/\log T)}{s_T}
\]
belongs to $I$ is at least $10^{-16}\mathbb P(G\in I)$, where
$G$ is a standard real Gaussian. The constant is absolute; the
starting height may depend on the character and the interval.
\end{theorem}

\begin{proof}
Fix an auxiliary interval $J$ of Gaussian mass $p>0$. The five
interval limits give the same phase-discard margin as before.
Each strict sign within $X_\rho\in J$ has absolute $C_\rho$ mass
at least $pNL/2000000$, eventually.
Delete $|\cos\psi_\chi|<10^{-12}$. By the restricted count and
second moment its mass is at most
\[
 6\cdot10^{-12}
 \left(\#\{X_\rho\in J\}\sum_J|A_\rho|^2\right)^{1/2}
       \le (6\cdot10^{-12}\sqrt b+o(1))pNL.
\]
The unit endpoint strips cost $o(NL)$. Choose a fixed curvature
constant $K_J$ so large that
$6\sqrt{C_M}/K_J<p/16000000$. Formula
\eqref{eq:weighted-curvature-new} bounds the mass lost by
$M_\rho>K_JL|D_\rho|$ by this last quantity times $NL$.
The retained nonzero summands are simple and have separated centers,
as in Lemma~\ref{lem:large-derivative-selection}.

Apply the absolute logarithmic estimate
\eqref{eq:log-kernel-new}, its reversed mean
\eqref{eq:reversed-kernel-new}, and
\eqref{eq:prime-comparison-new} on these centers. For every fixed
$\epsilon>0$, only $o(N)$ retained centers have
\[
 |\log|L_\chi(\rho+\eta/L)|-\Re P_T(\rho)|
                                      >\epsilon\sigma_T.
\]
Indeed Chebyshev bounds the prime comparison exceptions and Markov
bounds the kernel exceptions. Choose $\eta$ fixed and small enough
for the curvature constant $K_J$. Taylor's formula gives
$L_\chi(\rho+\eta/L)=(\eta/L)D_\rho(1+e_\rho)$ with
$|e_\rho|\le1/8$. Its bounded logarithmic factors imply
\begin{equation}\label{eq:derivative-gaussian-transfer}
       Y_\rho:=\frac{\log|D_\rho|-\log L}{s_T}
                   =X_\rho+o(1)
\end{equation}
in counting measure on the retained set, because
$\sigma_T/s_T\to1$ and $J$ is bounded.
Any $o(N)$ exceptions have $o(NL)$ signed absolute mass, by
$|C_\rho|\le6|A_\rho|$ and the global second moment.
Each sign still has mass at least $pNL/4000000$.
The restricted squared mass is at most $72pNL^2$ eventually.
Cauchy's inequality therefore gives at least
\[
              \frac{p}{72\cdot4000000^2}N>10^{-16}pN
\]
centers of either sign after every deletion.
For the derivative assertion take an inner interval $J$ of $I$,
apply \eqref{eq:derivative-gaussian-transfer}, take the lower limit
in $T$ and then exhaust $I$. Gaussian endpoints have zero mass.

Finally the exact velocity formula and the exterior Euler bounds
give, uniformly on the dyadic interval,
\[
             \log|v_\rho|=\log T-\log|D_\rho|+O(1).
\]
The retained phase margin makes
$\log|\Re v_\rho|-\log|v_\rho|$ bounded as well.
The normalized logarithmic speeds are thus $-Y_\rho+o(1)$.
Apply the preceding argument to reflected inner intervals and use
the symmetry of the Gaussian distribution. The strict velocity sign
is the sign of $C_\rho$ at every retained center, proving all claims.
\end{proof}

\begin{theorem}\label{thm:anisotropic-continuation}
There is an absolute $c>0$ such that the same two sets of at least
$10^{-24}N_{\chi,T}$ centers in
Theorem~\ref{thm:common-deformation} have simple disjoint branches on
\[
 \Omega_T=\{u\in\mathbb C:
                 (1+T|u|)e^{12T|\Im u|}\le c e^{s_T}\}.
\]
The sections are holomorphic on the interior and continue locally
across every boundary point. Their ordinates remain in $(T,2T)$.
In particular the real interval $|u|\le c e^{s_T}/(2T)$ is included
eventually. The strict-side conclusion holds for $0<u\le c_0/T$.

More generally, for each fixed $A>0$ there are positive constants
$c_A,c'_A$ such that each strict initial sign has at least
\[
        10^{-17}\mathbb P\{G\in[A+1,A+2]\}N_{\chi,T}
\]
simple centers with these continuation and ordinate properties on
\[
                  (1+T|u|)e^{12T|\Im u|}\le c_A e^{2As_T},
\]
and with strict sides on $0<u\le c'_A/T$.
All constants are independent of the fixed character.
\end{theorem}

\begin{proof}
For $\Re s>0$ and $|u|<1/2$, the identity
\begin{equation}\label{eq:shift-difference-series}
 F_\chi(u,s)-L_\chi(s)
 =\sum_{n\ge1}\chi(n)n^{-s}\{(1+u/n)^{-s}-1\}
\end{equation}
is locally normally convergent. Each summand is
$O(|u|n^{-1-\Re s})$ on compact subsets. It equals the difference
of the original Dirichlet series in $\Re s>1$, and analytic
continuation proves the displayed identity. The logarithm is the
one on $\Re(1+u/n)>0$.

On a center disk $|s-\rho|\le h/L$, $h\le1/10$, put
$U=T|u|$, $v=T|\Im u|$ and $\sigma_s=\Re s$.
For $0\le t\le1$,
\[
 \log|1+tu/n|\ge-2|\Re u|/n,\qquad
 |\arg(1+tu/n)|\le2|\Im u|/n.
\]
Here $2/5\le\sigma_s\le3/5$ and $|s|,|\Im s|\le3T$ eventually.
The trivial difference estimate and the integral identity
\[
 (1+u/n)^{-s}-1=-\frac{su}n
                         \int_0^1(1+tu/n)^{-s-1}dt
\]
give
$|(1+u/n)^{-s}-1|\le3e^2e^{6v}\min(1,U/n)$.
Splitting the sum at $n=U$ yields $O(U)$ for $U\le1$ and
$O(U^{1-\sigma_s})$ for $U\ge1$.
Since $U<T$ and $|\sigma_s-1/2|\le h/L$, the last power differs
from $\sqrt U$ by a factor at most $e^h$.
Consequently the absolute constant $C_0=10^4$ suffices in
\begin{equation}\label{eq:anisotropic-perturbation}
 \sup_{|s-\rho|\le h/L}|F_\chi(u,s)-L_\chi(s)|
                \le C_0 e^{6T|\Im u|}\sqrt{1+T|u|}.
\end{equation}

For the original selected centers,
$\sigma_T-s_T=O_\chi((\log V)^2/\sqrt V)=o(1)$.
Their derivative floor therefore gives
$|D_\rho|h/L\ge d_0e^{s_T/2}$ with an absolute $d_0>0$,
eventually. Choose $c\le(d_0/(10C_0))^2$.
On $\Omega_T$, the perturbation in
\eqref{eq:anisotropic-perturbation} is at most
$|D_\rho|h/(10L)$, whereas the original Taylor remainder is at most
$K_0h|D_\rho|h/(2L)$. Rouch\'e with reference function
$D_\rho(s-\rho)$ gives exactly one zero counted with multiplicity
in each disk. It is simple. The strict boundary inequality persists
in a neighborhood of every parameter point, and uniqueness in each
fixed center disk patches the local sections. The regions are
eventually contained in $|u|<1/2$, so the proof applies throughout.
The disks already provide the distinctness and ordinate conditions.
The stated real interval lies in $\Omega_T$ for large $T$.

For the second assertion use the proof of
Theorem~\ref{thm:gaussian-speed-component} on $J=[A+1,A+2]$.
Remove the $o(N)$ centers where the transfer
\eqref{eq:derivative-gaussian-transfer} has error at least $1/2$.
The remaining signed mass and restricted squared mass still give
more than $10^{-16}\mathbb P(G\in J)N$ centers of each sign.
In particular the asserted smaller fraction has
\[
              |D_\rho|\ge Le^{As_T},\qquad
              M_\rho\le K_AL|D_\rho|.
\]
Choose their separated radius $h_A/L$ with $K_Ah_A\le1/100$.
With $c_A\le(h_A/(10C_0))^2$, the same Rouch\'e argument proves
the second region. The fixed phase margin and the local relative
velocity estimate in the preceding section give $c'_A/T$ after
choosing $c'_A$ small enough. No additional centers are discarded.
\end{proof}

For every fixed $B>0$, these regions include both the complex disk
$|u|\le B\sqrt{\log\log T}/T$ and the real interval
$|u|\le T^{-1}e^{B\sqrt{\log\log T}}$ for some positive fraction
of each initial sign. Indeed choose $A=12B+1$; the logarithms of the
two left sides are at most $12B\sqrt{\log\log T}+O_B(\log\log\log T)$
and $B\sqrt{\log\log T}+O(1)$, respectively, while
$2As_T=\sqrt2 A\sqrt{\log\log T}$.
The construction exhibits the same analytic zero sections on the
larger regions, with their real normal signs controlled on the
explicit interval supplied by the theorem.

\section{Uniformity for a small power of the conductor}
\label{sec:uniform-deformation}
The conductor can vary with the height after the errors in the moment
formulas and the deleted Euler factors are controlled uniformly.
The estimates below provide constants independent of the character
and conductor throughout the stated range.

\begin{theorem}\label{thm:uniform-deformation}
There exist absolute $\kappa,c,c_1,T_0>0$ such that, for $T\ge T_0$
and every primitive character of conductor $q\le T^\kappa$, each
strict initial normal-velocity sign has at least
$10^{-30}N_{\chi,T}$ distinct simple centers in $(T,2T)$.
Their same branches are simple, disjoint and have ordinates in
$(T,2T)$ throughout
\[
                  (1+T|u|)e^{12T|\Im u|}\le c_1e^{s_T},
              \qquad s_T=\sqrt{(\log\log T)/2}.
\]
The branches are holomorphic in the interior and continue locally
across its boundary. This region contains the complex disk
$|u|\le\sqrt{\log\log T}/(100T)$ and the real interval
$|u|\le c_1e^{s_T}/(2T)$ eventually. Their strict normal sides persist
for every positive real $u\le c/T$.

For each fixed finite interval $I$ of positive length, the lower limit
of the infimum over these characters of each signed counting
fraction in Theorem~\ref{thm:gaussian-speed-component} is at least
$10^{-30}\mathbb P(G\in I)$. For every fixed $A>0$ each initial
strict sign also has a possibly different set of at least
\[
               10^{-31}\mathbb P\{G\in[A+1,A+2]\}N_{\chi,T}
\]
centers with the same continuation and ordinate properties on
$(1+T|u|)e^{12T|\Im u|}\le c_Ae^{2As_T}$, and strict sides on
$0<u\le c'_A/T$. The constants and eventual thresholds in these
last assertions may depend on the fixed interval or on $A$, and are
uniform in the displayed conductor range.
\end{theorem}

\subsection{Uniform zero counts and moment errors}
\begin{proof}[Proof of Theorem~\ref{thm:uniform-deformation}]
Retain the single fixed $P,\nu$ used throughout and write
\[
 L=\log T,\quad V=\log L,\quad d=\tfrac12-2\nu>0,
 \qquad \ell=\frac{\log(qT)}L.
\]
All preliminary bounds below hold for $2\le q\le T$ and
$y\le T^\theta$, with fixed $\theta<1/2$.
The uniform Hadamard identity used to prove
\eqref{eq:uniform-character-growth}, evaluated at $2+it$, gives
\[
 \sum_\rho\frac{3/2}{(3/2)^2+(t-\gamma)^2}
            \ll\log(q(|t|+3)).
\]
Every term is positive. The local zero count is therefore
$O(\log(q(|t|+3)))$. Integration over $[T,2T]$ gives
\begin{equation}\label{eq:uniform-count-new}
 N_{\chi,T}=\frac T{2\pi}\{\log(qT)+O(1)\}
                              +O(\log(qT)\log T).
\end{equation}
Indeed the integrated gamma and conductor terms are
$\tfrac12T\log(qT)+O(T)$, and the Euler logarithmic derivative has
bounded integral by termwise integration of its absolutely convergent
series. The Poisson kernel integral differs from
$\pi\mathbf1_{(T,2T)}(\gamma)$ by at most
$C\{(1+|\gamma-T|)^{-1}+(1+|\gamma-2T|)^{-1}\}$ for
$|\gamma|\le4T$. The local count bounds this total error by
$O(\log(qT)\log T)$; dyadic tails give $O(\log(qT))$.
The same argument includes conductor one after adding the two
polynomial factors of its entire completion.

Each unit interval used for a top or bottom contour contains a height
at distance $\gg1/\log(q(|t|+3))$ from every ordinate. Subtracting
the completed partial fractions at $2+it$ from those at $\sigma+it$
gives $L_\chi'/L_\chi=O(\log^2(q(|t|+3)))$ there on every fixed
strip used above. Near terms use the distance bound; far differences
are $O((t-\gamma)^{-2})$ and are summed by the local count.
For $w=C/L\le1/10$, the Euler series at $1+w$ and the functional
equation at $-w$ also give the uniform convexity bound
\[
 |L_\chi(\sigma+it)|\ll L(qT)^{(1+w-\sigma)/2},
       \quad -w\le\sigma\le1+w,\quad T/2\le t\le3T.
\]
One direct maximum-principle proof applies it to
$L_\chi(s)(qT)^{(s-1-w)/2}e^{(s-it)^2}$.
The boundary estimates are absolute, and the Gaussian removes the
horizontal edges. For zeta, first multiply by $(s-1)/(s+1)$.
Slightly wider strips and Cauchy circles of radius $c/L$ give fixed
shift and derivative versions.

The exact main terms in \eqref{eq:dirichlet-three-residue},
\eqref{eq:dirichlet-discrete}, and \eqref{eq:bilinear-first-shift},
and the imaginary-part bound in \eqref{eq:real-bilinear}, consequently
have the following common uniform error:
\begin{equation}\label{eq:uniform-moment-error-new}
                         O_\epsilon(qT^{1/2+\epsilon}y).
\end{equation}
We verify every source of the conductor factor. Opposite $L$ factors
on a high horizontal edge have product $O(L^C\sqrt{qT})$ by the
preceding convexity bound. Their coefficient-polynomial product is
$O_\epsilon(y^{1+\epsilon}L^C)$, and the good-height logarithmic
derivative costs $O(L^2)$. A bounded bottom and its low occurrences
cost $O_\epsilon(qyT^\epsilon)$; the same strip proof applies with
height replaced by three. Endpoint restoration and $O(1/t)$ gamma
errors cost $O_\epsilon(\sqrt{qT}\,yT^\epsilon)$.

In the stationary transform the critical point is
$t=2\pi m/(qk)$ and the exact scalar is $\tau(\bar\chi)/q$.
The pre-transform prefactor is $q^{c-1/2}k^{c-1}m^{-c}$,
where $c=1+O(L^{-1})$. After the first and lower endpoint errors in
the cited stationary-phase lemma, it becomes $k^{-1/2}m^{-1/2}$.
Their sum is
\[
 k^{-1/2}\sum_{m\le CqkT}\frac{\tau_R(m)}{\sqrt m}
                    \ll_\epsilon\sqrt{qT}(qkT)^\epsilon.
\]
The upper endpoint error is bounded by the same prefactor times
$\sum_{m\asymp qkT}(|m-qkU/(2\pi)|+qk\sqrt T)^{-1}\ll L$.
This includes the nearest integers. The nonstationary ranges and
the relative Stirling error give $O(\sqrt{qT}L^C)$ per outer
coefficient by integration by parts and the absolute Dirichlet norm.
Restoring critical points below the bottom cutoff costs
$O_\epsilon(\sqrt q\,yT^\epsilon)$. Summing proves the required
stationary error, with no conductor in the amplitude constant.

Finally write $d_1=(n,k)$, $Q=qk/d_1$ in the auxiliary-character
average. The exact factorization and Parseval estimate already
proved give the uniform Perron integrand bound
\[
 (d_1Q(|\Im s|+3)T)^\epsilon\tau_R(d_1)\sqrt{Q+y}.
\]
Take Perron height $T^{10}$, a half-integer temporary cutoff, and
shift to $1/2+\eta$ with $\eta$ small compared with the final
exponent. The cutoff lies between $T/(2\pi)$ and $CT^{5/2}$.
Truncate after the character combination, whose original multiplier
has modulus one and whose nonprincipal remainder has modulus at
most two. Thus no character-count loss enters the truncation.
The remaining vertical error is at most
\[
 T^\epsilon(qTk/d_1)^{1/2}\tau_R(d_1)\sqrt{qk/d_1+y}.
\]
Its two outer sums are bounded by
\begin{align*}
 q\sqrt T\,T^\epsilon
 \sum_{k\le y}|a_2(k)|\sum_{d_1\mid k}\frac{\tau_R(d_1)}{d_1}
                       &\ll_\epsilon q\sqrt T\,yT^\epsilon,\\
 \sqrt{qT}\sqrt y\,T^\epsilon
 \sum_{k\le y}\frac{|a_2(k)|}{\sqrt k}
       \sum_{d_1\mid k}\frac{\tau_R(d_1)}{\sqrt{d_1}}
                       &\ll_\epsilon\sqrt{qT}\,yT^\epsilon.
\end{align*}
The beneficial outer factor $q^{-1/2}$ has even been discarded.
For the principal residue cutoff, let $\mathcal F(s)$ be the
principal Dirichlet series in the preceding Perron integral.
Enclose all its possible shifted poles by $|s-1|=C_2/L$, choosing
the fixed $C_2$ larger than all scaled shift bounds. On this circle
the finite polynomial has a fixed logarithmic Dirichlet norm,
and the valuation bounds and principal Euler factors give
\[
 |\mathcal F(s)|\ll_\epsilon
       \tau_R(d_1)(d_1Q)^\epsilon L^{C_3}.
\]
The residue sum and its derivative are
\[
 \mathcal R(x)=\frac1{2\pi i}\oint\mathcal F(s)\frac{x^s}{s}\,ds,
 \qquad
 \mathcal R'(x)=\frac1{2\pi i}\oint\mathcal F(s)x^{s-1}\,ds.
\]
Since $x\ll T^{5/2}$, the factor $x^{s-1}$ is bounded on this
circle. After decreasing preliminary exponents, this gives
$\mathcal R'(x)\ll_\epsilon T^\epsilon\tau_R(d_1)$.
Changing $x$ to $\lfloor x\rfloor+1/2$ therefore has this same
bound; the outer divisor sums put it within
\eqref{eq:uniform-moment-error-new}. Fixed shift derivatives add
only logarithmic powers.
Pole cutoff changes and horizontal Perron edges lie within these
bounds. The principal Gauss scalar and all residue identities remain
exactly real, so the signed cancellation has the same error.
Fixed Cauchy derivatives add logarithmic powers only. This proves
\eqref{eq:uniform-moment-error-new}.

\paragraph{The primitive Gauss factor and the error term.}
The same proof gives the sharper common error
\begin{equation}\label{eq:retained-gauss-moment-error}
                         O_\epsilon(\sqrt{qT}\,yT^\epsilon),
                 \qquad 1\le q\le T,
\end{equation}
in \eqref{eq:dirichlet-three-residue},
\eqref{eq:dirichlet-discrete}, \eqref{eq:bilinear-first-shift}
and \eqref{eq:real-bilinear}, with the same fixed divisor envelopes,
length bound $y\le T^\theta$, $\theta<1/2$, and fixed shift circles.
Indeed, $|\tau(\bar\chi)/q|=q^{-1/2}$. Retaining this factor in
the two outer sums just estimated gives
\begin{align*}
 \sqrt{qT}\,T^\epsilon
 \sum_{k\le y}|a_2(k)|\sum_{d_1\mid k}\frac{\tau_R(d_1)}{d_1}
                    &\ll_\epsilon\sqrt{qT}\,yT^\epsilon,\\
 \sqrt T\sqrt y\,T^\epsilon
 \sum_{k\le y}\frac{|a_2(k)|}{\sqrt k}
       \sum_{d_1\mid k}\frac{\tau_R(d_1)}{\sqrt{d_1}}
                    &\ll_\epsilon\sqrt T\,yT^\epsilon.
\end{align*}
The preliminary exponents have again been decreased before absorbing
the fixed divisor and logarithmic factors.

Here are the remaining cutoff checks for this stronger bound.
The Perron horizontal edges have integrand bounded by
$T^\epsilon\tau_R(d_1)\sqrt{Q+y}\,x/H$, whereas the vertical
bound uses $T^\epsilon\tau_R(d_1)\sqrt{Q+y}\sqrt x$.
Since $x\ll T^{5/2}$ and $H=T^{10}$, the horizontal edges lie within
the same two outer bounds. Combining characters before truncation
gives the half-integer Perron error
$O(T^\epsilon d_1^\epsilon x/H)$. Its outer sum, with the Gauss
factor retained, is
$O_\epsilon(\sqrt q\,T^{1+\epsilon}y/H)$.

For the principal residue restoration, the displayed contour for
$\mathcal R'(x)$ gives
$\mathcal R'(x)\ll_\epsilon T^\epsilon\tau_R(d_1)$.
The principal character coefficient has modulus at most one by
Parseval. Consequently changing the half-integer cutoff back to
$x$ costs at most
\[
 q^{-1/2}T^\epsilon
 \sum_{k\le y}\frac{|a_2(k)|}{k}
                  \sum_{d_1\mid k}\tau_R(d_1)
                    \ll_\epsilon q^{-1/2}T^\epsilon,
\]
which is also within \eqref{eq:retained-gauss-moment-error}.
The factors produced by shifts of size $O(L^{-1})$ are bounded here;
fixed shift derivatives contribute only logarithmic powers.

The high horizontal edges, stationary-phase errors, gamma remainder
and high endpoint restorations were already bounded by
$O_\epsilon(\sqrt{qT}\,yT^\epsilon)$ before the weaker common
bound was taken. The bounded bottom and its low occurrences cost
$O_\epsilon(qyT^\epsilon)$, and $q\le\sqrt{qT}$ on $1\le q\le T$.
The restored low stationary points cost
$O_\epsilon(\sqrt q\,yT^\epsilon)$; the right-edge harmonic
off-diagonal costs $O_\epsilon(yT^\epsilon)$. Both satisfy the
stronger bound. This proves \eqref{eq:retained-gauss-moment-error},
including fixed Cauchy derivatives. The later conductor ranges and
choice of $\kappa$ continue to use \eqref{eq:uniform-moment-error-new}.

\subsection{Deleted Euler factors and the exact rescaling}
Let $E_q=\prod_{p\mid q}(1-p^{-1})$. Splitting the prime divisors
at $\log(q+3)$, and using $\sum_{p\mid q}\log p\le\log q$, gives
\[
 \sum_{p\mid q}\frac1p=O(1+\log\log\log(q+10)),\qquad
 E_q^{-1}\ll1+\log\log(q+3).
\]
For $q\le T$, $|w|\le C/\sqrt L$, split the logarithmic derivative
of the deleted factors at $p=L^2$. The small-prime part is $O_C(V)$,
because $p^{|\Re w|}\le e^{2CV/\sqrt L}$; the large-prime part is
at most $CL^{-2+2C/\sqrt L}\sum_{p\mid q}\log p=O_C(1)$.
Integration from zero to $w$ proves
\begin{equation}\label{eq:uniform-euler-new}
                Z_q(1+w)=\frac{E_q}{w}(1+O_C(|w|V))
\end{equation}
uniformly, with the residue interpretation at zero.
The local remainder factors in the Euler products are
$1+O(p^{-3/2})$ in a common neighborhood and equal one at zero.
Omitting any subset of primes therefore preserves their uniform
boundedness and convergence to one by a common summable majorant.
On the full Mellin lines, the restricted Euler series bound both
$Z_q$ and its reciprocal by $\zeta(\Re s)$ when $\Re s>1$.
The global bounds $L^3,L^6,L^2$ consequently acquire at most a
fixed factor $V^B$. The order-sixteen Mellin weights still give
tails $O(V^BL^{-5})$, $O(V^BL^{-2})$, and $O(V^BL^{-6})$.
On the remaining box \eqref{eq:uniform-euler-new} has relative
error $O(V/\sqrt L)=o(1)$. All $E_q$ powers cancel before physical
shift differentiation, exactly as in the fixed-character formulas.

Put $\mathcal L=\log(qT)=\ell L$. The actual same polynomial has
length $T^\nu=e^{(\nu/\ell)\mathcal L}$ and cutoff denominator
$\nu L=(\nu/\ell)\mathcal L$.
Changing every Mellin and physical shift variable from units $L^{-1}$
to units $\mathcal L^{-1}$ therefore gives exactly the earlier
kernels with length parameter $\nu/\ell$.
The gamma logarithm is $\mathcal L+O(1)$ and
$N\sim T\mathcal L/(2\pi)$ uniformly. A physical derivative
introduces $\ell$, and two introduce $\ell^2$ in the normalizations
used here. Thus the two actual constants are
\begin{equation}\label{eq:uniform-moment-constants-new}
                  a_\ell=\ell a(P,\nu/\ell),\qquad
                  b_\ell=\ell^2b(P,\nu/\ell).
\end{equation}
This is the precise rescaling of the finite kernels.
Their integrable majorants also prove continuity in $\ell$ near one.
The certified strict gap at one permits an absolute $\kappa_0>0$
such that, for $1\le\ell\le1+\kappa_0$,
\[
 0<a_\ell<2,\quad0<b_\ell<2,\quad
 a_\ell^2-\frac{b_\ell}{2(1-1/2500)}>\frac1{4000}.
\]
Choose once $0<\kappa<\min(1/10,d/4,\kappa_0)$.
For every fixed insertion degree,
$qT^{1/2+\epsilon}T^{2\nu+o(1)}=o(T)$ uniformly when $\epsilon$
is sufficiently small. The undifferentiated moment and Cauchy proof
of \eqref{eq:weighted-curvature-new} now have absolute constants
and one uniform eventual height on $q\le T^\kappa$.

\subsection{Uniform Gaussian restrictions and continuation}
The five moment limits of Theorem~\ref{thm:gaussian-direction-moments}
are uniform in this conductor range, with $a,b$ replaced by
$a_\ell,b_\ell$. Here are the remaining error checks.
Deleting prime divisors of $q$ changes $S_T$ by $O(1+\log V)$, so
\[
 S_T=V-(\log V)^2-2\log\log V+O(1+\log V),\qquad
                      \sigma_T-s_T=o(1)
\]
uniformly. The rational Landau contour has uniform good-height,
gamma and Euler bounds for $q\le T$, so its error, all diagonal
and repeated-prime counts, and the $T^{o(1)}$ off-diagonal errors
are uniform. The marked local Euler coefficients remain the same
four polynomial identities. Their central-box error is a fixed
power of $V$ times $d_TRV+z_T^{-1}$, with $R=d_T^{-1/2}$.
The intermediate box costs $O(V^{B+K}R^{-14})$ and the central
Laurent relative error is $O(RV/L)$, both $o(1)$. Formula
\eqref{eq:uniform-moment-error-new} handles all finite polynomial
insertions and physical shift derivatives.

For the even Hardy phase, retaining $q$ in the proof of
\eqref{eq:rational-phase-new} gives the same exact main coefficient
and the uniform error
\[
 O_\epsilon\left(T^{1/2+\epsilon}\sqrt{qb}
          +(q^{-1}a/b+q^{-1/2}\sqrt{a/b})L^2\right).
\]
The $q/\varphi(q)$ in its main term is $O(V)$.
After the exterior truncation and fixed prime insertions, the
main term is $O_k(T(1+V)^{k+1})$ and the errors are bounded by
$\sqrt{qT}T^{\epsilon+o(1)}$, $q^{-1}T^{1/4+o(1)}$ and
$q^{-1/2}T^{o(1)}$. They are $o(N)$ uniformly after fixing a small
$\epsilon$. For the signed moment choose
$\delta_0>0$ with $\kappa+\nu+\delta_0<1/2$ and use
\eqref{eq:uniform-moment-error-new} at length
$T^{\nu+\delta_0+o(1)}$; the exterior tail is removed by the
already uniform second and prime moments. Thus all five polynomial
limits are uniform. The interval conclusion follows from the same
Taylor-remainder and Fourier argument. Explicitly, a failure of
uniformity would give a sequence with a convergent subsequence of
$\ell\in[1,1+\kappa]$; the uniform moments and continuity in
\eqref{eq:uniform-moment-constants-new} give the required interval
limit on that subsequence, a contradiction.

The logarithm estimate \eqref{eq:log-kernel-new} is uniform: its
constant has no conductor, trivial zeros depend only on parity, and
its contour now has the uniform good-height bounds.
The coefficient comparison has squared norm
$O(1+(\log V)^2+\log V)=o(V)$ uniformly.
The reversed kernel mean is also uniform, since the ambient count
in $[T/2,3T]$ is $O(T\log(qT))=O(N)$ and the outer shell counts
are $O(2^jT(L+j))$. Thus every $o(N)$ derivative-selection
exception and its $o(NL)$ signed mass is uniform.

We finish with explicit margins. On $I=[1,2]$, of mass $p>1/24$,
the uniform moment gap gives
$\sum_{I,|\cos\psi|\ge1/50}|A_\rho|\ge(p/20000-o(1))NL$.
Hence total absolute $C_\rho$ mass is at least
$(p/3000000-o(1))NL$, and each sign has at least
$pNL/10000000$ eventually. We have
$\sum|A_\rho|^2\le3NL^2$ and $|C_\rho|\le6|A_\rho|$.
Take $\omega=10^{-11}$ and phase margin $10^{-14}$.
Endpoint strips cost $o(NL)$, the phase deletion at most
$6\sqrt3\,10^{-14}NL$, and choose $K_0$ so large that
$6\sqrt{C_M}/K_0<\omega/10$ for the curvature deletion.
The uniform logarithm argument gives the derivative floor
$L e^{\sigma_T/2-1}/(2\eta)$ on all but $o(N)$ of these separated
centers. Each sign still has mass $\omega NL$; its squared mass
is at most $108NL^2$. Thus it has at least
$\omega^2N/108>10^{-30}N$ centers.
All local deformation constants are absolute because their exterior
Dirichlet series have coefficients bounded by one. The proofs of
Theorems~\ref{thm:common-deformation} and
\ref{thm:anisotropic-continuation} therefore give the first
assertion with a uniform height and constants. The logarithm of the
left side on $|u|\le\sqrt V/(100T)$ is at most
$0.12\sqrt V+O(\log V)$, below $s_T+\log c_1$ eventually.

For an arbitrary fixed interval of mass $p>0$, repeat the restricted
selection with its initial sign mass $pNL/10000000$.
Choose its curvature cutoff so the loss is below $pNL/10^9$.
The uniform transfer of logarithmic derivatives leaves sign mass
$pNL/10^8$ with squared mass at most $144pNL^2$.
Cauchy gives at least $pN/(144\cdot10^{16})>10^{-30}pN$ centers.
Inner intervals and the exact reflected speed identities prove the
uniform distribution assertions. Applying this selection to
$[A+1,A+2]$ gives $|D_\rho|\ge Le^{As_T}$, after $o(N)$
deletions, and ample slack above the stated $10^{-31}$ fraction.
The same absolute difference estimate
\eqref{eq:anisotropic-perturbation} and the same relative direction
bound prove the final region and real-ray conclusions. This
completes every asserted uniformity statement.
\end{proof}

\section{The collision of a pair and its complementary divisor}
\label{sec:isolated-pair}

This section proves a local collision formula which keeps every complementary zero and the full real drift. Throughout, $D$ denotes differentiation in the spatial variable, and negative heat time is written $t=-a$, so $f_{-a}=e^{aD^2}f$.

\begin{lemma}[A centered canonical product]\label{lem:centered-product}
Let $G$ be real entire of order at most one, with real zeros and $G(b)\ne0$ for $b\in\mathbb R$. Put
\[
 L=G'(b)/G(b),\quad S=\sum_{G(\xi)=0}(b-\xi)^{-2},\quad
 R(h)=G(b+h)/(G(b)e^{Lh}),
\]
where the zeros are repeated according to multiplicity. Then
\[
 R(h)=\prod_{G(\xi)=0}
       \left(1+\frac h{b-\xi}\right)e^{-h/(b-\xi)},
 \qquad |R(h)|\le e^{S|h|^2/2}.
\]
Moreover, with $T_3=\sum_{G(\xi)=0}(b-\xi)^{-3}$,
\[
 R(h)=1-\tfrac12Sh^2+\tfrac13T_3h^3+E_4(h),\qquad
 |T_3|\le S^{3/2},\quad
 |E_4(h)|\le C S^2|h|^4e^{2S|h|^2},
\]
for an absolute constant $C$. Also
$|R(h)-1|\le8S|h|^2e^{2S|h|^2}$.
\end{lemma}
\begin{proof}
The genus-one Hadamard product has at most a linear exponential factor. Centering it at $b$ and matching the value and logarithmic derivative gives the product written above. For $u\in\mathbb C$, the inequality
$\log(1+v)\le v$ with $v=2\Re u+|u|^2\ge-1$ gives
$|(1+u)e^{-u}|\le e^{|u|^2/2}$; a zero factor satisfies the same inequality directly. Multiplying and taking the locally uniform limit proves the growth bound.

The expansion of $\log R$ at zero gives the quadratic and cubic coefficients. The inequality
$\sum |b-\xi|^{-3}\le(\sum |b-\xi|^{-2})^{3/2}$ proves the stated bound for $T_3$.
For $S>0$ put $v=h\sqrt S$. On $|v|=1$ the normalized function is bounded by $e^{1/2}$. Cauchy's coefficient bound and the geometric series bound the remainder after degree three by $2e^{1/2}|v|^4$ on $|v|\le1/2$. Outside that disk, the global growth estimate and the three explicit coefficients are bounded by a fixed multiple of $|v|^4e^{2|v|^2}$. The same argument after degree one gives the claimed bound with constant $8$. For $S=0$ the product is empty and both remainders vanish.
\end{proof}

\begin{theorem}[An isolated pair]\label{thm:isolated-pair}
Let $f$ be real entire of order at most one with real zeros, and let $b-\delta/2,b+\delta/2$ be adjacent simple zeros, where $\delta>0$. Define
\[
 G(z)=\frac{f(z)}{(z-b)^2-\delta^2/4},\quad
 L=G'(b)/G(b),\quad S=\sum_{G(\xi)=0}(b-\xi)^{-2},\quad
 \kappa=\delta^2S.
\]
There is an absolute $\kappa_0>0$ such that every pair with $\kappa<\kappa_0$ has exactly one collision of its two continued zeros during $0\le a\le\delta^2/4$. This collision is double, and
\[
 a_c=\delta^2\left(\frac18+\frac\kappa{32}+O(\kappa^2)\right),
 \qquad z_c=b-2a_cL+O(\delta\kappa^{3/2}).
\]
The constants in these estimates are absolute. The two zeros are real before $a_c$ and form a nonreal conjugate pair afterwards on the stated interval.
\end{theorem}
\begin{proof}
Set $a=\delta^2\tau$ and $z=b-2aL+\delta w$. The exact operator identity
\[
 e^{aD^2}(e^{Lh}p(h))
   =e^{Lh+aL^2}(e^{aD^2}p)(h+2aL)
\]
follows by expanding $(D+L)^2$, whose three terms commute. It also follows by completing the Gaussian square. With $V$ a standard real Gaussian it yields
\[
 \frac{f_{-\delta^2\tau}(b-2\delta^2\tau L+\delta w)}
 {\delta^2G(b)e^{\delta wL-\delta^2\tau L^2}}
 =J(\tau,w)
 =\mathbb E\bigl[(Y^2-1/4)R(\delta Y)\bigr],
 \quad Y=w+\sqrt{2\tau}V.
\]
All denominators are nonzero. The heat transform of an entire function of order at most one exists for every bounded complex heat parameter by Gaussian domination. The expectation above is even in $\sqrt\tau$ and therefore defines a holomorphic function of $\tau$ through zero.

Lemma~\ref{lem:centered-product} now gives, uniformly on fixed bounded complex $(\tau,w)$ sets,
\begin{align*}
 J(\tau,w)&=w^2+2\tau-\tfrac14-\tfrac\kappa2P_4(\tau,w)
             +\eta_3P_5(\tau,w)+O(\kappa^2),\\
 P_4(\tau,w)&=w^4+(12\tau-\tfrac14)w^2+12\tau^2-\tfrac12\tau,\\
 P_5(\tau,w)&=w^5+(20\tau-\tfrac14)w^3
                       +(60\tau^2-\tfrac32\tau)w,\\
 \eta_3&=\delta^3T_3/3,\qquad |\eta_3|\le\kappa^{3/2}/3.
\end{align*}
For completeness, the integrated remainder is bounded by a constant times
\[
 \kappa^2\mathbb E[(1+|V|)^6e^{C\kappa(1+V^2)}]=O(\kappa^2)
\]
for small absolute $\kappa$. Cauchy's formula gives the same remainder order for every fixed derivative on a smaller compact set. The correction of order $\kappa$ is even in $w$, and that of order $\kappa^{3/2}$ is odd.

For real $0\le\tau\le1/4$, the polynomial $w^2+2\tau-1/4$ has exactly two roots in $|w|<2$, and its modulus on the boundary is at least $15/4$. Rouch\'e's theorem therefore tracks exactly two roots of $J$ in the moving disk. Likewise $J_w=2w+O(\kappa)$ has exactly one zero there. Since $J_{ww}=2+O(\kappa)$, this critical point is a real analytic function $w_*(\tau)$. Evaluating $J_w$ at zero and using the parity just noted gives
$w_*(\tau)=O(\kappa^{3/2})$.

The critical value $D(\tau)=J(\tau,w_*(\tau))$ satisfies
\[
 D(\tau)=2\tau-\tfrac14
       -\tfrac\kappa2(12\tau^2-\tfrac12\tau)+O(\kappa^2),
 \qquad D'(\tau)=2+O(\kappa)>0.
\]
Its signs at $0$ and $1/4$ are opposite, so it has exactly one zero $\tau_c$. First $\tau_c=1/8+O(\kappa)$. At $\tau=1/8$ the expression $12\tau^2-\tau/2$ equals $1/8$, whence
$\tau_c=1/8+\kappa/32+O(\kappa^2)$.
At this point $J_{ww}\ne0$ and $D'\ne0$. The local quadratic factor thus has a simple discriminant zero. Positivity of $J_{ww}$ and $D'$ gives two real roots before the crossing and a conjugate pair after it. No further collision can occur because $J_w$ has only one zero and its critical value is strictly increasing. Returning to the original coordinates proves the result, including the spatial error. For $S=0$ all formulas reduce exactly to the heat transform of a quadratic times a linear exponential.
\end{proof}

The classical Lehmer-pair quantity of Csordas, Smith and Varga is
\[
 g=\sum_{\xi\ne b\pm\delta/2}
 \bigl((b-\delta/2-\xi)^{-2}+(b+\delta/2-\xi)^{-2}\bigr).
\]
Their Definition~1 and Theorem~1 \cite[pp.~109--110]{csv} use
$\delta^2g<4/5$ and give
$((1-5\delta^2g/4)^{4/5}-1)/(8g)\le\Lambda$.
Termwise symmetric expansion, justified by $\delta^2S\to0$, gives
$g=2S+O(\delta^2S^2)$. The first two terms of their bound are consequently
$-\delta^2/8-\delta^4S/32+O(\delta^6S^2)$, matching the negative of the collision time above. The theorem identifies the continued local pair and its unique crossing.

The two-angle example makes the normalization visible. For
$f(z)=2\cos(\delta/2)-2\cos(z-b)$ with $0<\delta<\pi$,
\[
 f_{-a}(z)=2\cos(\delta/2)-2e^{-a}\cos(z-b),\qquad
 a_c=-\log\cos(\delta/2)=\delta^2/8+\delta^4/192+O(\delta^6).
\]
The complementary periodic divisor has
$S=(2\sin^2(\delta/4))^{-1}-8/\delta^2=1/6+O(\delta^2)$, giving the same second coefficient.

\section{A common interval of backward survival}
\label{sec:survival}

We now turn the restricted pair-correlation theorem into a common backward interval for a large family of simple zeros. The local estimate depends on the entire complementary divisor through a positive inverse-square sum.

\begin{lemma}[One zero on a complex time disk]\label{lem:one-zero-survival}
Let $f$ be real entire of order at most one with real zeros, and let $x$ be a simple zero. Define
\[
 \ell_x=\frac{f''(x)}{2f'(x)},\qquad
 S_x=\sum_{f(y)=0,\,y\ne x}(x-y)^{-2}.
\]
For every $V>0$ with $S_x\le V$, the zero has a simple holomorphic continuation throughout $|a|\le2^{-16}/V$ under $e^{aD^2}f$. It obeys
\[
 z_x(-a)=x-2a\ell_x+w_x(a),\qquad
 |w_x(a)|<\frac1{32\sqrt V},
\]
and is real for real $a$ in the interval.
\end{lemma}
\begin{proof}
The centered product of Lemma~\ref{lem:centered-product} gives
$f(x+h)=f'(x)h e^{\ell_xh}R(h)$ and
$|R(h)-1|\le8V|h|^2e^{2V|h|^2}$. Conjugating the heat operator yields
\[
 \frac{f_{-a}(x-2a\ell_x+w)}{f'(x)e^{\ell_xw-a\ell_x^2}}
 =G_a(w)=\mathbb E[(w+\sqrt{2a}Z)R(w+\sqrt{2a}Z)],
\]
where $Z$ is standard Gaussian. As before, Gaussian domination and evenness in the square root give holomorphic dependence on $a$.
Set $\kappa=2^{-5}$, $c_0=2^{-16}$ and $r=\kappa/\sqrt V$. For $|w|=r$, write $v=\sqrt Vw$ and $\sigma=\sqrt{2aV}$. Then
\[
 \sqrt V|G_a(w)-w|
 \le8\mathbb E[|v+\sigma Z|^3e^{2|v+\sigma Z|^2}].
\]
On $|a|\le c_0/V$, use $|v|=\kappa$, $|\sigma|\le\sqrt{2c_0}$,
$(u+v)^3\le4(u^3+v^3)$, and
$2|v+\sigma Z|^2\le4\kappa^2+8c_0Z^2$.
The right side is at most
\[
 32e^{4\kappa^2}\bigl(
 \kappa^3\mathbb Ee^{8c_0Z^2}
 +(2c_0)^{3/2}\mathbb E|Z|^3e^{8c_0Z^2}\bigr).
\]
The elementary Gaussian integrals
$\mathbb Ee^{Z^2/4}=\sqrt2<2$ and
$\mathbb E|Z|^3e^{Z^2/4}=16/\sqrt{2\pi}<8$, together with
$e^{4\kappa^2}<2$, bound this by
$128\kappa^3+512(2c_0)^{3/2}<\kappa/4$.
Rouch\'e's theorem gives exactly one zero in $|w|<r$, counting multiplicity. It is consequently simple. Its implicit sections agree on overlaps by uniqueness and extend through the boundary of the parameter disk because the estimate is strict. Conjugation fixes the unique section on the real parameter interval. This proves all assertions.
\end{proof}

Let $N_H(X)$ count the positive zeros of $H_0$ up to $X$, with multiplicity, and put $M_T=N_H(2T)-N_H(T)$. The zero-location input and the Riemann--von Mangoldt formula give
$N_H(X)\sim X\log X/(4\pi)$.

Put
\[
 c_{\rm MT}=\frac32-\frac1{\sqrt2}\cot\frac1{\sqrt2}
            =0.6725007\ldots.
\]
This is the Montgomery--Taylor simple-zero proportion, whose smooth
test formulation is given in \cite[Lemma~3.2, pp.~13--15]{lamzouri-simple2026}.

\begin{theorem}[Common survival at the Montgomery--Taylor proportion]\label{thm:two-thirds-survival}
For every $0<\eta<c_{\rm MT}$ there exist $c_\eta>0$ and $T_\eta$ such that, for all $T\ge T_\eta$, at least $(c_{\rm MT}-\eta)M_T$ distinct simple zeros of $H_0$ in $(T,2T)$ have branches which remain distinct, real, simple, and in $(T,2T)$ throughout
\[
                  -\frac{c_\eta}{(\log T)^2}\le t\le0.
\]
Consequently, for every positive $a(T)$ with $a(T)(\log T)^2\to0$, the lower limiting proportion of initial centers having those properties on $[-a(T),0]$ is at least $c_{\rm MT}$.
\end{theorem}
\begin{proof}
We first state the precise correlation input. For each nonnegative even Schwartz function $g$ with $g(0)=1$ and Fourier support strictly inside $(-1,1)$, the degree-one case of Rudnick--Sarnak \cite[Theorems~1.1--1.2, pp.~4--5]{rudnick-sarnak}, using their positive-zero normalization and their multiplicity convention on pp.~1--3, gives
\[
 P(Y):=\sum_{\substack{0<x_j,x_k\le Y\\j\ne k}}
       g(F(x_j/2)-F(x_k/2))
       =I_gN_H(Y)+o(N_H(Y)),
\]
where
\[
 F(u)=\frac{u\log u}{2\pi},\qquad
 I_g=\int_{\mathbb R}g(u)
           \left(1-\left(\frac{\sin\pi u}{\pi u}\right)^2\right)du.
\]
The sum is over distinct indices, so repeated ordinates contribute.
There are admissible $g$ with $I_g<1-c_{\rm MT}+\eta/16$.
We verify the precise normalization of the Montgomery--Taylor test
used here. On $[-1/2,1/2]$ put
\[
 f_0(v)=\frac{\cos(\sqrt2v)}{\sqrt2\sin(1/\sqrt2)},
\]
and extend it by zero. It is nonnegative and has integral one.
For $g_0(u)=|\int f_0(v)e^{-2\pi iuv}\,dv|^2$, Plancherel and
Fourier inversion of the squared sinc give
\[
 I_{g_0}=\int f_0(v)^2\,dv-1+
       \iint |v-w|f_0(v)f_0(w)\,dv\,dw.
\]
The function $f_0(v)+\int|v-w|f_0(w)dw$ has second derivative
$f_0''+2f_0=0$ in the interior and is even, so it is constant.
At $v=1/2$ its value is
$1/2+\cot(1/\sqrt2)/\sqrt2=2-c_{\rm MT}$.
Integrating this identity against $f_0$ gives
$I_{g_0}=1-c_{\rm MT}$.
Approximate $f_0$ in $L^1\cap L^2$ by even nonnegative smooth
functions supported strictly inside $(-1/2,1/2)$, each normalized
to have integral one. Their squared Fourier moduli are nonnegative
Schwartz functions, have value one at zero and Fourier support
strictly inside $(-1,1)$. The $L^2$ convergence of the transforms
implies $L^1$ convergence of their squared moduli, hence convergence
of $I_g$, since the sinc weight is bounded. This proves the asserted
admissible approximation and the exact energy conversion.
Fix one sufficiently good smooth approximation. Choose $0<\delta<1$ with $g(u)\ge1-\eta/16$ for $|u|\le\delta$, and put $h=\delta/\log T$.

Call a zero occurrence in $(T,2T)$ isolated when every other occurrence has distance at least $h$. Such an occurrence is simple. Since
\[
 \frac{d}{dx}F(x/2)=\frac{\log(x/2)+1}{4\pi}\le\log T
 \quad(T-1\le x\le2T+1)
\]
for large $T$, every nonisolated occurrence has an ordered witness pair contributing at least $1-\eta/16$. Select one witness per first index; these ordered pairs are distinct. They all occur in
$P(2T+1)-P(T-1)$, whose asymptotic is $I_gM_T+o(M_T)$.
Thus the proportion of nonisolated occurrences is at most
\[
 \frac{I_g+o(1)}{1-\eta/16}
 <1-c_{\rm MT}+\frac\eta4
\]
eventually. At least $(c_{\rm MT}-\eta/4)M_T$ isolated centers remain.

We next control $S_x$ on most of these centers. For any fixed complementary occurrence $y\in(T/2,3T)$, the isolated centers on either side are mutually $h$-separated and, after omitting coincidence, are all at distance at least $h$ from $y$. Hence
\[
 \sum_{x\text{ isolated},\,x\ne y}(x-y)^{-2}
 \le2h^{-2}\sum_{k\ge1}k^{-2}<4h^{-2}.
\]
Summing this inequality over $y$ preserves every complementary multiplicity. Since $N_H(3T)\le4M_T$ eventually, the middle region contributes at most $16M_Th^{-2}$. The zero-count bound and partial summation show uniformly for $x\in(T,2T)$ that the zeros outside $(T/2,3T)$, including the negative zeros, contribute $O(\log T/T)$ to $S_x$. Therefore
\[
                 \sum_{x\text{ isolated}}S_x\le20M_Th^{-2}
\]
for large $T$. Put $B=80/(\eta\delta^2)$. Markov's inequality discards at most $\eta M_T/4$ centers with $S_x>B\log^2T$.

The bound $N_H(u+1)-N_H(u)=O(\log(u+3))$ follows directly from the positive Hadamard sum
\[
 \sum_\gamma\frac{3/2}{(3/2)^2+(v-\gamma)^2}
 =\Re\frac{\xi'}\xi(2+iv)=O(\log(v+3)).
\]
Here the last estimate uses the logarithmic derivative of the gamma factor and the absolutely convergent Euler series at real part $2$. Each occurrence with $|v-\gamma|\le1$ contributes at least $6/13$. Consequently the endpoint strips of width $\log T$ contain $O(\log^2T)=o(M_T)$ zeros. Remove them as well; at least $(c_{\rm MT}-\eta)M_T$ centers remain for large $T$.

Apply Lemma~\ref{lem:one-zero-survival} with the common value $V=B\log^2T$. To control the drift, the even canonical product gives
\[
 \ell_x=\frac1{2x}+
 \sum_{y>0,\,y\ne x}
       \left(\frac1{x-y}+\frac1{x+y}\right).
\]
The paired terms with $y<x/2$ or $y>2x$ sum absolutely to $O(\log T)$. In the middle region the terms $1/(x+y)$ have the same bound. The terms at distance below $1$ are at most
$(O(\log T)S_x)^{1/2}=O_B((\log T)^{3/2})$ by Cauchy's inequality. Summing the remaining unit-distance shells gives $O(\log^2T)$. Thus $|\ell_x|=O_B(\log^2T)$ uniformly on the selected set. The displacement on $|a|\le2^{-16}/V$ is $O_B(1)$, so the endpoint margin keeps each branch in $(T,2T)$.

Distinct selected branches cannot meet. A meeting at a simple zero would identify their local implicit sections, and the identity theorem on the connected common time disk would identify their different initial centers. Take $c_\eta=2^{-16}/B$.
Finally, for each fixed $\eta$, the condition $a(T)\log^2T\to0$ places $[-a(T),0]$ within the proved interval for all sufficiently large $T$. Taking the lower limit and then $\eta\downarrow0$ gives the final assertion.
\end{proof}

\subsection{Finite tests on the surviving positive-density packet}
\label{sec:fd-density-detection}

The common survival interval has an exact consequence for the finite
tests. Fix \(0<\eta<c_{\rm MT}\). For each sufficiently large \(T\),
choose the \(M\ge(c_{\rm MT}-\eta)M_T\) branches given by
Theorem~\ref{thm:two-thirds-survival}, and write them as
\(x_1(t),\ldots,x_M(t)\). Throughout
\[
 I_{T,\eta}=[-c_\eta/(\log T)^2,0]
\]
they are distinct simple real zeros in \((T,2T)\). Define
\[
 \mathcal F_t(w)=\frac{H_t(i\sqrt w)}{H_t(0)},\qquad
 P_{T,t}(w)=\prod_{j=1}^{M}\left(1+\frac{w}{x_j(t)^2}\right),
 \qquad \mathcal E_{T,t}(w)=\frac{\mathcal F_t(w)}{P_{T,t}(w)}.
\]
The first expression denotes its even-power series, so no square-root
branch is chosen. Every factor in the denominator divides the entire
numerator. Hence \(\mathcal E_{T,t}\) is real entire and has constant
coefficient one.

For any real series \(A(w)=\sum_{n\ge0}a_nw^n\), with \(a_0=1\)
and \(a_n=0\) for \(n<0\), define \(m_*(A)\) to be the least
order of a negative minor of \([a_{j-i}]_{i,j\ge0}\), and define
\(R_*(A)\) to be the least index window \(0,\ldots,R\) containing
such a minor. Define \(j_*(A)\) to be the least \(N\ge1\) for which
\[
 J_N(A;X)=\sum_{n=0}^{N}(N)_n a_nX^n
\]
has a nonreal zero. Each empty minimum is \(+\infty\). These
definitions also apply to the complementary quotient without imposing
an extra sign assumption on its coefficients.

\begin{corollary}[Detection on a positive-density surviving packet]
\label{cor:fd-density-detection}
For every \(t\in I_{T,\eta}\), the coefficient Toeplitz matrix of
\(P_{T,t}\) is totally nonnegative and every Jensen polynomial of
\(P_{T,t}\) has only real zeros. Thus
\[
 m_*(P_{T,t})=R_*(P_{T,t})=j_*(P_{T,t})=+\infty
\]
uniformly on this common interval. The exact factorization above gives
\[
 \begin{split}
 m_*(\mathcal F_t)&\ge m_*(\mathcal E_{T,t}),\\
 R_*(\mathcal F_t)&\ge R_*(\mathcal E_{T,t}),\\
 j_*(\mathcal F_t)&\ge j_*(\mathcal E_{T,t}).
 \end{split}
\]
All three comparisons use the same time and permit infinite values.
\end{corollary}
\begin{proof}
For \(c>0\), the Toeplitz matrix of \(1+cw\) has diagonal entries
one, entries \(c\) immediately above the diagonal, and all other
entries zero. In a minor with increasing row and column lists, a
nonzero determinant term cannot contain an inversion. Indeed, for
rows \(i_p<i_q\), nonzero entries require their assigned columns to
satisfy \(j_p\le i_p+1\le i_q\le j_q\). An inverted assignment
would contradict this inequality. Thus the only possible nonzero
term has the positive identity-permutation sign, proving total
nonnegativity.

Toeplitz convolution is exact multiplication:
\[
 T(AB)=T(A)T(B).
\]
For each fixed minor the intermediate indices range over a finite
set, because the matrices are upper triangular. Cauchy--Binet
therefore proves total nonnegativity for the finite product
\(P_{T,t}\). More generally, when every minor of \(T(B)\) of orders
at most \(r\) is nonnegative, the same is true of \(T(P_{T,t}B)\).
This proves the comparison for \(m_*\). For the finite window,
\[
 T_R(P_{T,t}B)=T_R(P_{T,t})T_R(B),
\]
where \(T_R\) uses precisely the indices \(0,\ldots,R\).
Applying the same finite Cauchy--Binet formula proves the comparison
for \(R_*\), with no change in the window.

For the Jensen test, reverse the polynomial and write \(D=d/dx\):
\[
 Q_N(A;x)=x^N J_N(A;1/x)=A(D)x^N.
\]
Only the first \(N+1\) terms of the formal differential operator
act nontrivially. Coefficient convolution now gives the exact map
\[
 Q_N(P_{T,t}B;x)
   =P_{T,t}(D)Q_N(B;x)
   =\prod_{j=1}^{M}\left(1+\frac{D}{x_j(t)^2}\right)Q_N(B;x).
\]
Each operator \(Q\mapsto Q+cQ'\), for \(c>0\), preserves real
roots. For a polynomial with distinct real roots this follows from
\(Q'/Q=\sum_r(x-r)^{-1}\), which decreases strictly from
\(+\infty\) to \(-\infty\) between consecutive roots and gives
one additional solution of \(Q'/Q=-1/c\) to the left of all roots.
Polynomials with repeated real roots follow by perturbing their
roots to distinct real ones and taking the coefficient limit;
real-rootedness is closed under that limit. Constants are fixed.
Starting from \(x^N\) therefore proves the assertion for the packet.
Starting from a real-rooted \(Q_N(B;x)\) proves it for the product.
Polynomial reversal preserves reality of nonzero roots and merely
adds or removes roots at zero; its leading coefficient here is one.
Thus \(J_N(B)\) real-rooted implies \(J_N(P_{T,t}B)\) real-rooted,
which proves the comparison for \(j_*\).
\end{proof}

This corollary transports the positive-density survival count to
infinite local detecting orders on one common interval. Its relation
to the global tests is the displayed product map: the complementary
quotient controls the lower bound for each global detecting index.
The question of divergent scaled global orders still asks for the
corresponding estimates for the full complementary factor as the
height and the backward time vary.

\section{A quantitative height for nonreal heat zeros}
\label{sec:height}
The following bound makes the existence of a nonreal zero at each small
negative time quantitative. Its analytic input is Dobner's uniform
approximation to the deformed Dirichlet series; the recurrence estimate
needed to place a zero at a finite height is proved here.

\begin{theorem}\label{thm:nonreal-height}
There are absolute constants $C,a_0>0$ such that, for $0<a<a_0$,
$H_{-a}$ has a simple zero $z_a$ satisfying
\[
 0<\Re z_a\le \exp\!\bigl(\exp(C/\sqrt a)\bigr),
 \qquad \Im z_a\le-\tfrac12.
\]
The complex conjugate of $z_a$ is another simple zero.
\end{theorem}

\begin{proof}
We use the normalization
\[
 \Xi_{-a}(s)=8H_{-a}(-i(2s-1)),\quad
 \gamma(s)=\tfrac12s(s-1)\pi^{-s/2}\Gamma(s/2),
 \quad \Xi_0(s)=\gamma(s)\zeta(s).
\]
Put
\[
 F_a(s)=\sum_{n\ge1}e^{-a(\log n)^2/4}n^{-s},\qquad
 J_a(s)=s+\frac a4\operatorname{Log}\frac{s}{2\pi},
 \quad
 \gamma_a(s)=\gamma(s)\exp\frac{(s-J_a(s))^2}{a}.
\]
The logarithm is principal in the upper half-plane. Dobner
\cite[Theorem~4, equations~(10)--(13)]{dobner} proves that absolute
constants $C_D,Y_D$ give
\begin{equation}\label{eq:dobner-height}
 \frac{\Xi_{-a}(J_a(s))}{\gamma_a(s)}
 =F_a(s)+O\bigl(C_D(\Im s)^{-1/5}e^{40/a}\bigr)
\end{equation}
uniformly for $0<a\le1$, $1/2\le\Re s\le2$ and $\Im s\ge Y_D$.
In particular, the threshold and error constant are independent of $a$.

The pole of the original zeta function produces the scaling function
\[
 E(w)=\int_0^\infty e^{-v^2/4-wv}\,dv.
\]
This function is entire and satisfies $|E(w)|\le2\sqrt\pi e^{|w|^2}$.
It has a zero: otherwise its entire logarithm $g$ has real part bounded
above by $O(1+|w|^2)$; Borel--Carath\'eodory and Cauchy's estimate imply
that $g$ is a polynomial of degree at most two. But
\[
 xE(x)=\int_0^\infty e^{-u-u^2/(4x^2)}\,du\longrightarrow1
 \quad(x\longrightarrow+\infty),
\]
so $\Re g(x)=-\log x+o(1)$, a contradiction. The integral is positive
on the real axis. Integration by parts gives $E'=2wE-2$, and thus every
zero of $E$ is nonreal and simple.

We shall also need an error estimate in this scaling. For
$s=1+\sqrt a\,w$, comparison of a sum with its integral yields
\begin{equation}\label{eq:heat-pole-profile}
 F_a(1+\sqrt a\,w)=a^{-1/2}E(w)+\mathcal E_a(w),
 \qquad |\mathcal E_a(w)|\le7,
\end{equation}
provided $\Re s\ge1/2$, $|s|\le2$ and $0<a\le1$. Indeed, for
$f_s(u)=e^{-a(\log u)^2/4}u^{-s}$ the error is at most
\[
 |f_s(1)|+\int_1^\infty|f_s'(u)|\,du
 \le 1+2\int_1^\infty u^{-3/2}du
       +\tfrac a2\int_1^\infty(\log u)u^{-3/2}du\le7.
\]
The substitution $v=\sqrt a\log u$ evaluates the integral exactly.

Fix a closed disk $D$ containing just one zero of $E$, and no real
point. Let $\mu=\min_{\partial D}|E|>0$, and choose $K\ge1$ with
$|w|\le K$ on $D$. Define
\[
 W_K=1+\int_0^\infty e^{-v^2/4+Kv}\,dv.
\]
Choose fixed $R>0$ and an integer $Q\ge2$ so that
\[
 \int_R^\infty e^{-v^2/4+Kv}\,dv\le\mu/64,
 \qquad 2\pi/Q\le\mu/(32W_K).
\]
For small $a$, put $N=\lceil e^{R/\sqrt a}\rceil$. The positive
majorant $e^{-a(\log u)^2/4}u^{-1+K\sqrt a}$ decreases for $u\ge1$;
therefore
\begin{align*}
 \sum_{n\ge1}e^{-a(\log n)^2/4}n^{-1+K\sqrt a}&\le W_K/\sqrt a,\\
 \sum_{n>N}e^{-a(\log n)^2/4}n^{-1+K\sqrt a}&\le\mu/(64\sqrt a).
\end{align*}
For any $Y>0$, pigeonhole the $Q^{N-1}+1$ points
\[
 \left(\left\{\frac{jY\log n}{2\pi}\right\}\right)_{2\le n\le N},
 \qquad 0\le j\le Q^{N-1},
\]
into $Q^{N-1}$ cubes of side $1/Q$. Their difference supplies
$\tau=kY$, $1\le k\le Q^{N-1}$, with
$|n^{-i\tau}-1|\le2\pi/Q$ for every $n\le N$. The last two estimates
then give
\begin{equation}\label{eq:heat-finite-recurrence}
 \sup_{w\in D}|F_a(1+\sqrt a\,w+i\tau)-F_a(1+\sqrt a\,w)|
 \le\frac{\mu}{16\sqrt a},\qquad
 Y\le\tau\le YQ^{N-1}.
\end{equation}

Take $Y=e^{250/a}$ and this value of $\tau$, and set
$s(w)=1+\sqrt a\,w+i\tau$. For small enough $a$, throughout $D$,
$3/4\le\Re s(w)\le5/4$ and $\Im s(w)\ge Y/2\ge Y_D$.
The function
\[
 B_a(w)=\Xi_{-a}(J_a(s(w)))/\gamma_a(s(w))
\]
is holomorphic near $D$. Equation~\eqref{eq:dobner-height} has error
at most $2^{1/5}C_De^{-10/a}\le\mu/(16\sqrt a)$. Combined with
\eqref{eq:heat-pole-profile} and \eqref{eq:heat-finite-recurrence}, this
shows on $\partial D$ that
\[
 |B_a-a^{-1/2}E|\le\frac{\mu}{8\sqrt a}+7
 <\frac{\mu}{\sqrt a}=|a^{-1/2}E|.
\]
Rouch\'e's theorem gives exactly one zero $w_a$ in $D$, with
multiplicity one. Write $s_a=s(w_a)$ and
$z_a=-i(2J_a(s_a)-1)$. Since $J_a'(s_a)=1+a/(4s_a)\ne0$,
the zero of $H_{-a}$ at $z_a$ is simple.

For small $a$, $|s_a|>2\pi$; hence
$\Re J_a(s_a)\ge3/4$ and $\Im z_a\le-1/2$.
Also $0<\Re z_a=2\Im J_a(s_a)\le3\tau$. Consequently
\[
 \log\Re z_a\le\log3+250/a+2(\log Q)e^{R/\sqrt a}
 \le e^{C/\sqrt a}
\]
for a fixed $C>R$ and all small enough $a$. Reality of $H_{-a}$ on
the real axis gives the conjugate zero.
\end{proof}

\section{Early collisions in the circular unitary ensemble}
\label{sec:cue-collisions}

We apply the local dynamics to the extreme gaps whose distribution was proved by Ben Arous and Bourgade. Let $U_N$ be Haar distributed in $U(N)$ and write its ordered eigenangles as
$0\le\theta_1<\cdots<\theta_N<2\pi$, with
$\theta_{i+N}=\theta_i+2\pi$. Define
\[
 Z_N(z)=\prod_{i=1}^N2\sin\frac{z-\theta_i}{2},\qquad
 Z_{N,-a}=e^{aD^2}Z_N,\qquad
 \delta_N=\min_i(\theta_{i+1}-\theta_i).
\]
The divisor of $Z_N$ is periodic, including when $Z_N$ itself is antiperiodic. Its first loss of reality occurs at
$A_N=\inf\{a>0:Z_{N,-a}\text{ has a nonreal zero}\}$.

\begin{theorem}[First collision and the early collision process]\label{thm:cue-collisions}
For $N\ge2$, almost surely $0<A_N<\infty$. Moreover,
\[
 \frac{8A_N}{\delta_N^2}=1+O_{\mathbb P}(N^{-1/3}),\qquad
 \lim_{N\to\infty}\mathbb P\left(
   \frac{8N^{8/3}A_N}{(72\pi)^{2/3}}>v\right)=e^{-v^{3/2}}
 \quad(v\ge0).
\]
With probability tending to one the first colliding pair starts at the smallest cyclic gap.
For each real double collision, let $a_{N,j}>0$ be its negative-time magnitude and $b_{N,j}$ its circular location. On bounded rescaled time intervals the point process
\[
 \sum_j\delta_{(8N^{8/3}a_{N,j}/(72\pi)^{2/3},\,b_{N,j})}
\]
converges to the Poisson process with intensity
\[
                  \frac32\sqrt v\,dv\,\frac{db}{2\pi}.
\]
With probability tending to one, every collision on such an interval is double, creates one nonreal conjugate pair, and that pair remains nonreal for the rest of the interval. In particular the number of nonreal pairs at rescaled time $v$ converges to $\operatorname{Poisson}(v^{3/2})$.
\end{theorem}

\begin{proof}
We first control the entire finite divisor. Write
$Z_N(z)=\sum_{k=0}^Nc_ke^{i(k-N/2)z}$. Its endpoint coefficients $c_0,c_N$ are nonzero. Heat multiplication changes $c_k$ to $c_ke^{-a(k-N/2)^2}$. Multiplication by $e^{iNz/2}$ gives a polynomial of degree $N$ in $w=e^{iz}$ with nonzero constant coefficient. On each compact interval of $a$, its roots lie in a compact annulus: the usual polynomial root bound applies to the polynomial and its reciprocal. A zero therefore cannot leave the real axis by escape to imaginary infinity at a finite time. A simple real zero stays real by conjugation and the implicit function theorem. The first loss of reality must occur at a real collision.

Before that time, the real angles $y_i(a)$ satisfy
\[
 y_i'(a)=-\sum_{j\ne i}\cot\frac{y_i(a)-y_j(a)}2.
\]
This follows by differentiating $Z_{N,-a}(y_i(a))=0$ and its sine product. On an empty adjacent arc of length $g=y_{i+1}-y_i$, the mutual terms give $-2\cot(g/2)$ to $g'$. Each other cotangent difference is negative before the overall minus sign, since the arc contains no pole. Thus
\[
          g'\ge-2\cot(g/2),\qquad (g^2)'\ge-8.
\]
Equivalently,
\[
                  (\cos(g/2))'\le\cos(g/2).
\]
Hence a closing gap has closing time at least $-\log\cos(g(0)/2)$.
A nonpositive initial cosine rules out closing by this inequality.
Since $\delta_N<\pi$ almost surely,
\[
                   A_N\ge-\log\cos(\delta_N/2)\ge\delta_N^2/8.
\]
Positivity follows. For even $N$, the constant Fourier coefficient is almost surely nonzero and eventually dominates, leaving no real zeros. For odd $N$, the nonzero frequencies $\pm1/2$ eventually dominate and leave one simple real zero per period. The relevant coefficients are nonzero real analytic functions of the angles: this can be checked when all angles tend to zero, where the product tends to $(2\sin(z/2))^N$. Their zero sets have Lebesgue measure zero, hence zero CUE probability. Uniform $C^1$ convergence after the appropriate rescaling proves the asserted eventual real-zero counts. This proves finiteness of $A_N$.

Put $\ell_N=N^{-7/6}$. With probability tending to one, no circular arc of length $\ell_N$ contains three eigenangles. Indeed, with
$v(x)=(1,e^{ix},\ldots,e^{i(N-1)x})/\sqrt{2\pi}$, the three-point intensity is
$\|v(x_1)\wedge v(x_2)\wedge v(x_3)\|^2$.
Successive divided differences extract $\prod_{i<j}(x_i-x_j)$, and the norms of $v,v',v''$ are respectively $O(N^{1/2}),O(N^{3/2}),O(N^{5/2})$. Hence
\[
 \rho_3(x_1,x_2,x_3)\le CN^9\prod_{i<j}|x_i-x_j|^2.
\]
Integrating over triples in a forward arc of length $\ell_N$ gives expected count at most $CN^9\ell_N^8=O(N^{-1/3})$. Every offending arc has such a triple with its first point as the initial endpoint. Markov's inequality proves the claim.

Ben Arous--Bourgade \cite[Theorem~1.1, p.~2650; Corollary~1.2, p.~2651]{ben-arous-bourgade} prove that the marked gaps $(N^{4/3}\delta,\theta)$ converge to a Poisson process with intensity
$u^2du/(24\pi)\,d\theta/(2\pi)$. In particular
\[
 \frac{N^{4/3}\delta_N}{(72\pi)^{1/3}}\Longrightarrow Y,
 \qquad \mathbb P(Y>y)=e^{-y^3},
\]
and the ratio of the second-smallest to smallest gap converges to a variable strictly larger than one almost surely.
On the no-triples event, once $\delta_N\le\ell_N/2$, every complementary periodic zero of the minimum pair is at distance at least $\ell_N/2$ from its midpoint. At most two zeros lie in each interval of length $\ell_N$, so summing the inverse-square shells gives $S_N\le C\ell_N^{-2}$. Consequently
$\kappa_N=\delta_N^2S_N=O_{\mathbb P}(N^{-1/3})$.
Theorem~\ref{thm:isolated-pair} gives a collision at
$a_{\min}=\delta_N^2(1+\kappa_N/4+O(\kappa_N^2))/8$.
The global lower bound and $A_N\le a_{\min}$ prove the ratio estimate. With probability tending to one, the square of every other initial gap divided by $8$ exceeds $a_{\min}$, by the joint extreme-gap limit. The differential inequality therefore identifies the first pair. Squaring the limit variable $Y$ proves the stated Weibull law.

To pass all early collisions, including those after the first one, fix $A>0$ and set
\[
 s_N=N^{-4/3},\quad T_N=As_N^2,\quad
 B_N=N^{1/24},\quad h_N=B_Ns_N.
\]
On the no-triples event, call every gap $\delta\le h_N$ a pair and every remaining point a singleton. For large $N$, $2h_N<\ell_N$, so the pairs are disjoint. A pair is at distance at least $\ell_N-h_N$ from every other zero. Singletons are mutually at distance at least $h_N$.

We record explicitly the uniform local estimate used for this partition. After removing a singleton at $b$, scale the exact conjugated Gaussian transform by $\sqrt T$, write $\tau=a/T$, and obtain
\[
 I(\tau,u)=\mathbb E[Y R(\sqrt T Y)]=u+O(TS),
 \quad Y=u+\sqrt{2\tau}V.
\]
For a pair use $R_0^2=\delta^2+T$, $d=\delta/R_0\le1$, and $\tau=a/R_0^2\le1$ to obtain
\[
 J(\tau,u)=\mathbb E[(Y^2-d^2/4)R(R_0Y)]
           =u^2+2\tau-d^2/4+O(R_0^2S).
\]
The bounds of Lemma~\ref{lem:centered-product} dominate every remainder and every fixed derivative on bounded complex sets, uniformly in $0<d\le1$. For sufficiently small $TS$ or $R_0^2S$, Rouch\'e on $|u|=2$ gives exactly one or two roots, respectively. The singleton is simple and real. In the pair case $J_u$ has one zero, $J_{uu}=2+O(R_0^2S)>0$ on the real disk, and the critical value has derivative $2+O(R_0^2S)>0$. Its initial value is negative because the initial two roots are $\pm d/2$. Thus the pair has at most one collision and cannot return to the real axis during the interval. When its collision lies in the interval, Theorem~\ref{thm:isolated-pair} gives its time, by uniqueness of continuation from the same initial two roots. This establishes the uniform estimate even when $\delta^2/T$ tends to zero.

For each pair, shell summation and the cotangent product give
\[
 S_{\rm pair}\le C\ell_N^{-2},\qquad
 |L_{\rm pair}|\le C\ell_N^{-1}\log(2/\ell_N).
\]
For each singleton they give
\[
 S_{\rm single}\le C(h_N^{-2}+\ell_N^{-2}),\qquad
 |L_{\rm single}|\le C(h_N^{-1}+\ell_N^{-1}\log(2/\ell_N)).
\]
The inverse-square estimates include all periodic copies. For the first logarithmic derivative use the finite cotangent sum, bounded by the harmonic sum of the circular distance shells; the two pair cotangents cancel at the midpoint. These estimates imply
\[
 (\delta^2+T_N)S_{\rm pair}=O(N^{-1/4}),\qquad
 T_NS_{\rm single}=O_A(B_N^{-2}+N^{-1/3}).
\]
Hence the uniform local estimates apply simultaneously to all clusters.
The moving pair disks have radii $O_A(h_N)=o(\ell_N)$ and center shifts $O_A(T_N\ell_N^{-1}\log(2/\ell_N))=o(\ell_N)$. The singleton disks have radius $O_A(s_N)$ and center shifts
$O_A(T_N/h_N+T_N\ell_N^{-1}\log(2/\ell_N))=o(s_N)$.
They remain disjoint modulo $2\pi$. Their root counts add to $N$, and the polynomial in $w$ has degree $N$ throughout, so these disks exhaust every root at every time in $[0,T_N]$.

For the pairs,
$a_c=\delta^2(1+O(N^{-1/4}))/8$ uniformly. Only gaps $\delta=O_A(s_N)$ can collide in this interval; their sharper interaction bound is $\delta^2S=O_A(N^{-1/3})$. Their collision positions differ from their initial left endpoints by $o(1)$. The early collision process is therefore the image of the extreme-gap process, up to uniformly vanishing displacements, under
\[
 (y,b)\longmapsto(y^2,b),\qquad
 y=N^{4/3}\delta/(72\pi)^{1/3}.
\]
The source intensity becomes $3y^2dy\,db/(2\pi)$ in this coordinate and then $(3/2)\sqrt v\,dv\,db/(2\pi)$ after squaring. Tightness of counts in bounded windows follows from the cited point-process convergence, and the limiting process has no atoms at deterministic boundaries. This justifies the displacement comparison and gives all asserted count limits.
\end{proof}

For zeta zeros, Theorem~\ref{thm:isolated-pair} already applies to any adjacent simple pair with small $\delta^2S$. An arithmetic version of the Poisson conclusion requires a marked extreme-gap limit and quantitative control of the complementary inverse-square sums for those pairs, together with exhaustion of early collisions. The exact restricted Fourier support in the Montgomery and Rudnick--Sarnak theorems controls the pair sums used in Section~\ref{sec:survival}; the empty-interval and extreme-gap quantities in the present proof are additional quantities to estimate. No arithmetic Poisson limit is proved here.

\section{Labels and gaps in the circular unitary model}
\label{sec:cue-labels}
We now compute a marked gap statistic in a finite-dimensional model whose
exterior phase is defined explicitly. The argument proves the joint limit
of that phase and the gap; a marginal equidistribution statement would
not suffice for the sign correlation.

Let $U_N$ be Haar distributed in $U(N)$, with angles
$0\le\theta_1<\cdots<\theta_N<2\pi$. Choose an index $I$ uniformly
and independently, after choosing the matrix. For fixed $r\in[0,1)$
and $a\in\mathbb R$, define
\[
 \phi_{N,r}(x)=\Im\sum_{l=1}^N\log(1-r e^{i(x-\theta_l)}),\qquad
 W_{N,i}=i-\frac{N\theta_i}{2\pi}-a-
                   \frac{\phi_{N,r}(\theta_i)}\pi,
 \quad \eta_{N,i}=\operatorname{sgn}\sin(\pi W_{N,i}).
\]
The logarithm is defined by its power series in $r$. Unwrap indices by
$\theta_{i+N}=\theta_i+2\pi$, and set
$D_{j,N}=N(\theta_{I+j}-\theta_I)/(2\pi)$.
Write $k(u)=\sin(\pi u)/(\pi u)$ and
\[
 K^0(x,y)=k(x-y)-k(x)k(y).
\]
This is the kernel of the sine process conditioned to have a point at
zero, with that point removed; this is its reduced Palm process.

\begin{theorem}\label{thm:cue-labels}
For every fixed positive integer $j$, the pair
$(W_{N,I}\bmod2,D_{j,N})$ converges in distribution to $(V,D_j)$,
where $V$ is uniform modulo two and independent of the $j$th positive
point $D_j$ of the reduced Palm sine process. If
$d_2(v)=\min_{m\in\mathbb Z}|v-2m|$, then
\[
 \lim_{N\to\infty}\mathbb P(\eta_{N,I}=1)=\tfrac12,
 \qquad
 \lim_{N\to\infty}\mathbb E(\eta_{N,I}\eta_{N,I+j})
       =\mathbb E\{1-2d_2(j-D_j)\}.
\]
These limits hold uniformly when $r$ ranges in a fixed compact subset
of $[0,1)$.
\end{theorem}

\begin{proof}
Use circle measure $d\theta/(2\pi)$ and let $P_N$ project onto
$1,e^{i\theta},\ldots,e^{i(N-1)\theta}$. Expanding both determinants
in the squared Vandermonde density and integrating coordinate by
coordinate proves
\begin{equation}\label{eq:cue-multiplier}
 \mathbb E\prod_l u(\theta_l)
       =\det(P_NuP_N|_{\operatorname{ran}P_N}).
\end{equation}
Conditioning the marked point to be $x$ replaces $P_N$ by
$Q_{N,x}=P_N-|v_x\rangle\langle v_x|$, where
$v_x(\theta)=P_N(\theta,x)/\sqrt N$. This follows from the Schur
complement of the marked row and column of the correlation determinant.
The marked point is uniform on the circle.

For a finite-rank orthogonal projection $Q$ and $|u|=1$, the positive
contraction $Qu^*(1-Q)uQ$ gives
\begin{equation}\label{eq:cue-leakage}
 |\det(QuQ)|^2
 \le\exp\{-\|(1-Q)uQ\|_{\rm HS}^2\}.
\end{equation}
For $P_N$, counting the Fourier basis vectors that leave its range
under multiplication by a Fourier mode gives
\[
 F_N(u):=\|(1-P_N)uP_N\|_{\rm HS}^2
       =\sum_{n\in\mathbb Z}\min(N,|n|)|\widehat u(n)|^2.
\]
A rank-one removal changes this quantity by at most one: for $P=Q+R$,
$\operatorname{rank}R=1$, the difference is
$-\|RuQ\|_{\rm HS}^2+\|(1-P)uR\|_{\rm HS}^2\in[-1,1]$.

Consider a nonconstant Fourier mode $e^{i\pi kW_{N,I}}$. If $k$ is
even, simultaneous rotation of the matrix leaves the exterior phase
and cyclic gaps invariant and averages $e^{-ikN\theta_I/2}$ to zero.
This conclusion continues to hold after multiplying by any bounded
function of the cyclic gaps. If $k$ is odd, condition on $x$.
Since $I=1+\#\{\theta_l\in(0,x):l\ne I\}$, the mode is, up to a
scalar of modulus one, the determinant in \eqref{eq:cue-multiplier}
with projection $Q_{N,x}$ and multiplier
\[
 u_x(\theta)=(-1)^{\mathbf1_{(0,x)}(\theta)}
     \exp\{-ik\Im\log(1-r e^{i(x-\theta)})\}.
\]
It has two jumps of absolute size two. Two integrations by parts on
its smooth pieces show
\[
 \widehat u_x(n)=\frac{J_0+J_xe^{-inx}}{2\pi in}+O_{x,k,r}(n^{-2}),
 \qquad |J_0|=|J_x|=2.
\]
For fixed $0<x<2\pi$, the cross sum
$\sum_{n\le N}e^{inx}/n$ is bounded. Hence
$F_N(u_x)\ge(4/\pi^2)\log N+O_{x,k,r}(1)$, and
\eqref{eq:cue-leakage} proves that the odd mode tends to zero.
Dominated convergence integrates over $x$. Bounds for the smooth
factor are uniform on every indicated range of $r$.

For the joint assertion fix $s>0$ and let $J$ be the arc of length
$2\pi s/N$ following $x$. In the determinant replace $u_x$ by
$w=u_x z^{\mathbf1_J}$, $|z|\le1$. Then
\[
 \|(w-u_x)Q_{N,x}\|_{\rm HS}^2
 \le4\operatorname{Tr}(Q_{N,x}\mathbf1_JQ_{N,x})\le4s.
\]
For a contraction multiplier $w$, the operator
$I-(QwQ)^*(QwQ)$ is the sum of $Q(1-|w|^2)Q$ and
$Qw^*(1-Q)wQ$. Its determinant therefore has the same bound as
\eqref{eq:cue-leakage} with exponent at least
\[
 \bigl(\|(1-Q)u_xQ\|_{\rm HS}-2\sqrt s\bigr)_+^2.
\]
It tends to infinity uniformly for $z$ on any circle $|z|=b<1$.
Cauchy coefficient extraction proves vanishing of every odd mixed
mode against $\mathbf1_{\#J=l}$, for fixed $l$. The event
$D_{j,N}>s$ is the union $\#J\le j-1$, so all its nonconstant
mixed modes vanish. The even modes already vanish by rotation.

Rescaling by $N/(2\pi)$, and removing a scalar phase, makes
$Q_{N,x}$ converge on compact intervals to $K^0$. Its count
generating functions converge: in the Fredholm series Hadamard's
bound gives a summable majorant $s^m/m!$. The limiting gap has no
atoms, since the intensity is continuous, and the finite gaps are
tight because cyclic summation gives $\mathbb ED_{j,N}=j$.
These facts and the mixed Fourier identities identify the joint
limit as the product measure in the theorem.

Finally,
\[
 \mathbb E\sup_x|\phi'_{N,r}(x)|
 \le\sum_{m\ge1}r^m\mathbb E|\operatorname{Tr}U_N^m|
 \le\sum_{m\ge1}r^m\sqrt m<\infty.
\]
The identity $\mathbb E|\operatorname{Tr}U_N^m|^2=\min(m,N)$ follows
by integrating the two-point determinant against the Fourier factor
$e^{im(\theta-\psi)}$.
Thus the phase change over a tight scaled gap tends to zero in
probability, and
\[
 W_{N,I+j}-W_{N,I}=j-D_{j,N}+o_{\mathbb P}(1).
\]
For uniform $V$ modulo two, the signs at $V$ and $V+v$ agree with
probability $1-d_2(v)$. Their discontinuity sets have zero limiting
measure, which completes the proof.
\end{proof}

\subsection{Distant labels}
A local central limit estimate supplies the fixed-frequency information
needed when the label separation grows.

\begin{theorem}\label{thm:palm-local-limit}
Let $f_j$ be the density of $D_j$, extended by zero on the negative
half-line, and let $\sigma_j=\sqrt{\log j}/\pi$, $j\ge2$. Then
\[
 \int_{\mathbb R}\left|\sigma_j f_j(j+\sigma_j y)
       -\frac{e^{-y^2/2}}{\sqrt{2\pi}}\right|dy\longrightarrow0.
\]
For every fixed $b\ne0$, $\mathbb Ee^{ib(D_j-j)}\to0$.
In particular the correlations in Theorem~\ref{thm:cue-labels}
tend to zero as $j\to\infty$, after taking its $N\to\infty$ limit.
\end{theorem}

\begin{proof}
The sine projection compressed to $(0,s)$ is a positive trace-class
contraction $A_s$. Its count generating function
$\det(I+(z-1)A_s)$ identifies the count as a sum of independent
Bernoulli variables with parameters the eigenvalues of $A_s$.
Direct integration gives
\[
 \operatorname{Tr}A_s=s,\quad
 \operatorname{Tr}(A_s-A_s^2)
 =s-2\int_0^s(s-u)k(u)^2du=\pi^{-2}\log s+O(1).
\]
Indeed $\int_0^\infty k(u)^2du=1/2$; the tail multiplied by $s$
is bounded, and $2\int_0^s uk(u)^2du=\pi^{-2}\log s+O(1)$.
Conditioning at zero and at $s$ removes two evaluation vectors.
The compressed two-point Palm operator $B_s$ therefore satisfies
$R=A_s-B_s\ge0$, $\operatorname{Tr}R\le2$, and
\[
 \operatorname{Tr}B_s=s+O(1),\qquad
 \operatorname{Tr}(B_s-B_s^2)=\pi^{-2}\log s+O(1).
\]
The variance difference is $\operatorname{Tr}R-
\operatorname{Tr}(R(A_s+B_s))$, of modulus at most two.

For completeness, a Bernoulli sum $X$ of mean $\mu$ and variance $v$
satisfies, uniformly over its parameters as $v\to\infty$,
\begin{equation}\label{eq:bernoulli-local}
 \sup_{l\in\mathbb Z}\left|\sqrt v\,\mathbb P(X=l)
  -(2\pi)^{-1/2}e^{-(l-\mu)^2/(2v)}\right|\longrightarrow0.
\end{equation}
Its centered characteristic function obeys
$|\psi(\theta)|\le e^{-2v\sin^2(\theta/2)}
\le e^{-2v\theta^2/\pi^2}$ for $|\theta|\le\pi$.
Taylor expansion of its logarithm gives
$\log\psi(u/\sqrt v)=-u^2/2+O(|u|^3/\sqrt v)$; the cubic
remainder is bounded by a constant times each Bernoulli variance.
Fourier inversion and dominated convergence, after extending the
scaled integrand by zero outside $[-\pi\sqrt v,\pi\sqrt v]$,
prove \eqref{eq:bernoulli-local}. Finite truncation proves the same
assertion for countable Bernoulli sums of finite mean.

The Palm Campbell formula gives
\[
 f_j(s)=(1-k(s)^2)\,
      \mathbb P^{0,s}\{N(0,s)=j-1\},\qquad s>0.
\]
It follows by counting a point at $s$ exactly when there are $j-1$
points between zero and $s$. Applying \eqref{eq:bernoulli-local} to
$B_s$ proves the stated density convergence uniformly on every
bounded interval of $y$. The one-point Palm count has the same
mean and variance estimates. Chebyshev at $s=j\pm M\sigma_j$
gives
\[
 \limsup_{j\to\infty}\mathbb P(|D_j-j|>M\sigma_j)\le C/M^2.
\]
There are infinitely many positive points almost surely: Chebyshev
at dyadic lengths and Borel--Cantelli prove this. Thus $f_j$ has
mass one, and the tail estimate upgrades compact convergence to
$L^1$ convergence. If $g_j(y)=\sigma_jf_j(j+\sigma_jy)$ and
$\varphi$ is the standard Gaussian density, then
\[
 \left|\mathbb Ee^{ib(D_j-j)}-e^{-b^2\sigma_j^2/2}\right|
 \le\|g_j-\varphi\|_1\longrightarrow0.
\]
Taking $b=\pi l$, $l\in\mathbb Z\setminus\{0\}$, proves Haar
convergence modulo two. The continuous function $1-2d_2$ has Haar
mean zero. Since both limiting signs are fair and their joint law
is invariant under flipping both signs, zero correlation makes
all four joint probabilities tend to $1/4$.
\end{proof}

\subsection{The adjacent probability and its certificate}
Let $K_s$ be the sine projection compressed to an interval of length
$s$, and let $E(s)=\det(I-K_s)$ be its empty-interval probability.

\begin{theorem}\label{thm:cue-agreement}
The adjacent agreement probability in Theorem~\ref{thm:cue-labels} is
\[
 A_{\rm adj}=1-2\sum_{n\ge1}(-1)^{n+1}E(n).
\]
The series converges absolutely, and rigorous bounds are
\[
 0.6665466018372107200719020708703
 <A_{\rm adj}<
 0.6665466018372107200719020748216.
\]
\end{theorem}

\begin{proof}
Differentiating the right endpoint of an empty interval and using
Campbell's formula gives $-E'(s)=\mathbb P^0(D_1>s)$, and hence
$E''=f_1$, distributionally. The probability of two points in an
added interval of length $h$ is $O(h^2)$, since the two-point
intensity is bounded by one; this justifies the first derivative.
At zero, $E(0)=1$ and $E'(0)=-1$. The eigenvalues of $K_s$ lie in
$[0,1]$ and sum to $s$, so $0\le E(s)\le e^{-s}$.
Put $g(s)=d_2(1-s)$. Then $g(0)=1$, $g'(0+)=-1$, and
$g''=2\sum_{n\ge1}(-1)^{n+1}\delta_n$ on $(0,\infty)$.
Twice integrating by parts gives
\[
 \mathbb E d_2(1-D_1)=\int_0^\infty g(s)E''(s)ds
      =2\sum_{n\ge1}(-1)^{n+1}E(n).
\]
The boundary terms at zero cancel. Those at infinity vanish because
$E$ decays exponentially and $E'(s)\to0$. Absolute convergence and
monotonicity of $E$ give the alternating upper and lower bounds.

Here is the analytic error estimate used for the numerical interval.
On $a,\xi\in[-1/2,1/2]$ the Fourier-feature kernel is
$T(a,\xi)=\sqrt s\,e^{2\pi is a\xi}$ and $K_s$ is unitarily
equivalent to $TT^*$. Let $T_p$ replace the exponential by its
Taylor polynomial of degree $p$. With
\[
 h=\pi s/2,\quad R=e^h h^{p+1}/(p+1)!,\quad
 \delta=sR(2+R),
\]
we have
\[
 \|TT^*-T_pT_p^*\|_1\le\delta,
 \qquad \operatorname{Tr}(T_pT_p^*)\le s(1+R)^2.
\]
For positive finite matrices, differentiating the determinant along
their positive interpolation, and using the adjugate bound
$\|\operatorname{adj}(I-C)\|\le\prod_l(1+\lambda_l(C))
\le e^{\operatorname{Tr}C}$, proves
\[
 |\det(I-A)-\det(I-B)|\le\|A-B\|_1
           e^{\max(\operatorname{Tr}A,\operatorname{Tr}B)}.
\]
Finite-rank approximation and the absolutely convergent
exterior-power series extend this to positive trace-class operators.
Thus the more generous error
$\delta\exp(1+s+s(1+R)^2)$ encloses the determinant error.

The finite calculation uses orthonormal Legendre polynomials
$L_n(u)=\sqrt{2n+1}\,P_n(2u)$ on $[-1/2,1/2]$. Their exact
moments, from Rodrigues' formula and integration by parts, are
\[
 M_{kn}:=\int_{-1/2}^{1/2}u^kL_n(u)du
 =\frac{\sqrt{2n+1}\,k!}
        {2^k(k-n)!!(k+n+1)!!}
\]
when $k\ge n$ and $k-n$ is even, and zero otherwise.
The matrix of the truncated feature map, divided by $\sqrt s$, is
\[
 T^{(p)}_{mn}=\sum_{k=0}^p
                 \frac{(2\pi is)^k}{k!}M_{km}M_{kn}.
\]
Its even block is real and its odd block purely imaginary, so its
Gram determinant is the product of two real determinants. This
provides a finite expression using rational numbers, square roots
and $\pi$, with no quadrature. Ball arithmetic at 180 decimal digits,
with $p=110$, evaluates $E(1),\ldots,E(7)$ and enlarges every ball
by the analytic error just proved. The alternating sums through
seven and six give, respectively, the asserted lower and upper
bounds. The program \texttt{certify\_cue\_agreement}, its complete
ball intervals and the exact command are available through the
workbench.
\end{proof}

\section{What the arithmetic gap theorems give for collisions}
\label{sec:arithmetic-collisions}
The isolated-pair theorem applies directly to a certified pair of zeta
zeros. Passing from that pairwise assertion to a statistical law requires
additional information about the complementary zeros. We make the
required quantities explicit and compare them with the gap theorems.

For consecutive distinct simple ordinates $\gamma_-<\gamma_+$, put
\[
 x_\pm=2\gamma_\pm,\qquad b=\gamma_-+\gamma_+,\qquad
 \delta=2(\gamma_+-\gamma_-),\qquad
 S=\sum_{\xi\ne x_-,x_+}(b-\xi)^{-2},\quad\kappa=\delta^2S.
\]
The sum runs through the full real zero divisor of $H_0$, including
negative zeros and multiplicities. These definitions give the exact
arithmetic specialization of Theorem~\ref{thm:isolated-pair}:
for the pairs in its stated range the collision time is
\[
 a_c=\frac{(\gamma_+-\gamma_-)^2}{2}
       +\frac{(\gamma_+-\gamma_-)^4}{2}S
       +O((\gamma_+-\gamma_-)^6S^2).
\]
The relation $x=2\gamma$ accounts for both factors of two. The moving
center is $b-2a_cG'(b)/G(b)+O(\delta\kappa^{3/2})$, where
$G(z)=H_0(z)/((z-b)^2-\delta^2/4)$.

The complementary sum admits a finite-height bound. For $R\ge1$,
write
\[
 S_R=\sum_{\substack{0<|\xi-b|\le R\\\xi\ne x_-,x_+}}(b-\xi)^{-2}.
\]
The zero count in every unit interval is $O(\log(|u|+3))$, including
multiplicity. Splitting the two tails into unit intervals and comparing
with $\int_R^\infty\log(b+u+3)u^{-2}du$ gives an absolute constant
$C$ such that
\begin{equation}\label{eq:pair-exterior-tail}
 0\le S-S_R\le C\frac{\log(b+R+3)}R
 \qquad(b\ge0,\ R\ge1).
\end{equation}
For example, integration by parts bounds that integral by
$\log(b+R+3)/R+\int_R^\infty[u(b+u+3)]^{-1}du$, and the last
integral is at most $1/R$. Thus interval enclosures for finitely many
zeros, together with an explicit unit-interval zero-count constant,
produce a rigorous upper bound for $\kappa$. The gaps enter the
calculation through the denominators in $S_R$; the tail estimate does
not replace them.

The small-gap theorems recently reached half the mean spacing.
Bondarenko and Heap \cite[Theorem~1, p.~1]{bondarenko-heap-half} obtain
\[
 \liminf_{n\to\infty}
       \frac{(\gamma_{n+1}-\gamma_n)\log\gamma_n}{2\pi}\le\frac12
\]
from the zero-location theorem. Their method uses resonance for a
quotient of shifted zeta functions. Durkan and Pearce-Crump
\cite[Theorems~1--2, pp.~1--2]{durkan-pearce-crump} establish the same
lower-limit bound for every fixed primitive Dirichlet character, and
for zeta a positive proportion of gaps smaller than
$(1/2+\epsilon)$ times the mean, for every fixed $\epsilon>0$.
Their count includes repeated ordinates. Lamzouri's
\cite[Theorem~1.1, p.~2]{lamzouri-small-gaps} positive-proportion bound
at $0.7892$ is obtained without a zero-location assumption.
These statements specify the present quantitative input at a fixed
multiple of the mean spacing.

There is an elementary obstruction to deducing $\kappa\to0$ from
a bound on that multiple alone. Choose $0<\lambda<1/2$ and consider
the entire function
\[
 f_\lambda(z)=\sin(\pi z/3)\sin(\pi(z-\lambda)/3)
                         \sin(\pi(z-2\lambda)/3).
\]
Its simple real zeros have mean spacing one and repeating gap pattern
$\lambda,\lambda,3-2\lambda$. For the adjacent pair $0,\lambda$,
the third zero $2\lambda$ alone gives
\[
 \delta^2 S\ge\lambda^2(2\lambda-\lambda/2)^{-2}=4/9.
\]
The fraction of gaps of length $\lambda$ is $2/3$.
The dilation $z\mapsto z/c$ multiplies every gap by $c$ and every
term of $S$ by $c^{-2}$, so it preserves $\delta^2S$ exactly.
Thus even arbitrarily small absolute gaps and a positive fraction
of small normalized gaps can coexist with a complementary interaction
bounded below. This example disproves that proposed implication of
those two scalar gap properties; it makes no claim to reproduce the
zeta divisor.

There is also a precise limitation to using the restricted correlation
identities alone. Lagarias and Rodgers
\cite[Theorem~3.1, p.~4, with the test class in Theorem~2.4, p.~3]{lagarias-rodgers}
construct a deterministic unit-density sequence with all gaps in
$\{1/2,1,3/2,\ldots\}$ whose translation-averaged correlation
statistics agree with the sine determinant for every order and every
Schwartz test in that class. On the hyperplane
$\sum\xi_i=0$, the allowed Fourier support is
$\sum|\xi_i|<2$, exactly the restricted range used here.
Every gap in that sequence is at least $1/2$, whereas a sine process
has positive probability of a gap in $(0,\epsilon)$ for every
$\epsilon>0$. The common restricted identities therefore do not
determine these small-gap events. This counterexample concerns the
specified statistics; it does not assert that the sequence has an
Euler product or satisfies the functional equation of zeta.

The next arithmetic question is consequently specific. Can one prove
that sufficiently many adjacent \emph{distinct simple} zeta pairs have
$\delta^2\sum_{\xi\ne x_-,x_+}(b-\xi)^{-2}\to0$, with quantitative
control of nearby triples and of the remaining singleton branches?
For an arithmetic collision point-process limit, can these estimates
be combined with a marked extreme-gap limit on the shrinking scale,
including its joint law and tightness? The proved local expansion
identifies the map from such pair data to collision data, but the
fixed-scale gap results quoted above leave these estimates open.

The recent half-gap work also gives a documented example of assistance
by OpenAI's ChatGPT in neighbouring research: Bondarenko and Heap
acknowledge exploratory discussions, calculations and numerics on p.~2
of their paper. Their theorem and proof remain credited to its named
authors. The present collision estimates use their theorem only in the
comparison of the arithmetic gap input.

\section{The lowest Weil eigenvector and the simplicity question}
\label{sec:weil-parity}
This section connects the lowest eigenvector in the programme of Connes,
Consani and Moscovici to the product over distinct zero ordinates, and then
to the multiplicities of the Riemann xi function.

Use the Fourier transform
$\widehat f(z)=\int_{\mathbb R}f(x)e^{-izx}\,dx$ and inner products linear
in their first argument. Write $\Gamma$ for the distinct real ordinates,
$m_\gamma$ for their multiplicities, and
\[
 q(f,g)=\sum_{\gamma\in\Gamma}
 m_\gamma\widehat f(\gamma)\overline{\widehat g(\gamma)}.
\]
The Riemann hypothesis, now provided by the input theorem, identifies
this expression with the positive Weil form. Its global completion is
$H=\ell^2(\Gamma,m)$ under sampling. The local completion $H_L$ of
$C_c^\infty(-L,L)$ embeds in $H$ and compactly in $L^2(-L,L)$.
The closed Weil operator $A_L$ has positive eigenvalues
$a_1(L)\le a_2(L)\le\cdots$, listed with multiplicity.
These local form and compact-resolvent facts are established in
\cite[\S3]{ccm2025}; the sampling realization and its unshifted inverse
compression are recorded in \cite[\S2.3 and Theorem 3.3]{suzuki}.
The form domain consists of zero-extended functions for which
$\int(1+\log(2+|t|))|\widehat f(t)|^2\,dt$ is finite. In this norm,
shrinking support by a dilation and then mollifying proves that compact
smooth functions are a form core. Dilation continuity follows first on
Schwartz functions, then by density and uniform boundedness for dilation
factors between $1/2$ and $1$.

Let $J_L:H_L\to L^2(-L,L)$ be inclusion, $\Pi_L:H\to H_L$ the
orthogonal projection, and $C$ multiplication by $\gamma^{-1}$ on $H$.
There is no zero at ordinate zero: the alternating-series identity gives
$(1-\sqrt2)\zeta(1/2)>0$. The order-one zero count gives
$\sum_\gamma m_\gamma/\gamma^2<\infty$, so $C$ is Hilbert--Schmidt.
Put $C_L=\Pi_LC\Pi_L$. Reflection decomposes $H_L$ into even and odd
parts, in which
\[
 C_L=\begin{pmatrix}0&X_L^*\\X_L&0\end{pmatrix},\qquad
 J_L^*J_L=\begin{pmatrix}T_{L,+}&0\\0&T_{L,-}\end{pmatrix}.
\]
Let $a_e(L)$ and $a_o(L)$ denote the minima in these two parts.
All the adjoints and projections in these formulas use $q$.

\begin{proposition}[The parity comparison]\label{prop:parity-comparison}
At a window with $a_e(L)<a_o(L)$ the lowest eigenvalue is simple and
even, and its eigenvector $e_L$ has nonzero integral. Put
\[
 \alpha_L=a_e(L),\quad M_{L,\pm}=\alpha_LT_{L,\pm},\quad
 S_{L,\pm}=I-M_{L,\pm},\quad
 \kappa_L=1-a_e(L)/a_o(L).
\]
On $H_L/\mathbb Ce_L$, with the form
$q-\alpha_L\langle\cdot,\cdot\rangle_2$, the compact primitive is
unitarily represented by the restriction off $e_L$ of
\[
 B_L=\begin{pmatrix}
 0&S_{L,+}^{1/2}X_L^*S_{L,-}^{-1/2}\\
 S_{L,-}^{-1/2}X_LS_{L,+}^{1/2}&0
 \end{pmatrix}.
\]
It is injective and selfadjoint. The reciprocals of its nonzero
eigenvalues are precisely the zeros of $\widehat e_L$, all of which
are real and simple. Moreover,
\begin{align}
 \|B_L-C_L\|_{\mathcal S_2}
 \le\sqrt2\bigl(&\kappa_L^{-1/2}
       \|M_{L,-}X_L\|_{\mathcal S_2}
       +\|X_LM_{L,+}\|_{\mathcal S_2}\bigr).
 \label{eq:parity-error}
\end{align}
Along every sequence $L\to\infty$ on which $\inf\kappa_L>0$,
\[
 \|B_L-C\|_{\mathcal S_2}\longrightarrow0,\qquad
 \frac{\widehat e_L(z)}{\widehat e_L(0)}\longrightarrow
 \Delta_\Gamma(z):=\prod_{\gamma>0,\ \gamma\in\Gamma}
              (1-z^2/\gamma^2)
\]
locally uniformly in $z$. Each distinct positive ordinate appears once
in this product.
\end{proposition}

\begin{proof}
Fix a window and write $\ell(f)=\int f$. Coercivity makes $\ell$ a
bounded functional on $H_L$. For a compact smooth mean-zero function,
\[
 C_Lf=\mathcal Pf,\qquad
 (\mathcal Pf)(x)=i\int_{-L}^x f(y)\,dy.
\]
Indeed its Fourier transform is $\widehat f(t)/t$, and its support
remains in the interval. This identity extends to every $f\in\ker\ell$:
approximate in the form norm by compact smooth functions, correct their
integrals with a fixed compact smooth function of integral one, and use
the boundedness of $C_L$. Their primitives also converge in
$H^1_0(-L,L)$, so the two limits agree in $L^2$.
Thus $DC_Lf=f$ on $\ker\ell$, where $D=-i\partial_x$, and
$C_LDg=g$ for compact smooth $g$.

An even bottom vector $u$ with integral zero would give a nonzero odd
primitive $v=C_Lu$. Its integral is zero by parity, so $w=C_Lv$ also
has compact support and $Dv=u$, $Dw=v$. The bottom eigenvector identity
and integration by parts give
\[
 q(v)=q(u,w)=\alpha_L\langle u,w\rangle_2
       =\alpha_L\langle Dv,w\rangle_2
       =\alpha_L\|v\|_2^2.
\]
This contradicts the strictly larger odd minimum. Two independent even
bottom vectors would have a nonzero mean-zero linear combination.
Consequently the bottom line is simple, even, and has nonzero integral.
Choose its $q$-unit vector $e$. On its orthogonal complement,
$q_\alpha=q-\alpha_L\langle\cdot,\cdot\rangle_2$ is an equivalent
Hilbert norm because $J_L^*J_L$ is compact and its largest eigenvalue
is simple.

Integration by parts for two mean-zero functions yields
$\langle C_Lf,g\rangle_2=\langle f,C_Lg\rangle_2$.
Hence $X_LM_{L,+}-M_{L,-}X_L$ vanishes on the even kernel of $\ell$.
Its value at $e$ is $S_{L,-}X_Le$, and therefore
\[
 X_LM_{L,+}-M_{L,-}X_L
   =S_{L,-}X_Le\,\frac{\ell}{\ell(e)}.
\]
Define $K_Lf=C_L(f-e\ell(f)/\ell(e))$. It kills $e$ and is the
compact primitive of the mean-zero expression in parentheses.
In quotient coordinates orthogonal to $e$, conjugation by
$S_{L,+}^{1/2}\oplus S_{L,-}^{1/2}$ gives the displayed $B_L$;
the commutator identity rewrites its lower-left block as
$S_{L,-}^{-1/2}X_LS_{L,+}^{1/2}$. The two blocks are adjoints.
No inverse of $S_{L,+}$ on its kernel is taken.

For injectivity, $C_LDg=g$ and the density of compact smooth $g$ first
show that $C_L$ is injective. If $K_Lf=ce$ and $c\ne0$, then
$e=C_Lg$ for a mean-zero $g$; injectivity and reflection make $g$ odd.
The primitive identity gives $De=g$. For compact smooth $h$,
\[
 q(g,h)=q(e,Dh)=\alpha_L\langle e,Dh\rangle_2
                   =\alpha_L\langle g,h\rangle_2.
\]
The form core extends this identity to $h=g$, contradicting the odd
minimum. Thus $c=0$, and injectivity of $C_L$ gives $[f]=0$.

Set $F=\widehat e$. It is even, entire of exponential type at most $L$,
and $F(0)\ne0$. The primitive identity gives
\[
 \widehat {K_Lf}(z)=
 \frac{\widehat f(z)-F(z)\ell(f)/\ell(e)}{z},\qquad z\ne0.
\]
At a zero of $F$, Fourier evaluation descends to a nonzero bounded
functional on the quotient. Its Riesz vector is an eigenvector with
eigenvalue $1/\overline z$, so selfadjointness makes $z$ real.
Conversely an eigenvector of the quotient with eigenvalue $\beta\ne0$
defines a distribution $\Lambda$ on $(-L,L)$ satisfying
$\Lambda(Dh)=\beta^{-1}\Lambda(h)$. Solving this constant-coefficient
first-order distributional equation gives
$\Lambda(h)=c\widehat h(\beta^{-1})$ with $c\ne0$. Since $\Lambda$
annihilates $e$, $F(\beta^{-1})=0$. This equation also shows the
eigenspace has dimension one. A double zero of $F$ would make evaluation
and its derivative descend simultaneously; differentiating the last
display gives a generalized eigenvector $w'$ satisfying
$(K_L-z^{-1})w'=-z^{-2}w$, where $w\ne0$ is the evaluation eigenvector.
Pairing with $w$ contradicts selfadjointness. All Fourier zeros are
therefore simple.

For the limiting comparison, take a compact smooth $h$ of integral one
and put $h_R(x)=R^{-1}h(x/R)$. Integration by parts gives
$q(h_R)\le R^{-2}\|h'\|_1^2\sum m_\gamma/\gamma^2$, while
$\|h_R\|_2^2=R^{-1}\|h\|_2^2$. Taking $R$ proportional to $L$ proves
$a_1(L)\to0$. Extend $M_L$ by zero to $H$. It is a positive contraction,
and for each fixed compact smooth $f$,
\[
 \|M_Lf\|_q^2\le q(M_Lf,f)=\alpha_L\|f\|_2^2\longrightarrow0.
\]
Density gives strong convergence $M_L\to0$. The increasing projections
$\Pi_L$ converge strongly to $I$, so Hilbert--Schmidt approximation gives
$C_L\to C$, $M_LC_L\to0$ and $C_LM_L\to0$ in
$\mathcal S_2$.

For $0\le x\le1-\kappa_L$,
\[
 (1-x)^{-1/2}-1
 =\frac{x}{\sqrt{1-x}(1+\sqrt{1-x})},\qquad
 |(1-x)^{1/2}-1|\le x.
\]
Apply these identities by functional calculus to the two blocks of
$B_L-C_L$. They give \eqref{eq:parity-error}, and the assumed positive
lower bound for $\kappa_L$ gives $B_L\to C$ in Hilbert--Schmidt norm
along the specified sequence.

Let $\beta_{L,j}>0$ be the positive eigenvalues of $B_L$ in decreasing
order. Min--max gives $\beta_{L,j}\to1/\gamma_j$ for every fixed $j$.
Reflection and Hilbert--Schmidt convergence also give
$\sum_j\beta_{L,j}^2\to\sum_j\gamma_j^{-2}$. Subtracting finitely
many convergent terms controls the remaining sums uniformly.
Hadamard factorization of the even exponential-type function $F$ gives
\[
 F(z)/F(0)=\prod_j(1-z^2\beta_{L,j}^2).
\]
For $|z|\le R$, the logarithms of sufficiently late factors have total
modulus at most $2R^2\sum_{j>J}\beta_{L,j}^2$. The uniform tail bound
and fixed-index convergence prove the asserted entire limit.
\end{proof}

\begin{proposition}[A subsequence criterion for simplicity]
\label{prop:liminf-rayleigh}
Let $K\in L^2(\mathbb R)$ be the inverse Fourier transform of $\Xi$,
and let $v_L$ be real even $L^2$-unit vectors in the local Weil form
domains with $v_L\to K/\|K\|_2$ in $L^2(\mathbb R)$.
The equality
\begin{equation}
 \liminf_{L\to\infty}\frac{q(v_L)}{a_2(L)}=0
 \label{eq:liminf-rayleigh}
\end{equation}
forces every nontrivial zeta zero to be simple. Equivalently, the
existence of a multiple zero forces a strictly positive lower limit
for this quotient. This applies to the original even prolate trial
of Connes, Consani and Moscovici.
\end{proposition}

\begin{proof}
Choose a sequence of windows along which
$r_L=q(v_L)/a_2(L)\to0$. Min--max gives $a_1\le q(v_L)<a_2$ eventually,
so the bottom line is simple. An odd bottom vector would be orthogonal
to the even trial and would give $q(v_L)\ge a_2$. Thus the bottom vector
$u_L$ is real and even. Choose its sign with positive overlap with $v_L$.
An eigenvector expansion gives
\[
 1-|\langle u_L,v_L\rangle_2|^2\le r_L,
 \qquad \|u_L-v_L\|_2^2\le2r_L.
\]
Hence $u_L\to K/\|K\|_2$. Furthermore
$a_e/a_o\le a_1/a_2\le r_L$, so the parity comparison applies along
this sequence and gives
$G_L=\widehat u_L/\widehat u_L(0)\to\Delta_\Gamma$ locally uniformly.
Plancherel gives $\widehat u_L\to\Xi/\|K\|_2$ in $L^2(\mathbb R)$.

Choose a small interval $I$ about zero on which $\Delta_\Gamma$ is
nonzero. With $c_L=\widehat u_L(0)$, the exact identity
\[
 c_L=\frac{\langle\widehat u_L,G_L\rangle_{L^2(I)}}
                {\|G_L\|_{L^2(I)}^2}
\]
shows that $c_L$ has a finite limit $c$. Thus
$c\Delta_\Gamma=\Xi/\|K\|_2$ almost everywhere on $I$, and continuity
and the identity theorem extend the equality to the whole plane.
The limit $c$ is nonzero because $\Xi(0)\ne0$, and evaluation at zero
gives $c=\Xi(0)/\|K\|_2$. Therefore
$\Delta_\Gamma=\Xi/\Xi(0)$. Since the factors of $\Delta_\Gamma$
use each ordinate once, every multiplicity is one.
\end{proof}

The estimate in the last proposition remains a question about the Weil
operator. It has the weaker subsequence quantifier displayed above;
neither positivity of the form nor the compact comparison proves it.
For the original prolate trial the available upper bound is
$q(v_{\log\lambda})\ll_\varepsilon
\lambda^\varepsilon(1-\nu_4)$, with
$1-\nu_4\asymp\lambda^9e^{-4\pi\lambda^2}$.
An upper bound for a trial does not provide the missing comparison with
$a_2$; the needed quantitative estimate is precisely the lower-limit
equality in the proposition.

\section{The prolate trials and their Rayleigh upper bounds}
\label{sec:prolate-trial}
This section defines the original trial in the spectral programme and
checks the upper estimate used in the Rayleigh criterion.
Andrews, Connes and La Gatta record a related upper estimate of Connes
and Kastoris for a Gevrey modification of the trial
\cite[Theorem~3.1, p.~14]{andrews-connes-lagatta2026}.
Their map first multiplies the raw prolate vector by an even cutoff in
an endpoint layer of width $\lambda^{-3}$, adds two fixed smooth bumps
to restore its value and integral at zero, applies the theta map, and
cuts off the lower logarithmic endpoint in a layer of width
$\lambda^{-4}$. They give a Rayleigh bound
$C(1-\mu_4)(\log\lambda)^3$ for that vector. In the notation below,
$1-\mu_4^2=(1+\mu_4)(1-\mu_4)$ exactly. We give the proof for the
original sharply truncated vector, whose exact definition is needed
in the Rayleigh criterion.
For raw variables use
$\mathcal Ff(v)=\int f(u)e^{-2\pi iuv}\,du$ and
$P_\lambda={\bf1}_{[-\lambda,\lambda]}$, with $L=\log\lambda$.
Let $\phi_{j,\lambda}$ be the real $L^2$-unit time-supported prolate
functions with Hermite index $j$, and write
\[
 P_\lambda\mathcal F\phi_{j,\lambda}
       =\mu_j\phi_{j,\lambda},\qquad
 \nu_j=\mu_j^2,\qquad d_j=\phi_{j,\lambda}(0).
\]
For $j=0,4$, the signs are chosen so that the limiting Hermite
functions have their usual real signs. For large $\lambda$ the two
$\mu_j$ are positive and $d_j\ne0$. Define
\[
 \psi_\lambda=-\frac34\left(
       \mu_4\frac{\phi_{0,\lambda}}{d_0}
       -\mu_0\frac{\phi_{4,\lambda}}{d_4}\right),\qquad
 k_\lambda(x)={\bf1}_{[-L,L]}(x)e^{x/2}
                  \sum_{n\ge1}\psi_\lambda(ne^x).
\]
The sum is finite for each $x$. The even normalized trial is
$v_L=(k_\lambda(x)+k_\lambda(-x))/(2\|k_\lambda^+\|_2)$,
where $k_\lambda^+$ denotes the even part.
This is the two-mode construction of
\cite[\S7, pp.~28--31]{ccm2025}, with its scalar fixed as follows.

The Hermite ratios at zero are
\[
 h_0(u)/h_0(0)=e^{-\pi u^2},\qquad
 h_4(u)/h_4(0)=
 \tfrac13(16\pi^2u^4-24\pi u^2+3)e^{-\pi u^2}.
\]
Their difference times $-3/4$ is
$h(u)=(4\pi^2u^4-6\pi u^2)e^{-\pi u^2}$.
Consequently
\[
 K(x)=e^{x/2}\sum_{n\ge1}h(ne^x),\qquad
 \widehat K(t)=\xi(1/2-it)=\Xi(t).
\]
For the transform identity, the Mellin integral of $h$ is
$\tfrac12s(s-1)\pi^{-s/2}\Gamma(s/2)$; multiply by $\zeta(s)$
and continue the absolutely convergent theta integral.
Poisson summation, $\mathcal Fh=h$, $h(0)=0$ and $\int h=0$ show
$K(-x)=K(x)$. Its displayed series gives decay faster than every
exponential at positive infinity, and hence at both ends.
Thus its Fourier integral is entire and $K\ne0$.

The fixed-index Hermite comparison in
\cite[Lemma~7.2, p.~29]{ccm2025} gives
$\sup_{|u|\le\lambda}|\psi_\lambda(u)-h(u)|\ll\lambda^{-2}$.
For completeness, the time-supported unit normalizations differ from
their Hermite limits in $L^2$ by $O(\lambda^{-3/2})$. Since the modes
$0,4$ are Fourier invariant in the limit, their leakage tends to zero;
this also gives $\mu_j\to1$ and the stated ratio comparison at zero.
Summing the uniform error over $n\le\lambda e^{-x}$ gives
$|k_\lambda(x)-K(x)|\ll\lambda^{-1}e^{-x/2}$ on $[-L,L]$, apart
from a Gaussian tail whose $L^2$ norm is exponentially small.
Integration yields
\begin{equation}
 \|k_\lambda-K\|_2=O(\lambda^{-1/2}),\qquad
 v_L\longrightarrow K/\|K\|_2.
 \label{eq:prolate-trial-limit}
\end{equation}
The raw mean is exactly zero: evaluation of the finite Fourier equation
at zero gives $\int_0^\infty\phi_{j,\lambda}=\mu_jd_j/2$, and these
terms cancel in $\psi_\lambda$.

\begin{proposition}[Leakage controls the original trial]
\label{prop:prolate-upper}
Put $\delta_\lambda=1-\nu_4$. For every $\varepsilon>0$,
\[
 q(k_\lambda)\ll_\varepsilon\lambda^\varepsilon\delta_\lambda,
 \qquad q(v_{\log\lambda})
                  \ll_\varepsilon\lambda^\varepsilon\delta_\lambda.
\]
The classical fixed-index prolate asymptotic gives
\[
 \delta_\lambda\sim
  \frac{2^{15}\sqrt2\pi^5}{3}\lambda^9e^{-4\pi\lambda^2},
\]
and consequently
$a_1(\log\lambda)\ll_\varepsilon
\lambda^{9+\varepsilon}e^{-4\pi\lambda^2}$.
\end{proposition}

\begin{proof}
We first prove the exterior estimates needed for the logarithmic tail.
Fix an even index $j$ and write $y=\mathcal F\phi_{j,\lambda}$,
$\delta_j=1-\nu_j$. Its differential eigenvalue $\chi_j$ obeys
$0\le\chi_j\ll_j\lambda^2$: a fixed-dimensional space of compact
smooth test functions proves the upper bound by min--max in the
prolate differential form. For $x>\lambda$, the Fourier continuation
satisfies
\[
 (py')'+Vy=0,\qquad
 p=x^2-\lambda^2,\qquad V=4\pi^2\lambda^2x^2-\chi_j.
\]
To justify the equation at the endpoint, extend the time-supported
function by zero. Multiplication of its distributional derivative by
$\lambda^2-x^2$ removes the endpoint delta terms; the remaining flux
also vanishes at both endpoints. Fourier transformation therefore gives
the equation globally, and $y$ is entire.
For $E=p|y'|^2+V|y|^2$, direct differentiation gives
\[
 (E/V)'=-\left(p'/V+pV'/V^2\right)|y'|^2\le0.
\]
Thus $|y(x)|\le|y(\lambda)|$. Integrating the equation from the
endpoint gives $|y'(x)|\ll\lambda^3|y(\lambda)|$ for
$\lambda\le x\le2\lambda$.

Here is an endpoint-to-leakage estimate with no differentiation of an
asymptotic formula. Put $c=2\pi\lambda^2$, $x=\lambda\cosh t$, and
$Z(t)=\sqrt{\sinh t}\,y(\lambda\cosh t)$. Then
\[
 Z''+\Omega Z=0,\qquad
 \Omega=c^2\cosh^2t-\chi_j-\tfrac14+(4\sinh^2t)^{-1}.
\]
For a fixed small $a>0$, set $t_0=a/c$ and $t_1=\operatorname{arcosh}2$.
The preceding derivative bound gives
$|y(\lambda\cosh t_0)|\ge|y(\lambda)|/2$.
On $[t_0,t_1]$, $\Omega$ lies between positive multiples of $c^2$ and
has one minimum, because
$\Omega'=\cosh t\,(4c^2\sinh^4t-1)/(2\sinh^3t)$.
The energy $\mathcal E=|Z'|^2+\Omega|Z|^2$ satisfies
$\mathcal E'=\Omega'|Z|^2$. On the decreasing portion use
$\mathcal E'\ge(\Omega'/\Omega)\mathcal E$; on the increasing
portion use $\mathcal E'\ge0$. The bounded ratio of the extrema gives
$\mathcal E(t_1)\gg c|y(\lambda)|^2$.

For $U(x)=xy(x)$ one obtains
\[
 U''+bU'+k^2U=0,\quad
 b=\frac{2\lambda^2}{x(x^2-\lambda^2)},\quad
 k^2=4\pi^2\lambda^2+
 \frac{4\pi^2\lambda^4-\chi_j-2\lambda^2/x^2}{x^2-\lambda^2}.
\]
On $x\ge2\lambda$, $k^2\ge4\pi^2\lambda^2$ and $(k^2)'<0$.
Hence $\mathcal E_U=|U'|^2+k^2|U|^2$ decreases.
At $2\lambda$ it is comparable to $\mathcal E(t_1)$, both being
comparable to
$\lambda^2|y'(2\lambda)|^2+\lambda^4|y(2\lambda)|^2$.
On $[2\lambda,3\lambda]$ the differential inequality
$\mathcal E_U'\ge(-2b+(k^2)'/k^2)\mathcal E_U$ shows that its
values are comparable. Choose a smooth cutoff $\theta$ in that
interval equal to one on its middle half, with
$|\theta'|\ll\lambda^{-1}$. Multiplication of the equation by
$\theta^2\overline U$ and integration by parts give
\[
 \int\theta^2|U'|^2
 =\int\theta^2k^2|U|^2+
  \operatorname{Re}\int(\theta^2b-2\theta\theta')U'\overline U.
\]
Young's inequality absorbs half the derivative integral and bounds
the rest by $O(\lambda^2\int|U|^2)$. Thus
$\lambda\mathcal E_U(2\lambda)
\ll\lambda^2\int_{2\lambda}^{3\lambda}|U|^2$, which yields
\[
 \delta_j\gg\lambda^{-3}\mathcal E_U(2\lambda)
                  \gg\lambda^{-1}|y(\lambda)|^2.
\]
The upper energy estimate at $2\lambda$, followed by its decrease,
also gives
\[
 |y(x)|\ll\lambda|y(\lambda)|/x,\qquad
 |y'(x)|\ll\lambda^4|y(\lambda)|/x\quad(x\ge\lambda).
\]
The interval $[\lambda,2\lambda]$ is covered by the earlier bounds.

Define the sharply cut logarithmic Fourier tail
$w_j(t)={\bf1}_{t>0}\sqrt\lambda e^{t/2}y(\lambda e^t)$.
Exactly $\|w_j\|_2^2=\delta_j/2$, while the exterior bounds give
\[
 |w_j(t)|\le A e^{-t/2},\quad |w_j'(t)|\le B e^{t/2},\quad
 A\ll\lambda\sqrt{\delta_j},\quad
 B\ll\lambda^5\sqrt{\delta_j}.
\]
For every fixed integer $m\ge0$ and $0<s<1/2$, a translation-difference
estimate gives
$\|t^mw_j\|_{H^s}^2\ll_{j,m,s}A^{2-2s}B^{2s}$.
Here are the details including the jump at zero. For a translation
$0<r<A/B$, split the positive half-line at
$R=\log(A/(Br))$. On $[0,R]$ integrate the derivative bound; on
$[R,\infty)$ use the amplitude bound. The squared difference,
including the interval crossing zero, is at most
$C_mABr(1+|\log(A/(Br))|)^{2m}$.
For $r\ge A/B$ it is at most $C_mA^2$.
Integration against $r^{-1-2s}\,dr$ gives the claim, since
$\int_0^1r^{-2s}(1+|\log r|)^{2m}\,dr<\infty$.
In particular $\|t^mw_j\|_{H^{1/4}}\ll\lambda^2\sqrt{\delta_j}$.

The exact $L^2$ norm improves the size bound. Split at $4\log\lambda$;
use the exact norm before the split and the amplitude bound afterwards.
This gives
$\|t^mw_j\|_2\ll\sqrt{\delta_j}(1+\log\lambda)^m$.
Fourier H\"older interpolation therefore proves, for $0<s<1/4$,
\begin{equation}
 \|t^mw_j\|_{H^s}
 \ll_{j,m,s}\sqrt{\delta_j}\lambda^{8s}
                     (1+\log\lambda)^{m(1-4s)}.
 \label{eq:prolate-log-tail}
\end{equation}
The two coefficients of $\psi_\lambda$ are bounded, and
$\delta_0\le\delta_4$. Thus the same estimate, with
$\sqrt{\delta_\lambda}$, applies to the logarithmic tail $w$ of
$\mathcal F\psi_\lambda$ and its first moment.

For clarity we describe the theta operation on this tail, including
the change of variables at the lower endpoint. For a raw function on
the positive half-line put
$\mathcal Uf(t)=e^{t/2}f(e^t)$. This is an isometry from
$L^2(0,\infty;du)$ to $L^2(\mathbb R;dt)$. In logarithmic coordinates
let $\mathcal E$ be the closed Fourier multiplier
$\zeta(1/2-i\tau)$. For compactly supported functions of bounded
variation, $(\mathcal E\mathcal Uf)(\log u)$ equals
\[
 Ef(u)=\sqrt u\sum_{n\ge1}f(nu)-u^{-1/2}\int_0^\infty f(v)\,dv.
\]
This follows by integrating from $u=\epsilon$ and using
$\sum_{n\le Y}n^{-s}=Y^{1-s}/(1-s)+\zeta(s)+O_s(Y^{-\Re s})$;
the first term cancels the subtracted integral, and the remaining
error tends to zero for $0<\Re s<1$.
The bound of Chandee and Soundararajan
\cite[Theorem~1.1, p.~1]{chandee-soundararajan2009} gives, under the
Riemann hypothesis,
$|\zeta(1/2+i\tau)|+|\zeta'(1/2+i\tau)|
\ll_b(1+|\tau|)^b$ for every $b>0$.
Here the derivative is obtained as follows. Their bound on the critical
line is smaller than every positive power of $1+|\tau|$. On the line
$\Re s=2$, the Dirichlet series is bounded. The three-lines theorem,
applied after cancelling the pole with $(s-1)/(s+1)$, gives the same
subpower bound throughout $1/2\le\Re s\le2$. The functional equation
then bounds $\zeta$ in $1/2-d\le\Re s\le1/2$ by an additional factor
$O_d((1+|\tau|)^d)$. Choose a fixed $d<\min(b/2,1/4)$, use exponent $b/2$
in the preceding subpower bounds, and apply Cauchy's formula on circles
of radius $d/2$ about $1/2+i\tau$. The bounded range of $\tau$ is
handled by compactness away from the pole. This proves the displayed
derivative estimate. Consequently, for arbitrarily small positive $s',b$,
\[
 \|\mathcal Ew\|_{H^{s'}}+\|t\mathcal Ew\|_{H^{s'}}
 \ll_{s',b}\|w\|_{H^{s'+b}}+\|tw\|_{H^{s'+b}}.
\]
The second term follows by differentiating the Fourier multiplier.
Write $\varphi=\mathcal F\psi_\lambda$, $g=P_\lambda\varphi$ and
$\eta=(1-P_\lambda)\varphi$. If
$\mathcal Mf(s)=\int_0^\infty f(u)u^{s-1}\,du$, the cosine Mellin
identity gives
\[
 \mathcal M\varphi(s)
 =2(2\pi)^{-s}\Gamma(s)\cos(\pi s/2)
                         \mathcal M\psi_\lambda(1-s).
\]
It is justified by inserting an Abel factor in the cosine integral
and then removing it: $\psi_\lambda$ is bounded and compactly
supported, and its bounded-variation Fourier transform is $O(1/u)$,
so the final integrals converge for $0<\Re s<1$.
The zeta functional equation says that the displayed gamma factor
times $\zeta(s)$ equals $\zeta(1-s)$. On $s=1/2-i\tau$, the theta
multiplier identity therefore gives, by Plancherel,
\[
 E\psi_\lambda(e^x)=E\varphi(e^{-x}).
\]
Here $E\varphi$ means the closed multiplier applied to
$\mathcal U\varphi$; no pointwise sum over the noncompact Fourier
tail is required. Both sides belong to $L^2(dx)$: the compact raw
terms have logarithmic Sobolev regularity of every order below
$1/2$, and the tail has the regularity already proved above.

Because $\int_0^\infty\psi_\lambda=0$, the full theta function
$F_\lambda(x):=E\psi_\lambda(e^x)$ equals the uncorrected sum in
the definition of $k_\lambda$; it vanishes for $x>L$.
For $y>\lambda$, the finite sum defining $Eg(y)$ is zero, and hence
\[
 E\varphi(y)=E\eta(y)-c_Qy^{-1/2},\qquad
 c_Q=\int_0^\infty g(u)\,du
       =-\tfrac38\mu_0\mu_4(\mu_0-\mu_4).
\]
For the last equality, on the time interval
$g=-\tfrac34\mu_0\mu_4(\phi_{0,\lambda}/d_0-
\phi_{4,\lambda}/d_4)$; the two integrals are $\mu_0/2$ and
$\mu_4/2$. Thus $|c_Q|\ll\delta_\lambda$.

The logarithmic tail previously denoted by $w$ is exactly
$\mathcal U\eta_\lambda$, where
$\eta_\lambda(u)=\sqrt\lambda\,\eta(\lambda u)$ on $u>0$.
The theta operation satisfies the exact dilation identity
$E\eta_\lambda(y)=E\eta(\lambda y)$, first on compact raw inputs
by substitution in the sum and integral, and then on its closed
multiplier domain. For the lower tail set $x=-L-t$, so that
$t>0$ and $e^{-x}=\lambda e^t$. The preceding identities give
\begin{equation}
 R(t):={\bf1}_{t>0}F_\lambda(-L-t)
 = {\bf1}_{t>0}\left(\mathcal Ew(t)
                    -c_Q\lambda^{-1/2}e^{-t/2}\right).
 \label{eq:prolate-lower-tail}
\end{equation}

For $0<s<s'<1/2$, sharp restriction obeys
$\|{\bf1}_{t>0}f\|_{H^s}\ll_{s,s'}\|f\|_{H^{s'}}$.
The extra Gagliardo term is a constant times
$\int_0^\infty|f(t)|^2t^{-2s}\,dt$; on $(0,1)$ apply Sobolev
embedding into $L^{2/(1-2s')}$ and H\"older, and elsewhere use $L^2$.
Combining this with \eqref{eq:prolate-log-tail}, with all indices
sufficiently small, proves
\[
 \|R\|_{H^s}+\|tR\|_{H^s}
       \ll_\sigma\sqrt{\delta_\lambda}\lambda^{8\sigma}
                          (1+\log\lambda).
\]

Finally the local zero count, with multiplicities, is
$\sum_{\gamma\in[j,j+1]}m_\gamma\ll\log(2+|j|)$.
For $f,tf\in H^s$ the one-dimensional $H^1$ pointwise bound for
$\widehat f$ on each unit interval and Plancherel give
\[
 \sum_\gamma m_\gamma|\widehat f(\gamma)|^2
      \ll_s\|f\|_{H^s}^2+\|tf\|_{H^s}^2.
\]
Indeed bound the supremum of $|\widehat f|^2$ by the integrals of
$|\widehat f|^2+|\widehat f'|^2$ on a slightly larger interval,
multiply by its zero count, and use
$\log(2+|\tau|)\ll_s(1+\tau^2)^s$.
The Fourier transform of $F_\lambda$ has factor
$\zeta(1/2-i\tau)$ and hence vanishes at every ordinate. Its upper
tail is zero. The bounds for $R,tR$ make $R$ integrable by
Cauchy--Schwarz, so $F_\lambda$ is integrable as well. Both sides of
the Fourier multiplier identity are continuous on the real line;
their $L^2$ equality therefore holds at each ordinate.
The exact substitution $x=-L-t$ in its lower tail gives
\[
 \widehat{k_\lambda}(\gamma)
       =-e^{i\gamma L}\widehat R(-\gamma).
\]
Reflection preserves the signed sampling sum, so
$q(k_\lambda)=\sum_\gamma m_\gamma|\widehat R(\gamma)|^2$.
This yields
\[
 q(k_\lambda)\ll_\sigma
       \lambda^{16\sigma}(1+\log\lambda)^2\delta_\lambda.
\]
Choose $16\sigma<\varepsilon$ and absorb the logarithm.
Reflection is a $q$-isometry, so taking the even part cannot increase
$q$, and its norm tends to $\|K\|_2>0$ by
\eqref{eq:prolate-trial-limit}. This proves the two Rayleigh bounds.

To state their exponential scale, Fuchs's fixed-index theorem
\cite[Theorem~1, p.~319]{fuchs1964} has parameter $c=2\pi\lambda^2$
and concentration eigenvalue $\nu_j$:
$1-\nu_j\sim(4\sqrt\pi\,8^j/j!)c^{j+1/2}e^{-2c}$.
Substitution of $j=4$ gives the constant in the proposition.
Min--max and the nonzero trial norm prove the asserted bound for $a_1$.
\end{proof}

\begin{proposition}[An odd prolate competitor]
\label{prop:odd-prolate-upper}
For the prolate modes of indices $2,6$, whose finite Fourier
eigenvalues are negative, define
\[
 \psi_\lambda^{\rm o}
   =\frac{15}{8}\left(\mu_6\frac{\phi_{2,\lambda}}{d_2}
                  -\mu_2\frac{\phi_{6,\lambda}}{d_6}\right),
 \qquad
 k_\lambda^{\rm o}(x)={\bf1}_{[-L,L]}(x)e^{x/2}
                   \sum_{n\ge1}\psi_\lambda^{\rm o}(ne^x).
\]
Its odd part has nonzero norm for sufficiently large $\lambda$.
Writing $v_L^{\rm o}$ for that part divided by its $L^2$ norm, one has
\[
 \|k_\lambda^{\rm o}-K'\|_2=O(\lambda^{-1/2}),\qquad
 q(v_L^{\rm o})\ll_\varepsilon\lambda^\varepsilon(1-\mu_6^2).
\]
In particular,
\begin{equation}
 a_2(\log\lambda)\ll_\varepsilon
                  \lambda^{13+\varepsilon}e^{-4\pi\lambda^2}
 \quad(\varepsilon>0).
 \label{eq:second-eigenvalue-upper}
\end{equation}
\end{proposition}

\begin{proof}
The fixed-index Hermite limit in \cite[equation~(7.10), p.~30]{ccm2025}
applies to indices $2,6$ and gives the ratios
\[
 \frac{h_2(u)}{h_2(0)}=(1-4\pi u^2)e^{-\pi u^2},\qquad
 \frac{h_6(u)}{h_6(0)}=
 \left(1-12\pi u^2+16\pi^2u^4-\frac{64}{15}\pi^3u^6\right)
 e^{-\pi u^2}.
\]
Since $\mu_2,\mu_6\to-1$, their displayed combination tends to
\[
 \frac{15}{8}\left(\frac{h_6}{h_6(0)}-\frac{h_2}{h_2(0)}\right)
  =(-15\pi u^2+30\pi^2u^4-8\pi^3u^6)e^{-\pi u^2}
  =\left(u\frac{d}{du}+\frac12\right)h(u).
\]
Differentiating the absolutely convergent theta series for $K$ shows
that the theta image of this last function is exactly $K'$.
The same uniform Hermite comparison and summation used in
\eqref{eq:prolate-trial-limit} give the asserted $L^2$ error.
The function $K'$ is odd and nonzero, because $K$ is even, nonconstant,
and tends to zero at infinity.  This proves the norm assertion.

The mean of $\psi_\lambda^{\rm o}$ is exactly zero by
$\int_0^\infty\phi_j=\mu_jd_j/2$.  Every exterior estimate in the
preceding proof was established for an arbitrary fixed even index.
It therefore applies to the two bounded coefficients of the modes
$2,6$, with leakage $\delta^{\rm o}=1-\mu_6^2$.
There is one scalar term in the Poisson tail to check: for
$g=P_\lambda\mathcal F\psi_\lambda^{\rm o}$,
\[
 \int_0^\infty g(u)\,du
   =\frac{15}{16}\mu_2\mu_6(\mu_2-\mu_6)
   =O(\delta^{\rm o}).
\]
Both eigenvalues approach $-1$, and
$|\mu_j+1|=(1-\mu_j^2)/(1-\mu_j)$, which proves this bound with
the sign and denominator explicit.  The exact reflected-tail
identity, Sobolev estimate, zeta multiplier, and sampling argument
from the preceding proof now give
$q(k_\lambda^{\rm o})\ll_\varepsilon
\lambda^\varepsilon\delta^{\rm o}$.  Taking the odd part decreases
the form and its normalization stays bounded away from zero.

The even trial and this odd trial are orthogonal both in $L^2$ and
for $q$: pair the two signed ordinates in the sampling formula for
the latter assertion.  The largest Rayleigh quotient on their
two-dimensional span is consequently the larger of their two
quotients.  Min--max and the fixed-index Fuchs formula give
\[
 a_2(\log\lambda)\ll_\varepsilon
       \lambda^\varepsilon(1-\mu_6^2),\qquad
 1-\mu_6^2\sim\frac{2^{22}\sqrt2\pi^7}{45}
                          \lambda^{13}e^{-4\pi\lambda^2}.
\]
This proves \eqref{eq:second-eigenvalue-upper}.
\end{proof}

The baseline power in the Rayleigh criterion is therefore $9$.
Precisely, the change of parameter between an absolute power $\beta$
and a gain $\alpha$ above the trial scale is $\beta=9+\alpha$:
\[
 \lambda^\beta e^{-4\pi\lambda^2}
  =\lambda^\alpha\bigl(\lambda^9e^{-4\pi\lambda^2}\bigr).
\]
Since $(1-\mu_4^2)/(\lambda^9e^{-4\pi\lambda^2})$ tends to a
positive constant, positive-power improvement over that leakage is
exactly a positive gain $\alpha$.  Equation
\eqref{eq:second-eigenvalue-upper} rules out every eventual lower
bound $a_2(\log\lambda)\ge c\lambda^\beta e^{-4\pi\lambda^2}$
with $\beta>13$: choose $0<\varepsilon<\beta-13$ and divide the two
bounds.  Thus the positive-gain range still open by this comparison
is $9<\beta\le13$, equivalently $0<\alpha\le4$.

\section{Theta lifts and one-mode coercivity}
\label{sec:theta-coercivity}

\subsection{The theta image and its polynomial hierarchy}

Use the Fourier convention $\widehat f(t)=\int_{\mathbb R}f(x)e^{-itx}\,dx$. Put

\[
 h(u)=(4\pi^2u^4-6\pi u^2)e^{-\pi u^2},\qquad
 (\Theta\psi)(x)=e^{x/2}\sum_{n\ge1}\psi(ne^x),\qquad K=\Theta h.
\]

The integral defining the Mellin transform of $h$ gives, exactly,

\begin{equation}
M_h(s):=\int_0^\infty h(u)u^{s-1}\,du
   =\frac12s(s-1)\pi^{-s/2}\Gamma(s/2).
 \label{eq:theta-hier-1}
\end{equation}

Indeed the two Gaussian moments are
$2\pi^{-s/2}\Gamma(s/2+2)$ and
$-3\pi^{-s/2}\Gamma(s/2+1)$; the Gamma recurrence gives \eqref{eq:theta-hier-1}.
For $\operatorname{Re}s>1$, the Mellin sum is absolutely convergent and gives
$\zeta(s)M_h(s)=\xi(s)$. Its Fourier restriction at
$s=1/2-it$, continued by the functional equation, is
$\widehat K(t)=\Xi(t)$. The exact Poisson relation for the self-Fourier function $h$, with $h(0)=\int h=0$, makes $K$ even.

Let $E=u\partial_u+1/2$. Termwise differentiation, justified by Gaussian decay at infinity, gives the exact intertwining

\begin{equation}
\partial_x\Theta\psi=\Theta E\psi.
 \label{eq:theta-hier-2}
\end{equation}

Thus $K^{(j)}=\Theta E^jh$ and
$\widehat{K^{(j)}}(t)=(it)^j\Xi(t)$. These profiles all have zero sampling vector under the actual map

\begin{equation}
\mathcal S f=(\sqrt{m_\gamma}\,\widehat f(\gamma))_{\gamma\in\Gamma}.
 \label{eq:theta-hier-3}
\end{equation}

Here the map is defined on rapidly decreasing smooth profiles for which the displayed sum converges. Its squared norm is the extended Weil seminorm. Compact cutoffs converge in that seminorm by the sampling bound proved below; consequently this map agrees with the completion of compact test functions. In particular, these nonzero physical profiles have image zero under \eqref{eq:theta-hier-3}. 

As $x\to+\infty$, with $u=e^x$, the first theta summand dominates, while the sum of all other terms and each fixed derivative is exponentially smaller. Hence

\begin{equation}
K^{(j)}(x)=4\pi^2(-2\pi)^j
       u^{9/2+2j}e^{-\pi u^2}(1+O_j(u^{-2})).
 \label{eq:theta-hier-4}
\end{equation}

The relative $u^{-2}$ error comes from the lower polynomial in $h$ and differentiation; terms $n\ge2$ have the additional exponential factor $e^{-3\pi u^2}$ times a fixed polynomial. Integration after the change of variable $u=e^x$ gives

\begin{equation}
\int_L^\infty|K^{(j)}(x)|^2\,dx
  \sim4\pi^3(2\pi)^{2j}
         \lambda^{7+4j}e^{-2\pi\lambda^2},\qquad\lambda=e^L.
 \label{eq:theta-hier-5}
\end{equation}

This follows from
$\int_\lambda^\infty u^a e^{-2\pi u^2}du
\sim\lambda^{a-1}e^{-2\pi\lambda^2}/(4\pi)$.

\subsection{Division at an actual zero: the exact first-order lift}

Fix an actual positive ordinate $\gamma$, of its actual multiplicity $m_\gamma$, and put $s_\gamma=1/2-i\gamma$. Define

\begin{equation}
(R_\gamma K)(x)
   =-i e^{i\gamma x}\int_x^\infty e^{-i\gamma y}K(y)\,dy.
 \label{eq:theta-hier-6}
\end{equation}

Differentiation gives
$(-i\partial_x-\gamma)R_\gamma K=K$.
The identity $\widehat K(\gamma)=0$ also gives the equivalent left-tail formula
$R_\gamma K=i e^{i\gamma x}\int_{-\infty}^x e^{-i\gamma y}K(y)dy$.
Both ends therefore decay faster than every exponential, and Fourier transformation yields

\begin{equation}
\widehat{R_\gamma K}(t)=\frac{\Xi(t)}{t-\gamma},
 \label{eq:theta-hier-7}
\end{equation}

with the entire extension at $\gamma$.

The corresponding raw theta input can be written directly:

\begin{equation}
\psi_\gamma(u)
  =-i u^{-s_\gamma}\int_u^\infty h(v)v^{s_\gamma-1}\,dv.
 \label{eq:theta-hier-8}
\end{equation}

It satisfies

\begin{equation}
(-iE-\gamma)\psi_\gamma=h,
 \qquad \Theta\psi_\gamma=R_\gamma K.
 \label{eq:theta-hier-9}
\end{equation}

The first identity follows by differentiation. The theta series and its derivative converge locally normally because the upper tail of \eqref{eq:theta-hier-8} is Gaussian. Equation \eqref{eq:theta-hier-2} shows that its theta image satisfies the same first-order equation as \eqref{eq:theta-hier-6}, and both decay at positive infinity. Their difference is a multiple of $e^{i\gamma x}$; positive-tail decay forces that multiple to be zero. This proves the second identity.

The two raw endpoint expansions are

\begin{equation}
\psi_\gamma(u)=-iM_h(s_\gamma)u^{-s_\gamma}+O_\gamma(u^2)
       \quad(u\downarrow0),
 \qquad
 \psi_\gamma(u)\sim\frac{-i}{2\pi u^2}h(u)
       \quad(u\to\infty).
 \label{eq:theta-hier-10}
\end{equation}

For the first, split the integral in \eqref{eq:theta-hier-8} at zero and use $h(v)=O(v^2)$. Its leading coefficient is nonzero by \eqref{eq:theta-hier-1}: $s_\gamma$ is neither $0$ nor $1$, and the Gamma function has no zeros. For the second, integrate the Gaussian tail once, or apply the elementary Laplace-tail asymptotic to each polynomial term. Thus the raw lift has a nonzero $u^{-1/2+i\gamma}$ term at zero and its squared modulus is not integrable there. The theta image nevertheless has the two rapidly decreasing tails proved directly by \eqref{eq:theta-hier-6}.

The exact Mellin relation behind this cancellation is

\begin{equation}
M_{\psi_\gamma}(s)=\frac{M_h(s)}{i(s-s_\gamma)},\qquad
 \zeta(s)M_{\psi_\gamma}(s)=\frac{\xi(s)}{i(s-s_\gamma)}.
 \label{eq:theta-hier-11}
\end{equation}

The first holds initially in the strip $\operatorname{Re}s>1/2$, where the endpoint integrations converge. The second holds initially for $\operatorname{Re}s>1$ by the absolutely convergent theta sum. The zero of $\zeta$ at $s_\gamma$ cancels the displayed raw Mellin pole. Equations \eqref{eq:theta-hier-6}--\eqref{eq:theta-hier-11} prove the relation directly, including the domain and the coefficient of the singular term.

\subsection{Repeated and paired divisors}

For $1\le r\le m_\gamma$, repeated application of \eqref{eq:theta-hier-6} gives

\begin{equation}
F_{\gamma,r}(x)=\frac{(-i)^r e^{i\gamma x}}{(r-1)!}
   \int_x^\infty (y-x)^{r-1}e^{-i\gamma y}K(y)\,dy,
 \qquad
 \widehat F_{\gamma,r}(t)=\frac{\Xi(t)}{(t-\gamma)^r}.
 \label{eq:theta-hier-12}
\end{equation}

The moment identities $\Xi^{(j)}(\gamma)=0$, $0\le j<r$, convert \eqref{eq:theta-hier-12} to the corresponding left-tail integral, proving decay at both ends. Its actual sampling norm is

\begin{equation}
q(F_{\gamma,r})=
 \begin{cases}
 0,&r<m_\gamma,\\
 m_\gamma\left|\Xi^{(m_\gamma)}(\gamma)/m_\gamma!\right|^2,
       &r=m_\gamma.
 \end{cases}
 \label{eq:theta-hier-13}
\end{equation}

This is obtained simply by evaluating \eqref{eq:theta-hier-12}'s entire transform at every signed ordinate: all values vanish except the final possible value at $\gamma$.

For the even real profiles define

\begin{equation}
\widehat G_{\gamma,r}(t)=\frac{\Xi(t)}{(t^2-\gamma^2)^r},
 \qquad 1\le r\le m_\gamma.
 \label{eq:theta-hier-14}
\end{equation}

The transform is real and even on the real axis, so its inverse is real and even. Applying $R_\gamma^rR_{-\gamma}^r$ gives this profile, with the same two-tail argument. At the two removed zeros the exact values are equal, so

\begin{equation}
q(G_{\gamma,r})=
 \begin{cases}
 0,&r<m_\gamma,\\
 2m_\gamma\left|
 \dfrac{\Xi^{(m_\gamma)}(\gamma)}{m_\gamma!(2\gamma)^{m_\gamma}}
                  \right|^2,&r=m_\gamma.
 \end{cases}
 \label{eq:theta-hier-15}
\end{equation}

For each application of $R_\beta$, the positive-tail Laplace integral multiplies its input's leading term by $-i/(2\pi u^2)$. Therefore

\begin{equation}
F_{\gamma,r}(x)\sim
   4\pi^2\left(\frac{-i}{2\pi}\right)^r
       u^{9/2-2r}e^{-\pi u^2},
 \qquad
 G_{\gamma,r}(x)\sim
   4\pi^2\frac{(-1)^r}{(2\pi)^{2r}}
       u^{9/2-4r}e^{-\pi u^2}.
 \label{eq:theta-hier-16}
\end{equation}

In particular $G_{\gamma,1}(x)\sim-e^{x/2}e^{-\pi e^{2x}}$, while $K(x)\sim4\pi^2e^{9x/2}e^{-\pi e^{2x}}$. The ratio is exactly asymptotic to $-e^{-4x}/(4\pi^2)$. Correspondingly,

\begin{equation}
\begin{aligned}
 \int_L^\infty |F_{\gamma,r}|^2
   &\sim4\pi^3(2\pi)^{-2r}\lambda^{7-4r}e^{-2\pi\lambda^2},\\
 \int_L^\infty |G_{\gamma,r}|^2
   &\sim4\pi^3(2\pi)^{-4r}\lambda^{7-8r}e^{-2\pi\lambda^2}.
\end{aligned}
 \label{eq:theta-hier-17}
\end{equation}

The first paired raw lift is explicitly
$\psi_{\gamma,\mathrm{pair}}=(\psi_\gamma-\psi_{-\gamma})/(2\gamma)$. Its Mellin transform is

\begin{equation}
\frac{M_h(s)}{-(s-1/2)^2-\gamma^2},
 \label{eq:theta-hier-18}
\end{equation}

and its theta image is $G_{\gamma,1}$. The nonzero endpoint terms at zero have exponents $-1/2\pm i\gamma$. Their squared integral diverges: after $u=e^{-v}$, the diagonal terms have a positive constant mean, while the cross term has a bounded oscillatory integral. For a repeated paired divisor, the raw Mellin poles have order $r$, giving nonzero polynomial factors in $\log u$ of degree $r-1$ at those two exponents. The same conclusion follows because the diagonal integral grows like $V^{2r-1}$, while integration by parts bounds the oscillatory cross term by $O(V^{2r-2})$. Thus the precise theta lifts carry explicitly quantified endpoint singularities; applying the square-integrable prolate estimates to them requires an additional truncation estimate for those terms.

\subsection{What direct compact truncation proves}

The local zero-count bound, including multiplicities, is
$\#\{\gamma:n\le\gamma<n+1\}\ll\log(2+|n|)$. Here is a direct justification of the bound used. For large positive $n$, apply Jensen's formula to $\zeta(s)$ on the disk of radius $3$ centered at $2+i(n+1/2)$; the disk of radius $5/2$ contains the whole zero strip with ordinates in $[n,n+1]$. The value at the center has modulus at least $1/\zeta(2)$, by the absolutely convergent reciprocal Dirichlet series. Euler--Maclaurin summation at an integer of size $n$, retaining the Bernoulli terms through order $4$, bounds the modulus throughout the outer disk by $O(n^2)$: on $-1\le\operatorname{Re}s\le5$, the finite sum is $O(n^2)$, and the remainder is bounded by a constant times $|s|^4 n^{-\operatorname{Re}s-3}=O(n^2)$. Jensen's formula gives the claimed logarithmic count. Conjugation treats negative ordinates, and the finitely many remaining intervals are absorbed in the constant. Unit-interval Sobolev evaluation therefore gives

\begin{equation}
q(f)\ll\int_{\mathbb R}\log(2+|t|)
      (|\widehat f(t)|^2+|\widehat f'(t)|^2)\,dt.
 \label{eq:theta-hier-19}
\end{equation}

For $\kappa\ge1$, use
$\log(2+|t|)\le\log(2+\kappa)+t^2/\kappa^2$. Parseval gives the concrete physical bound

\begin{equation}
q(f)\ll \log(2+\kappa)(\|f\|_2^2+\|xf\|_2^2)
     +\kappa^{-2}(\|f'\|_2^2+\|(xf)'\|_2^2).
 \label{eq:theta-hier-20}
\end{equation}

Choose a fixed smooth transition and an even cutoff $\chi_L$ supported in $[-L,L]$, equal to one on $[-L+\lambda^{-2},L-\lambda^{-2}]$, with $|\chi_L'|\ll\lambda^2$. Apply \eqref{eq:theta-hier-20}, with $\kappa=\lambda^2$, to $(1-\chi_L)f$. Equations \eqref{eq:theta-hier-4}, \eqref{eq:theta-hier-16}, and the corresponding differentiated estimates show that a profile with leading positive-tail exponent $u^a e^{-\pi u^2}$ obeys

\begin{equation}
q((1-\chi_L)f)\ll_f
      (\log\lambda)^3\lambda^{2a-2}e^{-2\pi\lambda^2}.
 \label{eq:theta-hier-21}
\end{equation}

The transition starts at $L-\lambda^{-2}$; replacing $\lambda$ by $\lambda e^{-\lambda^{-2}}$ changes the exponential by a bounded factor. The derivative of the cutoff has exactly the $\lambda^2$ size absorbed by $\kappa^{-2}$, and the factor $x$ contributes at most $O(L)$ in these tail integrals. This proves \eqref{eq:theta-hier-21}, including the transition interval.

Consequently,

\begin{equation}
q(\chi_L K^{(j)})\ll_j
      (\log\lambda)^3\lambda^{7+4j}e^{-2\pi\lambda^2}.
 \label{eq:theta-hier-22}
\end{equation}

For every profile \eqref{eq:theta-hier-12} or \eqref{eq:theta-hier-14}, the seminorm triangle inequality gives
$|\sqrt{q(\chi_L f)}-\sqrt{q(f)}|\le\sqrt{q((1-\chi_L)f)}$; substitute the exact values \eqref{eq:theta-hier-13}, \eqref{eq:theta-hier-15} and exponents from \eqref{eq:theta-hier-16}. In particular the zero entries in \eqref{eq:theta-hier-15} yield the upper scale
$(\log\lambda)^3\lambda^{7-8r}e^{-2\pi\lambda^2}$ for their compact approximants. These identities verify the tail hierarchy. The proved compact-cutoff bound has exponential $e^{-2\pi\lambda^2}$; the sharper original prolate construction has $e^{-4\pi\lambda^2}$. The exact maps above relate the profiles, but \eqref{eq:theta-hier-21} does not give a prolate-scale lower bound.

\subsection{Exact removal of the stated trial mode}

Let $v_L$ be the actual normalized even prolate trial, and define

\begin{equation}
b_L=\inf\{q(f):f\in\operatorname{Dom}q_L,
       \|f\|_2=1,\ \langle f,v_L\rangle_2=0\}.
 \label{eq:theta-hier-23}
\end{equation}

The exact comparison is

\begin{equation}
a_1(L)+a_2(L)-q(v_L)\le b_L\le a_2(L).
 \label{eq:theta-hier-24}
\end{equation}

For the lower inequality, every admissible $(v_L,f)$ is an orthonormal pair. Ky Fan's trace inequality gives $q(v_L)+q(f)\ge a_1+a_2$. For the upper inequality, the span of the first two eigenvectors meets $v_L^\perp$ nontrivially, and the Rayleigh quotient on that span is at most $a_2$. In particular $a_2-q(v_L)\le b_L\le a_2$, since $a_1\ge0$. This explicitly relates removal of the stated theta/prolate mode to the desired denominator; it needs no separate assertion that the trial aligns with the first eigenvector.

The remaining lower estimate concerns the actual restriction
\eqref{eq:theta-hier-23}:
\[
 b_L\ge c\lambda^{9+\alpha}e^{-4\pi\lambda^2},
 \qquad L=\log\lambda,
\]
with fixed $c>0$ and $\alpha>0$ along an unbounded sequence of
windows.  This estimate remains open.  Equation~\eqref{eq:theta-hier-24}
places $b_L$ within the already controlled trial energy of $a_2(L)$.
Equations~\eqref{eq:theta-hier-8} and \eqref{eq:theta-hier-18} identify
the raw endpoint terms that an argument using these theta lifts must
control at the displayed scale.

\subsection{A finite lower certificate for the second Weil eigenvalue}
\label{sec:finite-weil-lower}

The following bound isolates the remaining arithmetic inequality in a
finite matrix, with an explicit error for the discarded directions.
The geometric explicit formula on $[-L,L]$ is
\begin{equation}
 q(f)=2\operatorname{Re}\bigl(F(i/2)\overline{F(-i/2)}\bigr)
       +\frac1{2\pi}\int_{\mathbb R}\Psi_L(t)|F(t)|^2\,dt,
 \label{eq:weil-symbol-lower}
\end{equation}
where $F(t)=\int_{-L}^Lf(x)e^{-itx}\,dx$ and
\[
 \Psi_L(t)=\operatorname{Re}\psi\left(\frac14+\frac{it}{2}\right)
       -\log\pi-
       \sum_{\log n<2L}\frac{2\Lambda(n)}{\sqrt n}\cos(t\log n).
\]
Here $\psi=\Gamma'/\Gamma$ and $\Lambda$ is the von Mangoldt
function.  The strict prime-power cutoff is valid because the
autocorrelation vanishes at its support endpoints.  This representation
is also the starting point of Zhu's finite reduction
\cite[\S\S2--4, pp.~7--12]{zhu2026}; the estimate below uses a direct
polynomial projection error.

Put
\[
 \mathcal A_L=\sum_{\log n<2L}\frac{2\Lambda(n)}{\sqrt n},\qquad
 T=2\pi\exp(\mathcal A_L+b+1),\qquad b>0.
\]
On $L^2([-L,L])$ write $p_\pm(x)=e^{\pm x/2}$,
$g_t(x)=e^{itx}$, and $u\otimes v:f\mapsto\langle f,v\rangle u$.
Define the compact selfadjoint operator
\[
 B_T=p_+\otimes p_-+p_-\otimes p_+
       +\frac1{2\pi}\int_{-T}^{T}
                     (\Psi_L(t)-b)g_t\otimes g_t\,dt.
\]
Let $P_N$ be the $L^2$ orthogonal projection onto polynomials of degree
at most $N$, expressed for example in the normalized Legendre basis,
and let $M_{L,T,N}=bI+P_NB_TP_N$ act on that finite-dimensional space.
For $N\ge1$, set
\[
 \begin{split}
 W_{L,T,b}&=\frac1{2\pi}\int_{-T}^{T}|\Psi_L(t)-b|\,dt,\\
 R_N&=e^{TL}\frac{(TL)^{N+1}}{(N+1)!},\\
 \varepsilon_{L,T,N}
   &=4LR_NW_{L,T,b}
        +8Le^L\frac{(L/2)^{N+1}}{(N+1)!}.
 \end{split}
\]

\begin{proposition}[A finite second-eigenvalue lower bound]
\label{prop:finite-weil-lower}
The second eigenvalue of the windowed Weil form satisfies
\begin{equation}
 a_2(L)\ge
       \min\{b,\lambda_2(M_{L,T,N})\}
           -\varepsilon_{L,T,N}.
 \label{eq:finite-weil-lower}
\end{equation}
The entries of $M_{L,T,N}$ and an upper bound for its error term use
only the displayed finite prime-power sum and integrals over finite
intervals.
\end{proposition}

\begin{proof}
We first justify the frequency cutoff, including its numerical
constant.  The elementary Euler summation formula applied to
$(x+z)^{-1}$ gives, for $\operatorname{Re}z>0$,
\[
 \psi(z)=\log z-\frac1{2z}
        +\int_0^\infty\frac{\{x\}-1/2}{(x+z)^2}\,dx.
\]
For example, apply integration by parts on each unit interval, sum up
to an integer $M$, and take the limit in
$\psi(z)=\lim_{M\to\infty}(\log M-\sum_{j=0}^{M-1}(j+z)^{-1})$.
At $z=1/4+it/2$, $t\ne0$, the integral in absolute value is at most
\[
 \frac12\int_0^\infty
     \frac{dx}{(x+1/4)^2+t^2/4}
       \le\frac{\pi}{2|t|}.
\]
Together with $|1/(2z)|\le1/|t|$ and
$\log|z|\ge\log(|t|/2)$, this proves
\[
 \Psi_L(t)\ge\log\frac{|t|}{2\pi}
                     -\frac3{|t|}-\mathcal A_L.
\]
The right side increases with $|t|$.  At our choice of $T$ it is
$b+1-3/T>b$, since $T\ge2\pi e$.  Thus
$\Psi_L(t)\ge b$ outside $[-T,T]$.

Parseval's identity in \eqref{eq:weil-symbol-lower} now gives the
quadratic-form inequality
\[
 q(f)\ge b\|f\|_2^2+\langle B_Tf,f\rangle.
\]
It first holds on the compact smooth core and then on the closed form
domain.  Notice that the two signs of the pole term are both present;
on even functions it is positive and on odd functions it is negative.

The degree-$N$ Taylor polynomial approximates $g_t$ uniformly on
$[-L,L]$ with error at most
$e^{|t|L}(|t|L)^{N+1}/(N+1)!$.  Orthogonal projection is the best
$L^2$ approximation, so
\[
 \|(I-P_N)g_t\|_2\le\sqrt{2L}R_N,\qquad
 \|g_t\|_2=\sqrt{2L}\quad (|t|\le T).
\]
Writing $r=(I-P_N)g_t$ and using
\[
 g_t\otimes g_t-P_Ng_t\otimes P_Ng_t
      =r\otimes g_t+P_Ng_t\otimes r
\]
shows that the operator norm of this difference is at most $4LR_N$.
Its integral contribution is therefore at most $4LR_NW_{L,T,b}$.

The same approximation with $t=\pm i/2$ gives
\[
 \|p_\pm\|_2\le\sqrt{2L}e^{L/2},\qquad
 \|(I-P_N)p_\pm\|_2
       \le\sqrt{2L}e^{L/2}\frac{(L/2)^{N+1}}{(N+1)!}.
\]
Expand each of the two mixed rank-one terms just as above.  Their
total error is at most the second summand of
$\varepsilon_{L,T,N}$.  We have proved
\[
 \|B_T-P_NB_TP_N\|\le\varepsilon_{L,T,N}.
\]
On the orthogonal complement of the polynomial space,
$bI+P_NB_TP_N$ equals $bI$.  Its second min--max value is consequently
$\min\{b,\lambda_2(M_{L,T,N})\}$.  The min--max principle applied to
the form inequality and this operator norm estimate proves
\eqref{eq:finite-weil-lower}.
\end{proof}

There is a stronger bound which keeps the coupling quadratic in its
size.  In the notation above put
\[
 \begin{split}
 a&=\lambda_2(M_{L,T,N}),\qquad
 r=R_N,\qquad s=\frac{(L/2)^{N+1}}{(N+1)!},\\
 c&=b-2Lr^2W_{L,T,b}-4Le^Ls^2,\\
 z&=2LrW_{L,T,b}+4Le^Ls.
 \end{split}
\]

\begin{proposition}[The coupling retained in the lower certificate]
\label{prop:finite-weil-block}
For $N\ge1$,
\begin{equation}
 a_2(L)\ge\frac{a+c-\sqrt{(a-c)^2+4z^2}}2.
 \label{eq:finite-weil-block}
\end{equation}
For a target $u<\min\{a,c\}$, the right side exceeds $u$ precisely
when $(a-u)(c-u)>z^2$.
\end{proposition}

\begin{proof}
Write $Q_N=I-P_N$.  The same bounds on $g_t$ give
\[
 \|Q_N(g_t\otimes g_t)Q_N\|\le2Lr^2,
 \qquad
 \|Q_N(g_t\otimes g_t)P_N\|\le2Lr.
\]
The two pole ranks contribute at most $4Le^Ls^2$ to the first block
and $4Le^Ls$ to the second.  Thus the tail of $bI+B_T$ is bounded
below by $cI$, and its mixed block has norm at most $z$.
Exclude one unit eigenvector belonging to the smallest eigenvalue of
$M_{L,T,N}$.  This defines a codimension-one subspace of the full
Hilbert space, and on it the head block is bounded below by $aI$.
For $f=x+y$ in that subspace, with $x=P_Nf$ and $y=Q_Nf$,
\[
 q(f)\ge a\|x\|_2^2+c\|y\|_2^2-2z\|x\|_2\|y\|_2.
\]
The minimum of this two-variable quadratic form on
$\|x\|_2^2+\|y\|_2^2=1$ is the smaller eigenvalue of
$\left(\begin{smallmatrix}a&-z\\-z&c\end{smallmatrix}\right)$.
The codimension-one min--max characterization proves
\eqref{eq:finite-weil-block}.  Positivity of this matrix after
subtracting $uI$ is exactly the stated pair of positive diagonal
conditions and determinant inequality.  No positivity of $a$ or $c$
is needed for the lower bound itself.
\end{proof}

For computations the Taylor bounds can be replaced by smaller
Legendre tails.  Define, whenever its denominator is positive,
\[
 \rho_N(v)^2=
 \frac{(2N+3)v^{2N+2}}{((2N+3)!!)^2}
 \left(1-\frac{v^2}{(2N+3)(2N+5)}\right)^{-1}.
\]
Both lower bounds remain valid with
$r=\rho_N(TL)$ and $s=\rho_N(L/2)$, using
$\varepsilon=4LrW_{L,T,b}+8Le^Ls$ in
\eqref{eq:finite-weil-lower}.  To prove this, the normalized Legendre
basis $p_n(x)=\sqrt{(2n+1)/(2L)}P_n(x/L)$ has Fourier coefficients
\[
 \langle g_t,p_n\rangle
        =\sqrt{2L(2n+1)}\,i^n j_n(tL).
\]
Rodrigues' formula, integrated by parts $n$ times, gives
\[
 j_n(v)=\frac{v^n}{2^{n+1}n!}
                  \int_{-1}^1e^{ivu}(1-u^2)^n\,du.
\]
The beta integral yields
$|j_n(v)|\le |v|^n/(2n+1)!!$ for real $v$; for imaginary $v$ the
same bound acquires at most the factor $e^{|v|}$.
The successive squared terms
$(2n+1)v^{2n}/((2n+1)!!)^2$ have ratio
$v^2/((2n+1)(2n+3))$.  Summing from $n=N+1$ bounds
$\|(I-P_N)g_t\|_2^2$ by $2L\rho_N(TL)^2$.
For $p_\pm$, the corresponding bound is
$2Le^L\rho_N(L/2)^2$.  These are exactly the two estimates used in
the block proof.

The remaining positive-power comparison can now be specified without
an unspecified approximation error.  In the parameters $L=\log\lambda$
and $\beta=9+\alpha$, its finite arithmetic quantity is
\[
 \min\{b,\lambda_2(M_{L,T,N})\}-\varepsilon_{L,T,N}.
\]
A lower bound of size $\lambda^{9+\alpha}e^{-4\pi\lambda^2}$
with one $\alpha>0$ along an unbounded sequence of windows is the
inequality still needed in this reduction.  The prolate upper estimate
then makes the corresponding Rayleigh ratios tend to zero along that
sequence.  No such matrix inequality for an unbounded sequence is
proved here.  The large value of $T$ records the cost of bounding the
prime sum by its total mass; a sharper estimate must use cancellation
or the way the translations act on functions supported in the window.

\section{Multiplicity from Suzuki's inverse compressions}
\label{sec:suzuki-probes}
Suzuki's norm convergence supplies a second finite procedure for recovering
each fixed multiplicity, through the mass of a spectral projection.
His Theorem~3.3 identifies the local inverse with a compression of $C$;
equation~(4.1) proves convergence in Hilbert--Schmidt norm, and
Remark~4.1 gives resolvent convergence away from the ordinates
\cite[pp.~10--12]{suzuki}. His normalized spectral function
converges to $\Delta_\Gamma$, as Theorem~1.2 states. The following
multiplicity recovery is a consequence of that convergence and the
sampling norm. We credit the convergence mechanism to Suzuki.

Let $V_N\subset H$ be the continuous piecewise-linear functions supported
on $[-N,N]$, with mesh $2^{-N}$, zero endpoint values, and integral zero.
They form nested finite-dimensional spaces. Their union is dense in $H$.
Indeed, the fixed-window form is continuous in the $H^1$ norm, so dyadic
interpolation approximates every compact smooth mean-zero function.
Subtract a fixed unit-integral tent times the integral error to keep that
condition exact. General compact smooth functions are approximated in
$q$ by mean-zero ones: for $h\in C_c^\infty$ with $\int h=1$, the
rescaling $h_R(x)=R^{-1}h(x/R)$ satisfies
\[
 q(h_R)\le R^{-2}\|h'\|_1^2\sum_\gamma m_\gamma/\gamma^2\longrightarrow0,
\]
so $f-(\int f)h_R\to f$. Let $P_N$ be the $q$-orthogonal projection
onto $V_N$. Consequently
\begin{equation}
 C_N=P_NCP_N\longrightarrow C\quad\hbox{in }\mathcal S_2.
 \label{eq:finite-compression-limit}
\end{equation}
This is Suzuki's compact-compression argument applied to the indicated
finite spaces.

For a basis $b_1,\ldots,b_d$, all finite data are Weil pairings:
\[
 G=(q(b_j,b_i))_{i,j},\quad
 H_N=(q(\mathcal Pb_j,b_i))_{i,j},\quad
 c=(q(\phi,b_i))_i.
\]
Here $\mathcal P$ is the compact primitive defined above. The matrix
$S_N=G^{-1/2}H_NG^{-1/2}$ represents $C_N$, and $G^{-1/2}c$ represents
$P_N\phi$. The explicit formula uses only prime powers with logarithm
at most twice the largest support radius. Positivity of $G$ permits
certified computation of its inverse square root, by increasing precision
until a positive lower eigenvalue enclosure is obtained. Each finite
search terminates; no uniform condition number is claimed.

\begin{remark}[Fixed multiplicity recovery from Suzuki's convergence]
\label{rem:suzuki-multiplicity}
Fix a compact probe $\phi$ with $\widehat\phi(t)\ne0$ for real $t$.
For a fixed ordinate $\gamma$, let $E_{N,\gamma}$ be the finite spectral
projection converging to the isolated eigenvalue $1/\gamma$ of $C$, and
let $\theta_{N,\gamma}\to1/\gamma$ be its eigenvalue. Then
\[
 \frac{\|E_{N,\gamma}P_N\phi\|_q^2}
      {|\widehat\phi(1/\theta_{N,\gamma})|^2}
           \longrightarrow m_\gamma.
\]
Nearest-integer rounding therefore recovers $m_\gamma$ for every
sufficiently large $N$, with the threshold allowed to depend on $\gamma$.
\end{remark}

\begin{proof}
In the sampling realization, the unit coordinate vector at $\gamma$ has
value $m_\gamma^{-1/2}$ at that coordinate. Thus the squared mass of
$\phi$ on its eigenspace is
$m_\gamma|\widehat\phi(\gamma)|^2$. Norm convergence of $C_N$ gives
norm convergence of the spectral projection by the resolvent identity
on a small circle about $1/\gamma$. The limiting projection has rank
one; projections within distance less than one have equal finite rank,
so the corresponding finite eigenvalue is eventually simple.
Moreover $E_{N,\gamma}P_N\phi=E_{N,\gamma}\phi$.
The numerator therefore converges to the displayed coordinate mass,
and continuity and nonvanishing of the denominator prove the assertion.
A convergent sequence with integer limit eventually lies less than
$1/2$ from that integer.
\end{proof}

We now quantify the two errors in this division. The probe
$\phi=h*\widetilde h$, where
$h(x)=e^{-x}{\bf1}_{[0,1]}(x)$ and $\widetilde h(x)=h(-x)$, is real,
even, supported on $[-1,1]$, and belongs locally to $H^1_0$.
Writing $r=e^{-1}$, its transform is exactly
\[
 p(t)=\widehat\phi(t)
 =\left|\frac{1-e^{-1-it}}{1+it}\right|^2
 =\frac{1+r^2-2r\cos t}{1+t^2}.
\]
In particular
\begin{equation}
 p(t)\ge\frac{(1-r)^2}{1+t^2},\qquad
 |(\log p)'(t)|\le C_0:=1+\frac{2r}{(1-r)^2}.
 \label{eq:probe-bounds}
\end{equation}
The first inequality follows from $|1-re^{-it}|\ge1-r$. For the
second, differentiate the numerator and denominator; their logarithmic
derivatives have absolute values at most $2r/(1-r)^2$ and
$2|t|/(1+t^2)\le1$, respectively.

\begin{proposition}[A certified spectral rounding bound]
\label{prop:spectral-rounding}
Let $\lambda=1/\gamma$ and
$d=\operatorname{dist}(\lambda,\sigma(C)\setminus\{\lambda\})$,
where zero is included in the spectrum. Let $\widetilde C_N$ be a
selfadjoint finite approximation supported on $V_N$, with
$\|\widetilde C_N-C\|\le\varepsilon<d/4$. Let
$\widetilde a_N$ approximate $P_N\phi$ with error at most $\eta$,
and put $B=\|\phi\|_q$.
The spectral point $\widetilde\lambda$ enclosed by
$|z-\lambda|=d/2$ is simple, and its projection $\widetilde E$ satisfies
\[
 \|\widetilde E-E\|\le
 \delta:=\frac{2\varepsilon}{d-2\varepsilon},\qquad
 |\widetilde\gamma-\gamma|\le
 \xi:=\frac{\varepsilon\gamma^2}{1-\varepsilon|\gamma|},
 \quad\widetilde\gamma=1/\widetilde\lambda.
\]
The computed mass $u=\|\widetilde E\widetilde a_N\|_q^2$ satisfies
\[
 |u-m_\gamma p(\gamma)^2|
       \le\delta B^2+2B\eta+\eta^2=:\mathcal E.
\]
For $|\gamma|\le T$, the resulting multiplicity division has error
\begin{equation}
 \left|\frac{u}{p(\widetilde\gamma)^2}-m_\gamma\right|
 \le
 \frac{(1+(T+\xi)^2)^2}{(1-r)^4}\mathcal E
       +m_\gamma\bigl(e^{2C_0\xi}-1\bigr).
 \label{eq:spectral-rounding-error}
\end{equation}
Adding the final certified arithmetic error gives a rounding certificate
as soon as the sum is strictly less than $1/2$.
For a certificate the unknown $m_\gamma$ on the right may be replaced
by any certified upper bound. One always has the probe bound
$m_\gamma\le B^2(1+T^2)^2/(1-r)^4$; a bound for the total zero count
in a certified packet is another choice.
\end{proposition}

\begin{proof}
On the circle of radius $d/2$, the resolvent of $C$ has norm at most
$2/d$. The Neumann series bounds that of $\widetilde C_N$ by
$1/(d/2-\varepsilon)$. Integrate their resolvent identity around the
circle; its length is $\pi d$, giving the bound $\delta$.
Since $\varepsilon<d/4$, this bound is less than one. Both contour
projections are orthogonal, so their ranges have the same dimension,
namely one. The eigenvalue perturbation bound
$|\widetilde\lambda-\lambda|\le\varepsilon$ follows from the
selfadjoint spectral distance estimate, or the min--max principle.
Since $d\le|\lambda|$ and $\varepsilon<d/4$, taking reciprocals is
legitimate and gives the bound $\xi$.

Because $\widetilde E$ is supported on $V_N$,
$\widetilde EP_N\phi=\widetilde E\phi$. The difference of the two
exact probe masses is at most $\delta B^2$. Replacing $P_N\phi$ by
$\widetilde a_N$ contributes at most $2B\eta+\eta^2$ by expanding
the squared norm. This proves the mass estimate without requiring a
separate convergence modulus for $P_N\phi$.
Divide that estimate by $p(\widetilde\gamma)^2$ and use
\eqref{eq:probe-bounds}. Finally, integrating the logarithmic derivative
bound gives
$|p(\gamma)^2/p(\widetilde\gamma)^2-1|
\le e^{2C_0|\gamma-\widetilde\gamma|}-1$.
The claimed division error follows.
\end{proof}

\begin{proposition}[A local residual certificate for the probe]
\label{prop:local-probe-residual}
Let $v\in V_N$ have $q(v)=1$, and put
\[
 \mu=\langle Cv,v\rangle_q,\qquad
 \rho=\|(C-\mu)v\|_q.
\]
Assume $|\mu|>\rho$, and let $\lambda$ be the unique spectral point
in $[\mu-\rho,\mu+\rho]$.  Let $d>0$ be a lower bound for its
distance to the rest of $\sigma(C)$, including zero, with
$\rho<d/2$.  Define
\[
 \delta_v=\frac{\rho}{d-\rho},\qquad
 \omega=\min\left\{\rho,
             \frac{\rho^2(d-\rho)}{d(d-2\rho)}\right\},\qquad
 \xi_v=\frac{\omega}{|\mu|(|\mu|-\omega)}.
\]
Then, writing $E_\lambda$ for the spectral projection and
$\widetilde\gamma=1/\mu$, $\gamma=1/\lambda$, one has
\[
 \|v\otimes v-E_\lambda\|\le\delta_v,
 \qquad |\mu-\lambda|\le\omega,
 \qquad |\widetilde\gamma-\gamma|\le\xi_v.
\]
The finite probe mass $u=|\langle\phi,v\rangle_q|^2$ therefore obeys
\begin{equation}
 \left|\frac{u}{p(\widetilde\gamma)^2}-m_\gamma\right|
 \le\frac{(1+(T+\xi_v)^2)^2}{(1-r)^4}\delta_v B^2
             +m_\gamma(e^{2C_0\xi_v}-1)
 \quad (|\gamma|\le T),
 \label{eq:local-probe-rounding}
\end{equation}
with $r,B,C_0$ as in \eqref{eq:probe-bounds} and
Proposition~\ref{prop:spectral-rounding}.  The same upper bounds for
$m_\gamma$ and the same strict $1/2$ rounding threshold apply.
\end{proposition}

\begin{proof}
The spectral theorem gives
$\operatorname{dist}(\mu,\sigma(C))\le\rho$.  Since $|\mu|>\rho$,
the enclosed spectral point is nonzero.  Write
$v=\sum_j a_je_j$ in the orthonormal coordinate basis of $C$, and
let $w=|a_\lambda|^2$.  Every other spectral point obeys
$|\lambda_j-\mu|\ge d-\rho$.  Consequently
\[
 1-w\le\frac{\rho^2}{(d-\rho)^2}.
\]
The distance between the two rank-one orthogonal projections is
$\sqrt{1-w}$, which gives the first inequality.  The Rayleigh
identity yields
\[
 (\mu-\lambda)w
      =\sum_{j\ne\lambda}(\lambda_j-\mu)|a_j|^2.
\]
Cauchy--Schwarz bounds the absolute value on the right by
$\rho\sqrt{1-w}\le\rho^2/(d-\rho)$.  Dividing by
$w\ge1-\rho^2/(d-\rho)^2>0$ proves the second expression in the
definition of $\omega$; spectral inclusion already gives the first.
Taking reciprocals proves the bound $\xi_v$.

The projection estimate gives
$|u-\|E_\lambda\phi\|_q^2|\le\delta_v B^2$.
In the sampling space that exact mass is $m_\gamma p(\gamma)^2$.
Division and the logarithmic derivative bound for $p$ give
\eqref{eq:local-probe-rounding} exactly as in the preceding proof.
\end{proof}

All data in this residual can be computed from finite Weil pairings.
In addition to $G,H_N,c$, form the primitive Gram matrix
\[
 J_N=(q(\mathcal Pb_j,\mathcal Pb_i))_{i,j}.
\]
For a nonzero coefficient vector $x$, take
$v=\sum_jx_jb_j/(x^*Gx)^{1/2}$.  Then the exact formulas are
\begin{equation}
 \mu=\frac{x^*H_Nx}{x^*Gx},\qquad
 \rho^2=\frac{x^*J_Nx}{x^*Gx}-\mu^2,\qquad
 u=\frac{|x^*c|^2}{x^*Gx}.
 \label{eq:residual-arithmetic}
\end{equation}
Thus no separate approximation to $P_N\phi$ is needed.  The primitive
of a compact function with integral zero has the same compact support,
so the explicit formula for $J_N$ has a finite prime-power cutoff too.
Certified interval errors in \eqref{eq:residual-arithmetic} and in the
final divided scalar are included in the rounding inequality.

For translated tents this new Gram matrix has a direct scalar recipe.
Let $t$ be the tent centered at zero and let
$b_i(x)=t(x-ih)-t(x)$.  Define the even function
\[
 A_2(d)=\sum_{\gamma\in\Gamma}
      m_\gamma|\widehat t(\gamma)|^2
                      \frac{1-\cos(\gamma d)}{\gamma^2}.
\]
Expanding the two translation factors gives
\[
 (J_N)_{ij}=A_2(ih)+A_2(jh)-A_2((i-j)h).
\]
The inverse Fourier transform of
$(1-\cos(\gamma d))/\gamma^2$ is
$\tfrac12(|d|-|x|)_+$.  Therefore the geometric Weil functional
computes $A_2(d)$ from the compact test
$\tfrac12(|d|-|x|)_+*(t*\widetilde t)$.
It is piecewise polynomial of degree five, with explicitly known
breakpoints.  This identifies the additional finite integrals required
by a residual certificate.

For comparison with the original piecewise-polynomial probe, set
$\phi_0(x)=(1-|x|)_++(\sqrt2-|x|)_+$. Its transform obeys the
following fully explicit bounds:
\[
 p_0(t)=\frac{2(1-\cos t)+2(1-\cos(\sqrt2t))}{t^2},\quad p_0(0)=3,
\]
\begin{equation}
 p_0(t)\ge
 \begin{cases}
 12/\pi^2,&|t|\le\pi/\sqrt2,\\[2pt]
 \displaystyle\frac{16}{3t^2
     (\sqrt2|t|/\pi+(\sqrt2+1)/2)^2},&|t|>\pi/\sqrt2.
 \end{cases}
 \label{eq:two-tent-bound}
\end{equation}
Here is the arithmetic proof. By evenness take $t>0$, and choose
integers $n,m$ nearest to $t/(2\pi),\sqrt2t/(2\pi)$, respectively.
Put $e_1=t/(2\pi)-n$, $e_2=\sqrt2t/(2\pi)-m$. Above the stated
threshold these integers are not both zero. The nonzero integer
$m^2-2n^2$ has absolute value at least one, while
$|m+\sqrt2n|\le D:=\sqrt2t/\pi+(\sqrt2+1)/2$. Therefore
$|\sqrt2e_1-e_2|=|m-\sqrt2n|\ge1/D$, and
$e_1^2+e_2^2\ge1/(3D^2)$. Concavity gives
$\sin(\pi|e_j|)\ge2|e_j|$, since $|e_j|\le1/2$. Using
$1-\cos(2\pi e_j)=2\sin^2(\pi e_j)$ proves the second bound.
For $0<t\le\pi/\sqrt2$, apply the same sine inequality directly
to $t/2$ and $\sqrt2t/2$; it gives the first bound.
The value at zero follows by continuity. Thus division by $p_0^2$
has an explicit loss of order $T^8$, while the exponential
self-correlation probe has the order $T^4$ bound in
\eqref{eq:spectral-rounding-error}.

\subsection{A common bounded height and the remaining quantifiers}
This comparison identifies the separation data that both contour and
spectral certification must control. Let $T$ exceed the first ordinate,
and let $s(2T)$ be a certified positive lower bound for the distances
between distinct positive ordinates at most $2T$. For $0<\gamma\le T$,
\begin{equation}
 \operatorname{dist}\bigl(\gamma^{-1},
               \sigma(C)\setminus\{\gamma^{-1}\}\bigr)
 \ge\min\{s(2T)/(2T^2),\,1/(2T)\}.
 \label{eq:reciprocal-separation}
\end{equation}
For another ordinate $0<\eta\le2T$, the identity
$|\gamma^{-1}-\eta^{-1}|=|\gamma-\eta|/(\gamma\eta)$ proves the
first bound. Ordinates above $2T$, negative ordinates, and the spectral
point zero have distance at least $1/(2T)$.
This is the exact conversion from an ordinate gap to the spectral gap.
The same certified ordinate gaps permit disjoint disks of radius less
than one third of the gap in the contour calculation.

The interval selection in Proposition~\ref{prop:local-probe-residual}
can also use these gaps directly.  Write $d_T$ for the right side of
\eqref{eq:reciprocal-separation}.  A positive Rayleigh value with
$\mu-\rho\ge1/T$ and $2\rho<d_T$ has a spectral point in its
residual interval, by the spectral theorem.  That point corresponds to
an ordinate at most $T$, and the gap bound makes it the unique point
in the interval.  Thus the local residual test needs no additional
unknown spectral projection.  For example, an ordinate separation
exceeding $6\cdot10^{-10}$ through height $90$ yields
$d_{45}>1.48\cdot10^{-13}$.  Rounding still requires the separately
certified primitive Gram residual and the divided mass error in
\eqref{eq:local-probe-rounding}; separation alone does not bound them.

The common range consists of the heights for which the contour's
boundary and Taylor errors are certified and the right side of
\eqref{eq:spectral-rounding-error} or
\eqref{eq:local-probe-rounding}, uniformly over the enclosed
ordinates, is less than $1/2$. Suzuki's convergence supplies eventual
correctness at each fixed ordinate, and hence simultaneously below
each fixed height, because there are only finitely many such ordinates.
The stopping index can depend on the height; this convergence statement
does not specify a certified stopping index or assert that the recovered
integers equal one. An all-height certificate of simplicity must supply
an unbounded sequence of heights, finite packets covering all ordinates
below each height, certified errors less than $1/2$, and recovered
integers all equal to one. The contour boundary, Taylor and
leading-coefficient bounds, or the spectral residual and separation
bounds, supply the corresponding finite error tests. A finite
successful list recovers the multiplicities included in that list.

One exact identity also measures the unresolved global mass. The
Hilbert--Schmidt orthogonal projection $C\mapsto P_NCP_N$ gives
\[
 \|C-C_N\|_{\mathcal S_2}^2
      =\|C\|_{\mathcal S_2}^2-\operatorname{tr}(C_N^2).
\]
Here and below $\Gamma$ includes both signs of every distinct ordinate.
On the other hand the Hadamard product gives
$-\Xi''(0)/\Xi(0)=\sum_\gamma m_\gamma/\gamma^2$.
Consequently the computable decreasing quantities
\[
 \mathcal D_N=-\frac{\Xi''(0)}{\Xi(0)}
                      -\operatorname{tr}(C_N^2)
\]
satisfy the exact formula
\begin{equation}
 \mathcal D_N=
 \|C-C_N\|_{\mathcal S_2}^2+
           \sum_{\gamma\in\Gamma}\frac{m_\gamma-1}{\gamma^2}.
 \label{eq:trace-multiplicity-defect}
\end{equation}
The sum in \eqref{eq:trace-multiplicity-defect} is over signed ordinates,
and equals $2\sum_{\gamma>0}(m_\gamma-1)/\gamma^2$.
Their limit is the inverse-square weighted multiplicity excess.
Thus $\mathcal D_N\to0$ is another equivalent formulation of simplicity;
the identity proves the equivalence and does not establish that limit.
To verify the orthogonal-projection identity, choose an orthonormal
basis beginning with one for $V_N$: the squared Hilbert--Schmidt norm
is the sum of squared matrix entries, and compression retains exactly
the entries in the finite square. Increasing $N$ increases this square
norm. The weighted sum for $\Xi''(0)$ is obtained by differentiating
the paired product at zero, where absolute convergence is supplied by
$\sum m_\gamma/\gamma^2<\infty$.

\begin{proposition}[A trace certificate for simplicity below a height]
\label{prop:trace-height}
Let $\mathcal D_N$ be the preceding trace defect and let
$E(T)=\sum_{0<\gamma\le T}(m_\gamma-1)$, with the sum over distinct
positive ordinates. For every $N$ and $T>0$,
\begin{equation}
 E(T)\le\frac{T^2}{2}\mathcal D_N.
 \label{eq:trace-height}
\end{equation}
In particular, a certified upper bound
$\mathcal D_N\le U<2/T^2$ certifies simplicity of every zero with
$0<\gamma\le T$. When $U>0$, it certifies every height
$T<\sqrt{2/U}$; $U=0$ certifies simplicity at every height.
The following conditions are equivalent:
\[
 \text{all nontrivial zeros are simple},\qquad
 \mathcal D_N\longrightarrow0,\qquad
 \text{for every }T>0\text{ some }N\text{ has }\mathcal D_N<2/T^2.
\]
\end{proposition}

\begin{proof}
Every summand of the multiplicity excess and the squared norm in
\eqref{eq:trace-multiplicity-defect} is nonnegative. Restrict its signed
sum to $0<|\gamma|\le T$ and pair the two signs. Since
$\gamma^{-2}\ge T^{-2}$ there,
\[
 \mathcal D_N\ge2\sum_{0<\gamma\le T}
       \frac{m_\gamma-1}{\gamma^2}\ge\frac{2E(T)}{T^2}.
\]
The quantity $E(T)$ is a nonnegative integer, so an upper bound strictly
less than one forces it to vanish. This proves the certificate.
The decreasing limit of $\mathcal D_N$ is the weighted excess in
\eqref{eq:trace-multiplicity-defect}; it vanishes precisely when every
nonnegative integer $m_\gamma-1$ vanishes. This proves the first
equivalence. A sequence tending to zero eventually falls below each
positive threshold $2/T^2$. Conversely, existence of an $N$ below each
such threshold, together with monotonicity, makes the limit zero.
\end{proof}

The finite certificate uses the compression matrices and the constant
$S=-\Xi''(0)/\Xi(0)$. For instance, certified bounds
$S\le S^+$ and $\operatorname{tr}(C_N^2)\ge L_N$ give the usable upper
bound $U=S^+-L_N$. It needs neither an isolating circle for each
reciprocal eigenvalue nor division by an individual probe value.
The trace certificate supplies an upper bound for the aggregate
multiplicity excess. The local contour minors give the multiplicities
inside a chosen contour through their successive vanishing orders,
and the probe gives the multiplicity at an isolated spectral coordinate.
The exact formulas above relate these outputs to the same local
integers $m_\gamma$. A multiple zero above the requested height can
keep the trace defect large, so failure of its finite inequality says
nothing by itself about simplicity below that height. No decay estimate
for $\mathcal D_N$ is established here.

\subsection{An executed arithmetic probe certificate through height twenty}
\label{sec:executed-probe}

The local estimate can be executed with arithmetic separation data.
Take $N=4$, so $V_4$ has support $[-4,4]$, mesh $h=1/16$, and
dimension $126$.  The probe and its transform are
\[
 \phi(x)=\frac{e^{-|x|}-e^{|x|-2}}2\,\mathbf1_{|x|\le1},
 \qquad
 p(t)=\frac{1+e^{-2}-2e^{-1}\cos t}{1+t^2}.
\]
All the quantities in the following calculation are obtained from
finite arithmetic Weil pairings and completed moments at zero.

First, the completed moments themselves give an adequate local gap.
List the distinct positive eigenvalues of $C$ as
$\lambda_1>\lambda_2>\cdots>0$, and write
$S_k/2=\sum_jm_j\lambda_j^{2k}$.  For a positive Rayleigh value
$\mu$, one has $\lambda_1\ge\mu$.  The level-four calculation gives
\[
 b_7:=\bigl(S_7/2-\mu^{14}\bigr)^{1/14}
       <0.048230126,
 \qquad \mu-b_7>0.022515334.
\]
For every $j\ge2$,
\[
 \lambda_j^{14}\le S_7/2-\lambda_1^{14}
                        \le S_7/2-\mu^{14},
\]
so $\lambda_j\le b_7$.  The negative eigenvalues and zero have
distance at least $\lambda_1$ from $\lambda_1$.  Thus $d=9/400$
is a certified lower bound for the full spectral gap at $\lambda_1$.
Moreover, $\mu>1/20>b_7$ proves that the packet of positive ordinates
through height $20$ contains exactly one distinct point.
This moment inequality also implies $m_1=1$: the contrary inequality
$m_1\ge2$ would give $S_7/2\ge2\mu^{14}$, contradicting $b_7<\mu$.
The additional probe computation below verifies the projection-mass
rounding mechanism on these same arithmetic data.

\begin{proposition}[A completed local probe calculation]
\label{prop:executed-local-probe}
For a certified unit vector $v\in V_4$, the quantities in
Proposition~\ref{prop:local-probe-residual} and
\eqref{eq:residual-arithmetic} have the following enclosures.
\begin{center}
\begin{tabular}{ll}
\hline
Quantity & Certified enclosure or bound\\
\hline
$q(\phi)$ & $(9.98158838,9.98158839)\,10^{-5}$\\
$\mu$ & $(0.0707454607,0.0707454609)$\\
$\rho$ & $(0.0003983098,0.0003983100)$\\
$d$ & $9/400$ is a strict lower bound\\
$\delta_v$ & $<0.018022$\\
$\xi_v$ & $<0.001435$\\
$u/p(1/\mu)^2$ & $(0.9996487647,0.9996487649)$\\
Total multiplicity error & $<0.080978$\\
\hline
\end{tabular}
\end{center}
The resulting interval contains exactly the integer one.  This
executes the spectral-probe test for every positive ordinate through
height $20$.
\end{proposition}

\begin{proof}
We describe the arithmetic test functions, the certified linear
algebra, and the rounding inequality.  The primitive Gram matrix
$J_N$ is evaluated by the degree-five test preceding
\eqref{eq:two-tent-bound}.  For the additional probe pairings, put
$r=e^{-1}$, $g(x)=e^{-|x|}/2$, and $T_sg(x)=g(x+s)$.  Direct
substitution on the intervals separated by $-1,0,1$ gives
\[
 \phi=\bigl((1+r^2)I-r(T_1+T_{-1})\bigr)g.
\]
Splitting its convolution integral at $0$ and $x$ gives
$g*g=k$, where $k(x)=(1+|x|)e^{-|x|}/4$.  Consequently
\begin{align*}
 \phi*\phi={}&(1+4r^2+r^4)k
   -2r(1+r^2)(T_1+T_{-1})k\\
   &\quad+r^2(T_2+T_{-2})k.
\end{align*}
This is the compact test for $q(\phi)$.

For the central tent $t_h=(1-|x|/h)_+$, take the causal primitive
$A_\phi(x)=\int_{-\infty}^x(x-u)\phi(u)\,du$.  The distributional
identity $t_h''=(\delta_{-h}-2\delta_0+\delta_h)/h$ proves
\[
 K(x):=(\phi*t_h)(x)
   =\frac{A_\phi(x+h)-2A_\phi(x)+A_\phi(x-h)}h.
\]
The evenized test for $q(\phi,t_h(\cdot-jh)-t_h)$ is therefore
$K(x+jh)+K(x-jh)-2K(x)$.  Every piece of these tests is a sum
$F(x)=\sum_{a=-1}^1(c_a+d_ax)e^{ax}$, with exact dyadic
breakpoints.  Their supports have radius at most $5$; the test for
$q(\phi)$ has radius $2$, and the primitive Gram tests have radius
at most $8$.  Thus every prime sum is finite.

For completeness, the removable singularity in the first
archimedean integration interval is cancelled by the exact identity
\[
 \frac{F(x)-F(0)}x
   =\sum_{a=-1}^1
        \bigl(c_a a\operatorname{exprel}(ax)+d_a e^{ax}\bigr),
 \qquad
 \operatorname{exprel}(z)={}_1F_1(1;2;z).
\]
The constant multiplying the original test value $F(0)/2$ is
\[
 C_h=-\log\pi-\gamma_E
   +\int_0^h\frac{2(e^{-2x}-e^{-x/2})}{1-e^{-2x}}\,dx
   -\log(1-e^{-2h}).
\]
After this constant is extracted, the first-cell remainder is
\[
 -\int_0^h
    \frac{e^{-x/2}}{2\operatorname{exprel}(-2x)}
                      \frac{F(x)-F(0)}x\,dx.
\]
The other cells have no singular endpoint.  Polynomial-exponential
antiderivatives evaluate the pole terms.  Certified complex ball
integration evaluates the remaining meromorphic callbacks, and
each prime logarithm is separated from the piece boundaries before
its corresponding expression is evaluated.

A midpoint congruence preconditions the positive Gram matrix; the
full outward congruence residual certifies positivity.  An approximate
eigenvector then proposes coefficients, which are frozen as exact
dyadic numbers before certification.  Writing these coefficients
$x=a+ib$ and $H_N=-iB$ with real antisymmetric $B$, the exact formulas
used are
\[
 \begin{split}
 n&=a^TGa+b^TGb,\qquad
 \mu=2a^TBb/n,\\
 \rho^2&=(a^TJ_Na+b^TJ_Nb)/n-\mu^2,\qquad
 u=((a^Tc)^2+(b^Tc)^2)/n.
 \end{split}
\]
All these expressions are evaluated with outward balls at $120$
decimal digits, and repeated at $160$ digits.  The finite input
Gram and primitive pairings, candidate coefficients, pairing
enclosures, and strict tests are recorded with source hashes in the
external certificates.

The residual interval lies entirely above $b_7$, so its spectral
point is $\lambda_1$.  The bounds above give $\rho<d/2$.
Proposition~\ref{prop:local-probe-residual} therefore applies.
The exact evaluated denominator $p(1/\mu)^2$ can be used in its
division step.  Since
\[
 m_1\le\frac{q(\phi)e^{2C_0\xi_v}}{p(1/\mu)^2}<4,
\]
integrality permits the a priori upper bound $m_1\le3$ even without
using the direct moment implication $m_1=1$.  The final error bound
computed in balls is
\[
 \frac{\delta_vq(\phi)}{p(1/\mu)^2}
             +3\bigl(e^{2C_0\xi_v}-1\bigr)<0.080978.
\]
Together with the enclosure for the divided mass, this lies
strictly inside $(1/2,3/2)$ and proves the claimed rounding.
The arithmetic packet count already proved that it covers all
positive ordinates through height $20$.
\end{proof}

At level three the same bound for the logarithmic error exceeds
$0.755$, so the rounding interval does not isolate an integer.  No probe certificate for the full packet through height
$45$ is asserted by the present computation.

\section{A complete finite arithmetic test below a height}
\label{sec:fd-arithmetic}
The even trace powers of Suzuki's compressions give a finite arithmetic certificate for every bounded packet that consists of simple zeros.
Write \(\Xi(z)=\xi(1/2+iz)\), and let \(P_N\) be the orthogonal
projection in the Weil Hilbert space onto the mean-zero dyadic
piecewise-linear space \(V_N\) supported on \([-N,N]\).
Let \(C_N=P_NCP_N\), where \(C\) has the single eigenvalue
\(1/\gamma\) at each distinct signed ordinate \(\gamma\).
Its coordinate norm has weight \(m_\gamma\).
For \(k\ge1\) define
\[
 S_k=-2k[z^{2k}]\log\frac{\Xi(z)}{\Xi(0)},\qquad
 D_{N,k}=S_k-\operatorname{tr}(C_N^{2k}),\qquad
 H_{N,k}=\left(\frac2{D_{N,k}}\right)^{1/(2k)}.
\]
For a height \(T>0\), let
\[
 E(T)=\sum_{0<\gamma\le T}(m_\gamma-1).
\]
This is the multiplicity excess in the positive packet. It is an
integer and vanishes exactly when that packet is simple.

\begin{theorem}[Finite arithmetic certificates and the first multiple zero]
\label{thm:fd-arithmetic-complete}
For every \(N,k\ge1\), the number \(D_{N,k}\) is positive and
\[
 D_{N,k}\ge\frac{2E(T)}{T^{2k}}\qquad(T>0).
\]
For each fixed \(k\), it decreases with \(N\) and
\[
 \lim_{N\to\infty}D_{N,k}
   =2\sum_{\gamma>0}\frac{m_\gamma-1}{\gamma^{2k}}.
\]
If \(\gamma_*\) denotes the least positive multiple ordinate,
with value \(+\infty\) when there is none, then
\[
 \sup_{N,k\ge1}H_{N,k}=\gamma_*.
\]
Consequently, for each fixed \(T>0\),
\[
 \begin{split}
 &\text{every zero with }0<\Im\rho\le T\text{ is simple}\\
 &\hspace{16mm}\Longleftrightarrow\
 \text{there are finite }N,k\text{ with }D_{N,k}<2/T^{2k}.
 \end{split}
\]
The exact relation to the Jensen limit is
\[
 \lim_{a\downarrow0}a\,j_*(-a)
   =\frac12\left(\sup_{N,k\ge1}H_{N,k}\right)^2.
\]
\end{theorem}

\begin{proof}
The paired Hadamard product gives
\[
 S_k=2\sum_{\gamma>0}\frac{m_\gamma}{\gamma^{2k}}.
\]
The trace of \(C^{2k}\) is
\(2\sum_{\gamma>0}\gamma^{-2k}\), since each distinct signed
coordinate is one-dimensional. Compression decreases each
singular value: the approximation-number characterization gives
\(s_j(P_NCP_N)\le s_j(C)\). Summing their \(2k\)-th powers
therefore proves
\[
 D_{N,k}\ge
 2\sum_{\gamma>0}(m_\gamma-1)\gamma^{-2k}
 \ge 2E(T)T^{-2k}.
\]
The first inequality is strict before subtracting the excess:
\(C_N\) has finite rank and \(C\) has infinitely many nonzero
singular values. Thus \(D_{N,k}>0\).

The spaces \(V_N\) are nested and dense in the Weil space, as
proved in the compression construction. The min--max principle
makes each positive eigenvalue of \(C_N\) nondecreasing with \(N\);
applied to \(-C_N\), it makes the absolute values of the negative
eigenvalues nondecreasing. Adding zero eigenvalues when necessary,
their even-power sum is nondecreasing. This proves monotonicity
of \(D_{N,k}\).

Suzuki's compression argument gives \(C_N\to C\) in
Hilbert--Schmidt norm. To pass to higher powers, use
\[
 A^{2k}-B^{2k}
 =\sum_{j=0}^{2k-1}A^j(A-B)B^{2k-1-j}.
\]
In each summand, one of the powers besides \(A-B\) contains
at least one factor. Put that factor in Hilbert--Schmidt norm,
put \(A-B\) in the same norm, and bound the remaining powers
in operator norm. The ideal inequality
\(\|UV\|_1\le\|U\|_2\|V\|_2\) proves convergence in trace norm,
since the operator and Hilbert--Schmidt norms remain bounded.
The asserted limit follows.

For the supremum, list the numbers \(1/\gamma\) belonging to
multiple positive ordinates, each repeated \(m_\gamma-1\) times,
and denote this nonnegative sequence by \(v\). It belongs to
\(\ell^2\). Monotonicity in \(N\) gives
\[
 \sup_N H_{N,k}=\|v\|_{2k}^{-1},
\]
with the right side infinite for the empty sequence.
For a nonempty sequence, its largest entry is
\(v_{\max}=1/\gamma_*\), and
\[
 v_{\max}^{2k}\le\sum_jv_j^{2k}
 \le v_{\max}^{2k-2}\sum_jv_j^2.
\]
Taking \(2k\)-th roots shows
\(\|v\|_{2k}\to v_{\max}\). Moreover
\(\|v\|_{2k}\ge v_{\max}\) for every \(k\). These two facts prove
the supremum identity, including the empty case.

A strict defect bound makes the nonnegative integer \(E(T)\)
less than one and hence zero. Conversely, when \(E(T)=0\),
discreteness of the zeros gives \(\gamma_*>T\); the supremum
identity then supplies a finite pair with \(H_{N,k}>T\).
Finally \(x_*=2\gamma_*\) in Theorem~\ref{thm:fd-limits}, so
its Jensen formula is exactly the last displayed identity.
\end{proof}

This theorem supplies an exact relation among the spectral finite
data, bounded-height simplicity, and the Jensen detecting constant.
At every fixed height the strict test has a terminating search
precisely for a simple packet: enumerate \(N,k\) and increase
the certified precision until a strict inequality is found.
The assertion that this search terminates for every height is
the all-simple question.

\subsection{Operator control and the size of a finite certificate}

The same defect controls the operator approximation, without further
matrix entries. All operator norms below are for the Weil Hilbert norm.

\begin{proposition}[An operator bound from an even trace defect]
\label{prop:fd-even-operator}
For every \(N,k\ge1\),
\[
 \|C-C_N\|\le\frac2{\sqrt3}\,D_{N,k}^{1/(2k)}.
\]
The constant \(2/\sqrt3\) is sharp in the elementary selfadjoint
block estimate used in this proof.
\end{proposition}
\begin{proof}
Write \(P=P_N\), \(Q=I-P\), and let \(v_j\) be an orthonormal
eigenbasis for the finite compression on \(P\mathcal H\).
Scalar convexity, applied to the spectral measure of \(C\) at
\(v_j\), gives
\[
 |\langle Cv_j,v_j\rangle|^{2k}
       \le\langle |C|^{2k}v_j,v_j\rangle.
\]
After summing, and using \(S_k\ge\operatorname{tr}|C|^{2k}\),
we obtain
\[
 D_{N,k}\ge\operatorname{tr}(Q|C|^{2k}Q).
\]
For a unit vector \(v\in Q\mathcal H\), the power-mean
inequality for its spectral measure gives
\[
 \|Cv\|^{2k}\le\langle |C|^{2k}v,v\rangle
       \le\operatorname{tr}(Q|C|^{2k}Q)\le D_{N,k}.
\]
Thus \(\delta=\|CQ\|\le D_{N,k}^{1/(2k)}\).

Relative to \(P\mathcal H\oplus Q\mathcal H\), write
\[
 C-C_N=\begin{pmatrix}0&B^*\\B&D\end{pmatrix}.
\]
The column bound is \(BB^*+D^2\le\delta^2I\).
For a unit vector \((x,y)\), put \(a=\|x\|\),
\(b=\|y\|\), and \(w=\|Dy\|\). Then
\begin{align*}
 |\langle(C-C_N)(x,y),(x,y)\rangle|
 &\le2a\sqrt{\delta^2b^2-w^2}+bw\\
 &\le\delta\sqrt{4a^2b^2+b^4}
   =\delta\sqrt{4b^2-3b^4}
   \le\frac{2\delta}{\sqrt3}.
\end{align*}
The second inequality is the ordinary two-dimensional
Cauchy--Schwarz inequality. The last polynomial attains its
maximum at \(b^2=2/3\). Taking the supremum proves the assertion.
Sharpness of this block estimate follows from the scalar blocks
\(B=\delta\sqrt{2/3}\), \(D=\delta/\sqrt3\), whose omitted
column has norm \(\delta\) and whose resulting block matrix has
norm \(2\delta/\sqrt3\).
\end{proof}

Let \(\gamma_1\le\gamma_2\le\cdots\) now list the positive
ordinates with multiplicity. This counted list is used only in the
following exact size comparison.

\begin{proposition}[Dimension and moment order]
\label{prop:fd-certificate-size}
Put \(d_N=\dim V_N=2r_N\). For every \(k\ge1\),
\[
 D_{N,k}\ge2\sum_{j>r_N}\gamma_j^{-2k}.
\]
Consequently a strict certificate \(D_{N,k}<2/T^{2k}\)
requires \(d_N\ge2n(T)\), where \(n(T)\) counts positive
zeros through \(T\) with multiplicity. For each fixed \(N\),
\[
 \lim_{k\to\infty}H_{N,k}=\gamma_1.
\]
\end{proposition}
\begin{proof}
The rank is at most \(d_N\), and compression decreases every
singular value. Hence \(\operatorname{tr}(C_N^{2k})\) is at
most the sum of the largest \(d_N\) even powers of the singular
values of \(C\). The counted signed ordinate list contains each
such singular value at least as many times as the distinct
coordinate list. Its largest \(2r_N\) even powers sum to
\(2\sum_{j\le r_N}\gamma_j^{-2k}\). Subtraction from \(S_k\)
proves the tail bound. If \(n(T)>r_N\), its first omitted
term is at least \(2/T^{2k}\), contradicting the strict
certificate. This proves the dimension requirement.

Write \(\lambda_1=1/\gamma_1=\|C\|\) and
\(\beta_N=\|C_N\|\). We claim \(\beta_N<\lambda_1\).
Equality would give a unit vector in \(V_N\) attaining an
extreme Rayleigh value of \(C\). Positivity of
\(\lambda_1I\mp C\) would make it a \(C\)-eigenvector with
eigenvalue \(\pm\lambda_1\). The compact-primitive identity
then gives, as equality in the Weil space,
\[
 i\int_{-N}^{x}f(y)\,dy=\pm\lambda_1 f(x).
\]
The injective compact-support embedding makes this an equality
of ordinary functions. Differentiating gives a first-order
constant-coefficient equation, with \(f(-N)=0\), and therefore
\(f=0\). This contradicts its unit norm and proves the claim.

The summable counted reciprocal-square sequence gives
\(S_k^{1/(2k)}\to\lambda_1\). Moreover,
\[
 0\le\frac{\operatorname{tr}(C_N^{2k})}{S_k}
 \le\frac{d_N}{2}
       \left(\frac{\beta_N}{\lambda_1}\right)^{2k}
 \longrightarrow0.
\]
It follows that \(D_{N,k}^{1/(2k)}\to\lambda_1\), which is
exactly the last assertion.
\end{proof}

The dimension requirement and the fixed-space limit quantify two
constraints on the actual finite search. Increasing the moment
order alone in one space eventually returns its height threshold
to the first ordinate. The complete family in
Theorem~\ref{thm:fd-arithmetic-complete} permits the space and the
moment order to vary together.

\subsection{Evaluation using the arithmetic side of Weil's formula}

We now evaluate these certificates without inserting zero ordinates.
For \(h=f*\widetilde g\), where
\(\widetilde g(x)=\overline{g(-x)}\), the explicit formula is
\begin{align}
 q(f,g)
 ={}&\int_{\mathbb R}h(x)(e^{x/2}+e^{-x/2})\,dx
       -(\log\pi+\gamma_{\!E})h(0)\notag\\
 &+\int_0^\infty
   \frac{2e^{-2x}h(0)-e^{-x/2}(h(x)+h(-x))}
        {1-e^{-2x}}\,dx\notag\\
 &-\sum_{n\ge2}\frac{\Lambda(n)}{\sqrt n}
             \bigl(h(\log n)+h(-\log n)\bigr).
 \label{eq:fd-arithmetic-weil}
\end{align}
Here \(\gamma_{\!E}\) is Euler's constant and \(\Lambda(n)\)
is the von Mangoldt function. This normalization follows from
\cite[(3.2)--(3.7), pp.~5--6]{ccm2025} by \(u=e^x\).
Their constant is \(-(\log(4\pi)+\gamma_{\!E})h(0)\), with
\(2e^{-x}h(0)\) in the numerator. The identity
\[
 2\int_0^\infty\frac{e^{-2x}-e^{-x}}{1-e^{-2x}}\,dx=-\log4
\]
proves equality with~\eqref{eq:fd-arithmetic-weil}.
Only prime powers \(n<e^{2N}\) enter on the support in question.

Put \(h_N=2^{-N}\) and
\[
 t_j(x)=\max\{1-|x-jh_N|/h_N,0\},\qquad
 b_j=t_j-t_0,
\]
where \(-N/h_N<j<N/h_N\) and \(j\ne0\).
These \(2N2^N-2\) functions form the stated basis of \(V_N\):
all interior hats have the same integral \(h_N\), so their
coefficient sum is exactly the mean-zero condition.
The correlation of the central hat is
\[
 b(x)=(t_0*t_0)(x)
 =\frac1{6h_N^2}\sum_{\ell=0}^4
       (-1)^\ell\binom4\ell
             (x+(2-\ell)h_N)_+^3.
\]
It can be checked by integration on the two intervals:
\[
 b(x)=
 \begin{cases}
 2h_N/3-x^2/h_N+|x|^3/(2h_N^2),&|x|\le h_N,\\
 (2h_N-|x|)^3/(6h_N^2),&h_N\le|x|\le2h_N,\\
 0,&|x|\ge2h_N .
 \end{cases}
\]
Let \(W\) denote the linear functional in
\eqref{eq:fd-arithmetic-weil}, and define real kernels
\[
 g(d)=W(b(\,\cdot-d)),\qquad
 a(d)=W\left(\int_0^d b(\,\cdot-u)\,du\right).
\]
The first kernel is even and the second odd. In the Fourier
variable,
\[
 a(d)=\sum_{\gamma\in\Gamma}
      m_\gamma|\widehat t_0(\gamma)|^2
                   \frac{\sin(\gamma d)}{\gamma}.
\]
The identity follows by integrating the absolutely convergent
correlation series; equivalently it follows from the compact
primitive, after taking the basis differences.
Consequently the Gram and primitive matrices are
\[
 \begin{split}
 G_{ij}&=g((j-i)h_N)-g(jh_N)-g(ih_N)+g(0),\\
 B_{ij}&=a((j-i)h_N)-a(jh_N)+a(ih_N),\\
 H_{ij}&=q\left(i\int_{-N}^{\,\cdot}b_j,b_i\right)=-iB_{ij}.
 \end{split}
\]
Thus \(G\) is real symmetric and \(B\) real antisymmetric, and
\[
 \operatorname{tr}(C_N^{2k})
       =(-1)^k\operatorname{tr}\bigl((G^{-1}B)^{2k}\bigr).
 \label{eq:fd-arithmetic-trace}
\]
This evaluates the invariant trace without a matrix square root.

\subsection{Certified examples}

The following are strict rational height certificates from
\eqref{eq:fd-arithmetic-weil}. Each displayed upper bound \(U\)
encloses \(D_{N,k}\), and each row satisfies the exact rational
inequality \(U<2/T^{2k}\).
\[
\begin{array}{c|r|r|c|c|c}
 N&\dim V_N&\text{prime powers}&k&U&T\\ \hline
 1&2&5&9&3.946530144\,10^{-21}&14.134\\
 2&14&24&3&2.539105111\,10^{-8}&20.704\\
 3&46&98&3&5.419648073\,10^{-9}&26.782\\
 4&126&465&2&7.880330238\,10^{-7}&39.913\\
 5&318&2532&2&1.318023575\,10^{-7}&62.413\\
 6&766&15040&3&2.309270917\,10^{-12}&97.632
\end{array}
\]
In particular, this finite arithmetic calculation certifies
simplicity of every nontrivial zero with
\(0<\Im\rho\le97.632\).
The largest prime power in that calculation is \(162751\).
Proposition~\ref{prop:fd-even-operator} also gives the strict
operator estimate \(\|C-C_6\|<0.0133\): raising this bound
to the sixth power reduces its verification to the exact rational
inequality \(64(2.309270917\cdot10^{-12})<27(0.0133)^6\).

Here is the numerical error control. On each dyadic interval,
\(h(x)+h(-x)\) is a rational polynomial of degree at most four.
The polar integrals are polynomial-exponential antiderivatives.
The remaining integrals have the form
\[
 I_{\ell,r}=\int_{\ell h_N}^{(\ell+1)h_N}
       \frac{x^r e^{-x/2}}{1-e^{-2x}}\,dx,\qquad 0\le r\le4.
\]
At the origin the constant term is cancelled before evaluation.
Writing
\(\operatorname{exprel}(z)=(e^z-1)/z={}_1F_1(1;2;z)\),
the regular integrand for \(r\ge1\) is
\[
 \frac{x^{r-1}e^{-x/2}}{2\operatorname{exprel}(-2x)}.
\]
The constant-term integrand is
\[
 -\frac32 e^{-x/2}
       \frac{\operatorname{exprel}(-3x/2)}
            {\operatorname{exprel}(-2x)}.
\]
The contribution of \(2h(0)e^{-2x}/(1-e^{-2x})\) beyond the
first interval is exactly
\(-h(0)\log(1-e^{-2h_N})\).
All integration callbacks are meromorphic on their complex
domains; interval division detects their possible poles.

Outward ball arithmetic at \(120\) decimal digits encloses all
integrals, prime terms, and matrix entries. To certify positive
definiteness, an approximate \(LDL^T\) factor supplies an
invertible preconditioner \(R\); the decisive verified inequality is
\[
 \sum_{i,j}|(RGR^T-I)_{ij}|<1.
\]
It implies positive definiteness because the left side bounds
the operator norm of the real symmetric error.
For all six rows it is less than \(2\cdot10^{-94}\).
The trace powers are evaluated after the exact congruence
\((G,B)\mapsto(RGR^T,RBR^T)\), whose first matrix is close
to identity. This controls the error propagation through high
powers. The scalar moments use the ball Taylor series of the
completed zeta function at \(1/2\), followed by its formal
logarithm; no zero ordinate is an input.

The \href{\workbenchurl}{\emph{Arithmetic Weil compression}
program and certificate} give the full enclosures, exact
commands, and verified matrix residuals. The arithmetic uses
Arb's outward interval operations and rigorous integration
\cite{Johansson2017arb,Johansson2018integration}.


\subsection{Removing a certified simple packet}
The coordinate representation also records exactly what remains
after a finite certified simple packet has been removed.

\begin{proposition}\label{prop:finite-simple-deflation}
Let $S$ be a finite set of positive simple zero ordinates, and
write $\pm S$ for its signed counterpart. Delete these coordinates
from the Weil form and from the reciprocal multiplication operator:
\[
 q_S(f,g)=q(f,g)-\sum_{u\in\pm S}
                  \widehat f(u)\overline{\widehat g(u)},
 \qquad C_S=C|_{\{e_u:u\in\pm S\}^{\perp}}.
\]
For $k\ge1$, put $S_k^S=S_k-2\sum_{u\in S}u^{-2k}$. Then
\begin{equation}\label{eq:finite-deflation-excess}
 S_k^S-\operatorname{tr}C_S^{2k}
 =S_k-\operatorname{tr}C^{2k}
 =2\sum_{\gamma>0}(m_\gamma-1)\gamma^{-2k}.
\end{equation}
In the same coefficient basis as $G_N,H_N$, set
$e_u=(\widehat b_1(u),\ldots,\widehat b_d(u))$. The corresponding
matrices satisfy
\begin{align*}
 G_N^S&=G_N-\sum_{u\in\pm S}e_u^*e_u,\qquad
 H_N^S=H_N-\sum_{u\in\pm S}u^{-1}e_u^*e_u,\\
 TH_N^S-G_N^S
   &=(TH_N-G_N)-\sum_{u\in\pm S}(T/u-1)e_u^*e_u.
\end{align*}
At every $T>0$ with $S\subset(0,T)$, both the counted positive zero count and the
distinct positive coordinate count decrease by $|S|$, so their
difference remains $E(T)$.
\end{proposition}

\begin{proof}
In the weighted sampling realization, deleting the coordinates is
the orthogonal projection onto their complement. The image of the
dense test space is dense in that complement: approximate any vector
there in the original space and apply the projection. Thus the
completion of the projected tests, modulo their null space, is
exactly the displayed space for $C_S$. Each deleted coordinate has
weight one because its zero is simple. Its contribution to both
$S_k$ and $\operatorname{tr}C^{2k}$ is $u^{-2k}$. Subtracting
these finitely many terms in the absolutely convergent sums proves
\eqref{eq:finite-deflation-excess}.

For a coefficient column $c$, evaluation is $e_uc$, and its
squared modulus is $c^*e_u^*e_uc$. The reciprocal operator
multiplies that coordinate by $u^{-1}$. These two facts prove
the first two matrix identities; their stated linear combination
proves the third. Deleting a positive simple zero subtracts one
from each of the two positive counts, which proves the last claim.
\end{proof}

Finite deletion can change the approximation error of a chosen
compression. Formula~\eqref{eq:finite-deflation-excess} shows that
the intrinsic multiplicity excess itself is preserved exactly.
Iteration with a height-dependent deleted packet requires an
estimate for the resulting compression error and for the cross
terms with that growing packet. The positive-density argument
above does not supply those estimates for the successive remainders.

\subsection{Quantitative separation from the same finite data}
\label{sec:fd-arithmetic-separation}

The arithmetic certificate also gives a separation bound on the entire
certified packet. Let \(d_N=\dim V_N\), and let
\(\alpha_N>0\) be the minimum Rayleigh quotient of \(q\) against
the \(L^2\) norm on \(V_N\). Thus
\[
 q(f,f)\ge\alpha_N\|f\|_{L^2(-N,N)}^2\qquad(f\in V_N).
\]
Such a finite bound is supplied by the Gram matrix: if \(M_N\) is the
\(L^2\) Gram matrix in the hat-difference basis, then
\[
 \alpha_N\ge
 \frac{1}{\bigl(\sum_{i,j}|(G^{-1})_{ij}|\bigr)
                 \operatorname{tr}M_N},\qquad
 \operatorname{tr}M_N=\frac{(4d_N-2)2^{-N}}3.
\]
The right-hand side therefore gives an admissible finite lower bound.
Every generic hat difference has squared \(L^2\) norm
\(4\cdot2^{-N}/3\).
The two hats adjacent to the central hat each reduce that value by
\(2^{-N}/3\), proving the trace formula.

\begin{proposition}[A finite separation inequality]
\label{prop:fd-arithmetic-separation}
For every \(N,k\ge1\), \(T>0\), and every two distinct occurrences
in the zero multiset with \(0<\Im\rho\le T\), their
ordinates satisfy
\[
 |\gamma_1-\gamma_2|^2\ge
 \frac{3\alpha_N}{N^3}
       \left(1-\frac{D_{N,k}T^{2k}}2\right).
 \label{eq:fd-separation-bound}
\]
The finite arithmetic data above give the strict bounds
\[
 \alpha_6>7\cdot10^{-17},\qquad
 |\gamma_1-\gamma_2|>6\cdot10^{-10}
       \quad(0<\gamma_1,\gamma_2\le90),
\]
where the two zeros in the second statement are two distinct
occurrences in the multiset, so this bound also excludes a repeated ordinate.
For the reciprocal spectral values, define
\[
 d(T)=\inf\left\{
  \left|\frac1\gamma-\frac1{\gamma'}\right|:
  0<|\gamma|\le T,\ \gamma'\in\Gamma,\ \gamma'\ne\gamma
 \right\}.
\]
Then the same calculation gives \(d(45)>1.48\cdot10^{-13}\).
\end{proposition}

\begin{proof}
Let \(e_\gamma\) be the unit vector of the coordinate indexed by
the distinct signed ordinate \(\gamma\), and put
\(p_N(\gamma)=\|P_Ne_\gamma\|_q^2\).
Choose an orthonormal eigenbasis \(u_j\) of the finite compression.
Scalar convexity of \(x\mapsto x^{2k}\), applied to the spectral
measure of \(C\) in each \(u_j\), gives
\[
 \langle Cu_j,u_j\rangle_q^{2k}
       \le\langle C^{2k}u_j,u_j\rangle_q.
\]
Adding these inequalities and expanding the positive trace in the
coordinate basis proves
\[
 D_{N,k}\ge
 2\sum_{\gamma>0}\frac{m_\gamma-1}{\gamma^{2k}}
 +\sum_{\gamma\in\Gamma}
             \frac{1-p_N(\gamma)}{\gamma^{2k}}.
 \tag{*}
\]
All summands are nonnegative. If a zero in the packet is multiple,
the first sum already gives \(D_{N,k}\ge2/T^{2k}\), and the
asserted separation inequality has a nonpositive right-hand side.
It remains to consider two distinct simple positive ordinates.

Put \(\delta=|\gamma_1-\gamma_2|\) and
\(v_+=(e_{\gamma_1}-e_{\gamma_2})/\sqrt2\).
For \(f\in V_N\), its Fourier transform obeys
\[
 \begin{split}
 |\widehat f(\gamma_1)-\widehat f(\gamma_2)|
 &\le\delta\int_{-N}^N|x f(x)|\,dx\\
 &\le\delta\left(\frac{2N^3}3\right)^{1/2}\|f\|_2.
 \end{split}
\]
The coordinates at these two ordinates have multiplicity one.
Dividing the squared estimate by two, and taking the supremum over
\(q(f,f)=1\), yields
\[
 \|P_Nv_+\|_q^2\le\frac{\delta^2N^3}{3\alpha_N}.
\]
The identical estimate holds for the negative contrast
\(v_-=(e_{-\gamma_1}-e_{-\gamma_2})/\sqrt2\).
In each two-dimensional coordinate space, extend the contrast
to an orthonormal basis. Positivity of \(I-P_N\) then shows
\[
 \sum_{\gamma\in\{\pm\gamma_1,\pm\gamma_2\}}
      (1-p_N(\gamma))
 \ge2\left(1-\frac{\delta^2N^3}{3\alpha_N}\right).
\]
Each of these four ordinates has modulus at most \(T\).
Substitution into \((*)\) and rearrangement proves the inequality.

The Gram lower bound follows from
\(\lambda_{\min}(G)\ge(\sum_{i,j}|(G^{-1})_{ij}|)^{-1}\)
and \(\|M_N\|\le\operatorname{tr}M_N\).
The stored interval bound for the inverse matrix at level six
certifies that the resulting ratio exceeds \(7\cdot10^{-17}\).
Use the already verified rational upper bound
\(D_{6,3}<2.309270917\cdot10^{-12}\).
The remaining numerical comparison is the exact rational inequality
\[
 (6\cdot10^{-10})^2<
 \frac{3(7\cdot10^{-17})}{6^3}
 \left(1-\frac{(2.309270917\cdot10^{-12})90^6}{2}\right).
\]
It proves the claimed separation without a new matrix computation.

Finally consider a signed ordinate \(\gamma\) with
\(0<|\gamma|\le45\). An ordinate of the same sign and modulus
at most \(90\) has reciprocal distance greater than
\((6\cdot10^{-10})/(45\cdot90)\).
An ordinate of the same sign and modulus greater than \(90\)
has reciprocal distance greater than \(1/90\).
An opposite-sign ordinate has reciprocal distance greater than
\(1/45\). Taking the minimum proves
\[
 d(45)\ge\frac{6\cdot10^{-10}}{4050}
          >1.48\cdot10^{-13}.
\]
The \href{\workbenchurl}{\emph{Arithmetic separation} certificate}
records these scalar checks and the hash of their input interval data.
\end{proof}

\section{Transporting a finite trace certificate to detection growth}
\label{sec:trace-detection}

This subsection converts the finite trace defect into a lower bound for
the global detection thresholds.  Write $h_{m,1},\ldots,h_{m,m}$ for the
zeros of the monic probabilists' Hermite polynomial $\operatorname{He}_m$,
and put $D_m=\operatorname{diag}(h_{m,1},\ldots,h_{m,m})$.  Its spectrum
is symmetric, $\operatorname{tr}D_m=0$, and
$\operatorname{tr}D_m^2=m(m-1)$.  For a real diagonal matrix $X$, let
$R(X)$ be the difference between its largest and smallest entries and set
\[
 \mathcal O_m(X)=\int_{U(m)}e^{i\operatorname{tr}(D_mUXU^*)}\,dU,
 \qquad
 \mathfrak r_m=\inf\{R(X):\mathcal O_m(X)<0\},
\]
where Haar measure has mass one.  The symmetry of $D_m$ makes
$\mathcal O_m(X)$ real.

\begin{proposition}[An orbital radius bound]
\label{prop:orbital-unimodal-bound}
For every $m\ge2$,
\begin{equation}
 \mathfrak r_m^2\ge\frac{4\pi^2(m+1)}{3m^2},
 \qquad (m-1)\mathfrak r_m^2\ge\pi^2.
 \label{eq:orbital-unimodal-radius}
\end{equation}
The constant $\pi^2$ in the second inequality is attained for $m=2$.
\end{proposition}

\begin{proof}
We first prove the elementary inequality
\begin{equation}
 \frac{\sin u}{u}\ge\frac12-\frac{u^2}{2\pi^2}
 \quad (u\in\mathbb R),
 \label{eq:sinc-parabola}
\end{equation}
with the continuous value at zero.  The difference $g(u)$ between the
two sides is even.  On $[0,\pi]$, write
$h(u)=\sin u-u\cos u-u^3/\pi^2$.  Then
$g'(u)=-h(u)/u^2$ and
$h'(u)=u^2(\sin u/u-3/\pi^2)$.  The function $\sin u/u$ decreases
strictly from one to zero on this interval, since the derivative of
$u\cos u-\sin u$ is $-u\sin u$.  Thus $h$ first increases and then
decreases.  Its values at both endpoints are zero, so $h\ge0$.
Consequently $g$ decreases to $g(\pi)=0$.  On $[\pi,3\pi/2]$,
\[
 g''(u)=\left(-\frac1u+\frac2{u^3}\right)\sin u
              -\frac{2\cos u}{u^2}+\frac1{\pi^2}>0,
\]
and $g'(\pi)=0$.  Finally, for $u\ge3\pi/2$,
$g'(u)\ge u/\pi^2-1/u-1/u^2>0$.  The last inequality follows because
$u^3/\pi^2-u-1$ increases there and its initial value is
$15\pi/8-1>0$.  This proves \eqref{eq:sinc-parabola}.

Let $Y$ have an even probability density $p$ which is nonincreasing on
the positive half-line.  The Stieltjes measure
$\nu(da)=-2a\,dp(a)$ represents this law as a mixture of uniform laws on
$[-a,a]$; explicitly,
\[
 p(y)=\int_{a\ge |y|}\frac{\nu(da)}{2a},\qquad
 \int\nu(da)=1,
 \qquad
 \mathbb E Y^2=\frac13\int a^2\nu(da).
\]
These identities follow by Stieltjes integration by parts, first for
step functions and then by monotone approximation.  An atom at zero
can be included as a uniform law with $a=0$.  Integrating
\eqref{eq:sinc-parabola} therefore gives
\begin{equation}
 \mathbb E\cos Y\ge\frac12-\frac{3\mathbb E Y^2}{2\pi^2}.
 \label{eq:unimodal-cosine-bound}
\end{equation}

We apply this to $Y=\operatorname{tr}(D_mUXU^*)$.  We give the exact
map that makes its density unimodal.  For distinct entries
$x_1<\cdots<x_m$ of $X$, let $y^{(k)}_1\le\cdots\le y^{(k)}_k$ be the
eigenvalues of the upper left $k\times k$ corner of $UXU^*$.
Olshanski's corner-measure formula and its iterated form
\cite[Proposition 3.1 and Corollary 3.2, pp.~6--7]{olshanski2013}
send Haar measure to uniform Lebesgue measure on the convex
Gelfand--Tsetlin polytope
\[
 \{(y^{(1)},\ldots,y^{(m-1)}):
             y^{(1)}\prec\cdots\prec y^{(m-1)}\prec x\}.
\]
Here $\prec$ means interlacing.  If $s_k=\sum_jy^{(k)}_j$, $s_0=0$,
and $s_m=\operatorname{tr}X$, the $k$th diagonal entry of $UXU^*$ is
$s_k-s_{k-1}$.  Thus the map to $Y$ is precisely the linear functional
\[
 (y^{(1)},\ldots,y^{(m-1)})
       \longmapsto\sum_{k=1}^m h_{m,k}(s_k-s_{k-1}).
\]
For $m=2$ its density is uniform on an interval.  For $m>2$, the
Brunn--Minkowski inequality applied to parallel sections of this
convex polytope says that a positive power of the density is concave
on its support.  Conjugating $D_m$ by the permutation that reverses its
eigenvalues sends $Y$ to $-Y$ without changing its law.  The density is
therefore even, and concavity makes it nonincreasing on the positive
half-line.  Repeated entries of $X$ follow by approximation from
distinct entries: the orbital integrals and their second moments are
continuous in $X$.  Scalar $X$ gives $Y=0$ directly.

For completeness, its variance can be computed without an orbital
integration formula.  The squared absolute entries of one Haar row
are uniform on the probability simplex.  Their second moments are
$2/[m(m+1)]$ on the diagonal and $1/[m(m+1)]$ off the diagonal, as
follows by dividing independent exponential variables by their sum.
Hence, with $A=UXU^*$ and
$V_X=\operatorname{tr}X^2-(\operatorname{tr}X)^2/m$,
\[
 \operatorname{Var}(A_{11})=\frac{V_X}{m(m+1)},\qquad
 \operatorname{Cov}(A_{11},A_{22})
       =-\frac{V_X}{m(m+1)(m-1)}.
\]
The second identity also follows from permutation symmetry and the
constant value of $\operatorname{tr}A$.  Using the two trace identities
for $D_m$ now gives
\[
 \mathbb E Y^2=\frac{m}{m+1}V_X
                 \le\frac{m^2}{4(m+1)}R(X)^2.
\]
To verify the last inequality, put the entries of $X$ in $[b,b+R]$.
Averaging $(x_j-b)(b+R-x_j)\ge0$ shows that their variance is at most
$(\bar x-b)(b+R-\bar x)\le R^2/4$.

Equation \eqref{eq:unimodal-cosine-bound} proves that
$\mathcal O_m(X)\ge0$ whenever
$R(X)^2\le4\pi^2(m+1)/(3m^2)$.  Taking the infimum proves the first
inequality in \eqref{eq:orbital-unimodal-radius}; the second follows
from $(4/3)(1-m^{-2})\ge1$.  For $m=2$, $D_2$ has entries $-1,1$
and $Y$ is uniform on $[-R(X),R(X)]$, so
$\mathcal O_2(X)=\sin R(X)/R(X)$ and $\mathfrak r_2=\pi$.
\end{proof}

Let $L_{\rm T}=\lim_{a\downarrow0}\sqrt a\,m_*(-a)$ and
$L_{\rm J}=\lim_{a\downarrow0}a\,j_*(-a)$ be the extended limits in
the finite-detection theorems, with $H_0(x)=\Xi(x/2)/8$.  Their formulas
give
\[
 L_{\rm T}=\inf_{\gamma>0:\,m_\gamma\ge2}
             \frac{\gamma\mathfrak r_{m_\gamma}}{\sqrt2},
 \qquad
 L_{\rm J}=\frac{\gamma_*^2}{2},
\]
where $\gamma_*$ is the least positive multiple ordinate; an empty
infimum is $+\infty$.

\begin{corollary}[Finite trace bounds for global detection]
\label{cor:trace-global-threshold}
For every finite compression with $\mathcal D_N>0$,
\[
 L_{\rm T}\ge\frac{\pi}{\sqrt{\mathcal D_N}},
 \qquad L_{\rm J}\ge\frac1{\mathcal D_N}.
\]
When $\mathcal D_N=0$, both limits are infinite.
\end{corollary}

\begin{proof}
The signed trace-defect identity gives
$\mathcal D_N\ge2(m_\gamma-1)/\gamma^2$ at every positive multiple
ordinate.  Therefore
\[
 \frac{\gamma^2\mathfrak r_{m_\gamma}^2}{2}
       \ge\frac{(m_\gamma-1)\mathfrak r_{m_\gamma}^2}{\mathcal D_N}
       \ge\frac{\pi^2}{\mathcal D_N}.
\]
Taking the infimum proves the first assertion.  At the least multiple
ordinate the same defect identity gives
$\gamma_*^2/2\ge(m_{\gamma_*}-1)/\mathcal D_N
\ge1/\mathcal D_N$.  A zero trace defect forces every multiplicity
excess to vanish, and the finite-detection formulas then give both
infinite limits.
\end{proof}

\section{Integer certificates and quantitative arithmetic rank}
\label{sec:fd-inertia-density}

The finite trace test admits an exact integer version involving the
Gram and primitive matrices. We first realize the distinct-zero density
in these matrices, and then approximate the individual spectral
coordinates of a large isolated simple family.

\subsection{The exact integer test}
For \(T>0\), let \(n(T)\) count positive zero ordinates at most
\(T\), with multiplicity, and let \(n_{\mathrm{dist}}(T)\) count
these ordinates without multiplicity. Let
\[
 \iota_N(T)=\operatorname{ind}_+(T H_N-G_N),
\]
where \(\operatorname{ind}_+\) is the number of strictly positive
eigenvalues of a finite Hermitian matrix, counted with multiplicity.
The matrices \(G_N,H_N\) are the arithmetic matrices of
Section~\ref{sec:fd-arithmetic}.

\begin{proposition}[A finite inertia certificate]
\label{prop:fd-inertia-complete}
At every height \(T>0\) with \(\Xi(T)\ne0\), the integers
\(\iota_N(T)\) increase with \(N\), are bounded above by
\(n_{\mathrm{dist}}(T)\), and eventually equal that number.
Consequently the packet through \(T\) is simple exactly when
\[
 \iota_N(T)\ge n(T)
\]
holds for one finite \(N\). Simplicity of every nontrivial zero is
equivalent to these finite inequalities at an unbounded sequence of
heights with \(\Xi(T)\ne0\).
\end{proposition}
\begin{proof}
Congruence by \(G_N^{-1/2}\) turns \(T H_N-G_N\) into
\(T C_N-I\) on \(V_N\). Sylvester's law of inertia therefore
identifies \(\iota_N(T)\) with the number of positive eigenvalues
of \(C_N\) greater than \(1/T\). The variational principle bounds
these by the positive eigenvalues of \(C\), and makes them
nondecreasing under the nested projections. The latter eigenvalues
are \(1/\gamma\), one per distinct positive ordinate.

Only finitely many exceed \(1/T\). None equals \(1/T\), by the
boundary condition. Operator-norm convergence \(C_N\to C\)
therefore places every eigenvalue above this threshold on the correct
side for all sufficiently large \(N\); the upper variational bound
excludes additional ones. This proves eventual equality.
The counted and distinct counts agree exactly for a simple packet,
which proves the finite equivalence. An unbounded sequence of simple
packets covers every positive ordinate. Conversely, when every zero
is simple, choose unbounded zero-free contour heights and apply the
finite equivalence at each one.
\end{proof}

This test uses no high-order Taylor coefficients. In an even--odd
basis of \(V_N\), both parity dimensions equal
\(r_N=N2^N-1\). Write
\[
 G_N=\begin{pmatrix}G_+&0\\0&G_-\end{pmatrix},\qquad
 H_N=\begin{pmatrix}0&iK\\-iK^{\mathsf T}&0\end{pmatrix},
\]
with real \(K\); its sign depends only on the order of the parity
basis. Eliminating the negative definite upper block in
\(T H_N-G_N\) gives the real symmetric test
\[
 \iota_N(T)=\operatorname{ind}_+
       \bigl(T^2K^{\mathsf T}G_+^{-1}K-G_-\bigr).
\]
This is an exact Schur-complement identity, not an approximation to
its eigenvalues.

\subsection{A quantitative rank bound from separated coordinates}
\label{subsec:strong-inertia-density}

The positive inertia counts distinct spectral coordinates. A selected
coordinate may have multiplicity greater than one: its sampling weight
is positive and increases the lower quadratic-form estimate. We obtain
the following quantitative count. Put
\[
 d_{\rm MT}=\frac{1+c_{\rm MT}}2
 =\frac54-\frac1{2\sqrt2}\cot\frac1{\sqrt2}
 =0.8362503\ldots.
\]
The constant \(c_{\rm MT}\) and its admissible test approximation
were proved in Section~\ref{sec:survival}.

\begin{theorem}[Quantitative arithmetic realization of distinct zeros]
\label{thm:strong-inertia-density}
For every \(0<\eta<d_{\rm MT}\), there are constants
\(A_\eta,T_\eta>0\) such that
\[
 \operatorname{ind}_+(T H_N-G_N)\ge(d_{\rm MT}-\eta)n(T),
 \qquad N=\lceil A_\eta\log T\rceil,\quad T\ge T_\eta.
\]
The matrices use only the arithmetic Weil pairings. Their dimension is
\(O_\eta(T^{A_\eta\log2}\log T)\), and their prime terms have
\(n<e^{2N}\le e^2T^{2A_\eta}\).
More precisely, for every admissible nonnegative pair test \(g\) used
in Section~\ref{sec:survival}, the same conclusion holds with any
constant smaller than \(1-I_g/2\), where
\[
 I_g=\int_{\mathbb R}g(u)
      \left(1-\left(\frac{\sin\pi u}{\pi u}\right)^2\right)du.
\]
\end{theorem}

\begin{proof}
We first construct enough mutually separated coordinates. Fix an even
nonnegative Schwartz function \(g\) with \(g(0)=1\), Fourier support
strictly inside \((-1,1)\), and \(I_g\) sufficiently close to \(1-c_{\rm MT}\).
The approximations in the proof of
Theorem~\ref{thm:two-thirds-survival} supply such tests. Choose
\(\epsilon,\delta>0\) so that
\(g(u)\ge1-\epsilon\) for \(|u|\le\delta\). The restricted
pair-correlation formula gives
\[
 \sum_{\substack{0<\gamma_j,\gamma_k\le T\\j\ne k}}
 g\bigl(F(\gamma_j)-F(\gamma_k)\bigr)
       =(I_g+o(1))n(T),\qquad F(v)=\frac{v\log v}{2\pi},
\]
with all labels and multiplicities retained. Choose fixed
\(0<\alpha<\beta<1\), with \(\alpha\) small and \(\beta\)
close to one. For large \(T\), \(F'(v)\le\log T\) throughout
\([\alpha T,\beta T]\). Thus every unordered pair of distinct
labels in this interval at distance less than
\(h=\delta/\log T\) contributes at least \(2(1-\epsilon)\)
to the displayed ordered sum. The number of such pairs is at most
\[
 \left(\frac{I_g}{2(1-\epsilon)}+o(1)\right)n(T).
\]
In the finite graph joining precisely these close pairs, delete one
endpoint of every edge. The retained labels are mutually
\(h\)-separated. The Riemann--von Mangoldt count and suitable fixed
choices of the parameters give a set
\[
 Z_T=\{\gamma_1,\ldots,\gamma_M\}\subset[\alpha T,\beta T],
 \qquad M\ge(d_{\rm MT}-\eta/2)n(T).
\]
Its ordinates are distinct. Other occurrences, including repeated
copies at a selected ordinate, remain in the sampling measure.
Keeping \(I_g\) fixed in this argument gives every count below
\((1-I_g/2)n(T)\), which proves the final quantitative dependence
once the approximation argument below is established.

Write
\[
 Q_T(f)=\sum_{\gamma\in\Gamma}m_\gamma
             (T/\gamma-1)|\widehat f(\gamma)|^2,
 \qquad
 \|f\|_T^2=\sum_{\gamma\in\Gamma}m_\gamma
             (1+T/|\gamma|)|\widehat f(\gamma)|^2.
\]
The second expression bounds the absolute sesquilinear form of the
first. If \(\gamma_0>0\) is the least absolute ordinate, then
\(\|f\|_T^2\le(1+T/\gamma_0)q(f)\).
The unit-interval counted zero bound gives
\begin{equation}\label{eq:strong-inertia-count-sums}
 \begin{split}
 n(T)&=O(T\log(3T)),\\
 \sum_{0<|\gamma|\le T}\frac{m_\gamma}{|\gamma|}
       &=O(\log^2(3T)),\\
 \sum_{|\gamma|>T}\frac{m_\gamma}{\gamma^2}
       &=O\bigl(\log(3T)/T\bigr).
 \end{split}
\end{equation}
These follow by summing over unit intervals, with the finite initial
range absorbed in the constants.

Take \(N=\lceil A\log T\rceil\), \(a=\lfloor N/2\rfloor\), and
\[
 g_a(x)=\frac1a(1-|x|/a)_+,
 \qquad
 K_a(t)=\widehat g_a(t)
       =\left(\frac{\sin(at/2)}{at/2}\right)^2.
\]
The integral of \(g_a\) is one. Define
\[
 d_j=K_a(\gamma_j),\qquad
 f_j(x)=e^{i\gamma_jx}g_a(x)-d_jg_a(x),\qquad
 Uc=\sum_{j=1}^M c_jf_j.
\]
Every \(f_j\) has integral zero and support in \([-a,a]\).
The selected sampling matrix is exactly \(I+R-dd^*\), with
\(R_{kj}=K_a(\gamma_k-\gamma_j)\) for \(k\ne j\) and zero diagonal.
The elementary bound
\(|K_a(t)|\le\min(1,4/(a^2t^2))\) and separation give
\[
 \|R\|\le\frac{4}{a^2}\,2h^{-2}\sum_{r\ge1}r^{-2},
 \qquad
 \|d\|^2\le\frac{16M}{a^4\alpha^4T^4}=o(1).
\]
Choose \(A\) sufficiently large that the selected matrix is bounded
below by \(1/2\) for every sufficiently large \(T\). Each selected
coefficient satisfies \(T/\gamma_j-1\ge1/\beta-1\), and its
multiplicity is at least one. Its total contribution to \(Q_T(Uc)\)
is therefore at least
\begin{equation}\label{eq:strong-inertia-positive}
 \frac{1/\beta-1}{4}\|c\|_2^2.
\end{equation}
All remaining coordinates in \((0,T)\) have nonnegative
coefficients and can be omitted from a lower bound.

The negative part is confined to \(\gamma<0\) or \(\gamma>T\).
For these ordinates and all selected \(\gamma_j\),
\[
 |\gamma-\gamma_j|\ge c_0(T+|\gamma|),
 \qquad c_0=\min\{\alpha,(1-\beta)/2\}>0.
\]
For the uncorrected columns \(K_a(\gamma-\gamma_j)\), the row sum
is at most \(CM/[a^2(T+|\gamma|)^2]\). With output weight
\(m_\gamma(1+T/|\gamma|)\), the column sum is at most
\[
 \frac C{a^2}\sum_\gamma
 \frac{m_\gamma(1+T/|\gamma|)}{(T+|\gamma|)^2}
       \le\frac{C\log^2(3T)}{a^2T},
\]
by \eqref{eq:strong-inertia-count-sums}. Weighted Schur therefore
bounds the squared operator norm of the entire negative sampling
part by
\begin{equation}\label{eq:strong-inertia-negative}
 \frac{CM\log^2(3T)}{a^4T^3}
       =O\left(\frac{\log^3(3T)}{a^4T^2}\right)=o(1).
\end{equation}
Here the weighted Schur inequality follows directly from
\(|\sum_j k_{\gamma j}c_j|^2
 \le(\sum_j|k_{\gamma j}|)
        (\sum_j|k_{\gamma j}||c_j|^2)\), followed by summation
with the output weights.
The mean correction is negligible in the full norm, since
\[
 \|g_a\|_T^2\le\frac{16}{a^4}
   \sum_\gamma\frac{m_\gamma(1+T/|\gamma|)}{\gamma^4}
      \le\frac{C(1+T)}{a^4},
 \qquad \|d\|\|g_a\|_T=o(1).
\]
Equations~\eqref{eq:strong-inertia-positive} and
\eqref{eq:strong-inertia-negative} imply
\begin{equation}\label{eq:strong-inertia-margin}
 Q_T(Uc)\ge c_\beta\|c\|_2^2
\end{equation}
with fixed \(c_\beta>0\), for all large \(T\).

We also need a bound on the full coefficient map. For any output
ordinate, at most two selected points lie within distance \(h\);
the remaining distances are bounded below by successive multiples of
\(h\). Thus the unweighted row sum of
\(|K_a(\gamma-\gamma_j)|\) is \(O(1)\), including arbitrarily
close unselected ordinates. For a fixed column, the counted zeros
within distance one contribute \(O(\log T)\). Outside that interval
and with \(|\gamma|\ge\alpha T/2\), the weight
\(1+T/|\gamma|\) is bounded, and unit-distance shells contribute
\[
 O\left(\frac1{a^2}\sum_{r\ge1}
                    \frac{\log(T+r)}{r^2}\right)=O(\log T).
\]
The remaining range \(|\gamma|<\alpha T/2\) contributes
\(O(\log^2(3T)/(a^2T))\), again by
\eqref{eq:strong-inertia-count-sums}. Weighted Schur and the mean
correction consequently give
\begin{equation}\label{eq:strong-inertia-fullnorm}
 \|U\|_{\mathbb C^M\to\|\cdot\|_T}=O(\sqrt{\log T}).
\end{equation}

Finally interpolate into the specified arithmetic space. Let
\(b=2^{-N}\), let \(I_b\) be nodal PL interpolation, and set
\[
 v_j=I_bf_j-g_a\int I_bf_j,\qquad U_Nc=\sum_jc_jv_j.
\]
The knots \(-a,0,a\) belong to the mesh, and \(g_a\) is an exact
mesh function. Hence \(v_j\in V_N\), with integral exactly zero.
Within the mesh cells,
\(\|f_j''\|_2\le CT^2/\sqrt a\). The derivative of a cell
interpolant is the average derivative; the one-dimensional Poincar\'e
inequality on each cell, followed by integration, gives
\[
 \|(I_bf_j-f_j)'\|_2\le CbT^2/\sqrt a,
 \qquad
 \|I_bf_j-f_j\|_2\le Cb^2T^2/\sqrt a.
\]
No derivative jump lies inside a mesh cell. The sampling Sobolev
bound, equivalently \eqref{eq:fd-sampling-sobolev} or its displayed
Plancherel proof, yields
\[
 q(I_bf_j-f_j)\le Ca b^2T^4.
\]
Also \(|\int I_bf_j|\le Cb^2T^2\), so the exact mean correction
satisfies the same bound. Therefore
\[
 \|U_N-U\|_{\mathbb C^M\to\|\cdot\|_T}
   \le C\sqrt{Ma}\,bT^{5/2}
   \le CbT^3\sqrt{a\log(3T)}.
\]
Choose \(A\) also so that \(A\log2>3\). Together with
\eqref{eq:strong-inertia-fullnorm}, this makes the norm of the
pulled-back quadratic-form difference tend to zero:
\[
 \|U_N^*Q_TU_N-U^*Q_TU\|
 \le(\|U_N\|_T+\|U\|_T)\|U_N-U\|_T=o(1).
\]
The positive margin in \eqref{eq:strong-inertia-margin} persists.
Thus the range of \(U_N\) is an \(M\)-dimensional positive subspace
for \(Q_T\) in \(V_N\). Sylvester's law identifies its finite
matrix as \(T H_N-G_N\), and proves the stated inertia bound.
The dimension and prime range follow directly from
\(\dim V_N=2N2^N-2\) and \(N\le A\log T+1\).
\end{proof}

\subsection{Approximation of isolated simple coordinates}

\begin{theorem}[Localized spectral coordinates in logarithmic support]
\label{thm:fd-inertia-density}
For every \(0<\eta<c_{\rm MT}\) and \(\epsilon>0\), there are
constants \(A,\alpha,\beta,C,T_0>0\), with \(0<\alpha<\beta<1\),
such that for every \(T\ge T_0\) there is a set of simple ordinates
\(Z_T=\{\gamma_1,\ldots,\gamma_M\}\subset[\alpha T,\beta T]\)
with \(M\ge(c_{\rm MT}-\eta)n(T)\) and a map
\(U_N:\mathbb C^M\to V_N\), \(N=\lceil A\log T\rceil\), satisfying
\[
 \|U_N-E\|\le\epsilon,
 \qquad
 \|CU_N-U_N\Lambda\|\le\frac{C}{T^2\log T}.
\]
Here \(E\) sends the standard basis to the selected unit spectral
coordinates in \(\ell^2(\Gamma,m)\), and
\(\Lambda=\operatorname{diag}(1/\gamma_j)\); both norms are operator
norms into the Weil completion. The same maps give
\(\iota_N(T)\ge(c_{\rm MT}-\eta)n(T)\).
The dimension and prime range have the bounds of
Theorem~\ref{thm:strong-inertia-density}.
\end{theorem}

\begin{proof}
We give the approximation estimates explicitly. The proof of
Theorem~\ref{thm:two-thirds-survival}, before its heat continuation
step, supplies isolated simple ordinates with the following stronger
data. After restricting to finitely many adjacent dyadic intervals,
there are constants
\[
 0<\alpha<\beta<1,\qquad \delta>0,\qquad B>0
\]
and, for every sufficiently large \(T\), a set
\(Z_T=\{\gamma_1,\ldots,\gamma_M\}\subset[\alpha T,\beta T]\)
with
\begin{equation}\label{eq:fd-selected-data}
 \begin{split}
 M&\ge(c_{\rm MT}-\eta/2)n(T),\\
 |\gamma-\gamma_j|&\ge h:=\delta/\log T
     \quad(\gamma\ne\gamma_j),\\
 S_j:=\sum_{\gamma\ne\gamma_j}
           \frac{m_\gamma}{(\gamma-\gamma_j)^2}
     &\le B(\log T)^2.
 \end{split}
\end{equation}
The second line includes every other signed zero occurrence.
Here is the passage from the dyadic statement. Use the isolation
construction with a smaller error parameter on the finitely many
blocks ending at \(\beta T\) and beginning above \(\alpha T\).
Choose \(\beta\) sufficiently close to one and the lower endpoint
sufficiently small that the Riemann--von Mangoldt count loses less
than the allowed fraction. The proof already removes strips of width
\(\log T\) at each block endpoint, so adjacent selected blocks are
separated as well. There are only finitely many blocks, and their
heights are fixed positive multiples of \(T\). Their gap and
inverse-square constants may thus be replaced by common constants.
Changing from the heat coordinate \(x=2\gamma\) divides the gaps by
two and multiplies the inverse-square sums by four. This proves
\eqref{eq:fd-selected-data} with the displayed zero-ordinate
normalization.

Let \(\varphi\) be the convolution of eight copies of
\(4\mathbf1_{[-1/8,1/8]}\). It has integral one, support
\([-1,1]\), and square-integrable derivatives through order three.
Its Fourier transform in the convention used here is
\[
 K(v)=\left(\frac{\sin(v/8)}{v/8}\right)^8,
 \qquad K(0)=1.
\]
For every nonzero real \(v\),
\begin{equation}\label{eq:fd-bump-decay}
 |K(v)|\le\frac{64}{v^2},\qquad
 |vK(v)|\le\frac{512}{v^2}.
\end{equation}
Indeed, \(|K(v)|\le\min\{1,(8/|v|)^8\}\); split at
\(|v|=8\) to obtain both inequalities.

Put \(\ell=\log T\), choose \(N=\lceil A\ell\rceil\), and
let \(L=N/2\). For each selected ordinate define compact functions
\[
 g_j(x)=L^{-1}e^{i\gamma_jx}\varphi(x/L),\qquad
 f_j(x)=\frac1{i\gamma_j}\,g_j'(x).
\]
They are supported on \([-L,L]\), and \(f_j\) has integral zero.
Their Fourier transforms and compact primitives satisfy
\[
 \widehat f_j(\gamma)
    =\frac{\gamma}{\gamma_j}K(L(\gamma-\gamma_j)),
 \qquad C f_j=\frac{g_j}{\gamma_j}.
\]
The second identity follows either from the compact primitive
\(i\int_{-N}^x f_j\) or from division of the Fourier transform by
\(\gamma\).

Use the unitary spectral coordinates
\(f\mapsto(\sqrt{m_\gamma}\widehat f(\gamma))_\gamma\).
Let \(E:\mathbb C^M\to\mathcal H\) place the standard basis
at the selected simple coordinates, and let
\(U:\mathbb C^M\to\mathcal H\) have columns \(f_j\).
Define two matrices, omitting their selected diagonal entries,
\[
 \mathcal K_{\gamma j}=\sqrt{m_\gamma}
       K(L(\gamma-\gamma_j)),\qquad
 \mathcal L_{\gamma j}=\sqrt{m_\gamma}
       (\gamma-\gamma_j)K(L(\gamma-\gamma_j)).
\]
For every row ordinate, the selected ordinates on either side are
\(h\)-separated, and every retained distance is at least \(h\).
Consequently
\[
 \sum_{j:\gamma_j\ne\gamma}(\gamma-\gamma_j)^{-2}
       \le2h^{-2}\sum_{r\ge1}r^{-2}<4h^{-2}.
\]
Apply the weighted Schur test with row weights
\(\sqrt{m_\gamma}\) and column weights one. The row and column
bounds for \(\mathcal K\) are respectively
\(256/(L^2h^2)\) and \(64B\ell^2/L^2\). For
\(\mathcal L\) they are \(2048/(L^3h^2)\) and
\(512B\ell^2/L^3\). Thus
\begin{equation}\label{eq:fd-bump-matrix-bounds}
 \|\mathcal K\|\le\frac{128\sqrt B}{\delta}
                 \frac{\ell^2}{L^2},\qquad
 \|\mathcal L\|\le\frac{1024\sqrt B}{\delta}
                 \frac{\ell^2}{L^3}.
\end{equation}
The Schur test follows directly by applying Cauchy's inequality in
each row with these weights and then summing the column bounds.
In particular its use here includes every complementary multiplicity.

Let \(\Lambda=\operatorname{diag}(1/\gamma_j)\). The exact
Fourier formulas give
\[
 U-E=\mathcal K+\mathcal L\operatorname{diag}(1/\gamma_j),
 \qquad C U-U\Lambda
       =-\mathcal L\operatorname{diag}(1/\gamma_j^2).
\]
Since \(L\ge A\ell/2\), these imply
\begin{align}
 \|U-E\|&\le\frac{512\sqrt B}{\delta A^2}
       +\frac{8192\sqrt B}{\delta\alpha A^3T\ell},
       \label{eq:fd-bump-isometry}\\
 \|CU-U\Lambda\|&\le
       \frac{8192\sqrt B}{\delta\alpha^2A^3T^2\ell}.
       \label{eq:fd-bump-residual}
\end{align}
These estimates do not use a lower bound for a zeta derivative.

Put \(g=1/\beta-1>0\), \(M_0=1/\alpha-1\), and
\(\varepsilon_0=\min\{1/4,g/(12M_0)\}\).
Choose \(A\ge8\) sufficiently large and then \(T\) sufficiently
large that \(\|U-E\|\le\min(\varepsilon_0,\epsilon/2)\).
It follows that
\(\|U^*U-I\|\le3\varepsilon_0\).
For a unit vector \(a\in\mathbb C^M\), write
\(R=CU-U\Lambda\). Taking real parts in the exact identity
for the quadratic form gives
\begin{align*}
 \langle(C-T^{-1})Ua,Ua\rangle
 &\ge\frac gT-
       \|U^*U-I\|\frac{M_0}{T}-\|U\|\|R\|\\
 &\ge\frac{g}{2T}
\end{align*}
for all sufficiently large \(T\), by
\eqref{eq:fd-bump-residual}. Also \(U\) is injective, since it
is within \(1/4\) of the isometry \(E\).

It remains to place this whole positive subspace in the prescribed
finite piecewise-linear space. We include the needed norm estimates.
For every \(f\in H^1(\mathbb R)\) supported on \([-N,N]\),
\(N\ge1\), we have
\begin{equation}\label{eq:fd-sampling-sobolev}
 \|f\|_q\le50N\|f\|_{H^1}.
\end{equation}
Here are explicit constants. The completion estimates in
Lemma~\ref{lem:fd-multiplicity-log}, now at an arbitrary real ordinate
\(v\), give
\[
 \Re\frac{\xi'}{\xi}(2+iv)
 <\frac{13}{5}+\frac12\log(|v|+2)
 <4\log(|v|+3).
\]
Each zero in \([k,k+1)\) contributes at least \(3/5\) to the
positive logarithmic-derivative sum at \(v=k+1/2\). The counted
zeros in this interval therefore number at most
\(10\log(|k|+4)\). On a unit interval,
\(\sup|F|^2\le2\int(|F|^2+|F'|^2)\): choose a point where
\(|F|^2\) is no larger than its integral and integrate \(F'\)
from that point. For \(F=\widehat f\), summing gives
\[
 q(f,f)\le60\int_{\mathbb R}(1+v^2)
                   (|F(v)|^2+|F'(v)|^2)\,dv,
\]
where we used \(\log(|v|+5)\le3(1+v^2)\).
Plancherel identifies the integral as
\[
 2\pi\{\|f\|_2^2+\|f'\|_2^2+
                  \|xf\|_2^2+\|(xf)'\|_2^2\}
 \le6\pi(1+N^2)\|f\|_{H^1}^2.
\]
Thus \(q(f,f)\le720\pi N^2\|f\|_{H^1}^2
<2500N^2\|f\|_{H^1}^2\), proving
\eqref{eq:fd-sampling-sobolev}.

Write \(b=2^{-N}\), and let \(I_b f\) be linear interpolation
on the mesh in \([-N,N]\). The cellwise Poincar\'e inequality gives
\[
 \|(I_bf-f)'\|_2\le b\|f''\|_2,\qquad
 \|I_bf-f\|_2\le b^2\|f''\|_2.
\]
Let \(v_0(x)=(1-|x|)_+\), which belongs to the mesh space and
has integral one. Correct the mean by defining
\[
 J_bf=I_bf-v_0\int_{-N}^{N}I_bf.
\]
For mean-zero \(f\),
\[
 \left|\int I_bf\right|\le\sqrt{2N}\,b^2\|f''\|_2.
\]
Since \(b\sqrt N\le1\) and
\(\|v_0\|_{H^1}=\sqrt{8/3}\), the preceding three estimates imply
\begin{equation}\label{eq:fd-interpolation-h1}
 \|J_bf-f\|_{H^1}\le5b\|f''\|_2.
\end{equation}
The corrected interpolant belongs to \(V_N\).

For our functions, direct differentiation gives
\[
 f_j''(x)=e^{i\gamma_jx}
 \left[-\frac{\gamma_j^2}{L}\varphi(x/L)
       +\frac{3i\gamma_j}{L^2}\varphi'(x/L)
       +\frac3{L^3}\varphi''(x/L)
       +\frac1{i\gamma_jL^4}\varphi'''(x/L)\right].
\]
For large \(T\), \(1\le\gamma_j\le T\) and \(L\ge1\),
so
\[
 \|f_j''\|_2\le C_\varphi T^2/\sqrt L,\qquad
 C_\varphi=\|\varphi\|_2+3\|\varphi'\|_2+
             3\|\varphi''\|_2+\|\varphi'''\|_2\le1458.
\]
For the last bound, each original box has \(L^1\) norm one and
\(L^2\) norm two, and its distributional derivative has total
variation eight. Young's convolution inequality therefore gives
\(\|\varphi^{(r)}\|_2\le2\cdot8^r\) for \(0\le r\le3\).
Let \(U_N\) have columns \(J_bf_j\). Combining
\eqref{eq:fd-sampling-sobolev} and
\eqref{eq:fd-interpolation-h1}, then using
\(M\le C T\log T\), gives
\begin{align*}
 \|U_N-U\|
 &\le\sqrt M\max_j\|J_bf_j-f_j\|_q\\
 &\le C_\eta\,2^{-N}T^{5/2}\log T
 \le C_\eta T^{5/2-A\log2}\log T=o(T^{-1}).
\end{align*}
Here \(A\ge8\) makes the final exponent strictly less than
\(-1\). The form difference is bounded by
\[
 \|U_N^*(C-T^{-1})U_N-U^*(C-T^{-1})U\|
 \le(\|C\|+T^{-1})
       (\|U_N\|+\|U\|)\|U_N-U\|=o(T^{-1}).
\]
Consequently the form \(q(Cf,f)-T^{-1}q(f,f)\) is positive
definite on the \(M\)-dimensional range of \(U_N\) for all large
\(T\). The finite variational principle gives
\(\iota_N(T)\ge M\). The same interpolation estimate is
\(o(T^{-2}/\log T)\), since \(A\ge8\). Thus
\(\|U_N-E\|\le\epsilon\) for large \(T\), and
\[
 \|CU_N-U_N\Lambda\|
 \le\|CU-U\Lambda\|+(\|C\|+\|\Lambda\|)\|U_N-U\|
 \le C/(T^2\log T).
\]
This proves the coordinate and residual assertions. Finally, \(N\le A\log T+1\) and the explicit arithmetic formula
prove the stated prime range and matrix dimension.
\end{proof}

This estimate applies at every sufficiently large height and fixes
the support and mesh needed to realize the proved density in the
arithmetic matrices. The integer deficit
\(n(T)-\iota_N(T)\) is still permitted to be of order \(n(T)\).
Its vanishing at unbounded heights is the full simplicity gate in
Proposition~\ref{prop:fd-inertia-complete}; the approximation estimates
above do not establish that vanishing for the complementary set.

\subsection{The exact trace integral and the absolute integer scale}

Define the nonnegative integer deficit
\[
 A_N(T)=n(T)-\iota_N(T).
\]
The trace and inertia certificates are related by the exact identity
\begin{equation}\label{eq:fd-trace-inertia-integral}
 D_{N,k}=4k\int_0^\infty A_N(t)t^{-2k-1}\,dt,
 \qquad k\ge1.
\end{equation}
To prove it, parity pairs the nonzero eigenvalues of \(C_N\)
as \(\pm\lambda_{N,j}\), with \(\lambda_{N,j}>0\). For a
positive number \(\lambda\),
\[
 \lambda^{2k}=2k\int_{1/\lambda}^{\infty}t^{-2k-1}\,dt.
\]
Apply this identity first to the counted reciprocal ordinates in
\(S_k\), and then to the positive compression eigenvalues in
\(\operatorname{tr}(C_N^{2k})\). Both integrals converge:
the counting functions vanish near zero and grow at most as
\(O(t\log(t+3))\) at infinity. Subtraction proves
\eqref{eq:fd-trace-inertia-integral}; choices at the isolated
threshold points do not change the integrals. At zero-free heights,
\(A_N(T)\) decreases to the multiplicity excess
\(n(T)-n_{\mathrm{dist}}(T)\).

There is also an exact finite-dimensional comparison with the
positive-kernel replacement of
Section~\ref{sec:fd-positive-surgery}. Its prescribed simple
ordinates satisfy \(0<u<v\): the replacement removes the pair
\(\pm u\) and adds one occurrence at each of \(\pm v\).
Use the same compact mean-zero basis to form both Gram matrices.
For \(f,g\) in its span, write \(F=\widehat f\), \(G=\widehat g\)
and \(\mathcal P f=i\int_{-N}^{x}f\). The exact change of the
form \(T H_N-G_N\) is
\begin{align}
 &\widetilde q(T\mathcal P f-f,g)-q(T\mathcal P f-f,g)
 \notag\\
 &\quad=\sum_{\sigma=\pm1}\left[
  \left(\frac{T}{\sigma v}-1\right)
            F(\sigma v)\overline{G(\sigma v)}
 -\left(\frac{T}{\sigma u}-1\right)
            F(\sigma u)\overline{G(\sigma u)}\right].
 \label{eq:fd-inertia-surgery}
\end{align}
This follows by applying the already proved four-evaluation change
of \(q\) and using
\(\widehat{\mathcal P f}(\gamma)=F(\gamma)/\gamma\).
It compares forms on their common test functions and makes no
identification of the completed Hilbert spaces.

For \(T>v\), the terms at \(+v\) and \(-u\) have positive
coefficients, while those at \(-v\) and \(+u\) have negative
coefficients. Thus the finite matrix change is a sum of at most two
positive rank-one forms minus at most two positive rank-one forms.
The variational principle, applied to the kernels of these forms,
gives
\[
 |\widetilde\iota_N(T)-\iota_N(T)|\le2.
\]
For zero-free \(T>v\), the total counted zero numbers are equal,
whereas the distinct count for the replacement is smaller by one.
The eventual inertia counts therefore give
\[
 \lim_N\widetilde A_N(T)=\lim_N A_N(T)+1.
\]
The sampling construction and compact projections also apply to the
replacement. Indeed, its unit-interval bounds are preserved by the
finite change. A compactly supported nonzero function has a Fourier
transform of exponential type and therefore only \(O(R)\) zeros in
a disk of radius \(R\), by Jensen's formula centered at a nonzero
value. The unchanged positive-density set of simple ordinates has
order \(R\log R\), so the modified form is strictly positive on
such test functions. To verify density, an orthogonal vector would
define a tempered distribution supported on the modified ordinates
whose inverse Fourier transform annihilates all compact mean-zero
tests. That inverse transform must be constant, so the original
distribution is supported at zero. There is no zero ordinate, hence
the orthogonal vector vanishes. For support in \([-N,N]\), the
four-evaluation update and Cauchy's inequality give
\[
 \widetilde q(f)\le q(f)+|F(v)|^2+|F(-v)|^2
 \le q(f)+4N\|f\|_2^2.
\]
The sharper bound in the proof of
\eqref{eq:fd-sampling-sobolev} was
\(q(f)\le720\pi N^2\|f\|_{H^1}^2\). Since
\(720\pi+4<2500\) and \(N\ge1\), the same displayed
\(50N\) norm bound holds for \(\widetilde q\).
Mean-corrected interpolation therefore shows that the union of the
finite spaces is dense as well. This justifies the eventual inertia count
for the replacement. In particular the positive-density lower estimate
is preserved after changing its error parameter, since the finite
change is bounded by two and \(n(T)\to\infty\).

Equations \eqref{eq:fd-trace-inertia-integral} and
\eqref{eq:fd-inertia-surgery} specify the absolute scale of the
remaining estimate. The positive-density approximation controls a
fraction of the integer deficit. The displayed finite replacement
preserves that fractional conclusion while changing its eventual
integer value by one, and its change on the arithmetic side is exactly
the oscillatory Weil correction already computed. An estimate proving
zero deficit must therefore control this integer scale for the actual
prime form.

\section{A uniform arithmetic lower bound from the cosine pair test}
\label{sec:cosine-lower-rate}

The preceding arithmetic rank construction can also be used in a
uniform sampling estimate. We retain the exact constant
\[
 d=d_{\rm MT}=\frac54-\frac1{2\sqrt2}\cot\frac1{\sqrt2},
 \qquad p_{\rm MT}=d^{-1}=1.1958141\ldots.
\]
Mori proved the lower rate with exponent \(6/5\) by a uniform
interval version of Montgomery's pair argument and quantitative
sampling \cite[Theorem~4.3, pp.~11--13]{mori-rh-alone2026}.
We give the cosine-weight variant, including the uniform errors and
the precise sampling reduction. The zero-location input throughout
is the established theorem already used for the Weil sampling identity.

\begin{theorem}[An all-window lower rate]
\label{thm:cosine-lower-rate}
There are constants \(C,L_0>0\) such that, for \(L\ge L_0\),
\[
 a_1(L)\ge
 \exp\{-C e^{2p_{\rm MT}L}L(1+\log L)\}.
\]
The same lower bound holds for \(b_L\) and \(a_2(L)\).
Equivalently, with \(\mu=e^{2L}\), the exponent is
\(C\mu^{p_{\rm MT}}\log\mu(1+\log\log\mu)\).
\end{theorem}

\subsection{The uniform pair estimate}

Put \(c=1/\sqrt2\), and let
\[
 f_0(v)=\frac{\cos(\sqrt2v)}{\sqrt2\sin c}
              \mathbf1_{[-1/2,1/2]}(v),\qquad
 w=f_0*\widetilde f_0,
\]
where \(\widetilde f_0(v)=f_0(-v)\). Write
\[
 A=\int f_0^2=w(0),\qquad
 B=\int |v|w(v)\,dv.
\]
Then \(w\) is even and nonnegative, is supported in \([-1,1]\),
\(\int w=1\), and \(0\le w\le A\) by Cauchy's inequality.
The exact energy calculation in Section~\ref{sec:survival} gives
\begin{equation}\label{eq:cosine-energy}
 A+B=\frac12+c\cot c=3-2d.
\end{equation}
We shall use the elementary bounds
\begin{equation}\label{eq:cosine-coarse}
 1\le A\le\frac{51}{50},\qquad
 0<B\le\frac13,\qquad
 \frac56<d\le\frac78.
\end{equation}
Here \(A\ge1\) follows from \(\int f_0=1\). The variance of
\(f_0\) under uniform measure on its support is \(A-1\), and is
at most one quarter of its squared range. That range is
\(\tan(c/2)/\sqrt2\). Since
\(\tan(c/2)\le (c/2)/(1-c^2/8)=8/(15\sqrt2)\),
we get \(A-1\le4/225<1/50\).
The Taylor bounds
\(\cos c\le73/96\), \(\sin c\ge11c/12\) show
\(c\cot c\le73/88<5/6\). Thus \(A+B<4/3\),
which gives \(B<1/3\) and \(d>5/6\).
Conversely \(c\cot c\ge\cos c\ge3/4\), giving \(d\le7/8\).

Let \(J\) be any bounded interval, and define
\[
 \ell=|J|,\qquad H=\max\{2,\sup_{t\in J}|t|\},
 \qquad L_J=\log(2+H).
\]
Let \(N_J\) count all zero labels in \(J\), including multiplicity,
and let \(P_h(J)\) count unordered pairs of distinct labels at distance
less than \(h\); equal ordinates are allowed. There is an absolute
constant \(C_0\) such that, whenever
\(\kappa\) is sufficiently large,
\(L_J\ge\kappa\), \(\ell\ge e^\kappa\), and
\(h=\kappa^{-3/2}\),
\begin{equation}\label{eq:cosine-uniform-pair}
 N_J+2P_h(J)
 \le\frac{\ell}{2\pi}
       \left(\frac{A L_J^2}{\kappa}+B\kappa\right)
          (1+C_0/\kappa).
\end{equation}
All constants are independent of the position of \(J\).

Here are the arithmetic inputs and the proof of this uniformity.
For \(x\ge1\) set
\[
 \mathcal L_x(t)=2\sum_\gamma
           \frac{m_\gamma x^{i\gamma}}{1+(t-\gamma)^2}.
\]
Montgomery's explicit-formula decomposition is
\(\mathcal L_{e^u}=\mathcal A_u+\mathcal B_u+\mathcal E_u\),
where
\[
 \mathcal A_u(t)=-e^{iut}\sum_{n\ge1}
       \frac{\Lambda(n)}{n^{1/2+it}}
                  \min\{ne^{-u},e^u/n\},
 \qquad
 \mathcal B_u(t)=e^{-u+iut}
           \left(-\frac{\zeta'}\zeta(-1/2+it)\right).
\]
For an interval of length \(r\) with
\(\log(|t|+2)\le V\), the following bounds hold with absolute
constants:
\begin{equation}\label{eq:cosine-Montgomery-input}
 \begin{split}
 \int|\mathcal A_u(t)|^2dt
    &\le r(u+C)+Ce^u(u+1),\\
 |\mathcal B_u(t)|&\le e^{-u}(V+C),\\
 \int_{\mathbb R}|\mathcal E_u(t)|^2dt&\le Ce^u.
 \end{split}
\end{equation}
These precise interfaces, with the decomposition derived from
Landau's formula, are proved in
\cite[Lemmas~3.3--3.5, pp.~6--7]{mori-rh-alone2026}.
For clarity, the mean-value step uses the Montgomery--Vaughan
weighted Hilbert inequality on frequencies \(\log n\), whose local
spacing is at least \(1/(2n)\). The coefficient square sums satisfy
\(\sum|a_n(e^u)|^2\le u+C\) and
\(\sum n|a_n(e^u)|^2\le Ce^u(u+1)\), by partial summation
from the prime number theorem. The functional equation gives the
second bound, and the pole and trivial-zero terms give the third.
No close-pair hypothesis is used in these bounds.

Enlarge \(J\) by \(Y=\ell^{1/4}\) on each side, obtaining
\(J^+\). Put \(\ell^+=\ell+2Y\) and
\(L^+=\log(2+H+Y)\). On functions of \((u,t)\) define the
squared Hilbert norm
\[
 \|F\|_w^2=\frac2\kappa\int_0^\kappa w(u/\kappa)
                       \int_{J^+}|F(u,t)|^2dt\,du.
\]
The prime term in \eqref{eq:cosine-Montgomery-input} gives
\[
 \|\mathcal A\|_w^2
       \le \ell^+(B\kappa+C)+Ce^\kappa.
\]
Indeed \((2/\kappa)\int_0^\kappa u w(u/\kappa)du=B\kappa\),
and
\((2/\kappa)\int_0^\kappa w(u/\kappa)e^u(u+1)du
 \le 2Ae^\kappa\).
Similarly,
\[
 \|\mathcal B\|_w^2
       \le\frac{A\ell^+(L^++C)^2}{\kappa},
 \qquad
 \|\mathcal E\|_w^2\le Ce^\kappa/\kappa.
\]
It is essential to integrate the prime--archimedean cross term
before applying an inequality to the full \(u\)-norm. Cauchy's
inequality in \(t\) alone bounds its absolute value by
\[
 \frac{C(L^++C)\sqrt{\ell^+}}{\kappa}
 \int_0^\kappa e^{-u}
   \left(\sqrt{\ell^+(u+C)}+e^{u/2}\sqrt{u+1}\right)du
 \le\frac{C\ell^+(L^++C)}{\kappa}.
\]
The integrals on \((0,\infty)\) are finite absolute constants.

Let
\(Q=\ell(A L_J^2/\kappa+B\kappa)\).
Since \(H\ge\ell/2\),
\(L^+-L_J\le2\ell^{-3/4}\) and
\(\ell^+/\ell=1+2\ell^{-3/4}\).
Using \(L_J\ge\kappa\) and \(\ell\ge e^\kappa\), the preceding
bounds yield
\[
 \|\mathcal A+\mathcal B\|_w^2
       \le Q(1+C/\kappa),\qquad
 \|\mathcal E\|_w^2\le CQ/\kappa^2.
\]
The triangle inequality therefore gives
\begin{equation}\label{eq:cosine-average-upper}
 \|\mathcal L\|_w^2\le Q(1+C/\kappa).
\end{equation}

Expanding the square gives the exact positive pair average
\[
 \|\mathcal L\|_w^2
 =4\sum_{\gamma,\gamma'}m_\gamma m_{\gamma'}
 \widehat w_1\bigl(\kappa(\gamma-\gamma')\bigr)
 \int_{J^+}\frac{dt}
 {(1+(t-\gamma)^2)(1+(t-\gamma')^2)},
\]
where \(\widehat w_1(v)=\int w(y)e^{ivy}dy
=|\int f_0(y)e^{ivy}dy|^2\ge0\).
The double sum converges absolutely on the bounded integration
interval. For \(|v|\le\kappa h=\kappa^{-1/2}\),
\(\widehat w_1(v)\ge1-v^2/2\), since \(w\) is a probability
density supported in \([-1,1]\). For \(\gamma,\gamma'\in J\),
the Cauchy-kernel integral is at least
\[
 \frac{2\pi}{4+(\gamma-\gamma')^2}-\frac{2}{3Y^3}.
\]
Thus every retained ordered pair at distance less than \(h\),
including equal labels, contributes at least
\(2\pi(1-C/\kappa)\). Discarding all other nonnegative terms
and using \eqref{eq:cosine-average-upper} proves
\eqref{eq:cosine-uniform-pair}.

\subsection{Packing with the cosine constant}

The uniform interval count from Riemann--von Mangoldt is
\[
 N_J\ge\ell\left(\frac{L_J}{2\pi}-b_0\right)-c_0 L_J
\]
for absolute constants \(b_0,c_0\). Its proof averages \(\log|t|\)
on \(J\), with endpoint errors \(O(L_J)\); this gives a constant
loss in the average logarithm even for intervals containing zero
\cite[Lemma~2.2, pp.~7--8]{mori-sharp2026}.
Let \(N_0=\ell L_J/(2\pi)\), \(\beta=\kappa/L_J\), and
\(\rho=A/\beta+B\beta\). Then
\[
 N_J\ge N_0(1-C/\kappa),\qquad
 N_J+2P_h(J)\le\rho N_0(1+C/\kappa).
\]
In the close-pair graph, deletion of one endpoint per edge gives
an independent set of size at least \(N_J-P_h(J)\). Ordering its
vertices uniformly at random and retaining each vertex earlier than
all its neighbors gives expected size
\(\sum_v(1+\deg v)^{-1}\ge N_J^2/(N_J+2P_h(J))\).
Therefore the maximal size \(p_h(J)\) of an \(h\)-separated
subset is bounded below by the larger of these two expressions.

The exact scalar comparison needed here is
\begin{equation}\label{eq:cosine-scalar}
 \max\{3/2-\rho/2,1/\rho\}\ge d\beta
                     \qquad(0<\beta\le1).
\end{equation}
For \(\beta\ge3/5\), the difference in its first branch is
\[
 \frac{(1-\beta)((3-A)\beta-A)}{2\beta}\ge0,
\]
by \(A\le51/50\). For \(\beta\le3/5\),
\[
 \rho^{-1}=\frac\beta{A+B\beta^2}
       \ge\frac{50}{57}\beta>\frac78\beta\ge d\beta.
\]
In the first range \(\rho\) is bounded above and \(\beta\)
is bounded away from zero, so the preceding relative errors cost
only \(C/\kappa\) relative to \(d\beta\). In the second range
the reciprocal branch carries a relative \(O(1/\kappa)\) error
directly. We obtain, uniformly in the position of \(J\),
\begin{equation}\label{eq:cosine-separated}
 p_h(J)\ge\frac{d\kappa}{2\pi}\ell(1-C_1/\kappa),
 \quad L_J\ge\kappa,\quad\ell\ge e^\kappa,
 \quad h=\kappa^{-3/2}.
\end{equation}

\subsection{Quantitative sampling and the exact window map}

We recall the quantitative sampling consequence in the form needed
here, and specify its classical input. If a \(q\)-separated set
\(\Lambda\subset\mathbb R\) satisfies
\[
 \#(\Lambda\cap I)\ge(1+\epsilon)|I|-D
       \quad\hbox{for every bounded interval }I,
 \qquad0<q,\epsilon\le1,
\]
then its lower sampling constant on \(PW_\pi\) satisfies
\begin{equation}\label{eq:cosine-VH}
 \inf_{0\ne F\in PW_\pi}
 \frac{\sum_{x\in\Lambda}|F(x)|^2}{\|F\|_2^2}
 \ge\exp\{-C(D+1)(1+\log(1/q)+\log(1/\epsilon))\}.
\end{equation}
Here is the reduction to Voss's theorem. Put
\(r=(D+1)/(2(1+\epsilon))\) and \(\rho_0=1+\epsilon/2\).
For any finite set of integers, the union of the intervals
\([n/\rho_0-r,n/\rho_0+r]\) contains enough points of
\(\Lambda\) for an injective matching: a component containing
\(m\) centers has length at least \((m-1)/\rho_0+2r\), and the
count lower bound is at least \(m\). Finite Hall matching and a
locally finite diagonal limit give points \(x_n\in\Lambda\)
with \(|x_n-n/\rho_0|\le r\).

Set \(t_n=\rho_0x_n\),
\(M=\lfloor\rho_0r\rfloor+1\), and
\(\delta_0=\tfrac12\min(1,\rho_0q)\).
Then \(|t_n-n|<M\) and the separation is strictly greater than
\(\delta_0\). For \(F\in PW_\pi\),
\(G(x)=F(x/\rho_0)\) has type \(\tau=\pi/\rho_0<\pi\)
and \(\|G\|_2^2=\rho_0\|F\|_2^2\).
The sampling sequence is square summable: integrating the fundamental
theorem of calculus over the disjoint intervals of length
\(\delta_0\) centered at \(t_n\) gives
\[
 \sum_n|G(t_n)|^2\le
 \frac{2}{\delta_0}\|G\|_2^2+\delta_0\|G'\|_2^2
 \le(2/\delta_0+\delta_0\tau^2)\|G\|_2^2.
\]
Voss's published quantitative bound
\cite[Theorem~2 and Corollary~1, pp.~195--196]{voss1999} is
\[
 \sum_n|G(t_n)|^2\ge\frac{H_V^2}{2M+1}\|G\|_2^2,
\]
where
\[
 H_V=\delta_0^{2M}
 \left(\frac1{\pi-\tau}+\frac{21}{10}\right)^{-4M}
 \left(\frac2\pi\right)^{2M-1}
 \frac{M}{3e^{4M}}\left(4+\frac1M\right)^{-4M}.
\]
Since \(M=O(D+1)\), \(\delta_0\ge q/2\), and
\(\pi-\tau\ge\pi\epsilon/3\), taking logarithms proves
\eqref{eq:cosine-VH}. This reproduces the Voss--Hall reduction of
\cite[Proposition~2.4, p.~8]{mori-sharp2026}, with strict endpoint
parameters and with only an absolute constant retained.

There is one \(h\)-separated set of zero ordinates \(\Gamma_h\)
whose count in every bounded interval is at least \(p_h(J)-1\).
Indeed, left-to-right greedy selection on a finite interval is
optimal there. Its restriction to a subinterval loses at most one
point compared with starting the greedy selection at that
subinterval. Letting the finite endpoints tend to infinity along a
locally finite diagonal subsequence gives the asserted set.

Now put \(s=2L=\log\mu\), and choose
\[
 \kappa=\frac{s+1+2\pi}{d}+C_1+1,
 \qquad h=\kappa^{-3/2}.
\]
In particular \(d(\kappa-C_1)=s+1+2\pi+d\).
For sufficiently large \(s\), equation~\eqref{eq:cosine-separated}
gives
\(p_h(J)\ge((s+1)/(2\pi)+1)|J|\) whenever
\(L_J\ge\kappa\) and \(|J|\ge e^\kappa\).
Write \(r_s=(s+1)/(2\pi)\) and \(D=2r_s e^\kappa\).
On these long intervals the one-point greedy loss is paid by the
extra density one. On shorter intervals with \(L_J\ge\kappa\),
the possible deficit is at most \(r_s e^\kappa\). Intervals
with \(L_J<\kappa\) lie inside \((-e^\kappa,e^\kappa)\),
so their possible deficit is at most \(D\). Consequently
\[
 \#(\Gamma_h\cap J)\ge r_s|J|-D
                      \quad\hbox{for every bounded interval }J.
\]
Scale to \(\Lambda=s\Gamma_h/(2\pi)\). This set has separation
\(q=sh/(2\pi)\), satisfies the sampling count with
\(\epsilon=1/s\), and has
\[
 D=O(s\mu^{1/d}),\qquad
 1+\log(1/q)+\log(1/\epsilon)=O(1+\log s).
\]
Equation~\eqref{eq:cosine-VH} bounds its sampling constant below by
\(\exp\{-Cs\mu^{1/d}(1+\log s)\}\).

For the final normalization, take a compact smooth physical function
\(f\) supported in \((-L,L)\), and set
\(F(x)=\widehat f(\pi x/L)\). Then \(F\in PW_\pi\), and
Plancherel gives \(\|F\|_2^2=2L\|f\|_2^2=s\|f\|_2^2\).
The Weil sampling formula, with positive multiplicities,
gives exactly
\[
 q(f)\ge\sum_{\gamma\in\Gamma_h}|\widehat f(\gamma)|^2
       =\sum_{x\in\Lambda}|F(x)|^2.
\]
The compact smooth core and the variational characterization now
prove Theorem~\ref{thm:cosine-lower-rate} for \(a_1(L)\), after
absorbing the harmless factor \(s\) into the constant.
Since \(a_1(L)\le b_L\le a_2(L)\), its other assertions follow.

This lower rate concerns the Weil form at every
sufficiently large window. Its exponent satisfies
\(1<p_{\rm MT}<6/5\). The required denominator estimate in
Section~\ref{sec:theta-coercivity} has exponential
\(e^{-4\pi e^{2L}}\) and a positive polynomial gain. The bound
proved here does not supply that estimate: inserting it into the
prolate trial bound leaves the positive term
\(C e^{2p_{\rm MT}L}L(1+\log L)\) in the exponent of the resulting
upper ratio bound. No vanishing ratio is asserted from this rate.

\subsection{Conditioning and the accuracy of the arithmetic input}
\label{subsec:fd-arithmetic-precision}

The all-window coercivity estimate also gives a quantitative accuracy
budget for the finite trace family. This concerns the accuracy of the
matrix entries and the scalar Taylor trace; it makes no assertion
about the number of machine operations needed to obtain those entries.

\begin{proposition}[Arithmetic matrix conditioning and input accuracy]
\label{prop:fd-arithmetic-precision}
Put \(h=2^{-N}\), \(d_N=2N/h-2\), and let \(G_N\) be the
Gram matrix in the hat-difference basis used above. Then
\[
 \lambda_{\min}(G_N)\ge \frac h3 a_1(N),\qquad
 \lambda_{\max}(G_N)\le21250N^3h^{-2},
\]
and therefore
\[
 \operatorname{cond}_2(G_N)
 \le\frac{63750N^3h^{-3}}{a_1(N)}.
\]
In particular, the uniform lower rate of
Theorem~\ref{thm:cosine-lower-rate} gives
\[
 \log\operatorname{cond}_2(G_N)
 =O\bigl(e^{2N/d_{\rm MT}}N(1+\log N)\bigr).
\]

For every \(N,k\ge1\) and \(T\ge2\), the exact trace defect
\(D_{N,k}\) can be enclosed to radius less than \(T^{-2k}\)
using Hermitian rational approximations of the entries of \(G_N,H_N\)
and a rational approximation of \(S_k\), with absolute input
accuracy \(2^{-p}\), where
\[
 p=O\bigl(e^{2N/d_{\rm MT}}N(1+\log(N+1))
                  +k\log T+\log(k+1)\bigr).
\]
The central finite trace may be evaluated exactly over the rationals;
any numerical evaluation instead requires its own rounding bound.
The constants in this asymptotic accuracy statement are absolute.
\end{proposition}

\begin{proof}
Let \(c\in\mathbb C^{d_N}\) be the coefficients in the
hat-difference basis. In the full interior-hat basis its coefficient
vector \(v\) consists of these entries together with the central
entry \(-\sum_jc_j\). Thus
\[
 \|c\|_2^2\le\|v\|_2^2\le(d_N+1)\|c\|_2^2.
\]
The ordinary \(L^2\) mass matrix for the interior hats has diagonal
\(2h/3\) and neighboring entries \(h/6\). Its eigenvalues lie
between \(h/3\) and \(h\), by the elementary quadratic-form
bound on its off-diagonal terms. The derivative Gram matrix has
diagonal \(2/h\), neighboring entries \(-1/h\), and norm at
most \(4/h\). Consequently the corresponding function \(f\)
satisfies
\[
 \|f\|_2^2\ge\frac h3\|c\|_2^2,
 \qquad
 \|f\|_{H^1}^2\le(h+4/h)(d_N+1)\|c\|_2^2.
\]
Since \(h\le1/2\) and \(d_N+1\le2N/h\), the latter upper
bound is at most \((17/2)Nh^{-2}\|c\|_2^2\).
Coercivity gives \(q(f)\ge a_1(N)\|f\|_2^2\), while
\eqref{eq:fd-sampling-sobolev} gives
\(q(f)\le2500N^2\|f\|_{H^1}^2\). These prove the two
eigenvalue bounds and their stated ratio. Taking logarithms and
using the proved lower rate gives the conditioning estimate.

Here is a precise perturbation bound for the trace. Write \(d=d_N\)
and let \(g_0>0\) be any proved lower bound for
\(\lambda_{\min}(G_N)\). Suppose the Hermitian central matrices
\(G_0,H_0\) have entrywise absolute errors at most \(\varepsilon\).
Set
\[
 \tau=\frac{d\varepsilon}{g_0},\qquad
 B=G_N^{-1/2}H_NG_N^{-1/2},
\]
\[
 E=G_N^{-1/2}(G_0-G_N)G_N^{-1/2},\qquad
 F=G_N^{-1/2}(H_0-H_N)G_N^{-1/2}.
\]
Then \(\|E\|,\|F\|\le\tau\). For \(\tau\le1/4\),
the approximate Gram matrix is positive definite, and its generalized
eigenvalues are those of the Hermitian matrix
\[
 \widetilde B=(I+E)^{-1/2}(B+F)(I+E)^{-1/2}.
\]
The previously certified first packet in
Section~\ref{sec:executed-contour} gives \(\|C\|<1/14\):
its unique first positive heat zero lies in \((28,29)\), and the
heat coordinate is \(x=2\gamma\). Hence \(\|B\|<1/14\).
The elementary functional-calculus bound gives
\[
 \|\widetilde B-B\|
 \le\frac{\tau(1+\|B\|)}{1-\tau}\le2\tau.
\]
For completeness, with \(R=(I+E)^{-1/2}\), the contribution of
\(F\) is at most \(\tau/(1-\tau)\). Also
\(\|R-I\|(\|R\|+1)\le\tau/(1-\tau)\), which bounds
\(RBR-B\) by \(\|B\|\tau/(1-\tau)\).
Both \(B\) and \(\widetilde B\) have norm at most one.
Order their eigenvalues increasingly and apply the variational
principle to obtain a displacement at most \(2\tau\) for each.
The derivative of \(x^{2k}\) has absolute value at most \(2k\)
on \([-1,1]\). Therefore
\[
 \left|\operatorname{tr}(\widetilde B^{2k})
              -\operatorname{tr}(B^{2k})\right|
 \le4kd\tau
 =\frac{4kd^2\varepsilon}{g_0}.
\]
This uses the relative Hermitian pencil, so the trace error does not
raise the Euclidean condition number of \(G_N\) to the power \(2k\).

Choose
\[
 \varepsilon\le\frac{g_0T^{-2k}}{16kd^2},
 \qquad
 |S_{k,0}-S_k|\le\frac14T^{-2k}.
\]
These choices imply \(\tau\le1/4\) and bound the total defect
error by \(T^{-2k}/2\). The exact central value is
\[
 S_{k,0}-\operatorname{tr}\bigl((G_0^{-1}H_0)^{2k}\bigr),
\]
because this matrix has the same generalized eigenvalues as
\(\widetilde B\). The matrices can be chosen with rational real
and imaginary parts, preserving Hermitian symmetry, and this central
value is then an exact rational number.

Finally the lower rate supplies, after enlarging an absolute
constant to cover the finite initial range, a bound
\[
 g_0=\frac h3
 \exp\bigl[-C e^{2N/d_{\rm MT}}N(1+\log(N+1))\bigr].
\]
Taking binary logarithms of the two required accuracies, and using
\(\log d_N=O(N+\log(N+1))\), gives the stated value of \(p\).
This proves an input-accuracy bound. It does not estimate the
separation of the true value \(D_{N,k}\) from a certificate
threshold, which is a separate mathematical quantity.
\end{proof}

The bound fixes a concrete numerical issue in the two-parameter
search. Matrix conditioning can be controlled from an established
arithmetic lower rate, and finite trace values can be enclosed on
the required \(T^{-2k}\) scale. The remaining unbounded-height
question is an upper estimate for the true defects crossing
\(2T^{-2k}\); numerical accuracy alone does not supply that estimate.

\section{M\"obius inversion and the exact Mellin pole terms}
\label{sec:mobius-endpoint}

The theta map in Section~\ref{sec:theta-coercivity} has an explicit
arithmetic inverse. This inverse identifies the endpoint quantities
controlled by the all-window lower bound and the quantities still
present in the codimension-one restriction.

Write \(\mu_{\mathrm M}(n)\) for the M\"obius function. For a compact
smooth function \(f\) on the physical line, define
\begin{equation}\label{eq:mobius-map}
 (\mathfrak M f)(u)
  =\sum_{n\ge1}\mu_{\mathrm M}(n)(nu)^{-1/2}f(\log(nu)),
 \qquad u>0.
\end{equation}
The sum is finite for each \(u\). If \(f\) is supported in
\((-L,L)\), then \(\mathfrak M f\) vanishes for \(u>e^L\).
With
\[
 (\Theta\psi)(x)=e^{x/2}\sum_{n\ge1}\psi(ne^x),
 \qquad
 (\operatorname{Mel}\psi)(s)=\int_0^\infty\psi(u)u^{s-1}\,du,
\]
the exact inverse and Mellin identities are
\begin{equation}\label{eq:mobius-inverse}
 \Theta\mathfrak M f=f,
 \qquad
 (\operatorname{Mel}\mathfrak M f)(s)
    =\frac{\widehat f(i(s-1/2))}{\zeta(s)}
       \quad(\Re s>1).
\end{equation}
Indeed, both compositions expand with coefficient
\(\sum_{n\mid k}\mu_{\mathrm M}(n)\), which is one for \(k=1\)
and zero otherwise. For \(\Theta\mathfrak M f\) the double sum is
finite. Conversely, for a continuous \(\psi\) decreasing faster
than every power at infinity,
\begin{equation}\label{eq:mobius-other-inverse}
 \mathfrak M\Theta\psi=\psi.
\end{equation}
For each fixed \(u>0\), its rearrangement is justified by
\(\sum_{k\ge1}d(k)|\psi(ku)|<\infty\); no estimate at \(u=0\)
is needed for this pointwise assertion. To prove the Mellin identity,
substitute \(v=nu\) in every term. For \(\sigma=\Re s>1\),
the sum of the absolute integrals is bounded by
\[
 \left(\sum_{n\ge1}n^{-\sigma}\right)
 \int_0^\infty |f(\log v)|v^{\sigma-3/2}\,dv<\infty.
\]
The remaining Dirichlet series is \(1/\zeta(s)\), and the remaining
integral is \(\widehat f(i(s-1/2))\).

\subsection{Weighted physical Plancherel on every positive shift}

For every \(\epsilon>0\), the same compact smooth \(f\) satisfies
the exact identity
\begin{equation}\label{eq:mobius-physical-plancherel}
 \int_0^\infty |\mathfrak M f(u)|^2u^{2\epsilon}\,du
 =\frac1{2\pi}\int_{\mathbb R}
 \left|\frac{\widehat f(t+i\epsilon)}
 {\zeta(1/2+\epsilon-it)}\right|^2dt<\infty.
\end{equation}
Here the zero-location theorem supplies the smoothed estimate
\begin{equation}\label{eq:mobius-smoothed}
 \mathfrak M f(u)=O_{f,\delta}(u^{-1/2-\delta})
                         \quad(u\downarrow0)
\end{equation}
for every fixed \(\delta>0\). We give a contour proof to make the
precise input explicit.

First, for fixed \(0<\delta<1/2\),
\begin{equation}\label{eq:mobius-reciprocal-growth}
 |\zeta(\sigma+it)|^{-1}
 \le C_\delta(2+|t|)^{C_\delta},
            \qquad 1/2+\delta\le\sigma\le2.
\end{equation}
The elementary continuation
\[
 \zeta(s)=\frac{s}{s-1}
           -s\int_1^\infty\{x\}x^{-s-1}\,dx
                   \qquad(\Re s>0, s\ne1)
\]
follows by integrating the integer counting function. It bounds
\(|\zeta(s)|\) by \(C(2+|t|)\) on the disks below when
\(|t|\ge4\). Take the disk centered at \(2+it\) with radius
\(R=3/2-\delta/2\). It lies in \(\Re s>1/2\), has no zero
by the zero-location theorem, and contains no pole. Choose the
analytic branch of \(\log\zeta\) whose central value is given
by the absolutely convergent Euler product and hence is bounded.
Its real part is at most \(C+\log(2+|t|)\). The
Borel--Carath\'eodory inequality on the smaller disk of radius
\(r=3/2-\delta\) gives
\(|\log\zeta(s)|\le C_\delta\log(2+|t|)\) there.
This proves \eqref{eq:mobius-reciprocal-growth} for \(|t|\ge4\).
On the remaining compact set \(1/\zeta\) is continuous, including
its removable zero at \(s=1\), which proves the bound everywhere.

On \(\Re s=2\), ordinary Mellin inversion applies to
\(\mathfrak M f\). One direct justification is that
\(x\mapsto e^{2x}\mathfrak M f(e^x)\) is smooth, vanishes for
\(x>L\), and, with all its derivatives, is \(O_f(e^x)\) as
\(x\to-\infty\), by the absolute finite sum in
\eqref{eq:mobius-map}. Fourier inversion therefore gives
\[
 \mathfrak M f(u)=\frac1{2\pi i}
 \int_{2-i\infty}^{2+i\infty}
       \frac{\widehat f(i(s-1/2))}{\zeta(s)}u^{-s}\,ds.
\]
Repeated integration by parts in the compact physical variable
makes \(\widehat f(i(s-1/2))\) decrease faster than every power
of \(|\Im s|\), uniformly on
\(1/2+\delta\le\Re s\le2\).
Equation~\eqref{eq:mobius-reciprocal-growth} thus permits shifting
the contour to \(\Re s=1/2+\delta\); the horizontal integrals
tend to zero and no pole is crossed. The new vertical integral
is absolutely convergent. Taking its absolute value proves
\eqref{eq:mobius-smoothed}.

The bound for every \(\delta>0\) makes the Mellin integral
holomorphic in \(\Re s>1/2\). Analytic continuation from
\eqref{eq:mobius-inverse} identifies it there with the same
quotient. For fixed \(\epsilon>0\), choose
\(0<\delta<\epsilon\) in \eqref{eq:mobius-smoothed}; together
with the upper support bound this gives square integrability
with weight \(u^{2\epsilon}\). Finally apply ordinary Fourier
Plancherel to
\(x\mapsto e^{(1/2+\epsilon)x}\mathfrak M f(e^x)\).
The substitution \(u=e^x\), followed by \(t=-\Im s\), gives
exactly \eqref{eq:mobius-physical-plancherel}.

The pole-order mechanism behind M\"obius mean squares is classical;
it underlies the weak Mertens criterion
\(\int_1^X(M(t)/t)^2dt=O(\log X)\), discussed with its
zero-multiplicity consequences in
\cite[preprint \S1, equations~(1.14)--(1.15), p.~3]{goldston-suriajaya2020}.
Here its exact compact smoothing is visible directly. Set
\[
 M(x)=\sum_{n\le x}\mu_{\mathrm M}(n),\qquad
 \mathfrak m(y)=e^{-y/2}M(e^y),
\]
where \(M(x)=0\) for \(0<x<1\). For every real \(Y\),
\begin{equation}\label{eq:mobius-mertens-smoothing}
 g_f(Y):=e^{-Y/2}\mathfrak M f(e^{-Y})
 =\int_{\mathbb R}\mathfrak m(Y+z)
                  \bigl(f(z)/2-f'(z)\bigr)\,dz.
\end{equation}
Indeed, Stieltjes partial summation gives
\[
 \mathfrak M f(u)
 =-u^{-1/2}\int_0^\infty M(v)v^{-3/2}
          \bigl(f'(\log(vu))-f(\log(vu))/2\bigr)\,dv.
\]
The boundary terms vanish, since the test is compactly supported
away from zero and infinity in \(v\). Substitute \(u=e^{-Y}\)
and \(v=e^{Y+z}\) to obtain
\eqref{eq:mobius-mertens-smoothing}. Consequently the left side of
\eqref{eq:mobius-physical-plancherel} is exactly
\(\int_{\mathbb R}e^{-2\epsilon Y}|g_f(Y)|^2dY\).
This identifies the smoothing and all its normalization factors;
no weak Mertens mean-square estimate is assumed here.
Minkowski's inequality gives the further exact norm comparison
\[
 \|e^{-\epsilon Y}g_f(Y)\|_2
 \le\left(\int_{\mathbb R}|f(z)/2-f'(z)|e^{\epsilon z}\,dz\right)
       \|e^{-\epsilon Y}\mathfrak m(Y)\|_2,
\]
with extended nonnegative norms allowed on the right. This follows
by translating \(Y+z\) in each norm in
\eqref{eq:mobius-mertens-smoothing}. A bound of order
\(\epsilon^{-1}\) for the squared Mertens norm is the additional
mean-square estimate at issue; the fixed-shift contour argument
above supplies no such uniform estimate.

\subsection{Local Abel masses at a zeta zero}

Let \(\gamma>0\) be a zeta ordinate, put
\(\rho=1/2-i\gamma\), and write \(m=m_\gamma\).
For any nonzero entire function \(F\), suppose its order at
\(\gamma\) is \(k<m\), and put \(r=m-k\). The meromorphic
function
\[
 Q_F(s)=\frac{F(i(s-1/2))}{\zeta(s)}
\]
has a pole of order \(r\) at \(\rho\), with leading coefficient
\begin{equation}\label{eq:mobius-pole}
 a_{\gamma,F}
 =\frac{i^k m!F^{(k)}(\gamma)}{k!\,\zeta^{(m)}(\rho)}.
\end{equation}
Choose \(\delta>0\) so that the closed disk of radius \(2\delta\)
about \(\rho\) contains no other zero or pole of \(\zeta\).
The following transform-side Abel limit is exact:
\begin{equation}\label{eq:mobius-abel}
 \lim_{\epsilon\downarrow0}
 \frac{\epsilon^{2r-1}}{2\pi}
 \int_{\gamma-\delta}^{\gamma+\delta}
 \left|\frac{F(t+i\epsilon)}
 {\zeta(1/2+\epsilon-it)}\right|^2dt
 =c_r|a_{\gamma,F}|^2,
 \qquad
 c_r=\frac{\Gamma(r-1/2)}{2\sqrt\pi\,\Gamma(r)}.
\end{equation}
To check the coefficient and normalization, write
\(Q_F(s)=(s-\rho)^{-r}H(s)\), with \(H\) analytic and
\(H(\rho)=a_{\gamma,F}\). Taylor expansion of the numerator
uses \(i(s-\rho)\), giving \eqref{eq:mobius-pole}.
In the integral, set \(t=\gamma+\epsilon v\). Its normalized
integrand becomes
\[
 \frac1{2\pi}
 \frac{|H(\rho+\epsilon(1-iv))|^2}{(1+v^2)^r}
       \mathbf1_{|v|<\delta/\epsilon}.
\]
The numerator is uniformly bounded on the chosen disk, and
\((1+v^2)^{-r}\) is integrable. Dominated convergence applies.
Finally, the substitution \(y=v^2\) gives
\(\int_{\mathbb R}(1+v^2)^{-r}dv
 =\Gamma(1/2)\Gamma(r-1/2)/\Gamma(r)\).
This proves the full limit using only the local meromorphic function.
For compact smooth \(f\),
\eqref{eq:mobius-physical-plancherel} consequently gives
\begin{equation}\label{eq:mobius-physical-lower}
 \liminf_{\epsilon\downarrow0}\epsilon^{2r-1}
 \int_0^\infty|\mathfrak M f(u)|^2u^{2\epsilon}\,du
 \ge c_r|a_{\gamma,\widehat f}|^2.
\end{equation}
For any finite collection of ordinates with the same surviving
pole order \(r\), choosing disjoint local intervals gives the sum
of their right sides. This finite-packet inequality uses positivity;
it requires no interchange of an infinite zero sum with a limit.

In particular, at a simple zero and with \(F(\gamma)\ne0\),
\begin{equation}\label{eq:mobius-simple-mass}
 \lim_{\epsilon\downarrow0}\frac{\epsilon}{2\pi}
 \int_{\gamma-\delta}^{\gamma+\delta}
 \left|\frac{F(t+i\epsilon)}
 {\zeta(1/2+\epsilon-it)}\right|^2dt
 =\frac12\left|\frac{F(\gamma)}{\zeta'(\rho)}\right|^2.
\end{equation}
The corresponding local Weil weight is \(m_\gamma|F(\gamma)|^2\).
Thus their precise conversion at a simple zero multiplies the Abel
mass by \(2|\zeta'(\rho)|^2\). At higher order,
\eqref{eq:mobius-pole} displays the derivative of \(F\) and the
derivative of \(\zeta\) that enter each surviving pole.

\subsection{The common-divisor factor in the M\"obius norm}
\label{sec:mobius-gcd}

The squared norm also has an exact arithmetic expansion that separates
its common-divisor contribution. For a compact smooth \(f\), put
\[
 A_\epsilon(f)=\int_0^\infty
          |\mathfrak M f(u)|^2u^{2\epsilon}\,du,
 \qquad
 C_\epsilon(r)=\int_{\mathbb R}e^{2\epsilon y}
       f(y+r/2)\overline{f(y-r/2)}\,dy.
\]
Here \(\epsilon\) may be complex: in the first integral the absolute
square is fixed and the weight is \(u^{2\epsilon}\). If
\(\operatorname{supp}f\subset(-L,L)\), then
\(C_\epsilon(r)=0\) for \(|r|\ge2L\). On each compact set of
\(\epsilon\)'s the functions \(C_\epsilon\) are uniformly bounded.
For \(\Re\epsilon>1/2\), direct expansion gives
\begin{equation}\label{eq:mobius-double-sum}
 A_\epsilon(f)=\sum_{n,m\ge1}
  \frac{\mu_{\mathrm M}(n)\mu_{\mathrm M}(m)}
       {(nm)^{1/2+\epsilon}}
       C_\epsilon\bigl(\log(n/m)\bigr).
\end{equation}
To justify this identity before rearranging its terms, the absolute
integral of the term indexed by \((n,m)\) is bounded by
\(B_{f,\epsilon}(nm)^{-1/2-\Re\epsilon}\).
Indeed, set \(y=\log u+(\log n+\log m)/2\) in that integral.
The two test arguments become \(y+\tfrac12\log(n/m)\) and
\(y-\tfrac12\log(n/m)\), while the weight contributes
\((nm)^{-\epsilon}e^{2\epsilon y}\).
The sum of the stated bounds converges for \(\Re\epsilon>1/2\),
which proves both termwise integration and absolute convergence of
\eqref{eq:mobius-double-sum}.

Write \(n=da\), \(m=db\), with \((a,b)=1\). The exact identity
\[
 \mu_{\mathrm M}(da)\mu_{\mathrm M}(db)
 =\mu_{\mathrm M}(a)\mu_{\mathrm M}(b)\mu_{\mathrm M}(d)^2
           \mathbf1_{(d,ab)=1}
\]
holds prime by prime. A square divisor of \(a\), \(b\), or \(d\),
or a prime common to \(d\) and \(ab\), makes both sides zero.
In the remaining case, \(a,b,d\) are pairwise coprime and squarefree,
and the formula is multiplicativity of the M\"obius function.
Absolute convergence now permits summing first over \(d\), giving
\begin{equation}\label{eq:mobius-gcd-factor}
 A_\epsilon(f)=
 \frac{\zeta(1+2\epsilon)}{\zeta(2+4\epsilon)}R_\epsilon(f),
 \qquad \Re\epsilon>1/2,
\end{equation}
where
\begin{equation}\label{eq:mobius-coprime-series}
 R_\epsilon(f)=
 \sum_{\substack{a,b\ge1\\(a,b)=1}}
 \frac{\mu_{\mathrm M}(a)\mu_{\mathrm M}(b)}
      {(ab)^{1/2+\epsilon}}
 C_\epsilon\bigl(\log(a/b)\bigr)
 \prod_{p\mid ab}(1+p^{-1-2\epsilon})^{-1}.
\end{equation}
The inner common-divisor series used here is
\[
 \sum_{(d,ab)=1}\frac{\mu_{\mathrm M}(d)^2}{d^{1+2\epsilon}}
 =\prod_{p\nmid ab}(1+p^{-1-2\epsilon})
 =\frac{\zeta(1+2\epsilon)}{\zeta(2+4\epsilon)}
       \prod_{p\mid ab}(1+p^{-1-2\epsilon})^{-1}.
\]
These products converge absolutely in the stated half-plane.

The smoothed bound~\eqref{eq:mobius-smoothed} makes
\(A_\epsilon(f)\) holomorphic in \(\Re\epsilon>0\).
For a compact subset of this half-plane choose \(\delta>0\) smaller
than its minimum real part. Near zero the absolute integrand is then
bounded by a constant times
\(u^{-1+2(\Re\epsilon-\delta)}\), and each derivative with respect
to \(\epsilon\) adds a power of \(2\log u\). These functions
are integrable uniformly on that compact subset. At the upper endpoint
the integrand has compact support. Differentiation under the integral
therefore proves the asserted holomorphy. Since the two zeta factors
are finite and nonzero in this half-plane,
\begin{equation}\label{eq:mobius-coprime-continuation}
 R_\epsilon(f):=
 \frac{\zeta(2+4\epsilon)}{\zeta(1+2\epsilon)}A_\epsilon(f),
 \qquad \Re\epsilon>0,
\end{equation}
is the holomorphic continuation of
\eqref{eq:mobius-coprime-series}. The ordinary series representation
has been established for \(\Re\epsilon>1/2\).
For real \(\epsilon>0\), the continuation is nonnegative, and is
strictly positive when \(f\ne0\): the identity
\(\Theta\mathfrak M f=f\) makes \(\mathfrak M f\ne0\),
and its smoothness then makes its weighted squared norm positive.

The pole of zeta at one gives the precise normalization
\begin{equation}\label{eq:mobius-coprime-asymptotic}
 R_\epsilon(f)=2\zeta(2)\epsilon\bigl(1+O(\epsilon)\bigr)
                      A_\epsilon(f)
              \qquad(\epsilon\downarrow0).
\end{equation}
Combining this identity with the local pole calculation
\eqref{eq:mobius-physical-lower}, for the surviving pole order \(r\),
gives
\begin{equation}\label{eq:mobius-coprime-pole}
 \liminf_{\epsilon\downarrow0}
      \epsilon^{2r-2}R_\epsilon(f)
 \ge 2\zeta(2)c_r\left|a_{\gamma,\widehat f}\right|^2.
\end{equation}
The same argument sums the right sides over every finite collection
of ordinates with that pole order. Thus the extracted common-divisor
factor accounts for one power of \(\epsilon^{-1}\); the signed
coprime expression in \eqref{eq:mobius-coprime-series}, continued by
\eqref{eq:mobius-coprime-continuation}, carries the displayed remaining
pole contributions. A uniform upper bound for this continuation near
zero is the arithmetic estimate still to be proved.

\subsection{Completion factors and the divided theta kernels}

Set
\[
 \mathcal A(s)=\frac12s(s-1)\pi^{-s/2}\Gamma(s/2),
 \qquad \xi(s)=\mathcal A(s)\zeta(s),
 \qquad \Xi(z)=\xi(1/2+iz).
\]
The functional equation makes \(\Xi\) even, and hence
\(\xi(s)=\Xi(i(s-1/2))\). Comparing the first nonzero
Taylor coefficients at \(\rho\) gives
\begin{equation}\label{eq:mobius-completion}
 \mathcal A(\rho)\zeta^{(m)}(\rho)
       =i^m\Xi^{(m)}(\gamma),
 \qquad
 a_{\gamma,F}
       =i^{-r}\mathcal A(\rho)
          \frac{m!F^{(k)}(\gamma)}{k!\Xi^{(m)}(\gamma)}.
\end{equation}
Here \(\mathcal A(\rho)\ne0\).

For \(1\le j\le m\), consider the divided transform
\(F_j(z)=\Xi(z)/(z-\gamma)^j\) constructed in
Section~\ref{sec:theta-coercivity}. It has order \(m-j\) at
\(\gamma\), so the quotient of derivatives in
\eqref{eq:mobius-completion} cancels exactly. In fact the global
meromorphic identity is
\begin{equation}\label{eq:mobius-divided}
 Q_{F_j}(s)=\frac{\mathcal A(s)}{[i(s-\rho)]^j},
 \qquad a_{\gamma,F_j}=i^{-j}\mathcal A(\rho).
\end{equation}
The inverse Fourier functions here have the two rapidly decreasing
physical tails already proved in that section. Their raw theta
lifts are explicitly
\[
 \psi_j(u)=\frac{(-i)^j u^{-\rho}}{(j-1)!}
       \int_u^\infty\!\bigl(\log(v/u)\bigr)^{j-1}
                         h(v)v^{\rho-1}\,dv,
\]
where \(h(v)=(4\pi^2v^4-6\pi v^2)e^{-\pi v^2}\).
Fubini on \(\Re s>1/2\) and
\(\int_0^1 x^{s-\rho-1}(-\log x)^{j-1}dx
 =(j-1)!/(s-\rho)^j\)
give the Mellin quotient in \eqref{eq:mobius-divided}.
The repeated first-order intertwining proved earlier gives their
theta images as the physical divided kernels. Then
equation~\eqref{eq:mobius-other-inverse} identifies those lifts with
their M\"obius inverses. In particular, the nonzero raw endpoint
coefficient persists at every allowed division order.
For this noncompact divided family the identities used are the
pointwise inverse \eqref{eq:mobius-other-inverse} and the local Abel
formula \eqref{eq:mobius-abel}. The physical Plancherel identity
\eqref{eq:mobius-physical-plancherel} has been proved on the compact
smooth core, and retains that hypothesis in its applications here.

For \(1\le j<m\), the sampled norm of the physical
divided kernel is zero, whereas its normalized local Abel mass is
\(c_j|\mathcal A(\rho)|^2>0\). At the terminal order \(j=m\),
the sampled norm is
\(m|\Xi^{(m)}(\gamma)/m!|^2\), as previously proved.
These formulas give the complete conversion on this family;
in particular, the endpoint coefficient is determined even when
the sampled values vanish at the repeated ordinate.

\subsection{What the proved coercivity estimate controls}

Write
\[
 E_L=C e^{2p_{\rm MT}L}L(1+\log L),
 \qquad p_{\rm MT}=d_{\rm MT}^{-1}>1,
\]
with \(C\) chosen so that Theorem~\ref{thm:cosine-lower-rate}
gives \(\|f\|_2^2\le e^{E_L}q(f)\) on the compact smooth
core supported in \((-L,L)\). Fourier differentiation and
Cauchy's inequality give, for every integer \(k\ge0\),
\begin{equation}\label{eq:mobius-jet-bound}
 |\widehat f^{(k)}(\gamma)|^2
 \le \frac{2L^{2k+1}}{2k+1}\|f\|_2^2
 \le \frac{2L^{2k+1}}{2k+1}e^{E_L}q(f).
\end{equation}
Indeed, the squared norm of \(x\mapsto(-ix)^ke^{-i\gamma x}\)
on \((-L,L)\) is exactly \(2L^{2k+1}/(2k+1)\).
For \(F=\widehat f\) with order \(k<m\), substitution into
\eqref{eq:mobius-abel} proves the quantitative endpoint bound
\begin{equation}\label{eq:mobius-coercive-abel}
 c_r|a_{\gamma,F}|^2
 \le c_r
 \left(\frac{m!}{k!|\zeta^{(m)}(\rho)|}\right)^2
 \frac{2L^{2k+1}}{2k+1}e^{E_L}q(f).
\end{equation}
The zero-dependent derivative factor is explicit, and the same
exponent \(E_L\) is retained.

Orthogonality to the real even trial \(v_L\) does not impose a
zero value on these evaluations. To see this by an exact test,
take a nonnegative, nonzero, even smooth \(\chi\) compactly
supported in \((-L,L)\), and put \(f(x)=\chi(x)\sin(\gamma x)\).
Then \(f\) is odd, so \(\langle f,v_L\rangle_2=0\), while
\[
 \widehat f(\gamma)
 =-i\int_{-L}^L\chi(x)\sin^2(\gamma x)\,dx\ne0.
\]
The same explicit test has mean zero. Thus restricting to this
codimension-one complement leaves the pole term in
\eqref{eq:mobius-coercive-abel} present.

For comparison with the proved prolate numerator, put
\(\lambda=e^L\). The all-window lower bound gives
\[
 \frac{q(v_L)}{b_L}
 \le C_\varepsilon\lambda^{9+\varepsilon}
       \exp\{-4\pi\lambda^2+E_L\}.
\]
Since \(E_L/\lambda^2
 =C\lambda^{2(p_{\rm MT}-1)}L(1+\log L)\to\infty\),
this bound supplies no decaying ratio. The exact inversion and
endpoint transport above retain both the larger exponent and the
displayed derivative factors. The remaining arithmetic estimate
is still a lower bound for \(b_L\) on the scale
\(\lambda^{9+\alpha}e^{-4\pi\lambda^2}\), with a positive
\(\alpha\) in the range proved possible earlier, or a direct
vanishing estimate for the arithmetic trace defects. Neither follows
from \eqref{eq:mobius-coercive-abel} alone.

\subsection{The Nyman--Beurling comparison}
\label{subsec:mobius-nyman}

The Mellin inverse also gives a precise comparison with the
Nyman--Beurling approximation problem. Put
\[
 \chi(x)=\mathbf1_{(0,1)}(x),\qquad
 r_n(x)=\{1/(nx)\},\qquad
 D_N(s)=\sum_{n\le N}a_n n^{-s}.
\]
The coefficient map \(a\mapsto\chi+\sum_{n\le N}a_n r_n\)
satisfies
\begin{equation}\label{eq:nyman-mellin-map}
 \left\|\chi+\sum_{n\le N}a_n r_n\right\|_{L^2(0,\infty)}^2
 =\frac1{2\pi}\int_{\mathbb R}
  \frac{|1-\zeta(1/2+it)D_N(1/2+it)|^2}{1/4+t^2}\,dt.
\end{equation}
Indeed, for \(0<\Re s<1\), integration of the integer counting
function, followed by \(t=1/x\), gives
\[
 \int_0^\infty r_1(x)x^{s-1}\,dx=-\frac{\zeta(s)}s.
\]
One can see the boundary terms explicitly: the part \(0<t<1\)
of \(\int_0^\infty\{t\}t^{-s-1}\,dt\) is \(1/(1-s)\),
and the part \(t\ge1\) is \(1/(s-1)-\zeta(s)/s\).
They add to the displayed formula. Dilation gives the factor
\(n^{-s}\), whereas the Mellin transform of \(\chi\) is
\(1/s\). Mellin Plancherel proves \eqref{eq:nyman-mellin-map}.

Thus the squared distance \(d_N^2\), defined as the infimum in
\eqref{eq:nyman-mellin-map}, is a finite Gram minimum. Its entries
are
\[
 G_{nm}=\int_0^\infty r_n(x)r_m(x)\,dx,\qquad
 b_n=\int_0^1r_n(x)\,dx
       =\frac{1-\gamma_{\!E}+\log n}{n},
 \qquad d_N^2=1-b^TG^{-1}b.
\]
Here \(\gamma_{\!E}\) is Euler's constant. To evaluate \(b_n\),
substitute \(y=nx\); the integral on \((1,n)\) is \(\log n\),
and integration over the intervals
\((1/(j+1),1/j)\) gives
\(\int_0^1\{1/y\}\,dy=1-\gamma_{\!E}\).
The matrix \(G\) is positive definite. In a vanishing linear
combination of the \(r_n\), the jump at \(x=1\) first gives
\(a_1=0\). At \(x=1/j\), the jumping terms have indices
dividing \(j\); induction gives \(a_j=0\) for \(1\le j\le N\).
A function vanishing almost everywhere has zero jump because it
vanishes on both adjacent open intervals. Completing the square
now proves the formula for \(d_N^2\).
C\'aceres \cite[Sections~14, 16--17, pp.~50--52]{caceres-margin-v2}
also compares this finite Gram minimum with the zero-sum constant.
His linear coefficient is exactly \(c_n=-b_n\), and his Gram
matrix is the same fractional-part autocorrelation matrix; hence
\(1-c^TG^{-1}c=1-b^TG^{-1}b\).


Burnol's lower bound \cite[Theorem~1.3, p.~4; Section~5,
pp.~14--16]{burnol2002} yields
\begin{equation}\label{eq:nyman-multiplicity-excess}
 \liminf_{N\to\infty}(\log N)d_N^2
 \ge 2+\gamma_{\!E}-\log(4\pi)
   +\sum_{\rho\ {\rm distinct}}
       \frac{m_\rho(m_\rho-1)}{|\rho|^2}.
\end{equation}
The sum includes both signs of the ordinate. To specify the
comparison, each generator \(r_n\) in \eqref{eq:nyman-mellin-map}
is Burnol's fractional-part generator at parameter \(\theta=1/n\).
His continuous-parameter approximation space on \([1/N,1]\)
therefore contains this finite space, so its distance is at most
\(d_N\). His theorem gives the lower sum
\(\sum m_\rho^2/|\rho|^2\). The paired Hadamard product and the
functional equation give
\[
 \sum_\rho\frac{m_\rho}{|\rho|^2}
 =\frac{\xi'}\xi(1)-\frac{\xi'}\xi(0)
 =2+\gamma_{\!E}-\log(4\pi),
\]
where the zero-location theorem identifies
\(\rho(1-\rho)=|\rho|^2\). Subtracting this identity proves
\eqref{eq:nyman-multiplicity-excess}. The sums converge: the
local zero count and multiplicity bound give an upper bound
\(O((\log(2+|\gamma|))^2)\) for the squared multiplicities in
each unit interval, which is summable after division by
\(1/4+\gamma^2\).

The squared multiplicity in this formula comes from an explicit
jet Gram matrix. Burnol's normalized vectors at a zero of
multiplicity \(m\) have limiting matrix
\((1/(j+k+1))_{0\le j,k<m}\); its inverse has upper left entry
\(m^2\) \cite[Theorems~5.3--5.5, pp.~14--16]{burnol2002}.
This is the norm of the evaluation functional at zero on the
local jet space after the indicated rescaling.

The dependence on those jets can already be computed for one
double zero. In the Hardy space on the unit circle, let
\(P(z)=(1-z)^2\) and taper the coefficients of
\(P^{-1}=\sum_{n\ge0}(n+1)z^n\):
\[
 W_N(z)=\sum_{n=0}^N(1-n/N)(n+1)z^n,\qquad N\ge1.
\]
Multiplication gives the exact residual
\[
 1-P(z)W_N(z)=\frac2N\sum_{n=1}^Nz^n-z^{N+1},
 \qquad \|1-PW_N\|_{H^2}^2=1+\frac4N.
\]
The optimal degree-\(N\) multiplier has the smaller exact error
\begin{equation}\label{eq:nyman-double-example}
 \min_{\deg Q\le N}\|1-PQ\|_{H^2}^2
       =\frac{2(2N+5)}{(N+3)(N+4)}.
\end{equation}
To verify it, write the residual as a polynomial of degree at
most \(M=N+2\). Divisibility by \((1-z)^2\) says precisely
that the residual has value one and derivative zero at \(1\).
The two representing vectors have Gram matrix
\[
 \begin{pmatrix}
 M+1&M(M+1)/2\\
 M(M+1)/2&M(M+1)(2M+1)/6
 \end{pmatrix}.
\]
The squared minimum is the upper left entry of its inverse,
\(2(2M+1)/((M+1)(M+2))\), giving
\eqref{eq:nyman-double-example}. This is the double-zero case
of the polynomial approximation mechanism in
\cite[Section~2, pp.~4--6]{burnol2002}.

For zeta, Bettin--Conrey--Farmer use the logarithmically tapered
M\"obius polynomial
\(V_N(s)=\sum_{n\le N}(1-\log n/\log N)\mu_{\mathrm M}(n)n^{-s}\)
\cite[Theorem~1, p.~2; Section~3, pp.~5--9]{bettin-conrey-farmer2013}.
Their upper asymptotic uses a bound on the full sum of
\(|\zeta'(\rho)|^{-2}\). A multiple zero makes its summand
infinite, so finiteness of that input already includes simplicity
on each bounded packet. The local pole formulas above display
the corresponding higher-order terms before any such bound is
available. An upper estimate matching the first constant in
\eqref{eq:nyman-multiplicity-excess} has not been proved here.

\subsection{The rational sawtooth in the M\"obius Gram form}
\label{subsec:endpoint-sawtooth-gram}

The arithmetic quantity in the classical Vasyunin Gram formula can
be expressed by an ordinarily convergent series at every rational
argument. We record the convention, the exact coefficient map, and
the part of the resulting signed sum that remains to be estimated.
Define the periodic function
\[
 B(x)=\begin{cases}1-2\{x\},&x\notin\mathbb Z,\\
                   0,&x\in\mathbb Z,
       \end{cases}
 \qquad
 g_B(h/k)=\sum_{\ell\ge1}\frac{B(\ell h/k)}{\ell}
 \quad(h,k\in\mathbb N).
\]
The value of \(g_B\) depends on the rational number, so its value
at \(dh/dk\) is its value at \(h/k\).

\begin{lemma}[Rational sawtooth and the Vasyunin sum]
\label{lem:endpoint-rational-sawtooth}
The defining series for \(g_B(h/k)\) converges in its displayed
order. For coprime positive integers \(h,k\), put
\[
 V(h,k)=\sum_{r=1}^{k-1}\left\{\frac{rh}{k}\right\}
                       \cot\frac{\pi r}{k}.
\]
Then
\begin{equation}\label{eq:endpoint-g-vasyunin}
 g_B(h/k)=-\frac{\pi}{k}V(h,k).
\end{equation}
The empty sum for \(k=1\) has value zero.
\end{lemma}

\begin{proof}
Reduce the fraction first. For a reduced denominator \(k\), the
sequence \(B(\ell h/k)\) is periodic of period \(k\), and its
sum over one period is
\[
 \sum_{r=1}^{k-1}(1-2r/k)=0.
\]
Its partial sums are bounded, so summation by parts proves convergence.
For \(k>1\), grouping the first \(Mk\) terms into residue classes
and using \(\psi(z+M)-\psi(z)=\sum_{j=0}^{M-1}(j+z)^{-1}\)
gives
\[
 \sum_{\ell=1}^{Mk}\frac{B(\ell h/k)}{\ell}
 =\frac1k\sum_{r=1}^{k-1}B(rh/k)
                   \bigl(\psi(M+r/k)-\psi(r/k)\bigr).
\]
Here \(\psi=\Gamma'/\Gamma\). Since
\(\psi(M+r/k)=\log M+O_k(M^{-1})\), the zero period sum gives
\[
 g_B(h/k)=-\frac1k\sum_{r=1}^{k-1}B(rh/k)\psi(r/k).
\]
Pair \(r\) with \(k-r\). The identities
\(B(-x)=-B(x)\) and
\(\psi(1-z)-\psi(z)=\pi\cot(\pi z)\) yield
\[
 g_B(h/k)=\frac{\pi}{2k}
             \sum_{r=1}^{k-1}B(rh/k)\cot(\pi r/k).
\]
The sum of \(\cot(\pi r/k)\) over \(1\le r<k\) is zero
by the same pairing. Substitution of
\(B(rh/k)=1-2\{rh/k\}\) proves
\eqref{eq:endpoint-g-vasyunin}. For \(k=1\), every summand in
the defining series vanishes.
\end{proof}

Let \(\vartheta_n(x)=\{1/(nx)\}\), and write
\[
 G_{hk}=\int_0^\infty\vartheta_h(x)\vartheta_k(x)\,dx,
 \qquad \kappa=\log(2\pi)-\gamma_E.
\]
The integral is finite: the integrand is bounded near zero and is
\(1/(hkx^2)\) for all sufficiently large \(x\).
For coprime \(h,k\), Vasyunin's formula is
\begin{equation}\label{eq:endpoint-vasyunin-coprime}
 G_{hk}=\frac\kappa2\left(\frac1h+\frac1k\right)
       +\frac{k-h}{2hk}\log\frac hk
       -\frac{\pi}{2hk}\bigl(V(h,k)+V(k,h)\bigr).
\end{equation}
We use this classical identity with the precise normalization in
B\'aez-Duarte--Balazard--Landreau--Saias
\cite[Proposition~89, p.~39; definition in Section~4, p.~16]
{bbls-autocorrelation-2003}. Their autocorrelation is
\(\mathcal A(\alpha)=\int_0^\infty\{u\}\{\alpha u\}u^{-2}du\);
the substitution \(u=1/(hx)\) gives
\(G_{hk}=h^{-1}\mathcal A(h/k)\), which converts their formula
exactly into \eqref{eq:endpoint-vasyunin-coprime}.
Their Proposition~88, p.~38, also gives the series form below,
with their Bernoulli series \(\varphi_1=-g_B/2\).

\begin{proposition}[The signed Gram contribution]
\label{prop:endpoint-signed-gram}
For all positive integers \(h,k\),
\begin{equation}\label{eq:endpoint-gram-sawtooth}
 G_{hk}=\frac\kappa2\left(\frac1h+\frac1k\right)
       +\frac{k-h}{2hk}\log\frac hk
       +\frac{g_B(h/k)}{2h}+\frac{g_B(k/h)}{2k}.
\end{equation}
For arbitrary real coefficients \(a_1,\ldots,a_N\), define
\[
 S_0=\sum_{n\le N}a_n,\quad S_1=\sum_{n\le N}\frac{a_n}{n},
 \quad L_0=\sum_{n\le N}a_n\log n,\quad
 L_1=\sum_{n\le N}\frac{a_n\log n}{n}.
\]
Then the full quadratic form is
\begin{equation}\label{eq:endpoint-full-signed-gram}
 a^{\mathsf T}Ga=S_0(\kappa S_1+L_1)-S_1L_0
           +\sum_{h,k\le N}\frac{a_ha_k}{h}g_B(h/k).
\end{equation}
In particular, for \(N\ge2\) and
\(a_n=\mu(n)w_N(n)\), \(w_N(n)=1-\log n/\log N\),
the last contribution is exactly
\begin{equation}\label{eq:endpoint-mobius-sawtooth}
 \sum_{h,k\le N}\frac{\mu(h)\mu(k)}h
             w_N(h)w_N(k)g_B(h/k).
\end{equation}
\end{proposition}

\begin{proof}
For coprime \(h,k\), substitute
\eqref{eq:endpoint-g-vasyunin} into
\eqref{eq:endpoint-vasyunin-coprime}. For general \(h=da\),
\(k=db\), where \((a,b)=1\), the change of variables \(y=dx\)
gives \(G_{da,db}=d^{-1}G_{ab}\). Every term on the right of
\eqref{eq:endpoint-gram-sawtooth} has the same scaling:
the ratios and the values of \(g_B\) are unchanged, and each
coefficient is multiplied by \(d^{-1}\). This proves
\eqref{eq:endpoint-gram-sawtooth} for every pair. The formula
\eqref{eq:endpoint-vasyunin-coprime} is used only for coprime pairs;
no convention for an unreduced-denominator \(V\) is needed.

In the finite quadratic form, interchanging \(h,k\) in the last
term of \eqref{eq:endpoint-gram-sawtooth} combines the two
\(g_B\)-terms into \(\sum a_ha_k g_B(h/k)/h\).
The first term sums to \(\kappa S_0S_1\). The logarithmic term is
\[
 \frac12\sum_{h,k\le N}a_ha_k
           (h^{-1}-k^{-1})(\log h-\log k)
 =S_0L_1-S_1L_0.
\]
These identities prove \eqref{eq:endpoint-full-signed-gram}.
Substitution of the stated real coefficients proves
\eqref{eq:endpoint-mobius-sawtooth}; all sums in this step are finite.
\end{proof}

The linear coefficient has an equally explicit sign convention.
For \(b_n=\langle\mathbf1_{(0,1)},\vartheta_n\rangle\),
\[
 b_n=\frac{1-\gamma_E+\log n}{n},\qquad
 \left\|\mathbf1_{(0,1)}-\sum_{n\le N}a_n\vartheta_n\right\|_2^2
       =1-2b^{\mathsf T}a+a^{\mathsf T}Ga.
\]
Indeed \(u=1/(nx)\) reduces the integral for \(b_n\) to
\(n^{-1}\int_{1/n}^\infty\{u\}u^{-2}du\).
The part below one is \(\log n\), while the integral from one
to an integer \(M\) is \(\log M-H_M+1\), tending to
\(1-\gamma_E\). Thus the convention
\(c_n=(\gamma_E-1-\log n)/n\) in
\cite[Section~17, p.~52]{caceres-margin-v2} is exactly \(c=-b\),
and its expression \(1+2c^{\mathsf T}a+a^{\mathsf T}Ga\)
is the displayed physical norm.

The long-numerator sum displayed by Maier--Rassias
\cite[Theorem~2.1, p.~3]{maier-rassias-2018} has indexing range
\(k^D\le h<2k^D\), \(D\ge2\), hence \(h\ge k^2\).
Its transpose has \(k\ge h^2\). The signed form
\eqref{eq:endpoint-mobius-sawtooth} also contains the balanced
region \(h<k^2\), \(k<h^2\). This range comparison uses only
the displayed indexing range; no quantitative estimate from that
source is invoked here. The comparable pairs occurring in the
compact M\"obius correlation satisfy
\(\lvert\log(a/b)\rvert\le2L\). When
\(\min(a,b)>e^{2L}\), they satisfy both \(a<b^2\) and
\(b<a^2\): for example, \(a\ge b\) gives
\(a\le e^{2L}b<b^2\), and \(b\le a<a^2\).
An estimate for this balanced signed arithmetic contribution, and
the required uniform bound on the analytically continued
\(R_\epsilon(f)\) near \(\epsilon=0\), remain open in the present
argument. The coprime series defining \(R_\epsilon\) is used in
its proved absolute-convergence region \(\Re\epsilon>1/2\);
its extension to \(\Re\epsilon>0\) retains the analytic-continuation
meaning already established above.

\subsection{Burnol's local jets in the Weil completion}
\label{subsec:fd-burnol-weil}

Burnol's local vectors carry a full jet Gram matrix in physical
\(L^2\). The following computation describes their image under the
Weil sampling map and the term that remains when these vectors are
localized to a finite window.

Write \(\mathcal K=L^2((0,\infty),dt)\), and let
\[
 (\mathcal J f)(x)=e^{x/2}f(e^x).
\]
This is unitary from \(\mathcal K\) to \(L^2(\mathbb R,dx)\).
With the Fourier convention of this paper,
\(\widehat{\mathcal J f}(\nu)=\operatorname{Mel}f(1/2-i\nu)\)
where either side is defined. Burnol's unitary \(V\) on
\(\mathcal K\) has Mellin multiplier
\[
 v(s)=\left(\frac{s}{1-s}\right)^3
                \frac{\zeta(1-s)}{\zeta(s)}
 =\left(\frac{s}{1-s}\right)^3
       2(2\pi)^{-s}\Gamma(s)\cos\frac{\pi s}{2}.
\]
The second expression removes the apparent singularities at zeta
zeros, and \(|v(1/2+it)|=1\).

Fix a signed ordinate \(\gamma\), set \(\rho=1/2-i\gamma\),
and let \(0<\lambda\le e^{-1}\), \(R=\log(1/\lambda)\).
Let \(Q_\lambda\) be multiplication by
\(\mathbf1_{[\lambda,\infty)}\). Burnol's vectors have the
form
\[
 Y_{\rho,k}^{\lambda}=V^{-1}Q_\lambda\Phi_{\rho,k},
 \qquad
 X_k^\lambda=R^{-k-1/2}Y_{\rho,k}^{\lambda},
 \qquad k\ge0.
\]
Here \(\Phi_{\rho,k}\) is the boundary value of
\(V((\log(1/t))^k t^{-w}\mathbf1_{(0,1)}(t))\) as
\(w\to\rho\) from the left. The existence and physical
estimates are Burnol's Theorem~4.2 and Theorem~4.11
\cite[pp.~9--13]{burnol2002}. In his multiplicity application one
uses \(0\le k<m_\gamma\), the range for which the vectors have
the required orthogonality to the approximation space.

We use the following explicit consequences of that construction:
\begin{align}
 \operatorname{Mel}\Phi_{\rho,k}(s)
  &=\frac{k!v(s)}{(s-\rho)^{k+1}}
               &&(1/2<\Re s<1),\label{eq:fd-burnol-mellin}\\
 \Phi_{\rho,k}(t)
  &=\partial_\rho^k\bigl(v(\rho)t^{-\rho}\bigr)
       +r_{\rho,k}(t)
               &&(0<t\le1),\label{eq:fd-burnol-small}\\
 |r_{\rho,k}(t)|+|t r_{\rho,k}'(t)|&\le C_{\rho,k}t^2,
               &&(0<t\le1),\label{eq:fd-burnol-rem}\\
 \Phi_{\rho,k}(t)&=O_{\rho,k}((1+\log t)^2/t)
               &&(t\ge1).\label{eq:fd-burnol-large}
\end{align}
For clarity, the derivative bound in
\eqref{eq:fd-burnol-rem} follows from the convergent series
preceding Burnol's Theorem~4.11, followed by the two Hardy averages
in \(V=(1-M)^2U\). The displayed source series is stated for
\(k\ge1\). For \(k=0\), integration by parts in Burnol's
defining integral gives the boundary term
\(\sin(2\pi t)/(\pi t)\). Its constant term \(+2\) cancels
the constant \(-2\) from extending the remaining integral to
zero, so the same series starting at \(n=1\) holds. The
resulting formula for every \(k\ge0\) is
\[
 r_{\rho,k}(t)=2(-1)^k k!\sum_{n\ge1}
 \frac{(-1)^n(2\pi t)^{2n}}{(2n+1)!}
 \frac{2n}{(\rho+2n)^{k+1}}
 \left(\frac{2n}{2n+1}\right)^2.
\]
The series and its \(t\partial_t\) derivative converge absolutely
on \([0,1]\) and start with \(t^2\). The Mellin identity
\eqref{eq:fd-burnol-mellin} first follows from the multiplier
identity for \(\Re w<1/2\); the physical bounds permit passage
to \(w=\rho\) on each closed substrip of \(1/2<\Re s<1\).

\begin{proposition}[The leading Weil image of the local jets]
\label{prop:fd-burnol-weil}
The canonical Fourier transforms of \(\mathcal JX_k^\lambda\)
define vectors in the Weil completion. For each fixed \(\gamma\)
and \(k\),
\begin{equation}\label{eq:fd-burnol-coordinate}
 \left\|\mathcal JX_k^\lambda
       -\frac{\sqrt R}{k+1}\mathbf1_{\{\gamma\}}\right\|_q
       =O_{\gamma,k}(R^{-1/2}).
\end{equation}
The coordinate vector is in \(\ell^2(\Gamma,m)\), so its squared
norm is \(m_\gamma\). In particular,
\begin{equation}\label{eq:fd-burnol-weil-gram}
 R^{-1}q(\mathcal JX_j^\lambda,\mathcal JX_k^\lambda)
 \longrightarrow\frac{m_\gamma}{(j+1)(k+1)}.
\end{equation}
For every fixed finite coefficient vector \(a=(a_0,\ldots,a_{r-1})\)
with \(\sum_k a_k/(k+1)=0\),
\[
 q\left(\sum_{k<r}a_k\mathcal JX_k^\lambda\right)
       =O_{\gamma,r}(R^{-1})\|a\|_2^2.
\]
Its physical squared norm tends to
\(\sum_{j,k<r}a_j\overline{a_k}/(j+k+1)\), which is positive
for \(a\ne0\).
\end{proposition}

\begin{proof}
Put \(s_\nu=1/2-i\nu\), and define
\[
 E_\lambda(s)=\int_0^\lambda r_{\rho,k}(t)t^{s-1}\,dt.
\]
Integration by parts, using \eqref{eq:fd-burnol-rem}, gives on
the critical line
\begin{equation}\label{eq:fd-burnol-cutoff-error}
 |E_\lambda(s_\nu)|\le
                 C_{\rho,k}\lambda^{5/2}/|s_\nu|.
\end{equation}
Indeed the boundary term is \(\lambda^s r(\lambda)/s\), and
the remaining integral is \(-s^{-1}\int_0^\lambda r'(t)t^sdt\).
Both have the displayed bound.

Subtracting the lower tail from
\eqref{eq:fd-burnol-mellin} gives the canonical transform
\begin{align}
 F_{R,k}(\nu)
 &:={\widehat{\mathcal JY_{\rho,k}^{\lambda}}}(\nu)\notag\\
 &=\frac1{v(s_\nu)}\left\{
   \frac{k!v(s_\nu)}{(s_\nu-\rho)^{k+1}}
   -\partial_\rho^k
       \left(\frac{v(\rho)\lambda^{s_\nu-\rho}}
                         {s_\nu-\rho}\right)
   -E_\lambda(s_\nu)\right\}.
 \label{eq:fd-burnol-transform}
\end{align}
At \(\nu=\gamma\), the apparent singularity is removable.
One may obtain the value directly by integrating
\(Q_\lambda\Phi_{\rho,k}\). Equations
\eqref{eq:fd-burnol-small}--\eqref{eq:fd-burnol-large} give
\[
 \int_\lambda^\infty\Phi_{\rho,k}(t)t^{\rho-1}\,dt
 =\sum_{j=0}^k\binom kj v^{(j)}(\rho)
        \frac{R^{k-j+1}}{k-j+1}+O_{\rho,k}(1).
\]
Division by \(v(\rho)\) therefore yields
\begin{equation}\label{eq:fd-burnol-target}
 F_{R,k}(\gamma)=\frac{R^{k+1}}{k+1}+O_{\gamma,k}(R^k).
\end{equation}

For a distinct ordinate \(\nu\), the difference
\(|s_\nu-\rho|=|\nu-\gamma|\) is bounded below by a positive
constant depending on \(\gamma\). The derivatives of
\(\lambda^{s-\rho}/(s-\rho)\) are finite sums of
\[
 \lambda^{s-\rho}R^{\ell-j}j!(s-\rho)^{-j-1},
 \qquad 0\le j\le\ell\le k.
\]
On the critical line the exponential factor has modulus one.
Since \(|v(s_\nu)|=1\), equations
\eqref{eq:fd-burnol-cutoff-error} and
\eqref{eq:fd-burnol-transform} imply
\begin{equation}\label{eq:fd-burnol-other}
 |F_{R,k}(\nu)|\le
    C_{\gamma,k}\left(
       \frac{R^k}{|\nu-\gamma|}
           +\frac{\lambda^{5/2}}{|s_\nu|}\right).
\end{equation}
Both squared majorants are summable against the zero weights:
\[
 \sum_{\nu\ne\gamma}\frac{m_\nu}{(\nu-\gamma)^2}<\infty,
 \qquad
 \sum_\nu\frac{m_\nu}{|s_\nu|^2}<\infty.
\]
This follows from isolation of the fixed zero and the counted
unit-interval estimate. Multiplying
\eqref{eq:fd-burnol-target} and
\eqref{eq:fd-burnol-other} by \(R^{-k-1/2}\) proves
\eqref{eq:fd-burnol-coordinate}, provided the canonical transform
is interpreted in the sampling completion. We justify that
interpretation next.

For each fixed \(\lambda,\rho,k\),
\eqref{eq:fd-burnol-transform} is a continuous function of \(\nu\)
with a continuous derivative, including at the removable point.
The gamma expression for \(v\) gives
\[
 |v'(1/2-i\nu)|=O(\log(2+|\nu|)).
\]
For example this follows by logarithmic differentiation: the
digamma series, split at \(n=\lceil|\nu|\rceil+1\), bounds
the digamma term by \(O(\log(2+|\nu|))\), and the remaining
rational and tangent terms are bounded outside a fixed interval.
The derivative in \(\nu\) of the cutoff error satisfies the
same \((1+|\nu|)^{-1}\) bound with a fixed extra factor
\(R+1\), by the preceding integration by parts applied to
\(r(t)\log t\). Consequently, as \(|\nu|\to\infty\),
\[
 |F_{R,k}(\nu)|+|F_{R,k}'(\nu)|
       =O_{\lambda,\gamma,k}
          \left(\frac{\log(2+|\nu|)}{1+|\nu|}\right).
\]
Thus this transform belongs to
\[
 \mathcal E=\left\{F:\int_{\mathbb R}\log(3+|t|)
                    (|F(t)|^2+|F'(t)|^2)dt<\infty\right\}.
\]
The unit-interval sampling inequality is continuous on
\(\mathcal E\). Smooth frequency cutoffs and mollification show
that Schwartz functions are dense in this weighted Sobolev space.
Their inverse Fourier transforms are Schwartz; cutting these off
smoothly on the physical line then converges in the Schwartz
topology and hence in \(\mathcal E\). The resulting compact
smooth tests therefore converge to the displayed sampling vector
in Weil norm. This proves the stated completion-domain assertion.

Equation \eqref{eq:fd-burnol-weil-gram} and the coefficient
annihilation estimate now follow from
\eqref{eq:fd-burnol-coordinate} and the triangle inequality in
the completed Hilbert space.

For the physical Gram, unitarity of \(V\) and \(\mathcal J\)
reduces the calculation to the functions
\(Q_\lambda\Phi_{\rho,k}\). Their leading terms on
\((\lambda,1)\) are
\(v(\rho)t^{-\rho}(\log(1/t))^k\). The leading cross integral
is exactly
\[
 |v(\rho)|^2\int_\lambda^1
       (\log(1/t))^{j+k}\frac{dt}{t}
 =\frac{R^{j+k+1}}{j+k+1}.
\]
The lower logarithmic powers contribute
\(O(R^{j+k})\), and the remainder and upper tail contribute
lower order terms by
\eqref{eq:fd-burnol-rem}--\eqref{eq:fd-burnol-large}.
Normalization gives the Hilbert Gram limit, which is Burnol's
Theorem~5.3 in this notation \cite[p.~14]{burnol2002}.
Its positive definiteness follows directly from
\[
 \sum_{j,k<r}\frac{a_j\overline{a_k}}{j+k+1}
 =\int_0^1\left|\sum_{k<r}a_k t^k\right|^2dt>0
                         \quad(a\ne0).
\]
\end{proof}

This calculation gives an exact norm comparison: the normalized
physical jet Gram has full rank, while its leading Weil image is
the single value coordinate at the chosen ordinate. The calculation
uses no upper bound for a reciprocal zeta derivative. The
multiplicity restriction \(k<m_\gamma\) enters Burnol's
orthogonality theorem; the sampling estimates themselves apply to
every fixed \(k\).

\begin{proposition}[The physical window occupied by the jets]
\label{prop:fd-burnol-physical-window}
For each fixed \(\gamma\) and \(k\), put
\[
 p_{R,k}(x)=R^{-k-1/2}(-x)^k e^{i\gamma x}
                            \mathbf1_{[-R,0]}(x).
\]
Then
\begin{equation}\label{eq:fd-burnol-physical-window}
 \|\mathcal JX_k^{e^{-R}}-p_{R,k}\|_{L^2(\mathbb R)}
                  =O_{\gamma,k}(R^{-1/2}).
\end{equation}
For a fixed finite coefficient vector \(a\), set
\(P(t)=\sum_k a_k t^k\). Uniformly for \(0\le L\le R\),
\[
 \left\|\mathbf1_{[-L,L]}
             \sum_k a_k\mathcal JX_k^{e^{-R}}\right\|_2^2
       =\int_0^{L/R}|P(t)|^2dt+O_{\gamma,a}(R^{-1/2}).
\]
In particular, a window with \(L=o(R)\) retains vanishing
physical mass, although the full physical squared norm tends to
\(\int_0^1|P(t)|^2dt\).
\end{proposition}

\begin{proof}
The Fourier transform of the unnormalized polynomial indicator is
\[
 P_{R,k}(\nu)=\frac{k!}{(s_\nu-\rho)^{k+1}}
        -\partial_\rho^k\left(
                 \frac{\lambda^{s_\nu-\rho}}{s_\nu-\rho}\right).
\]
This follows by integrating \((-x)^k e^{i\gamma x}\) on
\([-R,0]\), with the removable value at \(\nu=\gamma\).
Subtract this expression from \eqref{eq:fd-burnol-transform}.
The result is
\[
 F_{R,k}(\nu)-P_{R,k}(\nu)
 =\frac1{v(s_\nu)}\partial_\rho^k
       \left\{\lambda^{s_\nu-\rho}
          \frac{v(s_\nu)-v(\rho)}{s_\nu-\rho}\right\}
       -\frac{E_\lambda(s_\nu)}{v(s_\nu)}.
\]
For each fixed \(j\), the divided difference
\[
 Q_j(\nu)=\partial_\rho^j
                 \frac{v(s_\nu)-v(\rho)}{s_\nu-\rho}
\]
has a removable singularity at \(\nu=\gamma\) and is bounded
on compact intervals. As \(|\nu|\to\infty\), the fact that
\(|v(s_\nu)|=1\), together with the fixed derivatives of
\(v\) at \(\rho\), gives \(Q_j(\nu)=O_{\gamma,j}(|\nu|^{-1})\).
Consequently \(Q_j\in L^2(\mathbb R,d\nu)\).
Differentiating the exponential factor yields powers of \(R\),
so the last display and \eqref{eq:fd-burnol-cutoff-error} imply
\[
 \|F_{R,k}-P_{R,k}\|_{L^2(d\nu)}
      \le \sum_{j=0}^k\binom kj R^{k-j}\|Q_j\|_2
                       +O_{\gamma,k}(e^{-5R/2})
      =O_{\gamma,k}(R^k).
\]
Divide by \(R^{k+1/2}\) and apply Fourier Plancherel to obtain
\eqref{eq:fd-burnol-physical-window}. A finite linear combination
has the same error order. Multiplication by
\(\mathbf1_{[-L,L]}\) is a contraction in \(L^2\), while direct
substitution \(t=-x/R\) gives exactly
\[
 \left\|\mathbf1_{[-L,L]}\sum_k a_kp_{R,k}\right\|_2^2
                      =\int_0^{L/R}|P(t)|^2dt.
\]
The squared-norm difference is bounded by the norm error times
the sum of the two norms, which is uniformly bounded. This proves
the assertion uniformly in \(L\).
\end{proof}

The remaining finite-window term is explicit. For a smooth physical
cutoff \(\chi_L\) and
\(F_R=\sum_{k<r}a_k\mathcal JX_k^\lambda\), with
\(\sum a_k/(k+1)=0\), the triangle inequality gives
\[
 q(\chi_LF_R)^{1/2}
 \le C_{\gamma,r}R^{-1/2}\|a\|_2
             +q((1-\chi_L)F_R)^{1/2}.
\]
For each fixed \(R\), smooth cutoffs approaching one converge in
the sampling domain after smoothing if necessary. One direct
justification uses their approximate-identity convolution in
frequency and the moderate weight \(\log(3+|t|)\); weighted
Young's inequality gives uniform boundedness, and density gives
strong convergence in \(\mathcal E\), also for the first
derivative. The following calculation gives the cutoff norm for a fixed
polynomial family at a fixed ordinate. A bound uniform in the
height and in the number of jet orders, at the scales needed for
the all-height criterion, remains unproved.

\subsection{The Weil norm of a coupled physical cutoff}
\label{subsec:burnol-coupled-cutoff}
The local vectors of
Proposition~\ref{prop:fd-burnol-weil} admit a quantitative
cutoff calculation for each fixed finite polynomial family.
Both endpoints of the physical polynomial enter the
uncut transform; the endpoint at zero determines the leading
sample after a shorter physical cutoff.

Fix $\gamma\in\Gamma$, put $\rho=1/2-i\gamma$, and retain
Burnol's multiplier
\[
 v(s)=\left(\frac{s}{1-s}\right)^3
       2(2\pi)^{-s}\Gamma(s)\cos\frac{\pi s}{2}.
\]
For a polynomial $P(t)=\sum_{k=0}^K b_kt^k$, define
\[
 H_R=\sum_{k=0}^K b_k\mathcal JX_k^{e^{-R}},
 \qquad
 p_R(x)=R^{-1/2}P(-x/R)e^{i\gamma x}
                    {\bf1}_{[-R,0]}(x).
\]
All constants below may depend on this fixed ordinate and
polynomial. Choose a real nonnegative
$\chi\in C_c^\infty((-1,1))$ equal to one near zero, and set
$\chi_L(x)=\chi(x/L)$.

\begin{proposition}\label{prop:burnol-coupled-cutoff}
Suppose $\int_0^1P(t)\,dt=0$, that $P$ has exact order $M\ge0$
at zero, and that it vanishes to order at least $M$ at one.
Write
\[
 a=\frac{P^{(M)}(0)}{M!}\ne0,\quad
 A_{\chi,M}=\int_0^1\chi(-t)t^Mdt,\quad
 B_{\chi,M}=\int_0^1\chi(-t)^2t^{2M}dt.
\]
Then $q(H_R)=O_{\gamma,P}(R^{-2M-1})$. For every
$L=L(R)\to\infty$ with $L/R\to0$,
\begin{align}
 q(\chi_LH_R)
 &\sim q((1-\chi_L)H_R)
 \sim m_\gamma|a|^2A_{\chi,M}^2
                  \frac{L^{2M+2}}{R^{2M+1}},
                 \label{eq:burnol-cutoff-q}\\
 \|\chi_LH_R\|_2^2
 &\sim |a|^2B_{\chi,M}
                  \frac{L^{2M+1}}{R^{2M+1}}.
                 \label{eq:burnol-cutoff-physical}
\end{align}
In particular,
\[
 \frac{q(\chi_LH_R)}{\|\chi_LH_R\|_2^2}
 \sim m_\gamma\frac{A_{\chi,M}^2}{B_{\chi,M}}\,L.
\]
\end{proposition}

\begin{proof}
We first record the endpoint cancellations in the
multiplier formula. Put $s_\nu=1/2-i\nu$ and
$d_\nu=s_\nu-\rho=-i(\nu-\gamma)$. Summing
\eqref{eq:fd-burnol-transform} against $b_kR^{-k-1/2}$
and applying Leibniz's rule gives, for $\nu\ne\gamma$,
\begin{align}
 \widehat H_R(\nu)
 =R^{-1/2}\bigg\{&
 \sum_{j=0}^K\frac{P^{(j)}(0)}{R^jd_\nu^{j+1}}\notag\\
 &-\frac{e^{-Rd_\nu}}{v(s_\nu)}
 \sum_{j=0}^K\frac{P^{(j)}(1)}{j!R^j}
       \partial_\rho^j\left(\frac{v(\rho)}{d_\nu}\right)
 \bigg\}+\mathcal R_R(\nu).
 \label{eq:burnol-endpoint-transform}
\end{align}
Indeed $\sum_{k\ge j}b_k\binom{k}{j}=P^{(j)}(1)/j!$,
while $b_jj!=P^{(j)}(0)$. The lower-tail remainder bound
\eqref{eq:fd-burnol-cutoff-error} gives
\begin{equation}\label{eq:burnol-poly-rem}
 |\mathcal R_R(\nu)|
 \le C_{\gamma,P}R^{-M-1/2}
           e^{-5R/2}/|s_\nu|.
\end{equation}
Here $b_k=0$ for $k<M$. The two endpoint sums in
\eqref{eq:burnol-endpoint-transform} start at $j=M$.
Isolation of the fixed ordinate and $|v(s_\nu)|=1$ therefore
give
\[
 |\widehat H_R(\nu)|
 \le C_{\gamma,P}R^{-M-1/2}|\nu-\gamma|^{-1}
      +|\mathcal R_R(\nu)|\quad(\nu\in\Gamma,\ \nu\ne\gamma).
\]
The squared majorants are summable against $m_\nu$.

At the target coordinate the direct integral calculation
used for \eqref{eq:fd-burnol-target} instead gives
\[
 \widehat H_R(\gamma)=
 \sqrt R\sum_{j=0}^K
 \frac{v^{(j)}(\rho)}{v(\rho)j!R^j}
             \int_0^1P^{(j)}(t)\,dt
          +O_{\gamma,P}(R^{-M-1/2}).
\]
The $j=0$ term vanishes by the integral constraint. For
$1\le j\le M$, the integral equals
$P^{(j-1)}(1)-P^{(j-1)}(0)=0$. Every remaining term is
$O(R^{-M-1/2})$. This proves $q(H_R)=O(R^{-2M-1})$.
It also identifies the two leading off-target endpoint terms:
for each fixed $\nu\ne\gamma$,
\[
 \widehat H_R(\nu)=
 \frac{R^{-1/2}}{-i(\nu-\gamma)}
 \left\{P(0)-P(1)\frac{v(\rho)}{v(s_\nu)}
                     e^{iR(\nu-\gamma)}\right\}
 +O_{\gamma,P,\nu}(R^{-3/2}).
\]

We next control multiplication by $\chi_L$ in the sampling
norm. Set $w(t)=\log(3+|t|)$ and
$\|F\|_w^2=\int w(t)|F(t)|^2dt$.
The counted unit-interval zero bound implies, for
$0<\epsilon\le1$,
\begin{equation}\label{eq:burnol-short-sampling}
 \sum_{\nu\in\Gamma}m_\nu|F(\nu)|^2
 \le C\left(\epsilon^{-1}\|F\|_w^2
                         +\epsilon\|F'\|_w^2\right).
\end{equation}
To prove it, apply the elementary one-dimensional Sobolev
inequality at $\nu$ on $[\nu-\epsilon,\nu+\epsilon]$.
At a point $t$, at most $C\log(3+|t|)$ counted ordinates
have intervals containing it. Summation and integration
give \eqref{eq:burnol-short-sampling}.

If $e$ has Fourier transform in $L^2(w\,dt)$, multiplication
by $\chi_L$ is convolution with
$k_L(t)=L\widehat\chi(Lt)/(2\pi)$ in frequency.
The inequality
$w(t+u)\le Cw(t)\log(3+|u|)$ and Minkowski's inequality give
\[
 \|k_L*F\|_w\le C_\chi\|F\|_w,\qquad
 \|k_L'*F\|_w\le C_\chi L\|F\|_w
 \qquad(L\ge1).
\]
Both follow by substituting $uL$ in the weighted $L^1$
norms of the Schwartz kernels. Taking $\epsilon=1/L$ in
\eqref{eq:burnol-short-sampling} proves the useful bound
\begin{equation}\label{eq:burnol-cutoff-sampling}
 q(\chi_Le)\le C_\chi L\|\widehat e\|_w^2.
\end{equation}
This explicitly uses weighted Fourier control; a physical
$L^2$ bound alone has not been substituted for a Weil bound.

The exact difference from the polynomial indicator is
\begin{equation}\label{eq:burnol-poly-error}
 \widehat H_R(\nu)-\widehat p_R(\nu)
 =\frac{R^{-1/2}e^{-Rd_\nu}}{v(s_\nu)}
 \sum_{j=M}^K\frac{P^{(j)}(1)}{j!R^j}\,Q_j(\nu)
       +\mathcal R_R(\nu),
\end{equation}
where
\[
 Q_j(\nu)=\partial_\rho^j
             \frac{v(s_\nu)-v(\rho)}{s_\nu-\rho}.
\]
Each singularity at $\gamma$ is removable, and
$Q_j(\nu)=O_{\gamma,j}(|\nu|^{-1})$ at infinity.
Consequently $Q_j\in L^2(w\,d\nu)$. On the real line the
two factors $e^{-Rd_\nu}$ and $v(s_\nu)^{-1}$ have absolute
value one. Thus
\[
 \|\widehat H_R-\widehat p_R\|_w
       \le C_{\gamma,P}R^{-M-1/2}.
\]
Equations \eqref{eq:burnol-cutoff-sampling} and Fourier
Plancherel now give
\begin{equation}\label{eq:burnol-two-cutoff-errors}
 q\bigl(\chi_L(H_R-p_R)\bigr)
       \le C_{\gamma,P,\chi}L R^{-2M-1},\qquad
 \|\chi_L(H_R-p_R)\|_2
       \le C_{\gamma,P}R^{-M-1/2}.
\end{equation}
These statements also place the compactly supported
functions in the completed form domain. For example,
physical mollification supported within the fixed support
margin converges in the weighted Fourier Sobolev norm
appearing in \eqref{eq:burnol-short-sampling}; it gives
compact smooth approximants in the same window.

It remains to compute the compact polynomial itself.
Write $g_{R,L}=\chi_Lp_R$. At the target,
\begin{align*}
 \widehat g_{R,L}(\gamma)
 &=R^{-1/2}L\int_0^1\chi(-t)P(Lt/R)\,dt\\
 &=aA_{\chi,M}\frac{L^{M+1}}{R^{M+1/2}}
                      \left(1+O_P(L/R)\right).
\end{align*}
At a distinct ordinate, integration by parts $M+1$ times
gives
\[
 |\widehat g_{R,L}(\nu)|
 \le C_{P,\chi}R^{-M-1/2}|\nu-\gamma|^{-M-1}
       \qquad(1\le L\le R/2).
\]
For completeness, all left endpoint terms vanish because
$\chi$ has compact support in $(-1,1)$. At zero the
derivatives of order less than $M$ vanish, the derivative
of order $M$ is $(-1)^MP^{(M)}(0)/R^M$, and the
$L^1$ norm of the derivative of order $M+1$ of
$\chi(x/L)P(-x/R)$ on $[-L,0]$ is $O_{P,\chi}(R^{-M})$.
The last bound follows term by term from Leibniz's rule,
using $P^{(j)}(t)=O_P(t^{M-j})$ for $j\le M$ and $L/R\le1/2$.
Therefore the off-target squared samples sum to
$O_{\gamma,P,\chi}(R^{-2M-1})$, whereas the target contributes
the leading expression in \eqref{eq:burnol-cutoff-q}.
The first error norm in \eqref{eq:burnol-two-cutoff-errors}
is smaller than this leading norm by
$O(L^{-M-1/2})$. This proves the asymptotic for
$q(\chi_LH_R)$.
The same leading sample has the exact global jet image
given by \eqref{eq:hb-evaluation-jet}:
\[
 (\mathcal U(\chi_LH_R))^{(m_\gamma-1)}(\gamma)
 =\frac{iF^{(m_\gamma)}(\gamma)}{\sqrt\pi}
                      \widehat{\chi_LH_R}(\gamma).
\]
Thus the coordinate responsible for the cutoff norm is
identified with its corresponding derivative evaluation,
with the multiplicity and normalization retained.

Direct substitution also gives
\[
 \|g_{R,L}\|_2^2
 =R^{-1}L\int_0^1\chi(-t)^2|P(Lt/R)|^2dt
 \sim |a|^2B_{\chi,M}\frac{L^{2M+1}}{R^{2M+1}}.
\]
The second error bound in \eqref{eq:burnol-two-cutoff-errors}
is smaller than its norm by $O(L^{-M-1/2})$, proving
\eqref{eq:burnol-cutoff-physical}. Finally
$\|H_R\|_q=O(R^{-M-1/2})$ is smaller than
$\|\chi_LH_R\|_q$ by $O(L^{-M-1})$.
The reverse triangle inequality proves that
$q((1-\chi_L)H_R)\sim q(\chi_LH_R)$.
\end{proof}

\paragraph{An exponential estimate for the local multiplier error.}
The multiplier gives a stronger version of the first
bound in \eqref{eq:burnol-two-cutoff-errors}. With the same
fixed data, for $1\le L\le R/2$,
\begin{equation}\label{eq:burnol-exp-local-error}
 \|\chi_L(H_R-p_R)\|_q
 \le C_{\gamma,P,\chi}R^{-M-1/2}
 \left\{L^{3/2}e^{-2(R-L)}
                    +\sqrt L\,e^{-5R/2}\right\}.
\end{equation}
Here and below $\|f\|_q=q(f)^{1/2}$.
We give the multiplier justification to keep the direction
of the physical translation explicit.

Let $\kappa_j$ be the inverse Fourier transform of
$Q_j(\nu)/v(s_\nu)$. For $x>0$,
\begin{equation}\label{eq:burnol-positive-kernel-decay}
 |\kappa_j(x)|+|\kappa_j'(x)|
                   \le C_{\gamma,j}e^{-2x}.
\end{equation}
Indeed,
\[
 \frac{Q_j(\nu)}{v(s_\nu)}
 =\frac{j!}{d_\nu^{j+1}}
  -\frac1{v(s_\nu)}
        \partial_\rho^j\left(\frac{v(\rho)}{d_\nu}\right).
\]
The apparent real pole cancels in the whole expression.
The reciprocal multiplier is holomorphic in
$0\le\Im\nu<5/2$: its apparent issue at $s=1$ is a zero
of order two, and its first pole above the critical line
is at $s=3$, corresponding to $\nu=5i/2$.
The gamma formula and Stirling's formula give
$|v(5/2-it)^{-1}|\le C(1+|t|)^{-2}$.
Shift the inverse Fourier contour to $\nu=t+2i$.
The vertical edge terms vanish because the whole integrand
before the Fourier exponential is $O(1/|t|)$ uniformly
across this fixed strip. On the shifted line the inverse
transform of $j!/d_\nu^{j+1}$ is zero for $x>0$.
One can check this without a contour closure from
\[
 \frac{j!}{(2-i(t-\gamma))^{j+1}}
     =\int_0^\infty u^je^{-2u}e^{i(t-\gamma)u}\,du:
\]
its inverse transform is supported on $x\le0$.
The remaining shifted-line integral has integrand
$O((1+|t|)^{-3})$, and its derivative in $x$ has an
integrable $O((1+|t|)^{-2})$ majorant. Its factor $e^{-2x}$
proves \eqref{eq:burnol-positive-kernel-decay}, with a
continuous first derivative for $x>0$.

The nonremainder part of \eqref{eq:burnol-poly-error}
therefore has inverse transform
\[
 R^{-1/2}e^{-iR\gamma}
       \sum_{j=M}^K\frac{P^{(j)}(1)}{j!R^j}
                                     \kappa_j(x+R).
\]
On $|x|\le L$ its value and first derivative are bounded
by $CR^{-M-1/2}e^{-2(R-L)}$.
After multiplication by $\chi_L$ its physical $H^1$ norm
is at most $C\sqrt L R^{-M-1/2}e^{-2(R-L)}$.
The compact-support sampling estimate
$\|f\|_q\le CL\|f\|_{H^1}$ follows from
\eqref{eq:burnol-short-sampling} with $\epsilon=1$,
Fourier Plancherel, and
$w(\nu)\le C(1+\nu^2)$ applied to $f$ and $xf$.
This proves the first term in
\eqref{eq:burnol-exp-local-error}. The remainder
\eqref{eq:burnol-poly-rem} has weighted Fourier norm at
most $CR^{-M-1/2}e^{-5R/2}$; applying
\eqref{eq:burnol-cutoff-sampling} proves its second term.

\paragraph{The two-jet family and the all-window lower rate.}
For $P(t)=1-2t$, $M=0$ and $a=1$. Thus, at
$L=c\log\log R$ with any fixed $c>0$,
\[
 q((1-\chi_L)H_R)\sim
 m_\gamma A_{\chi,0}^2
             \frac{c^2(\log\log R)^2}{R},
 \qquad
 \|\chi_LH_R\|_2^2\sim
 B_{\chi,0}\frac{c\log\log R}{R}.
\]
The same asymptotics hold for the exact target-canceling
combination
\[
 Z_R=\mathcal JX_0^{e^{-R}}-c_R\mathcal JX_1^{e^{-R}},
 \qquad
 c_R=
 \frac{\widehat{\mathcal JX_0^{e^{-R}}}(\gamma)}
      {\widehat{\mathcal JX_1^{e^{-R}}}(\gamma)}
     =2+O_\gamma(R^{-1}).
\]
In fact $q(Z_R)=O_\gamma(R^{-1})$, and all $M=0$ cutoff
error estimates above are uniform over bounded polynomial
coefficients. Its local polynomial
$1-c_Rt$ has the same nonzero constant term and the same
leading cutoff samples and physical mass. No exact
integral constraint is needed for that cutoff calculation;
the uncut bound follows explicitly from
\[
 \|Z_R\|_q\le\|H_R\|_q+|c_R-2|\,
                   \|\mathcal JX_1^{e^{-R}}\|_q
                =O_\gamma(R^{-1/2}),
\]
using Proposition~\ref{prop:fd-burnol-weil}.

The proved all-window estimate of
Theorem~\ref{thm:cosine-lower-rate} is
\[
 a_1(L)\ge
 \exp[-C e^{2L/d_{\rm MT}}L(1+\log L)].
\]
At $L=c\log\log R$ its right side is
\[
 \exp\!\left[-C(\log R)^{2c/d_{\rm MT}}\,
        c\log\log R\bigl(1+\log(c\log\log R)\bigr)\right].
\]
The Rayleigh quotient of the compact two-jet
function is asymptotic to
\[
 m_\gamma\frac{A_{\chi,0}^2}{B_{\chi,0}}\,c\log\log R.
\]
The cutoff Weil norm is asymptotic to a positive constant times
$(\log\log R)/\sqrt R$, whereas the uncut Weil norm is
$O_\gamma(R^{-1/2})$.
These formulas quantify the missing term for this fixed
family and window scale. They provide no bound uniform in
the ordinate or as the number of jet orders grows.
The original all-height integer-deficit estimate remains
unproved.


\subsection{Multiplication by the coordinate and the prime form}
\label{subsec:fd-prime-commutator}

There is an exact arithmetic form identity for physical
multiplication by the coordinate. It also identifies a domain issue
in an attempted comparison of local multiplicity traces with the
ordinary spectral trace.

Let \(D=-i\partial_x\) and \(Xf(x)=xf(x)\), initially on compact
smooth functions. For \(f,g\) in this class write
\(h=f*\widetilde g\), with
\(\widetilde g(x)=\overline{g(-x)}\), and let \(W\) be the
arithmetic Weil functional. Direct convolution gives
\[
 (Xf)*\widetilde g-f*\widetilde{Xg}=xh(x).
\]
Consequently
\begin{align}
 q(Xf,g)-q(f,Xg)
 &=W(xh)\notag\\
 &=\int_{\mathbb R}xh(x)(e^{x/2}+e^{-x/2})\,dx\notag\\
 &\quad-\int_0^\infty
 \frac{xe^{-x/2}(h(x)-h(-x))}{1-e^{-2x}}\,dx\notag\\
 &\quad-\sum_{n\ge2}\frac{\Lambda(n)\log n}{\sqrt n}
                  \bigl(h(\log n)-h(-\log n)\bigr).
 \label{eq:fd-prime-commutator}
\end{align}
The prime sum is finite because \(h\) has compact support.
The local term in the Weil formula vanishes since \(xh(x)\) is zero
at the origin. The remaining integral is regular there because
\(h(x)-h(-x)=O(x)\). On the spectral side the same identity is
\[
 q(Xf,g)-q(f,Xg)
 =i\sum_\gamma m_\gamma
        (\widehat f\,\overline{\widehat g})'(\gamma),
\]
as \(\widehat{Xf}=i\widehat f'\). Also
\([D,X]f=-if\) holds exactly on this test space.

\begin{proposition}[A nonzero residual for global multiplication]
\label{prop:fd-multiplication-domain}
The operator \(X\), defined on compact smooth mean-zero tests and
viewed as an operator into the full Weil completion, is not closable.
For every signed ordinate \(\gamma\) there are such tests \(f_R\)
with
\[
 \|f_R\|_q\longrightarrow0,
 \qquad
 Xf_R\longrightarrow i\mathbf1_{\{\gamma\}}
       \quad\hbox{in }\ell^2(\Gamma,m).
\]
The squared norm of the limiting vector is \(m_\gamma\).
\end{proposition}

\begin{proof}
Choose an even \(\varphi\in C_c^\infty(\mathbb R)\) of integral
one, and put \(K=\widehat\varphi\). Before correcting the mean,
define
\[
 g_R(x)=R^{-1}e^{i\gamma x}\varphi(x/R),\qquad
 f_R^{\,0}=(D-\gamma)g_R.
\]
Their Fourier transforms satisfy
\[
 F_R^{\,0}(z)=(z-\gamma)K(R(z-\gamma)),
\]
\[
 \widehat{Xf_R^{\,0}}(z)
 =i\bigl(K(R(z-\gamma))
        +R(z-\gamma)K'(R(z-\gamma))\bigr).
\]
At \(z=\gamma\) these values are zero and \(i\), respectively.
Every distinct ordinate is isolated, and the counted zero bound gives
\[
 S_\gamma=\sum_{\lambda\ne\gamma}
             \frac{m_\lambda}{(\lambda-\gamma)^2}<\infty.
\]
Indeed the sum is finite locally after removing the isolated
ordinate, and the tail converges by the unit-interval count
\(O(\log(|\lambda|+3))\).

Both \(K\) and \(K+vK'\) are Schwartz functions. Hence, for
\(R\ge1\) and \(d\ne0\),
\[
 |dK(Rd)|\le C R^{-2}|d|^{-1},
 \qquad
 |K(Rd)+RdK'(Rd)|\le C R^{-1}|d|^{-1}.
\]
The first estimate follows from the boundedness of \(v^2K(v)\),
and the second from the boundedness of \(v(K(v)+vK'(v))\).
Summing their squares against the zero weights proves
\[
 q(f_R^{\,0})\le C R^{-4}S_\gamma\longrightarrow0,
\]
\[
 \|Xf_R^{\,0}-i\mathbf1_{\{\gamma\}}\|_q^2
 \le C R^{-2}S_\gamma\longrightarrow0.
\]

For the mean correction choose one fixed
\(b\in C_c^\infty(\mathbb R)\) with \(\int b=1\), and set
\[
 f_R=f_R^{\,0}-F_R^{\,0}(0)b.
\]
There is no zero ordinate, so
\(F_R^{\,0}(0)=-\gamma K(-R\gamma)\to0\).
Both \(b\) and \(Xb\) have finite Weil norm. The correction
therefore vanishes in both displayed limits and makes
\(\int f_R=0\). The limiting graph pair is
\((0,i\mathbf1_{\{\gamma\}})\), whose second entry is nonzero.
This is exactly the obstruction to closability.
The same sequence can have its first physical moment zero as well.
Indeed \((F_R^{\,0})'(0)\to0\). Choose fixed compact smooth
functions \(b_0,b_1\) with moment pairs
\((\int b_j,\int xb_j)=(1,0),(0,1)\), respectively, and subtract
the two moments of \(f_R^{\,0}\) times these functions. Both
coefficients tend to zero, so both graph limits remain unchanged.
Then \(f_R\) and \(Xf_R\) are both mean-zero tests. For example,
one may take an even mean-one \(b_0\) and \(b_1=-b_0'\).
\end{proof}

The trace comparison can be stated without identifying two different
trace functionals. The unitary map
\[
 U:\ell^2(\Gamma,m)\longrightarrow\ell^2(\Gamma),
 \qquad (Ua)_\gamma=\sqrt{m_\gamma}\,a_\gamma
\]
conjugates every diagonal translation
\(a_\gamma\mapsto e^{-it\gamma}a_\gamma\), and every defined
diagonal function of \(D\), to the same diagonal operation. In
particular, for a spectral function \(v\ge0\) with
\(\sum_\gamma m_\gamma v(\gamma)<\infty\),
the ordinary Hilbert-space trace is
\[
 \operatorname{tr}v(D)=\sum_{\gamma\in\Gamma}v(\gamma),
\]
whereas the zero-divisor functional is
\[
 \tau_{\rm div}(v)=\sum_{\gamma\in\Gamma}m_\gamma v(\gamma).
\]
Their exact difference is
\(\sum_\gamma(m_\gamma-1)v(\gamma)\); for
\(v(\gamma)=\gamma^{-2k}\), it is the limiting trace defect
already proved above. Thus conjugating the commuting translation
operators does not change their ordinary trace or manufacture this
additional term.

The form commutator \eqref{eq:fd-prime-commutator} retains the
derivative evaluations and hence additional information about the
physical embedding. The proposition shows why passing directly from
that identity to a cyclic trace of a global closed multiplication
operator is invalid: the stated test operator has the explicit
nonzero graph residual above. A bound controlling these derivative
channels in a suitable growing-window trace comparison has not been
established here. The prime-translation identity remains exact, and
no conclusion about the impossibility of an arithmetic comparison
is drawn from this failed operator passage.

\subsection{Explicit global growth from certified heights}
\label{sec:fd-global-growth-numeric}

The arithmetic certificate also gives explicit leading constants
in the global detection laws. We first control the multiplicity at an
individual ordinate by an elementary positive logarithmic derivative.

\begin{lemma}[An explicit multiplicity bound]
\label{lem:fd-multiplicity-log}
Every nontrivial zero at positive ordinate \(\gamma\ge10\) has
multiplicity
\[
 m_\gamma<\frac34\log\gamma.
\]
\end{lemma}
\begin{proof}
The paired Hadamard product and the zero-location theorem give
\[
 \frac{2m_\gamma}{3}
 \le\Re\frac{\xi'}{\xi}(2+i\gamma).
\]
Indeed, the zero at \(1/2+i\gamma\) contributes exactly
\(m_\gamma/(3/2)\), and every other real part in the paired
logarithmic-derivative sum is positive.

For real \(y\), differentiating Euler's product for the gamma
function gives
\[
 \Re\psi(1+iy)=-\gamma_E+
       \sum_{n\ge1}\frac{y^2}{n(n^2+y^2)}.
\]
The summand, as a function of positive real \(n\), decreases.
Bounding its sum by its first term plus its integral from one to
infinity yields
\[
 \Re\psi(1+iy)
 \le-\gamma_E+\frac{y^2}{1+y^2}
                +\frac12\log(1+y^2)
 <1-\gamma_E+\frac12\log(1+y^2).
\]
The absolutely convergent Euler series gives
\[
 A_2=-\frac{\zeta'}{\zeta}(2)=\sum_{n\ge2}\frac{\Lambda(n)}{n^2}
 <\frac35.
\]
Here this last rational bound has a short finite certificate: summing
the \(193\) prime-power terms through \(1000\) and using
\(\Lambda(n)\le\log n\) gives
\[
 A_2\le\sum_{n\le1000}\frac{\Lambda(n)}{n^2}
             +\frac{\log1000+1}{1000}
 <0.57687<\frac35.
\]
The same scalar certificate verifies
\(\gamma_E\ge\sum_{j=1}^{6}j^{-1}-\log7>1/2\) and
\(\log(2\pi)>9/5\).

The completion formula, with \(y=\gamma/2\), now implies
\begin{align*}
 \Re\frac{\xi'}{\xi}(2+i\gamma)
 &<\frac3{\gamma^2}+\frac35-\frac12\log\pi
     +\frac12-\frac12\gamma_E
     +\frac14\log(1+\gamma^2/4)\\
 &=\frac12\log\gamma+\frac3{\gamma^2}+\frac{11}{10}
     -\frac12\{\gamma_E+\log(2\pi)\}
     +\frac14\log(1+4/\gamma^2)\\
 &<\frac12\log\gamma+\frac4{\gamma^2}-\frac1{20}
 <\frac12\log\gamma\qquad(\gamma\ge10).
\end{align*}
In the first line, the rational completion terms are bounded by
\(2/\gamma^2+1/\gamma^2\), and the Euler series has real part at
most \(A_2\). The final two inequalities use
\(\log(1+x)<x\) and \(4/\gamma^2\le1/25\).
Combining the first and last bounds proves the lemma.
\end{proof}

Write the extended global limits as
\[
 L_T=\lim_{a\downarrow0}\sqrt a\,m_*(-a)
     =\lim_{a\downarrow0}\sqrt a\,R_*(-a),\qquad
 L_J=\lim_{a\downarrow0}a\,j_*(-a).
\]

\begin{corollary}[Certified global constants and their comparison]
\label{cor:fd-global-growth-numeric}
For the actual zeta heat family,
\[
 L_T>1.65\cdot10^{12},\qquad L_J>4.5\cdot10^{24}.
\]
The arithmetic trace certificate through $97.632$ by itself gives
$L_T>216.88$ and $L_J>4766.003712$. The larger bounds use the
Platt--Trudgian simplicity certificate through $3\cdot10^{12}$
\cite[abstract, p.~792; preprint Section~2, pp.~2--3]{platt-trudgian2021}.
Moreover,
\[
 \frac{4\pi}{3}
     \sqrt{\frac{L_J}{\log\sqrt{2L_J}}}
 \le L_T\le
 \frac{\sqrt3\,\pi}{2}
     \sqrt{L_J\log\sqrt{2L_J}}.
\]
The two functions of \(L_J\) are interpreted as \(+\infty\) at
\(+\infty\).
\end{corollary}
\begin{proof}
The arithmetic certificate proves that the first multiple ordinate
\(\gamma_*\) exceeds \(97.632\), including the value infinity.
The exact limits in Theorem~\ref{thm:fd-limits}, with \(x=2\gamma\),
are
\[
 L_T=\min_{m_\gamma\ge2}\frac{\gamma\mathfrak r_{m_\gamma}}{\sqrt2},
 \qquad L_J=\frac{\gamma_*^2}{2}.
\]
Thus \(L_J>(97.632)^2/2=4766.003712\).

For multiplicities two and three, the exact radii
\(\mathfrak r_2=\pi\), \(\mathfrak r_3=2\pi/\sqrt3\) give a
Toeplitz contribution strictly greater than
\[
 \frac{97.632\,\pi}{\sqrt2}>216.88.
\]
For an integer multiplicity \(m\ge4\), the preceding lemma gives
\(\gamma>e^{4m/3}\). Proposition~\ref{prop:orbital-unimodal-bound}
therefore gives
\[
 \frac{\gamma\mathfrak r_m}{\sqrt2}
 >\pi\sqrt{\frac23}\,e^{4m/3}\frac{\sqrt{m+1}}m.
\]
The logarithmic derivative of the last function of real \(m\) is
\(4/3+1/(2(m+1))-1/m>0\) for \(m\ge4\). Its value at four is
\[
 \frac\pi4\sqrt{\frac{10}3}\,e^{16/3}>297.
\]
This covers every possible multiple zero and proves the stated
lower bound for \(L_T\).

For the comparison, the same orbital estimate implies
\(\mathfrak r_m>2\pi/\sqrt{3m}\). At every multiple ordinate,
Lemma~\ref{lem:fd-multiplicity-log} consequently gives
\[
 \frac{\gamma\mathfrak r_m}{\sqrt2}
 >\frac{2\sqrt2\,\pi}{3}
             \frac{\gamma}{\sqrt{\log\gamma}}.
\]
The function \(\gamma/\sqrt{\log\gamma}\) increases for
\(\gamma\ge10\), so the minimum is bounded below by its value
at \(\gamma_*\). For the upper bound, evaluate the same minimum
at the first multiple ordinate and use the previously proved
\(\mathfrak r_m\le\pi\sqrt{m-1}\). This gives
\[
 L_T\le\pi\sqrt{\frac38}\,
          \gamma_*\sqrt{\log\gamma_*}.
\]
Substituting \(\gamma_*=\sqrt{2L_J}\) proves both displayed
comparisons for finite \(\gamma_*\). The infinite case follows
from the exact simplicity equivalences in
Theorem~\ref{thm:fd-limits} and the stated extended interpretation.

Finally, the Platt--Trudgian certificate gives
\(\gamma_*>H:=3\cdot10^{12}\). Substitution into the exact
Jensen limit and the increasing lower comparison just proved gives
\[
 L_J>H^2/2=4.5\cdot10^{24},\qquad
 L_T>\frac{2\sqrt2\pi}{3}\frac{H}{\sqrt{\log H}}
                         >1.65\cdot10^{12}.
\]
For the last strict comparison it suffices to use
\(\sqrt2>1.414\), \(\pi>3.14\), \(\log H<28.8\), and
\(\sqrt{28.8}<5.367\); their substitution gives
\(2(1.414)(3.14)/5.367>1.65\). These inequalities are also
certified by exact rational and outward-ball arithmetic.
\end{proof}

All decimal comparisons in this subsection are strict outward-ball
inequalities in the accompanying scalar certificate. The smaller input height
is the arithmetic trace certificate, the larger is the cited
Platt--Trudgian packet, and the Euler sum
uses prime powers only through \(1000\). These are fixed lower bounds
and a quantitative comparison of the global constants. The question
whether the constants are infinite remains the all-height simplicity
question.

\subsection{An exact commutator with unbounded quotient norm}
The following calculation records why a shrinking parity gap requires
an estimate beyond strong convergence and compactness.
On two copies of $\ell^2(\mathbb N)$ let
$D=\operatorname{diag}(2^{-j})_{j\ge1}$ and
$C=\left(\begin{smallmatrix}0&D\\D&0\end{smallmatrix}\right)$.
For $n\ge1$ put $\tau=2^{-n}$, $\kappa=\tau^4$, $r=1-\kappa$, and
$v_{n+k}=2^{-(k+1)/2}$ for $k\ge0$, with all earlier coordinates zero.
Thus $\|v\|=1$ and $v\notin\operatorname{ran}D$. Define
\[
 H=\tau^4D^2,\quad
 a=\langle(I-H)^{-1}v,v\rangle^{-1},\quad
 b=\langle(rI-H)^{-1}Dv,Dv\rangle^{-1},
\]
\[
 M_+=H+a\,v\otimes v^*,\qquad
 M_-=H+b\,Dv\otimes(Dv)^*,\qquad S_\pm=I-M_\pm.
\]
The top eigenvalues of $M_+,M_-$ are respectively $1,r$, both simple.
Indeed, for $t>\|H\|$ and $c=\langle(tI-H)^{-1}u,u\rangle^{-1}$,
\[
 tI-H-c\,u\otimes u^*
 =(tI-H)^{1/2}(I-cw\otimes w^*)(tI-H)^{1/2},\quad
 w=(tI-H)^{-1/2}u,
\]
and $c\|w\|^2=1$. This proves positivity and the one-dimensional
kernel. Each $M_\pm$ is compact and strictly positive on nonzero
vectors because $H$ has positive diagonal entries.

Let $e,f$ be the unit normalizations of $(I-H)^{-1}v$ and
$(rI-H)^{-1}Dv$. The tail supports escape to infinity,
$a\le1$, and $b\|Dv\|^2\le r$, so $M_\pm\to0$ strongly.
Finite-rank approximation of $D$ yields
$\|M_-D\|_{\mathcal S_2}+\|DM_+\|_{\mathcal S_2}\to0$.
Nevertheless the quotient operator below becomes unbounded in norm.

The exact commutator is
\[
 DM_+-M_-D=Dv\otimes\ell^0,\qquad
 \ell^0(h)=a\langle h,v\rangle-b\langle h,D^2v\rangle.
\]
Evaluating at $e$ gives $S_-De=Dv\,\ell^0(e)$.
Since $S_-\ge\kappa I$ and $D$ is injective, $\ell^0(e)\ne0$.
Set $\ell=\ell^0/\ell^0(e)$ on the even copy and zero on the odd
copy. Then
\[
 DM_+-M_-D=S_-De\otimes\ell,\qquad
 D(I-e\otimes\ell)=S_-^{-1}DS_+.
\]
The quotient of $C(I-e\otimes\ell)$ by $\mathbb Ce$, with metric
$I-M$, is therefore represented isometrically by
\[
 B_n=\begin{pmatrix}0&S_+^{1/2}DS_-^{-1/2}\\
 S_-^{-1/2}DS_+^{1/2}&0\end{pmatrix}
\]
off its zero vector $e$. At each $n$ it is compact and selfadjoint.
It is injective on the quotient: the lower-left block has kernel
$\mathbb Ce$; the kernel of its adjoint requires
$DS_-^{-1/2}h\in\mathbb Ce$. Since
$e_j$ is proportional to $v_j/(1-\tau^4 2^{-2j})$, it does not belong
to $\operatorname{ran}D$, and hence $h=0$.

For the norm, put $f^0=Dv/\|Dv\|$. Exact geometric sums give
\[
 \|Dv\|^2=\frac47\tau^2,\quad
 \|D^2v\|^2=\frac{16}{31}\tau^4,\quad
 \langle Df^0,(I-v\otimes v^*)Df^0\rangle
       =\frac{72}{217}\tau^2.
\]
On the common tail, $\|D\|=\tau$ and $\|H\|=\eta=\tau^6$.
Since $r\ge15/16$ and $\eta\le1/64$, the resolvent geometric series
and normalization give
$\|f-f^0\|\le2(\eta/r)/(1-\eta/r)\le3\eta$.
Also $1-\eta\le a\le1$, so
$\|S_+-(I-v\otimes v^*)\|\le2\eta$ on that tail.
It follows that
\[
 \langle Df,S_+Df\rangle
 \ge\frac{72}{217}\tau^2-8\tau^8
 \ge\frac{359}{1736}\tau^2>\frac15\tau^2.
\]
Since $S_-f=\kappa f$, one obtains
\[
 \|B_n\|^2\ge\|B_n(0,f)\|^2
       =\kappa^{-1}\langle Df,S_+Df\rangle
       \ge\frac1{5\tau^2}\longrightarrow\infty.
\]
The construction realizes the same rank-one commutator equation and
compact-product limits appearing in the parity comparison. Its
quantified divergence shows that those displayed data alone do not
bound the quotient norm. Applying the parity method to the Weil
operator still requires control of its gap or of the weighted product
in \eqref{eq:parity-error}.

\section{Arithmetic identities at the multiplicity boundary}
\label{sec:boundary-multiplicity}
This section expresses boundary multiplicities through the Euler
product, local conductor counts, and the derivative companion of
$\Xi$. The formulas identify the exact maps between arithmetic
residues, Fourier evaluations, global jets, and spectral trace defects.
They also locate the boundary estimates still needed for simplicity.

\subsection{A positive Poisson sum}
For $a>0$ and real $t$, define
\[
 P_a(t)=a\Re\frac{\xi'}{\xi}(\tfrac12+a+it).
\]
The centered even Hadamard product gives the absolutely convergent
identity
\begin{equation}\label{eq:poisson-multiplicity}
 P_a(t)=\sum_{\gamma\in\Gamma}
             m_\gamma\frac{a^2}{a^2+(t-\gamma)^2}.
\end{equation}
To verify its normalization, differentiate each paired factor
$\bigl(1+(s-1/2)^2/\gamma^2\bigr)^{m_\gamma}$.
Its logarithmic derivative is the sum of the fractions at
$1/2+i\gamma$ and $1/2-i\gamma$; their real parts give the two
terms displayed above. The zero count makes the tails locally
uniformly convergent for $a>0$.

At any zero ordinate $\gamma_0$ the corresponding term is exactly
$m_{\gamma_0}$. Every other term decreases as $a$ decreases to
zero, and its value at one fixed positive $a$ is a summable
majorant. Consequently
\begin{equation}\label{eq:poisson-infimum}
 \inf_{a>0}P_a(\gamma_0)
 =\lim_{a\downarrow0}P_a(\gamma_0)=m_{\gamma_0}.
\end{equation}
Thus simplicity at this ordinate is equivalent to the existence
of $a>0$ with $P_a(\gamma_0)<2$.

For $a>1/2$, put $s=1/2+a+it$. The absolutely convergent Euler
series gives exactly the same quantity as
\begin{equation}\label{eq:poisson-euler}
 P_a(t)=a\left[
 \Re\left\{\frac1s+\frac1{s-1}+\frac12\psi(s/2)\right\}
 -\frac12\log\pi
 -\sum_{n\ge2}\frac{\Lambda(n)}{n^{1/2+a}}\cos(t\log n)
 \right].
\end{equation}
The distance from the line matters quantitatively. The argument
bound $S(t)=O(\log t/\log\log t)$, available from the zero-location
theorem and \cite[Theorem~1, p.~2]{carneiro-chandee-milinovich2015},
implies
\begin{equation}\label{eq:poisson-exterior-size}
 P_{1/2}(t)=\tfrac14\log(t/(2\pi))
                +O(\log t/\log\log t).
\end{equation}
Here are the details of the passage to this formula. Integrate
the kernel $(1/4)/(1/4+(u-t)^2)$ against the positive zero-counting
measure. On $t/2\le u\le3t/2$, the main density is
$(2\pi)^{-1}\log(u/(2\pi))$; replacing its logarithm by
$\log(t/(2\pi))$ costs $O(\log t/t)$.
The integral of the kernel on the whole line is $\pi/2$.
Integration by parts against the argument error costs
$O(\log t/\log\log t)$, because the kernel has bounded total
variation. The remaining positive and negative tails cost
$O(\log t/t)$ by the zero count. This proves
\eqref{eq:poisson-exterior-size}.
Since \eqref{eq:poisson-multiplicity} increases with $a$, every
$a\ge1/2$ gives a value exceeding $2$ at all sufficiently large
heights. The upper threshold in \eqref{eq:poisson-infimum}
therefore requires control closer to the line. The logarithmic
derivative estimates of Chirre--Gon\c calves
\cite[Theorems~1--2 and (1.1)]{chirre-goncalves2022}
give explicit bounds in their stated distance range; no upper bound
below $2$ at every zero follows from the estimates used here.

\subsection{The prime ramp and its individual coefficients}
For $y>1$, define the arithmetic function
\[
 G(y)=\sum_{n\le y}\Lambda(n)(y/n-1)-y\log y+(1+\gamma_E)y.
\]
Direct integration of each summand, absolutely in $\Re s>1$,
gives
\begin{equation}\label{eq:prime-ramp-mellin}
 \int_1^\infty G(y)y^{-s-1}\,dy
 =-\frac{\zeta'(s)}{s(s-1)\zeta(s)}
       -\frac1{(s-1)^2}+\frac{1+\gamma_E}{s-1}.
\end{equation}
The pole at $s=1$ cancels, by
$\zeta'/\zeta=-1/(s-1)+\gamma_E+O(s-1)$.
At zero the residue is $\log(2\pi)$. At a nontrivial zero
$\rho$ it is $-m_\rho/[\rho(\rho-1)]$, and at $-2n$ it is
$-1/[2n(2n+1)]$. Mellin inversion therefore gives
\begin{equation}\label{eq:prime-ramp-zero-sum}
 G(y)=\log(2\pi)
  +\sum_{\gamma\in\Gamma}
       \frac{m_\gamma y^{1/2+i\gamma}}{1/4+\gamma^2}
  -\sum_{n\ge1}\frac{y^{-2n}}{2n(2n+1)}.
\end{equation}
For completeness, apply Mellin inversion first to
\[
 A(y)=\sum_{n\le y}\Lambda(n)(y/n-1),\qquad
 \int_1^\infty A(y)y^{-s-1}\,dy
       =-\frac{\zeta'(s)}{s(s-1)\zeta(s)}.
\]
On every fixed line $\Re s>1$, this transform is
$O((1+|\Im s|)^{-2})$, so its inversion is absolutely convergent.
At $s=1$ its Laurent terms are
$(s-1)^{-2}-(1+\gamma_E)(s-1)^{-1}$; after multiplication
by $y^s$, their residue is
$y\log y-(1+\gamma_E)y$.
Truncate the inversion at ordinates a distance
$\gg1/\log H$ from all zeros in the neighboring unit intervals.
The partial fractions and the unit-interval zero count give
$\zeta'/\zeta=O(\log^2H)$ on the horizontal portions of a fixed
strip. Their integrals tend to zero after division by $s(s-1)$.
Shift the left line through negative odd integers; the functional
equation bounds its logarithmic derivative by
$O(\log(|s|+2))$ away from the trivial zeros, and $y^{\Re s}$
makes its contribution tend to zero. The residue sums converge
absolutely: for the nontrivial zeros this follows from
$\sum m_\gamma/(1+\gamma^2)<\infty$, and for the trivial zeros
it follows from the displayed geometric factor. This justifies
the successive contour limits for $A(y)$. Subtracting the
explicit $s=1$ contribution, exactly as in the definition of
$G(y)$, proves \eqref{eq:prime-ramp-zero-sum}.

Put $R(u)=e^{-u/2}G(e^u)$. Its difference from
\[
 \sum_{\gamma\in\Gamma}
       \frac{m_\gamma}{1/4+\gamma^2}e^{i\gamma u}
\]
is integrable on $u\ge0$, by \eqref{eq:prime-ramp-zero-sum}.
The series converges uniformly and absolutely. Termwise averaging,
followed by dominated convergence over its summable coefficients,
therefore proves
\begin{equation}\label{eq:prime-ramp-coefficient}
 \lim_{U\to\infty}\frac1U\int_0^U
            R(u)e^{-i\gamma_0u}\,du
       =\frac{m_{\gamma_0}}{1/4+\gamma_0^2}.
\end{equation}
This is an exact map from the prime expression to each multiplicity.
An upper bound below $2/(1/4+\gamma_0^2)$ for this coefficient,
at every zero ordinate, would be precisely the simplicity
inequality. That arithmetic estimate remains open here.
The global identity
\[
 \sum_{\gamma\in\Gamma}\frac{m_\gamma}{1/4+\gamma^2}
       =2+\gamma_E-\log(4\pi)
\]
follows by evaluating twice the centered logarithmic derivative
at $s=1$. Its individual-coefficient consequence permits a bound
growing quadratically with $\gamma_0$ after the factor
$1/4+\gamma_0^2$ is restored.

One can see directly why a positive boundary residue and a
square-root size bound permit arbitrary integer multiplicities
in an analytic model. Fix $\gamma>0$, let $D=\gamma^2+1/4$,
and put
\[
 F_m(s)=\bigl(s(s-1)+D\bigr)^m,\qquad
 k_m=m/D,\qquad m=1,2,\ldots.
\]
Since $F_m'(1)/F_m(1)=m/D$, exact cancellation gives
\[
 -\frac{F_m'(s)/F_m(s)-sF_m'(1)/F_m(1)+k_m(1-s)}{s(s-1)}
 =\frac1D\frac{F_m'(s)}{F_m(s)}.
\]
The right side is the Laplace transform, initially for
$\Re s>1/2$, of
\[
 g_m(u)=\frac{2m}{D}e^{u/2}\cos(\gamma u).
\]
Its two residues are $m/D>0$ and
$|g_m(\log y)|\le(2m/D)\sqrt y$, while $F_m$ has order $m$
at both prescribed boundary zeros. These explicit formulas
verify the compatibility at every positive integer $m$.
The prime identity \eqref{eq:prime-ramp-mellin} supplies the
additional arithmetic structure in the zeta problem.

\subsection{Local zero counts for growing conductor}
\label{subsec:conductor-local-count}
The accepted Dirichlet zero-location input gives a local obstruction
to approximation across the critical line that also allows
holomorphic normalizing factors.

\begin{proposition}\label{prop:conductor-local-count}
Fix $t_0\in\mathbb R$ and $r>0$. For a primitive nonprincipal
character $\chi$ of conductor $q$, let $n_\chi(t_0,r)$ count the
nontrivial zeros, with their analytic multiplicities, whose
ordinates satisfy $|\gamma-t_0|<r$. There is a finite constant
$C_{t_0,r}$, independent of $q$ and $\chi$, such that
\[
 n_\chi(t_0,r)\ge \frac r{12}\log q-C_{t_0,r}.
\]
Let $U$ be an open neighborhood of $s_0=1/2+it_0$, let
$q_j\to\infty$ be conductors of primitive nonprincipal
characters $\chi_j$, and let $a_j$ be holomorphic and zero-free
on $U$. Every locally uniform holomorphic limit on $U$ of
$a_j(s)L(s,\chi_j)$ vanishes identically on the connected
component of $U$ containing $s_0$.
\end{proposition}

\begin{proof}
For the Fourier convention
$\widehat h(v)=\int_{\mathbb R}h(t)e^{-itv}\,dt$, put
\[
 w(y)=(1-y^2)\left(\frac{\sin y}{y}\right)^4,\qquad
 h(t)=w((t-t_0)/r).
\]
Removable values are understood. On the real line,
$h\le{\bf1}_{\{|t-t_0|<r\}}$: inside the interval its value
lies between zero and one, outside it is nonpositive, and at
the two endpoints it is zero. Moreover $\|h\|_1\le4r$, since
$|w(y)|\le1$ for $|y|\le1$ and $|w(y)|\le y^{-2}$ otherwise.
It is entire of exponential type $4/r$ and has decay
$O_{r,t_0}((1+|\Re z|)^{-2})$ on every fixed horizontal strip.

The transform identity
\[
 \widehat{(\sin y/y)^2}(v)=\frac\pi2(2-|v|)_+
\]
follows by convolving the transforms in
$\sin y/y=\frac12\int_{-1}^1e^{iyu}\,du$.
Plancherel, applied also to the weak derivative of this
triangle, gives
\[
 \int_{\mathbb R}\left(\frac{\sin y}{y}\right)^4dy
   =\frac{2\pi}{3},\qquad
 \int_{\mathbb R}y^2\left(\frac{\sin y}{y}\right)^4dy
   =\frac\pi2.
\]
Consequently
\[
 \int_{\mathbb R}h(t)\,dt=\frac{\pi r}{6},\qquad
 \operatorname{supp}\widehat h\subset[-4/r,4/r].
\]
The support assertion follows because the fourth power has
transform supported on $[-4,4]$, and multiplication by
$1-y^2$ adds its second transform derivative without
enlarging the support.

Write $a=0$ or $1$ according as $\chi(-1)=1$ or $-1$, and
$\psi=\Gamma'/\Gamma$. The primitive Guinand--Weil formula
in this convention is
\begin{align*}
 \sum_\rho h((\rho-\tfrac12)/i)
 ={}&\frac{\log(q/\pi)}{2\pi}\int_{\mathbb R}h(t)\,dt\\
 &+\frac1{2\pi}\int_{\mathbb R}h(t)
       \Re\psi\left(\frac14+\frac a2+\frac{it}{2}\right)dt\\
 &-\frac1{2\pi}\sum_{n\ge2}\frac{\Lambda(n)}{\sqrt n}
   \left\{\chi(n)\widehat h(\log n)
       +\overline{\chi(n)}\widehat h(-\log n)\right\}.
\end{align*}
See \cite[Section~2.1, equation~(2.1), p.~3]{chirre-pereira-delaat2021};
that source uses the Fourier exponential $e^{-2\pi ixt}$,
so its transform arguments $\pm\log n/(2\pi)$ become the
arguments displayed here. The conductor factor follows
from $(q/\pi)^{(s+a)/2}$ in the completion; reflection by
the functional equation gives the conjugate-character
term. There is no principal-character pole term.

Every zero lies on the critical line by the accepted input.
The zero sum is absolutely convergent by the decay of $h$
and the standard zero count. Its value is at most
$n_\chi(t_0,r)$. Only $n\le e^{4/r}$ occur in the prime sum,
whose absolute value is at most
\[
 \frac{4r}{\pi}\sum_{2\le n\le e^{4/r}}
                       \frac{\Lambda(n)}{\sqrt n}.
\]
The two possible gamma integrals are bounded independently
of $q$ and $\chi$, because
$\Re\psi(\alpha+it/2)=O(1+\log(2+|t|))$ for
$\alpha=1/4,3/4$. For example, this estimate follows from
\[
 \Re\psi(\alpha+iy)=\psi(\alpha)+
 \sum_{n\ge0}\frac{y^2}
 {(n+\alpha)((n+\alpha)^2+y^2)}
\]
by splitting at $n\asymp |y|$. Thus the formula gives
$\sum_\rho h(\gamma)=(r/12)\log q+O_{t_0,r}(1)$, proving
the stated bound. The constant is uniform in the conductor;
the radius is fixed throughout.

For the limit assertion, let $g$ be such a limit and suppose
its restriction to the component containing $s_0$ is not
identically zero. Choose $R>0$ with
$\overline{D(s_0,R)}\subset U$ and with no zero of $g$ on
the boundary. Isolatedness of the zeros permits this choice.
Local uniform convergence and Rouch\'e's theorem give a
fixed finite number of zeros in this disk for
$a_jL(\,\cdot\,,\chi_j)$ when $j$ is large. The factors
$a_j$ preserve every zero order. The first assertion with
$r=R/2$ instead gives at least
$(R/24)\log q_j-C_{t_0,R/2}$ zero occurrences on a critical
segment inside the disk, which tends to infinity. This
contradiction proves the component-wise conclusion.
\end{proof}

\subsection{The derivative companion and its common divisor}
\label{subsec:hb-boundary}
Put $F(z)=\Xi(z)$ and retain the distinct signed ordinates
$\Gamma$ and multiplicities $m_\gamma$. The paired canonical
product, with $\sum_\Gamma m_\gamma/\gamma^2<\infty$, gives
\[
 -\Im\frac{F'(z)}{F(z)}
 =(\Im z)\sum_{\gamma\in\Gamma}
          \frac{m_\gamma}{|z-\gamma|^2}>0
 \qquad(\Im z>0).
\]
Hence $E=F+iF'$ satisfies the upper-half-plane
Hermite--Biehler inequality
\[
 |E(z)|^2-|E^\#(z)|^2
  =4|F(z)|^2\left(-\Im\frac{F'(z)}{F(z)}\right)>0,
 \qquad E^\#(z)=\overline{E(\bar z)}=F-iF'.
\]
Reflection exchanges $E$ and $E^\#$ and preserves their
real boundary norm.

For the de Branges norm
$\|H\|_{\mathcal H(E)}^2=\int_{\mathbb R}|H(x)/E(x)|^2dx$,
where both $H/E$ and $H^\#/E$ belong to the upper
half-plane Hardy space, the kernel is
\begin{align}
 K_E(w,z)
 &=\frac{\overline{F(w)}F'(z)-\overline{F'(w)}F(z)}
                {\pi(\bar w-z)}\notag\\
 &=\frac1\pi\sum_{\gamma\in\Gamma}m_\gamma
       \frac{F(z)}{z-\gamma}
       \overline{\frac{F(w)}{w-\gamma}}.             \label{eq:hb-kernel}
\end{align}
The first equality follows by expanding the usual numerator
$\overline{E(w)}E(z)-\overline{E^\#(w)}E^\#(z)$ and dividing
by $2\pi i(\bar w-z)$. Subtracting the two logarithmic
derivatives gives the second. The series converges locally,
with removable values filled in. In particular
\begin{equation}\label{eq:hb-laguerre}
 \pi K_E(x,x)=F'(x)^2-F(x)F''(x).
\end{equation}
Its equality case at a multiple zero is also the $n=1$
case of \cite[Theorem~1.3 and Section~4, pp.~3, 10--12]{sra-xi2026}.

Define the convergent paired product
\[
 G(z)=\prod_{\gamma>0}(1-z^2/\gamma^2)^{m_\gamma-1},
 \qquad A=F/G,\quad B=F'/G,\quad \widetilde E=A+iB.
\]
At an order-$m$ zero, write
$F(z)=c_\gamma(z-\gamma)^m+O((z-\gamma)^{m+1})$.
Then $G$ has order $m-1$ there, $B(\gamma)=mA'(\gamma)\ne0$,
and
\[
 F'^2-FF''=mc_\gamma^2(z-\gamma)^{2m-2}
                       +O((z-\gamma)^{2m-1}).
\]
Thus $\widetilde E$ has no real zero, $E=G\widetilde E$, and
\[
 K_E(w,z)=G(z)\overline{G(w)}K_{\widetilde E}(w,z),
 \qquad K_{\widetilde E}(\gamma,\gamma)
                =mA'(\gamma)^2/\pi.
\]
The spaces satisfy the isometric identity
$\mathcal H(E)=G\mathcal H(\widetilde E)$.
Indeed, an entire $H$ with $H/E\in L^2(\mathbb R)$ must
vanish to at least the order of each real zero of $E$:
an integer deficit of one produces a nonintegrable
$|x-\gamma|^{-2}$ singularity. Therefore $H/G$ is entire,
and its two Hardy quotients and norm agree by cancellation.
Multiplication by $G$ proves the reverse inclusion.

The model-space inner function and its angular data are
\[
 \Theta=\frac{F-iF'}{F+iF'}
       =\frac{\widetilde E^\#}{\widetilde E},\qquad
 \Theta(\gamma)=-1,\qquad \Theta'(\gamma)=-\frac{2i}{m_\gamma}.
\]
Here $K(\Theta)=H^2(\mathbb C_+)\ominus\Theta H^2(\mathbb C_+)$,
with the boundary $L^2(dx)$ norm, and $K_\Theta$ denotes its
reproducing kernel.
The last two identities follow by substituting
$F/F'=(z-\gamma)/m_\gamma+O((z-\gamma)^2)$.
In the boundary $L^2(dx)$ normalization,
\[
 K_\Theta(\gamma,\gamma)=\frac1{\pi m_\gamma},
 \qquad
 \mu_\Theta(\{\gamma\})=\frac{2\pi}{|\Theta'(\gamma)|}
                       =\pi m_\gamma.
\]
These are the Clark masses for parameter $-1$ in
\cite[Section~2.4, equation~(2.8), and Proposition~4.1,
pp.~6--7, 10--11]{suzuki-weil2025}.

\subsection{Sampling, jets, and the boundary limit}
The normalized model-space basis from
\cite[Proposition~4.1]{suzuki-weil2025} is
\[
 e_\gamma(z)=i\sqrt{\frac{m_\gamma}{\pi}}\,
                  \frac{F(z)}{E(z)(z-\gamma)},\qquad
 e_\gamma(\eta)=\frac{\delta_{\gamma\eta}}{\sqrt{\pi m_\gamma}}.
\]
The source uses $\mathcal F_+h(\gamma)=\int h(x)e^{i\gamma x}dx$.
For the present convention $\widehat h=\mathcal F_-h$, the exact
conversion is $\mathcal F_+(h(-\,\cdot\,))(\gamma)=\widehat h(\gamma)$.
Equivalently $\mathcal F_+h(-\gamma)=\widehat h(\gamma)$;
these are alternative unitary reflections since
$m_{-\gamma}=m_\gamma$.
The weighted sampling isometry and this basis yield the unitary
\[
 \mathcal Uh
 =E\sum_{\gamma\in\Gamma}\sqrt{m_\gamma}\widehat h(\gamma)e_\gamma
 =\frac{i}{\sqrt\pi}\sum_{\gamma\in\Gamma}
                m_\gamma\widehat h(\gamma)\frac{F(z)}{z-\gamma}.
\]
Cauchy--Schwarz and $\sum_\Gamma m_\gamma/\gamma^2<\infty$
give locally uniform convergence of this series, with each removable
value filled in. Cauchy's formula therefore permits derivative
evaluation term by term. At an order-$m$ zero every nontarget
summand vanishes to order $m$, and hence
\begin{equation}\label{eq:hb-evaluation-jet}
 (\mathcal Uh)^{(m-1)}(\gamma)
       =\frac{i}{\sqrt\pi}F^{(m)}(\gamma)\widehat h(\gamma),
 \qquad
 (\mathcal Uh/G)(\gamma)
       =\frac{imA'(\gamma)}{\sqrt\pi}\widehat h(\gamma).
\end{equation}
This explicitly identifies the Fourier evaluations with
the corresponding global jets.

The multiplication extension in
\cite[Section~6, p.~18]{suzuki-weil2025} has the distinct
eigenvalues $\gamma$, each once. Its inverse, transported by
this unitary, is $C=\operatorname{diag}(1/\gamma)$. Therefore
the counted logarithmic trace and the distinct spectral
trace satisfy the exact identity
\[
 S_k-\operatorname{tr}C^{2k}
   =\sum_{\gamma\in\Gamma}\frac{m_\gamma-1}{\gamma^{2k}}
   =\sum_{\gamma\in\Gamma}
           \frac{2/|\Theta'(\gamma)|-1}{\gamma^{2k}}.
\]
This is also the limit of the finite defects $D_{N,k}$.

For $\omega>0$ put $E_\omega(z)=F(z+i\omega)$ and let
$K_\omega$ be its de Branges kernel. Direct differentiation
and the logarithmic derivative give
\[
 K_\omega(x,x)=\frac{\omega}{\pi}|F(x+i\omega)|^2
  \sum_{\eta\in\Gamma}\frac{m_\eta}{(x-\eta)^2+\omega^2},
 \qquad \omega^{-1}K_\omega(w,z)\longrightarrow K_E(w,z)
\]
locally uniformly. At an order-$m$ zero the exact leading
scale is
\begin{equation}\label{eq:hb-shift-scale}
 K_\omega(\gamma,\gamma)
   =\frac{m|c_\gamma|^2}{\pi}\omega^{2m-1}
                           (1+O_\gamma(\omega^2)).
\end{equation}
For the error term, the remaining zero sum is bounded near
$\gamma$, and reality of the Taylor coefficients gives
$|F(\gamma+i\omega)|^2
=|c_\gamma|^2\omega^{2m}(1+O_\gamma(\omega^2))$.

Strict interior positivity has a precise boundary limitation.
For $U(z)=|E(z)|^2-|E^\#(z)|^2$, the fact that $F'$ also
has only real zeros gives
\[
 \begin{gathered}
 \Delta U=4\bigl(|E'|^2-|(E^\#)'|^2\bigr)\ge0\quad(\Im z>0),\\
 U(\gamma+iy)=4m|c_\gamma|^2y^{2m-1}+O(y^{2m}).
 \end{gathered}
\]
The real-zero assertion for $F'$ follows by differentiating
the locally convergent paired polynomial products for $F$,
using Rolle's theorem and then Hurwitz's theorem.
The inward derivative vanishes at a multiple zero. The
subharmonic sign does not give the positive-superharmonic
boundary minimum inference. Explicitly, $F(z)=z^2$ gives
$U(x+iy)=8y(x^2+y^2)>0$ in the upper half-plane and
$\Delta U=64y>0$, with zero inward derivative at the origin.
The positive-kernel surgery already proved in this paper
gives the same cubic behavior in \eqref{eq:hb-shift-scale}
at its imposed double ordinate.

For comparison with the arithmetic canonical construction,
\cite[equations~(2.4)--(2.13), Theorems~2.1--2.4,
pp.~6--9]{suzuki-hamiltonians2021} uses
$E_{\omega,\nu}=E_\omega^\nu$ with $\nu\omega>1$ for zeta,
and obtains a positive determinant-one Hamiltonian for
each fixed permitted pair under the accepted zero-location
input. At real points the power has the exact kernel effect
\[
 K_{\omega,\nu}(x,x)
       =\nu|E_\omega(x)|^{2\nu-2}K_\omega(x,x).
\]
The displayed source results retain their fixed-parameter
quantifiers. A positive linear lower bound for
$K_\omega(\gamma,\gamma)$ as $\omega\downarrow0$ at every
zero ordinate remains unproved here; by
\eqref{eq:hb-shift-scale} it is the original multiplicity
margin.


\section{The estimates that remain open}
The all-simplicity question is the unboundedness of the exact finite
heights in Theorem~\ref{thm:fd-arithmetic-complete}:
\[
 \sup_{N,k}\left(\frac2{S_k-\operatorname{tr}C_N^{2k}}\right)^{1/(2k)}
          =+\infty.
\]
The arithmetic computations establish this height bound through
$97.632$. The equivalent assertion at every height remains open.
The dimension and fixed-space limits in
Proposition~\ref{prop:fd-certificate-size} identify two constraints
on a quantitative proof: its dimension must reach twice the counted
packet size, and increasing the moment order in one fixed space
eventually returns the threshold to the first ordinate. The unresolved
estimate concerns the finite arithmetic defect with both indices
allowed to grow.

Equivalently, the nonnegative integer
\(n(T)-\operatorname{ind}_+(T H_N-G_N)\) must vanish for finite
indices at unbounded zero-free heights. The logarithmic-support
construction proves a lower inertia density \(d_{\rm MT}-\eta\)
at every sufficiently large height; the complementary integer deficit
is still uncontrolled. Equation~\eqref{eq:fd-trace-inertia-integral}
transports this assertion exactly to the trace family. The rank-four
comparison in \eqref{eq:fd-inertia-surgery} preserves the density
conclusion and changes the limiting integer deficit by one.

The prime form also has the exact multiplication identity
\eqref{eq:fd-prime-commutator}, whose spectral side consists of
derivative evaluations at the zeros. Proposition~\ref{prop:fd-multiplication-domain}
constructs compact mean-zero tests tending to zero in the Weil norm
whose products by the coordinate converge to a nonzero spectral
coordinate. Thus this test operator has no closed extension, and a
cyclic trace argument through that proposed extension cannot provide
the missing divisor trace. The explicit prime form identity remains
available; a quantitative comparison retaining its derivative terms is
still unproved.

The derivative-moment formulation has an equally precise absolute
threshold. On every bounded packet, the minimum of the polynomial
square sum with constant term one is the total length of the
multiple-zero fibers. Its value below $2$ certifies simplicity.
The contour formula and coefficientwise error budget in
Proposition~\ref{prop:fd-moment} specify a finite test. Leading
asymptotic derivative moments allow errors proportional to the
number of zeros, while this test asks for a certified absolute
margin below $2$. The exact positive-kernel replacement in
Section~\ref{sec:fd-positive-surgery} preserves the stated leading
moments, simple-zero proportions, and zero heat threshold, while
its divisor and prime-form corrections are explicit. Thus the
missing step in this route must use arithmetic information beyond
those transported asymptotic quantities. The boundary identities
\eqref{eq:poisson-infimum} and \eqref{eq:prime-ramp-coefficient}
express that information as two further exact scalar thresholds;
their required upper estimates remain unproved.

For the original prolate trial, the unresolved ratio is
$\liminf q(v_L)/a_2(L)=0$. The trial upper estimate has scale
$\lambda^{9+\varepsilon}e^{-4\pi\lambda^2}$, where
$\lambda=e^L$. The finite matrix bounds of
Section~\ref{sec:finite-weil-lower} give a concrete lower certificate
with a frequency tail and a quadratic coupling error. No arithmetic
estimate establishing a sufficient positive margin on an unbounded
sequence of windows has been proved. The comparison with the odd
trial places the remaining positive polynomial gain in
$0<\alpha\le4$, meaning the absolute scale
$\lambda^{9+\alpha}e^{-4\pi\lambda^2}$.
The all-window lower rate in Theorem~\ref{thm:cosine-lower-rate}
provides \(\exp\{-C\lambda^{2/d_{\rm MT}}\log\lambda
(1+\log\log\lambda)\}\) for all three variational quantities.
Its explicit exponent comparison with the trial leaves the required
ratio estimate open.
The compact-operator example proves that the exact commutator and
the indicated compactness estimates permit a divergent quotient
norm; every transfer from them to the desired denominator must
therefore include an additional quantitative estimate for this form.

M\"obius inversion gives the exact weighted physical norm in
\eqref{eq:mobius-physical-plancherel} and the local pole masses in
\eqref{eq:mobius-abel}. Its compact smoothing of normalized Mertens
is \eqref{eq:mobius-mertens-smoothing}. The fixed positive shifts
are controlled; an \(O(\epsilon^{-1})\) squared-norm estimate as
\(\epsilon\downarrow0\) is not established. The explicit odd test
in that section also shows that the one-mode orthogonality constraint
leaves these pole evaluations present. These formulas preserve the
higher-pole and derivative factors that a sharper arithmetic estimate
must control.

For Burnol's local jets, Proposition~\ref{prop:burnol-coupled-cutoff}
computes the discarded Weil norm for a fixed polynomial and a fixed
ordinate. Its compact Rayleigh quotient is asymptotic to
$m_\gamma(A_{\chi,M}^2/B_{\chi,M})L$. The estimate uniform in
height and in the number of jet orders, needed to apply this construction
to growing packets, remains open. The boundary-kernel identity
\eqref{eq:hb-shift-scale} gives another exact expression of the
same multiplicity: at each fixed ordinate,
$\lim_{\omega\downarrow0}K_\omega(\gamma,\gamma)/\omega$
is positive precisely when $m_\gamma=1$. A positive arithmetic
lower bound for this limit at every ordinate has not been proved.

Half balance of the two arithmetic Hurwitz directions also remains
open. Proposition~\ref{prop:periodic-gaussian-gate} gives the finite
Abel estimate that isolates its missing data: fixed odd phase Fourier
means and the mass near the integer phase boundaries. The explicit
phase-measure map proved there preserves the established even Fourier
and normalized Gaussian information while fixing one sign. The required
odd-mode and boundary estimates for the arithmetic labels remain
unproved.

The contour and probe tests have now been executed on the same
positive-ordinate packet through height $20$. The contour calculation
also certifies the eight-zero packet through height $45$.
For larger contours the leading-coefficient bound needs separation
of distinct centers, while its Taylor estimate needs a zero-free
complex-time boundary and a logarithmic-derivative bound there.
The probe needs a spectral isolation bound and a residual small
enough that its divided mass has error below $1/2$.
All these quantities enter explicit inequalities in
Sections~\ref{sec:local} and~\ref{sec:suzuki-probes}.
Suzuki's convergence gives eventual recovery at every fixed
ordinate, and hence eventual simultaneous recovery on each fixed
finite packet. It supplies neither the effective index used in the
present computation nor the conclusion that every recovered integer
equals one. The all-height task is to certify that last conclusion
for packets covering every positive ordinate.

For zeta collisions, the isolated-pair expansion is proved, including
the arithmetic coordinate conversion. The statistical questions
still require control of the exterior interaction
$\delta^2\sum_{\xi\ne x_-,x_+}(b-\xi)^{-2}$ on shrinking gaps,
of nearby triples, and of the remaining simple branches. The
fixed-scale gap results and the restricted correlation identities
have their exact comparisons in
Section~\ref{sec:arithmetic-collisions}; the displayed sine-product
and separated-sequence examples prove the limitations of those
particular statistics. A zeta Weibull or Poisson collision law is
not established by the present results.


\section*{Note on method and records}
The mathematics and text were produced by GPT-6 Astra (ChatGPT Pro mode). Separate instances of the same model independently re-proved specified statements; the records distinguish those events from proof reading, computation, and the use of published background. No human has checked this text. Claude Opus 5.5 at high effort checked the corrected Hurwitz sign, its phase-distribution comparison, and the adjacent CUE agreement formula, including independent numerical calculations for the first and last; it also checked the fixed-multiplicity consequence of inverse compression. Claude Fable 5.2 at maximum effort identified the attribution of that consequence to Suzuki. The records separate these checks from the independent derivations, complete proof readings, and rigorous computations accompanying the present text. We dedicate our text and programs to the public domain under CC0 1.0. The programs and certificates are available through the workbench, \url{\workbenchurl}; the permanent address of the records is \url{\recordurl}.

\begin{thebibliography}{99}
\setlength{\itemsep}{2pt plus 1pt minus 1pt}
\setlength{\parsep}{0pt}
\setlength{\parskip}{0pt}
\bibitem{alpoge} L. Alp\"oge, announcement of the proof of RH and Dirichlet GRH, 20 November 2026, pp.~1, 3--4. \url{https://alpo.ge/remark.pdf}.

\bibitem{lean} Lean FRO and Mathlib Initiative, \emph{Lean verification of the Riemann Hypothesis}, 21 November 2026, items 1--4. \url{https://leanprover-community.github.io/blog/posts/riemann-hypothesis/}.

\bibitem{waterloo} M. Rubinstein, \emph{Alp\"oge/Claude proof of the Generalized Riemann Hypothesis}, verification report, 21--23 November 2026. \url{https://www.math.uwaterloo.ca/~mrubinst/alpoge_claude_grh/}.

\bibitem{roy} M.-F. Roy, \emph{Subdiscriminant of symmetric matrices are sums of squares}, Dagstuhl Seminar Proceedings 05021 (2006), pp.~1--4, Proposition~1, pp.~1--2. \url{https://doi.org/10.4230/DagSemProc.05021.20}.

\bibitem{dlmf} F.~W.~J.~Olver and others, eds., \emph{NIST Digital Library of Mathematical Functions}, Sections~25.11, 25.15 and~5.5: equations~25.11.17, 25.15.2, 25.15.5--6, 25.15.9, 5.5.2 and~5.5.4. \url{https://dlmf.nist.gov/}.

\bibitem{suzuki} M. Suzuki, \emph{Spectral convergence associated with Weil's quadratic form}, version of 26 October 2026, Theorem~3.3, p.~10; equation~(4.1), p.~11; Remark~4.1, p.~12. \url{https://arxiv.org/abs/2610.32704}.

\bibitem{csv} G.~Csordas, W.~Smith and R.~S.~Varga, \emph{Lehmer pairs of zeros, the de Bruijn--Newman constant $\Lambda$, and the Riemann Hypothesis}, Constructive Approximation \textbf{10} (1994), 107--129. Definition~1 and Theorem~1, pp.~109--110. Author-hosted copy: \url{https://www.math.kent.edu/~varga/pub/paper_206.pdf}.

\bibitem{rudnick-sarnak} Z.~Rudnick and P.~Sarnak, \emph{Zeros of principal $L$-functions and random matrix theory}, Duke Mathematical Journal \textbf{81} (1996), 269--322, \url{https://doi.org/10.1215/S0012-7094-96-08115-6}. Theorems~1.1--1.2 on pp.~4--5 of the manuscript at \url{https://empslocal.ex.ac.uk/people/staff/mrwatkin/zeta/rudnick-sarnak.pdf}.

\bibitem{ben-arous-bourgade} G.~Ben Arous and P.~Bourgade, \emph{Extreme gaps between eigenvalues of random matrices}, Annals of Probability \textbf{41} (2013), 2648--2681, \url{https://doi.org/10.1214/11-AOP710}. Theorem~1.1, p.~2650, and Corollary~1.2, p.~2651. Author-hosted published version: \url{https://math.nyu.edu/~bourgade/papers/ExtremeGaps.pdf}.

\bibitem{dobner} A.~Dobner, \emph{A proof of Newman's conjecture for the extended Selberg class}, Acta Arithmetica \textbf{201} (2021), 29--62, \url{https://doi.org/10.4064/aa200603-23-7}; published-version preprint \url{https://arxiv.org/abs/2005.05142v2}, Theorem~4, equations~(10)--(13), pp.~12--13.

\bibitem{arias} J.~Arias de Reyna, \emph{On the distribution (mod 1) of the normalized zeros of the Riemann zeta-function}, Journal of Number Theory \textbf{153} (2015), 37--53, \url{https://doi.org/10.1016/j.jnt.2015.01.006}; \url{https://arxiv.org/abs/1403.7126v1}, Lemmas~2.2 and~2.5, p.~5, and Corollary~3.2, equation~(13), p.~7.

\bibitem{benli-elma-ng} K.~Benli, E.~Elma and N.~Ng,
\emph{A discrete mean value of the Riemann zeta function},
Mathematika \textbf{70} (2024), no.~4, e12281, 47 pp.
\href{https://doi.org/10.1112/mtk.12281}{doi:10.1112/mtk.12281}.
Theorem~1.1(i), equations~(9)--(14), pp.~3--4;
\href{https://arxiv.org/abs/2311.13554}{arXiv:2311.13554}.

\bibitem{heap-li-zhao} W.~Heap, J.~Li and J.~Zhao, \emph{Lower bounds for discrete negative moments of the Riemann zeta function}, Algebra \& Number Theory \textbf{16} (2022), 1589--1625, \url{https://doi.org/10.2140/ant.2022.16.1589}; \url{https://arxiv.org/abs/2003.09368v1}, Theorem~2, equations~(5)--(9), pp.~3--5.

\bibitem{durkan-pearce-crump} B.~Durkan and A.~Pearce-Crump, \emph{Small and large gaps between zeros of Dirichlet $L$-functions}, \url{https://arxiv.org/abs/2610.28527v1} (2026), Theorems~1--2, pp.~1--2.

\bibitem{lamzouri-small-gaps} Y.~Lamzouri, \emph{The distribution of small gaps between zeros of the Riemann zeta function}, \url{https://arxiv.org/abs/2611.14486v1} (2026), Theorem~1.1, p.~2.


\bibitem{bondarenko-heap-half} A.~Bondarenko and W.~Heap, \emph{Half-sized gaps between zeros of the Riemann zeta function}, \url{https://arxiv.org/abs/2610.26038v1} (2026), Theorem~1, p.~1; acknowledgments, p.~2.

\bibitem{lagarias-rodgers} J.~C.~Lagarias and B.~Rodgers, \emph{Higher correlations and the Alternative Hypothesis}, Quarterly Journal of Mathematics \textbf{71} (2020), 257--280; \url{https://arxiv.org/abs/1905.12123v3}, Theorem~2.4, p.~3, and Theorem~3.1, p.~4.

\bibitem{ccm2025} Alain Connes, Caterina Consani, Henri Moscovici, \emph{Zeta Spectral Triples}, 2025. \url{https://arxiv.org/abs/2511.22755}.


\bibitem{fuchs1964} W.~H.~J.~Fuchs, \emph{On the eigenvalues of an integral equation arising in the theory of band-limited signals}, Journal of Mathematical Analysis and Applications \textbf{9} (1964), 317--330, \url{https://doi.org/10.1016/0022-247X(64)90017-4}, theorem on p.~319.

\bibitem{chandee-soundararajan2009} Vorrapan Chandee, K. Soundararajan, \emph{Bounding {$|\zeta(1/2+it)|$} on the {Riemann} hypothesis}, 2009. \url{https://arxiv.org/abs/0908.2008}.


\bibitem{andrews-connes-lagatta2026} Ronnie Andrews, Alain Connes, Antonio La Gatta, \emph{The uncertainty operator, the {Weil} ground state and the zeros of zeta:  a numerical study and conditional positivity of the {Weil} quadratic form}, 2026. \url{https://arxiv.org/abs/2611.13129}.

\bibitem{Edrei1953} A. Edrei, \emph{Proof of a conjecture of Schoenberg on the generating function of a totally positive sequence},
Canad. J. Math. \textbf{5} (1953), 86--94.
\url{https://doi.org/10.4153/CJM-1953-010-3}.

\bibitem{Shi2026} P. Shi, \emph{Coefficient dynamics and consecutive Toeplitz minors for the de Bruijn--Newman family},
hexagon:2610.00521v1 (2026).
\url{https://hexagonmath.org/2610.00521v1}.

\bibitem{Zuber2018} J.-B. Zuber, \emph{Horn's problem and Harish-Chandra's integrals. Probability density functions},
Ann. Inst. Henri Poincar\'e D \textbf{5} (2018), 309--338.
\url{https://doi.org/10.4171/AIHPD/56}.

\bibitem{BB2009} J. Borcea and P. Br\"and\'en, \emph{P\'olya--Schur master theorems for circular domains and their boundaries},
Ann. of Math. \textbf{170} (2009), 465--492.
\url{https://annals.math.princeton.edu/2009/170-1/p14}.

\bibitem{HSW2021} E. Hasanalizade, Q. Shen and P.-J. Wong, \emph{Counting zeros of the Riemann zeta function},
arXiv:2107.06506 (2021).
\url{https://arxiv.org/abs/2107.06506}.

\bibitem{Carnicer2023} J. M. Carnicer, E. Mainar and J. M. Pe\~na,
\emph{Spherical Bessel functions and critical lengths},
Rev. R. Acad. Cienc. Exactas F\'is. Nat. Ser. A Mat. \textbf{117} (2023), article 29.
\url{https://doi.org/10.1007/s13398-022-01365-8}.

\bibitem{Novak2022} J. Novak, \emph{Topological expansion of oscillatory BGW and HCIZ integrals at strong coupling},
arXiv:2203.10746 (2022).
\url{https://arxiv.org/abs/2203.10746}.

\bibitem{LasserrePauwels2019} J.-B. Lasserre and E. Pauwels,
\emph{The empirical Christoffel function with applications in data analysis},
Adv. Comput. Math. \textbf{45} (2019), 1439--1468.
\url{https://doi.org/10.1007/s10444-019-09673-1}.
Preprint: \url{https://arxiv.org/abs/1701.02886}.

\bibitem{Claude2026} Claude (Anthropic), \emph{More than two thirds of the zeros of the Riemann zeta function are simple and on the critical line},
11 August 2026.
\url{https://assets.claude.com/95c246936988e43127bc6b2ceb7077c1dad2d68e.pdf}.

\bibitem{Etingof2011} P. Etingof, O. Golberg, S. Hensel, T. Liu, A. Schwendner,
D. Vaintrob and E. Yudovina,
\emph{Introduction to representation theory},
Student Mathematical Library 59, American Mathematical Society, 2011.
Author-hosted edition:
\url{https://math.mit.edu/~etingof/representationtheorybook.pdf}.

\bibitem{Faraut2007} J. Faraut, \emph{Infinite dimensional spherical analysis},
notes written by S. Matsumoto, Kyushu University, September 2007.
\url{https://webusers.imj-prg.fr/~jacques.faraut/Kyushu.Notes.JF.2007.pdf}.

\bibitem{frisk-degosson} H.~Frisk and S.~de Gosson,
\emph{On the motion of zeros of zeta functions} (2001),
\href{https://arxiv.org/abs/math-ph/0102007}{arXiv:math-ph/0102007}.
Section~2, equation~(3), PDF p.~2.

\bibitem{garunkstis-steuding} R.~Garunk\v stis and J.~Steuding,
\emph{Questions around the nontrivial zeros of the Riemann
zeta-function. Computations and classifications},
Mathematical Modelling and Analysis \textbf{16} (2011), 72--81.
\href{https://doi.org/10.3846/13926292.2011.560616}{doi:10.3846/13926292.2011.560616}.
Section~2, pp.~73--74.

\bibitem{gross-signs} B.~H.~Gross, \emph{Some remarks on signs in functional equations}, Ramanujan Journal \textbf{7} (2003), 91--93. Author manuscript, Section~1, equations~(1.3)--(1.6), p.~2: \url{https://people.math.harvard.edu/~gross/preprints/signs.pdf}.

\bibitem{Johansson2017arb} F. Johansson, \emph{Arb: efficient arbitrary-precision midpoint-radius interval arithmetic},
IEEE Trans. Comput. \textbf{66} (2017), no.~8, 1281--1292.
\url{https://doi.org/10.1109/TC.2017.2690633}.

\bibitem{Johansson2018integration} F. Johansson, \emph{Numerical integration in arbitrary-precision ball arithmetic},
arXiv:1802.07942 (2018).
\url{https://arxiv.org/abs/1802.07942}.

\bibitem{olshanski2013} G.~Olshanski, \emph{Projections of orbital measures, Gelfand--Tsetlin polytopes, and splines}, Journal of Lie Theory \textbf{23} (2013), 1011--1022; \url{https://arxiv.org/abs/1302.7116v1}, Proposition~3.1 and Corollary~3.2, pp.~6--7.


\bibitem{montgomery1973} H.~L.~Montgomery, \emph{The pair correlation of zeros of the zeta function}, Analytic Number Theory, Proceedings of Symposia in Pure Mathematics \textbf{24} (1973), 181--193, Corollary~2, p.~182. \url{https://public.websites.umich.edu/hlm/paircor1.pdf}.
\bibitem{conrey-ghosh-gonek1998} J.~B.~Conrey, A.~Ghosh and S.~M.~Gonek, \emph{Simple zeros of the Riemann zeta-function}, Proceedings of the London Mathematical Society \textbf{76} (1998), 497--522. \url{https://doi.org/10.1112/S0024611598000306}.
\bibitem{bui-heathbrown2013} H.~M.~Bui and D.~R.~Heath-Brown, \emph{On simple zeros of the Riemann zeta-function} (2013), theorem on p.~1. \url{https://arxiv.org/abs/1302.5018v1}.
\bibitem{debruijn1950} N.~G.~de Bruijn, \emph{The roots of trigonometric integrals}, Duke Mathematical Journal \textbf{17} (1950), 197--226, Theorem~13, p.~205. \url{https://doi.org/10.1215/S0012-7094-50-01720-0}.
\bibitem{newman1976} C.~M.~Newman, \emph{Fourier transforms with only real zeros}, Proceedings of the American Mathematical Society \textbf{61} (1976), 245--251, Theorem~3 and Remark~2, p.~247. \url{https://doi.org/10.2307/2041319}.
\bibitem{rodgers-tao2020} B.~Rodgers and T.~Tao, \emph{The de Bruijn--Newman constant is non-negative}, Forum of Mathematics, Pi \textbf{8} (2020), e6. \url{https://doi.org/10.1017/fmp.2020.6}. Corrected author revision: \url{https://arxiv.org/abs/1801.05914v5}, Theorem~1, p.~3.
\bibitem{polymath2019} D.~H.~J.~Polymath, \emph{Effective approximation of heat flow evolution of the Riemann $\xi$ function, and a new upper bound for the de Bruijn--Newman constant} (2019), Theorems~1.1--1.2, pp.~2--3. \url{https://arxiv.org/abs/1904.12438v2}.
\bibitem{griffin-ono-rolen-zagier2019} M.~Griffin, K.~Ono, L.~Rolen and D.~Zagier, \emph{Jensen polynomials for the Riemann zeta function and other sequences} (2019), Theorems~1--3, p.~2. \url{https://arxiv.org/abs/1902.07321v2}.
\bibitem{zhu2026} X.~Zhu, \emph{Weil positivity in compact windows: a finite reduction, certified two-sided bounds, and a Landau--Widom decay law}, version~2 (2 September 2026), Sections~2--4, pp.~7--12; Theorem~6.2, p.~16. \url{https://arxiv.org/abs/2608.24827v2}.

\bibitem{Csordas2015}
G. Csordas, \emph{Fourier transforms of positive definite kernels and the Riemann
\(\xi\)-function}, Comput. Methods Funct. Theory \textbf{15} (2015), 373--391.
\url{https://doi.org/10.1007/s40315-014-0105-8}.
Preprint: \url{https://arxiv.org/abs/1309.0055}.


\bibitem{flint-acb-calc} The FLINT team, \emph{acb\_calc: rigorous complex integration}, documentation, section ``Integration'', accessed 23 November 2026, \url{https://flintlib.org/doc/acb_calc/index.html}.
\bibitem{carneiro-chandee-milinovich2015} E.~Carneiro, V.~Chandee and M.~B.~Milinovich, \emph{A note on the zeros of zeta and $L$-functions}, Mathematische Zeitschrift \textbf{281} (2015), 315--332, \url{https://doi.org/10.1007/s00209-015-1485-9}; \url{https://arxiv.org/abs/1503.00955v1}, Theorem~1, p.~2.
\bibitem{chirre-goncalves2022} A.~Chirre and F.~Gon\c calves, \emph{Bounding the log-derivative of the zeta-function}, Mathematische Zeitschrift \textbf{300} (2022), 1041--1053, \url{https://doi.org/10.1007/s00209-021-02820-9}, Theorems~1--2 and equation~(1.1), pp.~1--2 of \url{https://arxiv.org/abs/2103.06237v2}.

\bibitem{alpoge-furman2026} L.~Alp\"oge and R.~Furman, \emph{More than two thirds of the zeta zeros are simple and on the critical line}, \url{https://arxiv.org/abs/2608.13637v2} (2026), Theorems~A--B, pp.~1--2; provenance on p.~1.

\bibitem{lamzouri-simple2026} Y.~Lamzouri, \emph{A new proof that more than $2/3$ of the zeros of the Riemann zeta function are simple and on the critical line}, \url{https://arxiv.org/abs/2609.02882v2} (2026), Theorem~1.1, p.~3; Lemma~3.2, pp.~13--15, and Remark~3.4, p.~15.
\bibitem{knausgard2026} K.~M.~Knausg\aa rd, \emph{More than 83.9\% of the zeros of the Riemann zeta function are distinct and more than 67.35\% are simple and on the critical line}, \url{https://arxiv.org/abs/2610.08965v1} (2026), Theorems~1.1 and~1.3, pp.~2--3; Section~8, pp.~32--37.
\bibitem{cobeli-zaharescu2026} C.~Cobeli and A.~Zaharescu, \emph{More than 8/9 zeros of the Riemann zeta function are simple or on the critical line}, \url{https://arxiv.org/abs/2610.23467v1} (2026), Theorem~1, p.~2.

\bibitem{mori-rh-alone2026}
R.~Mori, \emph{Windowed Weil forms under RH alone: an
\(e^{-\mu^{6/5+\varepsilon}}\) lower bound and a barrier from distinct
zeros}, version~1, 16 November 2026, Zenodo~24281606;
\url{https://zenodo.org/records/24281606}.
\bibitem{mori-sharp2026}
R.~Mori, \emph{The smallest eigenvalue of windowed Weil forms decays
like \(e^{-4\pi\mu}\) under RH and hypotheses on close pairs of zeros},
version~8, Zenodo~24246078;
\url{https://zenodo.org/records/24246078}.
\bibitem{voss1999}
J.~J.~Voss, On discrete and continuous norms in Paley--Wiener spaces
and consequences for exponential frames,
\emph{J. Fourier Anal. Appl.} \textbf{5} (1999), 193--201;
\url{https://doi.org/10.1007/BF01261609}.

\bibitem{goldston-suriajaya2020}
D.~A.~Goldston and A.~I.~Suriajaya, A singular series average and the
zeros of the Riemann zeta-function, \emph{Acta Arith.}
\textbf{200} (2021), 71--90; preprint
\url{https://arxiv.org/abs/2007.16099}.

\bibitem{platt-trudgian2021} D.~Platt and T.~Trudgian, \emph{The Riemann hypothesis is true up to $3\cdot10^{12}$}, Bulletin of the London Mathematical Society \textbf{53} (2021), 792--797, \url{https://doi.org/10.1112/blms.12460}; preprint \url{https://arxiv.org/abs/2004.09765v1}, Theorem~1 and Section~2, pp.~2--3.

\bibitem{burnol2002} J.-F.~Burnol, \emph{A lower bound in an approximation problem involving the zeros of the Riemann zeta function}, Advances in Mathematics \textbf{170} (2002), 56--70; \url{https://arxiv.org/abs/math/0103058v2}, Theorem~1.3, p.~4; Section~2, pp.~4--6; Theorems~5.3--5.5, pp.~14--16.
\bibitem{bettin-conrey-farmer2013} S.~Bettin, J.~B.~Conrey and D.~W.~Farmer, \emph{An optimal choice of Dirichlet polynomials for the Nyman--Beurling criterion}, \url{https://arxiv.org/abs/1211.5191v1}, Theorem~1, p.~2; Lemmas~2--3 and proof, pp.~6--9.

\bibitem{caceres-margin-v2} P.~C\'aceres, \emph{The Riemann hypothesis as a margin: a two-sided criterion verified to $2\cdot10^{29}$, the fourth-moment anatomy of M\"obius cancellation, and one constant seen from every side}, Hexagon~2610.00688v2 (2026), \url{https://hexagonmath.org/pdf/2610.00688v2}, Part~III, Sections~14, 16--17, pp.~50--52.

\bibitem{bbls-autocorrelation-2003}
L.~B\'aez-Duarte, M.~Balazard, B.~Landreau, and E.~Saias,
\emph{Sur l'autocorr\'elation multiplicative de la fonction
\guillemotleft\ partie fractionnaire\ \guillemotright},
arXiv:math/0306251v1 (2003),
\url{https://arxiv.org/abs/math/0306251v1};
Section~4, pp.~16--17, and Propositions~88--89, pp.~38--39.

\bibitem{maier-rassias-2018}
H.~Maier and M.~Th.~Rassias,
\emph{Explicit estimates of sums related to the Nyman--Beurling
criterion for the Riemann Hypothesis},
Journal of Functional Analysis \textbf{276} (2019), 3832--3857;
\url{https://arxiv.org/abs/1806.05070v1} (2018),
Theorem~2.1, p.~3, cited here only for its displayed summation range.

\bibitem{banks-simple-transform}
W.~D.~Banks, \emph{A variant of the Linnik--Sprind\v zuk theorem
for simple zeros of Dirichlet $L$-functions},
Indagationes Mathematicae \textbf{36} (2025), 1459--1475,
\url{https://doi.org/10.1016/j.indag.2025.05.007};
\url{https://arxiv.org/abs/2410.11605v2},
Section~2, p.~3, and Lemma~2.2, p.~4.

\bibitem{chirre-pereira-delaat2021}
A.~Chirre, V.~J.~Pereira J\'unior, and D.~de Laat,
\emph{Primes in arithmetic progressions and semidefinite programming},
\url{https://arxiv.org/abs/2005.02393v3} (2021),
Section~2.1, equation~(2.1), p.~3.

\bibitem{suzuki-weil2025}
M.~Suzuki, \emph{On the Hilbert space derived from the Weil distribution},
\url{https://arxiv.org/abs/2301.00421v3} (2025),
Sections~2.3--2.4, Proposition~4.1, and Section~6,
pp.~6--7, 10--11, 18.

\bibitem{suzuki-hamiltonians2021}
M.~Suzuki, \emph{Hamiltonians arising from $L$-functions in the Selberg class},
Journal of Functional Analysis \textbf{281} (2021), 109116,
\url{https://doi.org/10.1016/j.jfa.2021.109116};
\url{https://arxiv.org/abs/1606.05726v3},
equations~(2.4)--(2.13), Theorems~2.1--2.4, pp.~6--9.

\bibitem{sra-xi2026}
S.~Sra, \emph{On some positivity properties of the Riemann $\Xi$ function},
author version~v3 (2 November 2026),
\url{https://arxiv.org/abs/2611.02944v1};
\url{https://suvrit.github.io/papers/riem-xi/riem-xi-v3.pdf},
Theorem~1.3 and Section~4, pp.~3, 10--12.


\bibitem{soundararajan-moments2009}
K.~Soundararajan, \emph{Moments of the Riemann zeta function},
Annals of Mathematics \textbf{170} (2009), 981--993,
\url{https://annals.math.princeton.edu/wp-content/uploads/annals-v170-n2-p17-p.pdf},
Lemma~1 and its integrated form, pp.~985--987.

\bibitem{sono-dirichlet2013}
K.~Sono, \emph{Upper bounds for the integral moments of Dirichlet
$L$-functions}, Josai Mathematical Monographs \textbf{6} (2013),
109--124, \url{https://libir.josai.ac.jp/il/user_contents/02/G0000284repository/pdf/JOS-13447777-06_109.pdf},
equations~(2.2)--(2.6), pp.~111--112.

\end{thebibliography}
\end{document}
